\documentclass[11pt,twoside,openany]{book}
\newcommand{\FgaLang}{spanish}
\usepackage[a4paper,inner=31mm,outer=27mm,top=28mm,bottom=31mm,marginparwidth=17mm]{geometry}
\usepackage{fontspec}
\defaultfontfeatures{Ligatures=TeX,Scale=MatchLowercase}
\setmainfont{TeX Gyre Pagella}
\setsansfont{TeX Gyre Heros}
\setmonofont{TeX Gyre Cursor}
\usepackage{polyglossia}
\setdefaultlanguage{\FgaLang}
\usepackage{microtype}
\usepackage{graphicx}
\usepackage{mathtools}
\usepackage{amsthm}
\usepackage{unicode-math}
\setmathfont{Latin Modern Math}
\usepackage{bookmark}
\usepackage{hyperref}
\hypersetup{hidelinks,unicode=true}

\setcounter{secnumdepth}{3}
\setcounter{tocdepth}{2}
\numberwithin{equation}{section}

\newcommand{\src}[1]{%
  \phantomsection\hypertarget{src:#1}{}%
  \marginpar{\raggedleft\footnotesize\sffamily #1}%
}
\newcommand{\sid}[1]{\phantomsection\hypertarget{fga:#1}{}}
\newcommand{\Osh}{\underline{O}}
\newcommand{\Zsh}{\underline{\mathbf Z}}
\newcommand{\Ccat}{\mathcal C}
\newcommand{\Homsh}{\underline{\operatorname{Hom}}}
\newcommand{\Extsh}{\underline{\operatorname{Ext}}}
\newcommand{\hsh}{\underline{h}}
\newcommand{\fquote}[1]{\textquotedbl#1\textquotedbl}
\newcommand{\seminar}[1]{%
  \begin{flushleft}\sffamily SEMINARIO BOURBAKI\\#1\end{flushleft}%
}
\newcommand{\seminaren}[1]{%
  \begin{flushleft}\sffamily BOURBAKI SEMINAR\\#1\end{flushleft}%
}
\newcommand{\paperhead}[2]{%
  \begin{center}
    {\Large\bfseries #1\par}
    \vspace{0.6\baselineskip}
    #2
  \end{center}
  \vspace{0.8\baselineskip}
}

\theoremstyle{plain}
\newtheorem{proposition}{Proposición}[section]
\newtheorem{theoreme}[proposition]{Teorema}
\newtheorem{theorem}[proposition]{Theorem}
\newtheorem{corollaire}[proposition]{Corolario}
\newtheorem{corollary}[proposition]{Corollary}
\newtheorem{lemme}[proposition]{Lema}
\newtheorem{lemma}[proposition]{Lemma}

\theoremstyle{definition}
\newtheorem{definition}[proposition]{Definición}
\newtheorem{definitionen}[proposition]{Definition}

\theoremstyle{remark}
\newtheorem*{remarque}{Observación}
\newtheorem*{remark}{Remark}

\newcommand{\unit}[1]{%
  \chapter{#1}%
  \markboth{#1}{#1}%
}
\newcommand{\fsection}[2]{%
  \section*{#1. #2}%
  \addcontentsline{toc}{section}{#1. #2}%
}
\newcommand{\result}[2]{%
  \par\medskip\noindent\textbf{#1%
  \if\relax\detokenize{#2}\relax\else\space#2\fi. ---}\enspace
}

% Printed page numbers recur across different works and errata.
% Namespace only the invisible PDF destinations; preserve all body bytes.
\newcommand{\FgaUnitKey}{frontmatter}
\renewcommand{\src}[1]{%
  \phantomsection\hypertarget{src:\FgaUnitKey:#1}{}%
  \marginpar{\raggedleft\footnotesize\sffamily #1}%
}
\setlength{\emergencystretch}{3em}
% The only inherited boldsymbol use is the Greek mu of the group scheme.
% Select its Unicode bold-italic glyph instead of a nonexistent bold math font.
\renewcommand{\boldsymbol}[1]{\symbfit{#1}}
\newcommand{\seminares}[1]{\begin{flushleft}\sffamily SEMINARIO BOURBAKI\\#1\end{flushleft}}
\theoremstyle{plain}
\newtheorem{teorema}[proposition]{Teorema}
\newtheorem{corolario}[proposition]{Corolario}
\newtheorem{lema}[proposition]{Lema}
\theoremstyle{definition}
\newtheorem{definicion}[proposition]{Definición}
\theoremstyle{remark}
\newtheorem*{observacion}{Observación}
\hypersetup{pdftitle={Fundamentos de la geometría algebraica -- Edición española actual},pdfauthor={Alexandre Grothendieck},pdfsubject={Ocho exposiciones, sus erratas y los comentarios de 1962}}
\begin{document}
% BEGIN RIGHTS AND AI TRANSLATION NOTICE
\begingroup
\small
\noindent\textbf{Derechos, disponibilidad y autoría}

Los derechos sobre las obras originales y sobre las contribuciones protegidas a estas traducciones y transcripciones derivadas corresponden a sus respectivos titulares. Este proyecto no dedica esas obras o contribuciones al dominio público ni concede una licencia sobre ellas. Las menciones de derechos de terceros y las autorizaciones ya concedidas siguen siendo aplicables.

Estas ediciones se ofrecen gratuitamente, con fines no comerciales y exclusivamente educativos, bajo la presunción de que los respectivos titulares de derechos no se opondrían a que estas traducciones y transcripciones estuvieran disponibles gratuitamente junto a los originales franceses en NUMDAM (www.numdam.org/). Esta presunción no significa que se haya obtenido autorización; el acceso gratuito a los originales no autoriza por sí solo su traducción, adaptación o republicación.

Si usted es titular de derechos sobre algún material presentado aquí y se opone a su disponibilidad, puede abrir una incidencia en GitHub (github.com/KokunoYumeto/ega-fga-es/issues). El responsable de la publicación retirará el material afectado a petición del titular de los derechos.

La traducción al español está producida por IA. Las correcciones, la edición y la composición también están producidas por IA. El responsable de la publicación no tradujo el texto. Se conservan la autoría y las atribuciones de las obras originales y de las contribuciones heredadas; no se atribuyen a la IA. Los modelos documentados se identifican en cada edición.
\par
\endgroup
\clearpage
% END RIGHTS AND AI TRANSLATION NOTICE

\frontmatter
\begin{titlepage}
\centering
\vspace*{0.13\textheight}
{\Huge\bfseries Fundamentos de la\par geometría algebraica\par}
\vspace{1.5\baselineskip}
{\Large Alexandre Grothendieck\par}
\vfill
{\Large Edición española actual\par}
\vspace{\baselineskip}
{\large Ocho exposiciones, sus erratas y los comentarios de 1962\par}
\vspace{1.5\baselineskip}
{\normalsize Edición española\par}
\end{titlepage}
\chapter*{Nota sobre esta edición}
\addcontentsline{toc}{chapter}{Nota sobre esta edición}
Este volumen contiene la traducción española de las exposiciones 149, 182, 190, 195, 212, 221, 232 y 236, sus erratas y los comentarios de 1962.

La edición actual y la traducción según el impreso se ofrecen como textos distintos.

Las cifras marginales remiten a las páginas de las exposiciones originales.

Traducción, correcciones y composición producidas por IA: OpenAI Codex, GPT-6 Astra, esfuerzo Ultra.
La autoría de la obra original corresponde a Alexandre Grothendieck.

El PDF se acompaña de fuentes LaTeX editables y de los archivos necesarios para su compilación.

\tableofcontents
\mainmatter
\def\FgaUnitKey{149}\unit{Exposición 149 — Teoremas de dualidad para los haces algebraicos coherentes}
\sid{149}
\src{169}
\seminar{(Mayo de 1957)}
\paperhead{TEOREMAS DE DUALIDAD PARA LOS HACES ALGEBRAICOS COHERENTES}%
  {por Alexandre GROTHENDIECK}

Los resultados que siguen, inspirados en el \fquote{teorema de dualidad algebraica}
de Serre, fueron obtenidos en los inviernos de 1955 y 1956. Se demuestran muy
sencillamente con ayuda de resultados bastante elementales sobre la cohomología de los espacios
proyectivos [3] y mediante el uso intensivo del álgebra homológica de
Cartan--Eilenberg, en la forma de [2].

\fsection{1}{Los Ext de haces de módulos ({[2]}, caps. 3 y 4)}

Sea $X$ un espacio topológico provisto de un haz $\Osh$ de anillos con unidad
(no necesariamente conmutativos). Consideramos la categoría abeliana
$\Ccat^{\Osh}$ de los haces de $\Osh$-módulos, llamados también $\Osh$-Módulos.
Se sabe que todo objeto de esta categoría admite una resolución inyectiva, lo que
permite definir en ella los funtores $\operatorname{Ext}$ con las propiedades
formales bien conocidas. Más precisamente, para evitar confusiones,
denotamos por $\operatorname{Hom}_{\Osh}(X;A,B)$, o simplemente por
$\operatorname{Hom}(X;A,B)$, el grupo abeliano de los $\Osh$-homomorfismos de $A$
en $B$, mientras que $\Homsh_{\Osh}(A,B)$ denotará el haz de gérmenes
de homomorfismos de $A$ en $B$ ($A,B\in\Ccat^{\Osh}$). Para
$A\in\Ccat^{\Osh}$ fijo se definen funtores $h_A$ y $\hsh_A$ con valores, respectivamente,
en la categoría $\Ccat$ de los grupos abelianos y en la categoría
$\Ccat^{\Zsh}$ de los haces abelianos sobre $X$, mediante las fórmulas:
\begin{equation}\tag{1.1}
  h_A(B)=\operatorname{Hom}_{\Osh}(X;A,B),
  \qquad
  \hsh_A(B)=\Homsh_{\Osh}(A,B).
\end{equation}

Los funtores $h_A$ y $\hsh_A$ son funtores covariantes exactos por la izquierda,
de los cuales consideramos los funtores derivados derechos, denotados respectivamente por
$\operatorname{Ext}_{\Osh}^{p}(X;A,B)$ y $\Extsh_{\Osh}^{p}(A,B)$. Por tanto,
por definición,
\begin{equation}\tag{1.2}
\left\{
\begin{aligned}
\operatorname{Ext}_{\Osh}^{p}(X;A,B)
  &= (R^{p}h_A)(B)
   = H^{p}\bigl(\operatorname{Hom}_{\Osh}(X;A,C(B))\bigr),\\
\Extsh_{\Osh}^{p}(A,B)
  &= (R^{p}\hsh_A)(B)
   = H^{p}\bigl(\Homsh_{\Osh}(A,C(B))\bigr),
\end{aligned}
\right.
\end{equation}

donde el símbolo $R^{p}$ designa el paso a los funtores derivados derechos, y donde
$C(B)$ designa una resolución inyectiva arbitraria de $B$ en $\Ccat^{\Osh}$.
Denotemos por
\[
  \Gamma:\Ccat^{\Zsh}\longrightarrow\Ccat
\]
el funtor \fquote{secciones}. Recordemos que sus funtores derivados derechos se
denotan por
\[
  B\longmapsto H^{p}(X,B).
\]

\src{170}
\begin{equation}\tag{1.3}
  H^{p}(X,B)=(R^{p}\Gamma)(B)=H^{p}\bigl(\Gamma(C(B))\bigr).
\end{equation}
Evidentemente, se tiene $h_A=\Gamma\hsh_A$; por otra parte, se comprueba que $\hsh_A$
transforma objetos inyectivos en objetos acíclicos para $\Gamma$. Se concluye del
modo habitual:

\result{PROPOSICIÓN}{1} Para todo $\Osh$-Módulo $A$ existe un funtor
espectral cohomológico sobre $\Ccat^{\Osh}$, que converge al funtor graduado
$\bigl(\operatorname{Ext}_{\Osh}^{*}(X;A,B)\bigr)$ y cuyo término inicial es
\begin{equation}\tag{1.4}
  E_{2}^{p,q}(A,B)=H^{p}\bigl(X,\Extsh_{\Osh}^{q}(A,B)\bigr).
\end{equation}
De ella se deducen \fquote{homomorfismos de borde} y una sucesión exacta de cinco términos
que no escribiremos.

\result{COROLARIO}{1} Si $A$ es localmente isomorfo a $\Osh^{n}$, entonces se tienen
isomorfismos canónicos
\begin{equation}\tag{1.5}
  \operatorname{Ext}_{\Osh}^{p}(X;A,B)
  \overset{\sim}{\longleftarrow}
  H^{p}\bigl(X,\Homsh_{\Osh}(A,B)\bigr)
\end{equation}
(dados por homomorfismos de borde de la sucesión espectral). En particular, se tiene
un isomorfismo canónico
\begin{equation}\tag{1.6}
  \operatorname{Ext}_{\Osh}^{p}(X;\Osh,B)=H^{p}(X,B).
\end{equation}

Para utilizar estos resultados, es preciso saber explicitar los
$\Extsh_{\Osh}^{p}(A,B)$. Son funtores que se calculan localmente,
es decir, si $U$ es un abierto de $X$, se tiene
\[
  \Extsh_{\Osh}^{p}(A,B)|U
  =\Extsh_{\Osh|U}^{p}(A|U,B|U),
\]
como resulta del hecho de que la restricción a $U$ de un $\Osh$-Módulo inyectivo
es un $(\Osh|U)$-Módulo inyectivo. Además, para $x\in X$ fijo, se tienen homomorfismos funtoriales
\begin{equation}\tag{1.7}
  \Homsh_{\Osh}(A,B)_x
  \longrightarrow
  \operatorname{Hom}_{\Osh_x}(A_x,B_x),
\end{equation}
que se prolongan de manera única a un homomorfismo de funtores cohomológicos
en $B$:
\begin{equation}\tag{1.8}
  \Extsh_{\Osh}^{p}(A,B)_x
  \longrightarrow
  \operatorname{Ext}_{\Osh_x}^{p}(A_x,B_x).
\end{equation}

\result{PROPOSICIÓN}{2} Si $A$ es, en un entorno de $x$, isomorfo al
conúcleo de un homomorfismo $\Osh^{m}\longrightarrow\Osh^{n}$, entonces (1.8) es un isomorfismo en grado $p$ cuando $A$ admite localmente una resolución libre parcial hasta el grado $p+1$; es un isomorfismo para todo $p$ cuando $A$ es pseudocoherente. Una sola presentación finita da únicamente la comparación en grado cero.\footnote{El texto impreso dice « (1.7) »; la
presencia de $p$ y la frase anterior exigen « (1.8) ».} Esto ocurre, en
particular, si $A$ es un $\Osh$-Módulo coherente [3].

\src{171}
\result{PROPOSICIÓN}{3} Sea $L_{*}=(L_i)$ una resolución por la izquierda del
$\Osh$-Módulo $A$ por $\Osh$-Módulos cada uno de los cuales es localmente isomorfo a
un $\Osh^{n}$. Entonces $\Extsh_{\Osh}^{*}(A,B)$ se identifica con
$H^{*}\bigl(\Homsh_{\Osh}(L_{*},B)\bigr)$, y
$\operatorname{Ext}_{\Osh}^{*}(X;A,B)$ se identifica con la hipercohomología de $X$
con respecto al complejo $\Homsh_{\Osh}(L_{*},B)$.

La demostración es estándar: se considera el bicomplejo
$\Homsh_{\Osh}(L_{*},C(B))$, donde $C(B)$ es una resolución inyectiva de $B$, y los
homomorfismos naturales de $\Homsh_{\Osh}(L_{*},B)$ y de
$\Homsh_{\Osh}(A,C(B))$ en este último.

Para terminar, señalemos que los dos funtores $\operatorname{Ext}$ introducidos
en (1.2) no solo son funtores cohomológicos en $B$, sino también
bifuntores cohomológicos, covariantes en $B$ y contravariantes en $A$.

\fsection{2}{Ley de composición en los Ext}

Los resultados de este número se deben, independientemente, a CARTIER y a YONEDA; véase
una exposición de CARTIER [1] para más detalles. Sea $\Ccat$ una categoría abeliana.
Sean $K$ y $L$ dos objetos graduados de $\Ccat$. Denotamos por
$\operatorname{Hom}(K,L)$ el grupo abeliano graduado cuya componente de grado
$n$ está formada por los homomorfismos homogéneos de grado $n$ de $K$ en $L$ (es decir,
los sistemas $(u_i)$ de homomorfismos $K^{i}\longrightarrow L^{i+n}$).

Si $K$ y $L$ son complejos (con operadores diferenciales de grado $+1$, para
fijar ideas), se introduce en $\operatorname{Hom}(K,L)$ el operador
diferencial dado por
\begin{equation}\tag{2.1}
  \delta(u)=du+(-1)^{n+1}ud,\qquad\text{donde }n=\deg(u),
\end{equation}
que lo convierte en un complejo con operador diferencial de grado $+1$. Los ciclos de
grado $n$ son los operadores de grado $n$ que anticonmutan con $d$ (como
operadores homogéneos).\footnote{El texto impreso dice « con $u$ »; la ecuación
(2.1) exige « con $d$ ».} Se puede entonces considerar
$H^{*}\bigl(\operatorname{Hom}(K,L)\bigr)$; este es un invariante de los tipos
de homotopía de $K$ y de $L$, que podremos denotar por $H^{*}(K,L)$.

Si se tiene un tercer complejo $M$, la composición de los homomorfismos
define un acoplamiento
\[
  \operatorname{Hom}(K,L)\times\operatorname{Hom}(L,M)
  \longrightarrow\operatorname{Hom}(K,M)
\]
compatible con los operadores diferenciales, de donde, al pasar a la cohomología,
se obtienen acoplamientos
\begin{equation}\tag{2.2}
  H^{*}(K,L)\times H^{*}(L,M)\longrightarrow H^{*}(K,M),
\end{equation}
denotados por $(u,v)\longmapsto vu$. Estos acoplamientos satisfacen la evidente propiedad
de asociatividad. En particular, $H^{*}(K,K)$ es un anillo graduado asociativo
con unidad, y $H^{*}(K,L)$ (resp. $H^{*}(L,K)$) es un módulo graduado por la derecha
(resp. por la izquierda) sobre este anillo, etc. En grado $0$, (2.2) se reduce a la
composición de los morfismos de complejos. Finalmente, una sucesión exacta de
complejos
\[
  0\longrightarrow K'\longrightarrow K\longrightarrow K''\longrightarrow 0,
\]
\src{172}
tal que $K'^{i}$ se identifica con un sumando directo de $K^{i}$ para todo $i$,
da lugar a una sucesión exacta de complejos de grupos
$\operatorname{Hom}(K'',L)$, etc., y de ahí a un operador de conexión
\[
  H^{i}(K',L)\longrightarrow H^{i+1}(K'',L).
\]
Se definen de manera análoga los operadores de conexión relativos a una sucesión exacta en $L$. Los
acoplamientos (2.2) son compatibles, en el sentido usual, con estos operadores de conexión.

Supongamos ahora que $\Ccat$ es una categoría tal que todo objeto $A$
de $\Ccat$ admite una resolución inyectiva $C(A)$. Se comprueba entonces, utilizando
una de las numerosas variantes del teorema del bicomplejo, que
\[
  H^{*}\bigl(C(A),C(B)\bigr)
  =H^{*}\bigl(\operatorname{Hom}(C(A),C(B))\bigr)
\]
es canónicamente isomorfo a
\[
  H^{*}\bigl(\operatorname{Hom}(A,C(B))\bigr)
  =\operatorname{Ext}^{*}(A,B).
\]
Los operadores de conexión considerados anteriormente dan los homomorfismos de conexión de los
$\operatorname{Ext}$. Además, los acoplamientos (2.2) dan aquí
acoplamientos asociativos, compatibles con los homomorfismos de conexión:
\begin{equation}\tag{2.3}
  \operatorname{Ext}^{*}(A,B)\times\operatorname{Ext}^{*}(B,C)
  \longrightarrow\operatorname{Ext}^{*}(A,C).
\end{equation}
En particular, $\operatorname{Ext}^{*}(A,A)$ es un anillo graduado asociativo
con unidad, etc. (Se demuestra de manera análoga que los funtores
$\operatorname{Ext}$ operan en los funtores derivados de un funtor
cualquiera; no utilizaremos aquí este hecho.)

En el caso en que la categoría considerada sea la categoría $\Ccat^{\Osh}$ de los
$\Osh$-Módulos sobre $X$, se obtienen, por tanto, acoplamientos
\begin{equation}\tag{2.4}
  \operatorname{Ext}_{\Osh}^{p}(X;A,B)
  \times\operatorname{Ext}_{\Osh}^{q}(X;B,C)
  \longrightarrow\operatorname{Ext}_{\Osh}^{p+q}(X;A,C),
\end{equation}
que pueden explicitarse como se ha indicado. Por lo demás, los mismos
desarrollos, sustituyendo la categoría de los grupos abelianos por la
categoría de los haces abelianos sobre $X$, y los funtores $\operatorname{Hom}$
por los funtores $\Homsh$, definen también acoplamientos con las mismas
propiedades formales y, esta vez, de \fquote{naturaleza local}:
\begin{equation}\tag{2.5}
  \Extsh_{\Osh}^{p}(A,B)\times\Extsh_{\Osh}^{q}(B,C)
  \longrightarrow\Extsh_{\Osh}^{p+q}(A,C).
\end{equation}
Estos últimos se precisan observando que los homomorfismos (1.8) son
compatibles con los acoplamientos entre los $\operatorname{Ext}$.

\src{173}
Recordemos, por último, que existe también una estructura multiplicativa entre los funtores
$H^{p}(X,A)$ (producto cup). Se comprueba entonces que las sucesiones espectrales de la
proposición 1 son compatibles con las estructuras multiplicativas; más
precisamente, se tiene un acoplamiento de la sucesión espectral $E(A,B)$ y de la sucesión
espectral $E(B,C)$ hacia la sucesión espectral $E(A,C)$ que, en el límite,
se reduce al acoplamiento entre los $\operatorname{Ext}$ globales y, en el
término inicial, al acoplamiento deducido, en el segundo miembro de (1.4), del
producto cup y de los acoplamientos de los $\operatorname{Ext}$ locales. De ello resulta,
en particular, que los \fquote{homomorfismos de borde}
\begin{equation}\tag{2.6}
  \operatorname{Ext}_{\Osh}^{n}(X;A,B)
  \longrightarrow H^{0}\bigl(X,\Extsh_{\Osh}^{n}(A,B)\bigr)
\end{equation}
y
\begin{equation}\tag{2.7}
  H^{n}\bigl(X,\Homsh_{\Osh}(A,B)\bigr)
  \longrightarrow\operatorname{Ext}_{\Osh}^{n}(X;A,B)
\end{equation}
son compatibles con las estructuras multiplicativas. Por tanto, si nos restringimos a los
haces localmente isomorfos a un $\Osh^{m}$, esto explicita por completo la
composición entre los $\operatorname{Ext}$ globales mediante el producto cup, teniendo
en cuenta el isomorfismo (1.5).

\fsection{3}{Resultados de cohomología local}

Sea $A$ un anillo conmutativo con unidad provisto de un ideal $J$. Vamos a
definir, para un $A$-módulo variable $M$, homomorfismos funtoriales
\begin{equation}\tag{3.1}
\left\{
\begin{aligned}
\operatorname{Ext}_{A}^{p}(A/J,M)
  &\longrightarrow
  \operatorname{Hom}_{A}\!\left(\bigwedge^{p}J/J^{2},M\otimes A/J\right),\\
\operatorname{Tor}_{p}^{A}(A/J,M)
  &\longleftarrow
  \left(\bigwedge^{p}J/J^{2}\right)
  \otimes\operatorname{Hom}_{A}(A/J,M),
\end{aligned}
\right.
\end{equation}
donde los productos tensoriales y exteriores se toman sobre el anillo $A$; nótese,
por otra parte, que $J/J^{2}$ es en realidad un $A/J$-módulo, y que sus potencias
exteriores como $A$-módulo o como $A/J$-módulo son las mismas.
Definir los homomorfismos (3.1) equivale a definir, para todo
sistema $\underline{x}=(x_1,\ldots,x_p)$ de elementos de $J$, homomorfismos
$\varphi_{\underline{x}}$:
\begin{equation}\tag{3.2}
\left\{
\begin{aligned}
\varphi_{\underline{x}}:
  \operatorname{Ext}_{A}^{p}(A/J,M)&\longrightarrow M\otimes A/J,\\
\varphi_{\underline{x}}:
  \operatorname{Hom}_{A}(A/J,M)&\longrightarrow
  \operatorname{Tor}_{p}^{A}(A/J,M),
\end{aligned}
\right.
\end{equation}
de tal manera que se satisfagan las condiciones siguientes:

\src{174}
\begin{enumerate}
\item[\textup{i.}] $\varphi_{x_1,\ldots,x_p}$ depende de manera
$A$-multilineal alternada del sistema de los $x_i\in J$.
\item[\textup{ii.}] $\varphi_{x_1,\ldots,x_p}$ es nulo cuando alguno de los $x_i$ está
en $J^{2}$.
\end{enumerate}
En realidad, ii. resulta de i., pues $a\varphi_{\underline{x}}=0$ para $a\in J$,
como se ve observando que todos los módulos de (3.2) son anulados por
$J$.

Para definir $\varphi_{\underline{x}}$, consideremos el complejo
$K_{\underline{x}}$ cuyo $A$-módulo subyacente es $\bigwedge A^{p}$, y
cuyo operador diferencial es el producto interior $i_{\underline{x}}$ por
$\underline{x}$ considerado como forma lineal sobre $A^{p}$ de componentes
$x_1,\ldots,x_p$. El operador diferencial es de grado $-1$, los grados son
positivos, y la homología de este complejo en grado $0$ es
$A/(x_1A+\cdots+x_pA)$. Como los $x_i$ están en $J$, se deduce una
aumentación $K_{\underline{x},0}\longrightarrow A/J$. Así,
$K_{\underline{x}}$ aparece como un complejo libre aumentado, con módulo
de aumentación $A/J$. De ello se deducen, del modo habitual, homomorfismos
\[
\operatorname{Ext}_{A}^{*}\bigl(H_0(K_{\underline{x}}),M\bigr)
\longrightarrow
H^{*}\bigl(\operatorname{Hom}_{A}(K_{\underline{x}},M)\bigr)
\]
y
\[
\operatorname{Tor}_{*}^{A}\bigl(H_0(K_{\underline{x}}),M\bigr)
\longleftarrow
H_{*}\bigl(K_{\underline{x}}\otimes_A M\bigr),
\]
de donde, componiendo con los homomorfismos sobre los $\operatorname{Ext}$ y los
$\operatorname{Tor}$ deducidos del homomorfismo de aumentación
$H_0(K_{\underline{x}})\longrightarrow A/J$, se obtienen homomorfismos
\begin{equation}\tag{3.3}
\left\{
\begin{aligned}
\psi_{\underline{x}}:
\operatorname{Ext}_{A}^{*}(A/J,M)
&\longrightarrow H^{*}\bigl(\operatorname{Hom}_{A}(K_{\underline{x}},M)\bigr),\\
\psi_{\underline{x}}:
\operatorname{Tor}_{*}^{A}(A/J,M)
&\longleftarrow H_{*}\bigl(K_{\underline{x}}\otimes_A M\bigr).
\end{aligned}
\right.
\end{equation}
Ahora bien, se comprueba inmediatamente que, en grado máximo $p$, la cohomología de los segundos
miembros es $M/(x_1M+\cdots+x_pM)$ (resp. es el conjunto de los elementos de $M$
anulados por cada uno de los $x_i$). Como los $x_i$ están en $J$, se deducen
homomorfismos
\begin{equation}\tag{3.4}
\left\{
\begin{aligned}
H^{p}\bigl(\operatorname{Hom}_{A}(K_{\underline{x}},M)\bigr)
&\longrightarrow M\otimes A/J,\\
H_{p}\bigl(K_{\underline{x}}\otimes_A M\bigr)
&\longleftarrow\operatorname{Hom}_{A}(A/J,M).
\end{aligned}
\right.
\end{equation}
Componiendo los homomorfismos (3.3) y (3.4), se obtienen los homomorfismos (3.2)
que queríamos definir. La comprobación de i. es laboriosa, pero no
presenta dificultades.

\src{175}
\result{PROPOSICIÓN}{4} Sea $A$ un anillo conmutativo con unidad, y sea
$(x_1,\ldots,x_p)$ una sucesión de elementos de $A$ tal que, para $1\leq i\leq p$,
la imagen de $x_i$ en el cociente de $A$ por el ideal generado por
$(x_1,\ldots,x_{i-1})$ no sea divisor de $0$. Sea $J$ el ideal generado
por los $x_i$. Entonces $J/J^{2}$ es un $(A/J)$-módulo libre que tiene por base las
imágenes canónicas de los $x_i$, el complejo $K_{\underline{x}}$ es una resolución
libre de $A/J$, y para todo $A$-módulo $M$, los homomorfismos (3.1) en
grado $p$ son biyectivos. Lo mismo ocurre con los homomorfismos análogos
definidos para grados $i$ cualesquiera, siempre que $J.M=0$.

(El punto esencial, del cual resultan todos los demás, es la aciclicidad de
$K_{\underline{x}}$, que es un hecho bien conocido bajo las condiciones indicadas.)

\result{COROLARIO}{1} Sean $A$ y $J$ como arriba, y supongamos además que
$A$ es un álgebra afín regular de $\dim n$ sobre un cuerpo perfecto $k$, y
que $A/J$ es un álgebra afín regular. Denotemos por $\Omega^{1}(A)$,
$\Omega^{1}(A/J)$ los módulos de diferenciales de Kähler. Entonces se tiene un
isomorfismo canónico compatible con la localización,
\begin{equation}\tag{3.5}
  \operatorname{Ext}_{A}^{p}\bigl(\Omega^{n-p}(A/J),\Omega^{n}(A)\bigr)=A/J.
\end{equation}
En efecto, $\Omega^{n-p}(A/J)$ es un $(A/J)$-módulo proyectivo de rango $1$, y asimismo
$\Omega^{n}(A)$ es un $A$-módulo proyectivo de rango $1$,\footnote{El texto impreso
dice « rango $n$ »; los dos módulos de formas de grado máximo tienen rango
$1$, como exige también la identificación con $\bigwedge^{p}(J/J^{2})$ más abajo.}
por lo que el primer miembro es igual a
\[
  \operatorname{Ext}_{A}^{p}(A/J,A)
  \otimes\Omega^{n-p}(A/J)'\otimes\Omega^{n}(A)
\]
(donde el símbolo « $'$ » designa el $(A/J)$-módulo dual). El producto tensorial de los
dos últimos factores se identifica con $\bigwedge^{p}(J/J^{2})$, por lo que el conjunto
se identifica con
\[
  \operatorname{Ext}_{A}^{p}\!\left(A/J,\bigwedge^{p}(J/J^{2})\right),
\]
y, por la proposición, con
\[
  \operatorname{Hom}_{A}\!\left(\bigwedge^{p}J/J^{2},
  \bigwedge^{p}J/J^{2}\right),
\]
es decir, con $A/J$. En particular, en
$\operatorname{Ext}_{A}^{p}\bigl(\Omega^{n-p}(A/J),\Omega^{n}(A)\bigr)$ hay un
elemento distinguido, correspondiente al elemento unidad de $A/J$, llamado
\emph{clase fundamental del ideal $J$ en $A$}. (En realidad, puede
definirse bajo condiciones sensiblemente más generales.) El
corolario 1 puede darse en una forma más geométrica y más global:

\result{COROLARIO}{2} Sean $X$ una variedad no singular sobre un cuerpo
algebraicamente cerrado $k$, $Y$ una subvariedad cerrada no singular de $X$,
$\Osh_X$ el haz de anillos locales de $X$, y $\Osh_Y$ el haz de anillos
locales de $Y$ considerado como haz cociente de $\Osh_X$. Sean $n$ la
dimensión de $X$ y $n-p$ la de $Y$.

\src{176}
Sea $\Omega_X$ (resp. $\Omega_Y$) el haz de gérmenes de formas
diferenciales regulares sobre $X$ (resp. $Y$). Entonces se tienen isomorfismos
canónicos:
\begin{equation}\tag{3.6}
  \Extsh_{\Osh_X}^{p}\bigl(\Omega_Y^{n-p},\Omega_X^{n}\bigr)=\Osh_Y,
\end{equation}
o también
\begin{equation}\tag{3.6 bis}
  \Extsh_{\Osh_X}^{p}\bigl(\Osh_Y,\Omega_X^{n}\bigr)=\Omega_Y^{n-p}.
\end{equation}
La fórmula (3.6 bis) puede servir como definición de $\Omega_Y^{n-p}$ cuando $Y$
es una variedad singular. Más precisamente:

\result{PROPOSICIÓN}{5} Sea $Y$ un subconjunto algebraico de dimensión
$q=n-p$ en una variedad algebraica no singular $X$ de dimensión $n$. Sean
$F$ un haz algebraico coherente sobre $X$ cuyo soporte está contenido en $Y$, y $L$
un haz algebraico localmente libre sobre $X$. Entonces los haces
$\Extsh_{\Osh_X}^{i}(F,L)$ son nulos para $i<p$, mientras que para $i=p$ se tiene un
isomorfismo canónico
\begin{equation}\tag{3.7}
  \Extsh_{\Osh_X}^{p}(F,L)
  =\Homsh_{\Osh_X}\!\left(F,
  \Extsh_{\Osh_X}^{p}(\Osh_X/J,L)\right),
\end{equation}
donde $J$ designa un haz de ideales arbitrario sobre $X$ que anula $F$ y cuyo
conjunto de ceros es $Y$. En particular, si $F$ es un haz algebraico
coherente sobre $Y$, se tiene
\begin{equation}\tag{3.7 bis}
  \Extsh_{\Osh_X}^{p}(F,L)
  =\Homsh_{\Osh_Y}\!\left(F,\Extsh_{\Osh_X}^{p}(\Osh_Y,L)\right).
\end{equation}
Por último, siendo todavía $F$ un haz algebraico coherente sobre $Y$, los haces
\[
  E^{i}(F)=\Extsh_{\Osh_X}^{p+i}(F,\Omega_X^{n})
\]
no dependen de la manera de sumergir el espacio algebraico $Y$ en una variedad
no singular $X$.

Como la cuestión es local, se puede suponer que $X$ es afín y $L=\Osh_X$. Esto reduce
el problema a una cuestión de álgebra conmutativa, e incluso de álgebra local: si $A$ es un
anillo local regular y $M$ un $A$-módulo cuyo soporte tiene dimensión
$\leq q=n-p$, hay que demostrar que $\operatorname{Ext}_{A}^{i}(M,A)=0$ para $i<p$, y
\[
  \operatorname{Ext}_{A}^{p}(M,A)
  =\operatorname{Hom}_{A}\!\left(M,\operatorname{Ext}_{A}^{p}(A/J,A)\right),
\]
donde $J$ es un ideal cualquiera de « dimensión » $\leq q$ que anula $M$. Para el
primer punto se procede por inducción sobre $q$: un devissage inmediato reduce
al caso en que $M$ es de la forma $A/J$ y luego, sustituyendo $J$ por un ideal más
pequeño y utilizando la hipótesis de inducción y la sucesión exacta de los
$\operatorname{Ext}$, al caso en que $J$ está generado por un « sistema de parámetros »
como en la proposición 4, donde el resultado es inmediato. El resultado
anterior implica que, si $J$ es un ideal

\src{177}
fijo de « dimensión » $\leq q$, entonces el funtor contravariante
$E(M)=\operatorname{Ext}_{A}^{p}(M,A)$ sobre la categoría de los $A/J$-módulos es
exacto por la izquierda; además, transforma sumas directas en productos directos,
de donde resulta fácilmente que
\[
  E(M)=\operatorname{Hom}_{A}\bigl(M,E(A/J)\bigr).
\]
\footnote{El texto impreso dice $E(A)$; en la categoría considerada, el objeto
representante es $A/J$, y la igualdad requerida para (3.7) es $E(A/J)$.} Por último, la
última afirmación de la proposición 5 es más sutil y resulta de una
caracterización intrínseca de los $E^{i}(F)$ mediante un teorema de dualidad
local que no puede exponerse aquí.

\result{COROLARIO}{} Denotemos por $\omega_Y^{q}$ el haz
$\Extsh_{\Osh_X}^{p}(\Osh_Y,\Omega_X^{n})$. Entonces se tiene un isomorfismo
funtorial para los haces algebraicos coherentes $F$ sobre $Y$:
\begin{equation}\tag{3.8}
  \Extsh_{\Osh_X}^{p}(F,\Omega_X^{n})
  =\Homsh_{\Osh_X}(F,\omega_Y^{q}).
\end{equation}

\fsection{4}{Clase de cohomología asociada a una subvariedad}

En todo lo que sigue, $X$ designa un conjunto algebraico de dimensión $n$, definido
sobre un cuerpo $k$, que supondremos algebraicamente cerrado para simplificar. Salvo
en el n\textsuperscript{o} 6, se supone que $X$ es no singular. Denotamos por
$\Osh_X$ el haz de anillos locales sobre $X$, y por
\[
  \Omega_X^{*}=\bigoplus_p\Omega_X^{p}
\]
el haz de gérmenes de formas diferenciales sobre $X$.\footnote{El texto
impreso emplea un signo de unión; la estructura graduada es la suma directa
de los haces $\Omega_X^p$.} Si $Y$ es una parte cerrada de $X$, identificamos
los haces algebraicos coherentes sobre $Y$ con haces coherentes sobre $X$
nulos fuera de $Y$; en particular, así ocurre con $\Osh_Y$ y $\Omega_Y$.

\result{LEMA}{1} Sean $F$ un haz algebraico coherente sobre $X$ cuyo
soporte tiene dimensión $\leq n-p$, y $L$ un haz algebraico coherente
localmente libre sobre $X$. Entonces $\operatorname{Ext}_{\Osh_X}^{i}(X;F,L)$ es nulo
para $i<p$, y se tiene un isomorfismo canónico
\begin{equation}\tag{4.1}
  \operatorname{Ext}_{\Osh_X}^{p}(X;F,L)
  =H^{0}\bigl(X,\Extsh_{\Osh_X}^{p}(F,L)\bigr).
\end{equation}
Si $F$ es un haz algebraico coherente sobre una parte cerrada $Y$ de $X$ de
dimensión $\leq n-p$, se tiene un isomorfismo canónico
\begin{equation}\tag{4.1 bis}
  \Extsh_{\Osh_X}^{p}(F,L)
  =\Homsh_{\Osh_X}\!\left(
    F\otimes L'\otimes\Omega_X^{n},\omega_Y^{n-p}
  \right),
\end{equation}
donde $\omega_Y^{n-p}$ es el haz sobre $Y$ definido en el corolario de la
proposición 5 (que se identifica con $\Omega_Y^{n-p}$ si $Y$ es no singular).

La fórmula (4.1) es una consecuencia inmediata de la sucesión espectral de la
proposición 1 y de la proposición 5; en virtud de la fórmula (3.8), se puede
escribir

\src{178}
\begin{align*}
\Extsh_{\Osh_X}^{p}(F,L)
&=L\otimes(\Omega_X^{n})'\otimes
  \Extsh_{\Osh_X}^{p}(F,\Omega_X^{n})\\
&=L\otimes(\Omega_X^{n})'\otimes
  \Homsh_{\Osh_X}(F,\omega_Y^{q})\\
&=\Homsh_{\Osh_X}\!\left(
  F\otimes L'\otimes\Omega_X^{n},\omega_Y^{q}\right),
\end{align*}
donde $q=n-p$, de donde resulta la fórmula (4.1 bis).

Tomando, en particular, $F=\Osh_Y$, $L=\Omega_X^{p}$, se obtiene, teniendo en cuenta
que $\Omega_X^{n}\otimes(\Omega_X^{p})'=\Omega_X^{n-p}$, un isomorfismo
canónico
\begin{equation}\tag{4.2}
  \operatorname{Ext}_{\Osh_X}^{p}(X;\Osh_Y,\Omega_X^{p})
  =\operatorname{Hom}_{\Osh_X}
  (X;\Omega_X^{n-p},\omega_Y^{n-p}).
\end{equation}
Supongamos ahora que $Y$ es no singular, para simplificar, de modo que
$\omega_Y^{n-p}=\Omega_Y^{n-p}$. Se tiene un homomorfismo natural de
$\Omega_X^{n-p}$ en $\Omega_Y^{n-p}$, de donde resulta una sección canónica $s_Y$ del
haz $\Extsh_{\Osh_X}^{p}(\Osh_Y,\Omega_X^{p})$, que llamaremos, si
todas las componentes de $Y$ tienen dimensión $n-p$, \emph{sección
fundamental} del haz $\Extsh_{\Osh_X}^{p}(\Osh_Y,\Omega_X^{p})$. Esta
sección define, en virtud de (4.1), un elemento de
$\operatorname{Ext}_{\Osh_X}^{p}(X;\Osh_Y,\Omega_X^{p})$. Ahora bien, el homomorfismo
natural $\Osh_X\longrightarrow\Osh_Y$ define un homomorfismo
\[
  \operatorname{Ext}_{\Osh_X}^{p}(X;\Osh_Y,\Omega_X^{p})
  \longrightarrow
  \operatorname{Ext}_{\Osh_X}^{p}(X;\Osh_X,\Omega_X^{p})
  =H^{p}(X,\Omega_X^{p}).
\]
Se obtiene así un elemento de $H^{p}(X,\Omega_X^{p})$, denotado por $P_X(Y)$ y llamado
\emph{clase de cohomología de $Y$ en $X$}. Se deduce, por tanto, de la sección
$s_Y$ de $\Extsh_{\Osh_X}^{p}(\Osh_Y,\Omega_X^{p})$ mediante el diagrama
de homomorfismos siguiente:
\begin{equation}\tag{4.3}
\begin{aligned}
H^{p}(X,\Omega_X^{p})
&=\operatorname{Ext}_{\Osh_X}^{p}(X;\Osh_X,\Omega_X^{p})\\
&\longleftarrow\operatorname{Ext}_{\Osh_X}^{p}(X;\Osh_Y,\Omega_X^{p})
\xrightarrow{\ \sim\ }
H^{0}\!\left(X,\Extsh_{\Osh_X}^{p}(\Osh_Y,\Omega_X^{p})\right).
\end{aligned}
\end{equation}

Llamemos \emph{ciclo no singular de dimensión $n-p$} a un elemento del grupo
abeliano libre generado por las subvariedades irreducibles no singulares de
dimensión $n-p$ en $X$. Entonces la función $Y\longmapsto P(Y)$ se prolonga a
un homomorfismo del grupo de ciclos no singulares de dimensión $n-p$ sobre $X$
en el grupo $H^{p}(X,\Omega_X^{p})$.

Sean $Z^{n-p}$ y $Z'^{n-p'}$ ciclos no singulares de dimensiones $n-p$ y
$n-p'$; se dice que se \emph{intersecan transversalmente} si toda componente
de $Z$ interseca transversalmente toda componente de $Z'$. Entonces el ciclo $Z.Z'$ está
definido y es un ciclo no singular de dimensión $n-p-p'$. Dicho esto, se tiene el

\src{179}
\result{TEOREMA}{1} Si $Z^{n-p}$ y $Z'^{n-p'}$ son ciclos no singulares
que se intersecan transversalmente, entonces se tiene
\begin{equation}\tag{4.4}
  P_X(Z.Z')=P_X(Z).P_X(Z'),
\end{equation}
donde el producto del segundo miembro es el producto cup
\[
  H^{p}(X,\Omega_X^{p})\times H^{p'}(X,\Omega_X^{p'})
  \longrightarrow H^{p+p'}(X,\Omega_X^{p+p'}).
\]
(Se supone que $X$ es isomorfo a una parte localmente cerrada de un espacio
proyectivo.)

Esta última hipótesis sirve únicamente para poder concluir que todo
haz algebraico coherente sobre $X$ es cociente de un haz algebraico
coherente localmente libre (SERRE) y, por tanto, admite una resolución por la izquierda mediante
haces localmente libres. Para demostrar el teorema 1, se puede suponer que
$Z$ y $Z'$ son subvariedades irreducibles no singulares $Y$ y $Y'$ que se
intersecan transversalmente. Sea $L_{*}$ una resolución por la izquierda de $\Osh_Y$ mediante
haces localmente libres; entonces, en virtud de la proposición 3, el diagrama
de homomorfismos (4.3) se identifica con el diagrama
\[
R^{p}\Gamma(\Omega_X^{p})
\xleftarrow{\alpha}
\mathbb R^{p}\Gamma\!\left(\Homsh_{\Osh_X}(L_{*},\Omega_X^{p})\right)
\xrightarrow[\sim]{\beta}
\Gamma\!\left(H^{p}\bigl(\Homsh_{\Osh_X}(L_{*},\Omega_X^{p})\bigr)\right),
\]
donde $\Gamma$ es el funtor « grupo de secciones » sobre la categoría de los
haces abelianos sobre $X$, y $\mathbb R^{p}\Gamma$ designa
la hipercohomología de grado $p$ de dicho funtor, y $R^{p}\Gamma$ su
$p$-ésimo funtor derivado. Supondremos, para simplificar, que $L_0=\Osh_X$ y que
la aumentación $L_0\longrightarrow\Osh_Y$ es el homomorfismo natural (lo cual
es lícito); entonces $\alpha$ se deduce del homomorfismo de complejos
$\Osh_X\longrightarrow L_{*}$ (considerando $\Osh_X$ como un complejo concentrado
en grado $0$), teniendo en cuenta que $\mathbb R^{p}\Gamma(K)=R^{p}\Gamma(K_0)$ si $K$
es un complejo de haces concentrado en grado $0$. El homomorfismo $\beta$ es un
« homomorfismo de borde » bien conocido. Consideremos un diagrama análogo, relativo a
una resolución localmente libre $L'_{*}$ de $\Osh_{Y'}$, y consideremos el
diagrama conmutativo de acoplamientos:
\begin{equation}\tag{4.5}
\resizebox{\textwidth}{!}{$
\begin{array}{ccc}
R^{p}\Gamma(\Omega_X^{p})
&\longleftarrow
\mathbb R^{p}\Gamma\!\left(\Homsh_{\Osh_X}(L_{*},\Omega_X^{p})\right)
&\xrightarrow{\sim}
\Gamma H^{p}\!\left(\Homsh_{\Osh_X}(L_{*},\Omega_X^{p})\right)
\\[3pt]
\times&\times&\times\\[3pt]
R^{p'}\Gamma(\Omega_X^{p'})
&\longleftarrow
\mathbb R^{p'}\Gamma\!\left(\Homsh_{\Osh_X}(L'_{*},\Omega_X^{p'})\right)
&\xrightarrow{\sim}
\Gamma H^{p'}\!\left(\Homsh_{\Osh_X}(L'_{*},\Omega_X^{p'})\right)
\\[3pt]
\downarrow&\downarrow&\downarrow\\[3pt]
R^{p+p'}\Gamma(\Omega_X^{p+p'})
&\longleftarrow
\mathbb R^{p+p'}\Gamma\!\left(
\Homsh_{\Osh_X}(L_{*}\otimes L'_{*},\Omega_X^{p+p'})\right)
&\xrightarrow{\sim}
\Gamma H^{p+p'}\!\left(
\Homsh_{\Osh_X}(L_{*}\otimes L'_{*},\Omega_X^{p+p'})\right).
\end{array}
$}
\end{equation}

\src{180}
Los acoplamientos de las dos columnas de la derecha se deducen del acoplamiento de
complejos de haces
\[
\Homsh_{\Osh_X}(L_{*},\Omega_X^{p})
\times\Homsh_{\Osh_X}(L'_{*},\Omega_X^{p'})
\longrightarrow
\Homsh_{\Osh_X}(L_{*}\otimes L'_{*},\Omega_X^{p+p'}),
\]
que se define mediante el producto exterior
$\Omega_X^{p}\times\Omega_X^{p'}\longrightarrow\Omega_X^{p+p'}$.
El acoplamiento de la primera columna es el producto cup (relativo al producto
exterior). Afirmo que la última fila de (4.5) se identifica con el diagrama
de isomorfismos análogo a (4.3), donde $Y$ se sustituye por $Y\cap Y'$ y $p$ por
$p+p'$. Para ello basta demostrar que $L\otimes L'$ es una resolución
(evidentemente localmente libre) de $\Osh_{Y\cap Y'}$. En efecto, se tiene
\[
  H_0(L\otimes L')=\Osh_Y\otimes\Osh_{Y'}=\Osh_{Y\cap Y'}
\]
y
\[
  H_i(L\otimes L')=\operatorname{Tor}_{i}^{\Osh_X}(\Osh_Y,\Osh_{Y'})=0
  \qquad(i>0),
\]
por el hecho de que $Y$ y $Y'$ se intersecan transversalmente. El teorema 1 resulta
ahora de la fórmula:
\begin{equation}\tag{4.6}
  s_Y.s_{Y'}=s_{Y.Y'}
\end{equation}
(donde el producto del primer miembro es el de la última columna de (4.5)).
Esta fórmula (4.6), de naturaleza puramente local, se comprueba sin dificultad
tomando para $L_{*}$ y $L'_{*}$ las resoluciones consideradas en la proposición
4. Se demuestra asimismo (mediante una demostración más sencilla) que
$Z\longmapsto P_X(Z)$ es compatible con el producto cartesiano:
\begin{equation}\tag{4.7}
  P_{X\times X'}(Z\times Z')=P_X(Z)\otimes P_{X'}(Z')
\end{equation}
(fórmula válida si $Z$, resp. $Z'$, es un ciclo no singular sobre la variedad
no singular $X$, resp. $X'$, considerando $Z\times Z'$ como un ciclo no
singular sobre $X\times X'$). De (4.4) y (4.7) se deduce que $P_X(Z)$ es también
compatible con la operación « imagen inversa » por un morfismo de variedades no
singulares $f:X\longrightarrow X'$:
\begin{equation}\tag{4.8}
  P_X\bigl(f^{-1}(Z')\bigr)=f^{*}\bigl(P_{X'}(Z')\bigr),
\end{equation}
fórmula válida si $Z'$ es un ciclo no singular sobre $X'$ tal que $f$ sea
« transversal » a $Z'$, es decir, tal que la gráfica de $f$ sea transversal al ciclo
$X\times Z'$ en $X\times X'$.

\src{181}
\result{COROLARIO}{1} Sean $X$, $X'$ dos variedades no singulares
localmente cerradas en un espacio proyectivo, y supongamos que $X'$ es completa. Sean $U$
un ciclo no singular sobre $X\times X'$, $a$ y $b$ dos puntos de $X'$ tales que
$U$ interseca transversalmente los ciclos $X\times(a)$ y $X\times(b)$, y sean $Z$
y $Z'$ los ciclos no singulares sobre $X$ tales que
\[
  Z\times(a)=(X\times(a)).U,
  \qquad
  Z'\times(b)=(X\times(b)).U.
\]
Entonces se tiene
\[
  P_X(Z)=P_X(Z').
\]

En efecto, sea $f_a:X\longrightarrow X\times X'$ definido por
$f_a(x)=(x,a)$; entonces, en virtud de (4.8), se tiene la fórmula
$P(Z)=f_a^{*}(P(U))$, y análogamente $P(Z')=f_b^{*}(P(U))$. Ahora bien, utilizando la fórmula de
Künneth
\[
  H^{*}(X\times X',\Omega_{X\times X'}^{*})
  =H^{*}(X,\Omega_X^{*})\otimes H^{*}(X',\Omega_{X'}^{*})
\]
y el hecho de que $H^{0}(X',\Osh_{X'})$ se reduce a los escalares, se obtiene
fácilmente $f_a^{*}=f_b^{*}$, de donde se sigue la conclusión.

Para todo $x\in X$, $(x)$ es una subvariedad no singular de $X$ de
codimensión $n$, y define por tanto un elemento $\varepsilon_x$ de
$H^{n}(X,\Omega_X^{n})$. Si $X$ es una variedad proyectiva no singular, del
corolario 1 resulta que $\varepsilon_x$ no depende del punto $x$ elegido; se
denota por $\varepsilon_X$ y se llama \emph{clase fundamental} de
$H^{n}(X,\Omega_X^{n})$.

\result{OBSERVACIÓN}{} Para disponer de una teoría satisfactoria, habría que definir
$P_X(Z)$ para ciclos $Z$ cualesquiera y demostrar el teorema 1 para una
intersección propia de ciclos. (En el momento de escribir esta exposición, esto aún no se
ha hecho con toda generalidad.) Admitiendo que se haya hecho, el corolario 1
pasa a ser: si $Z$ y $Z'$ son dos ciclos algebraicamente equivalentes, entonces
$P_X(Z)=P_X(Z')$ (afirmación que no parece resultar de lo anterior, incluso si
$Z$ y $Z'$ son no singulares).

\fsection{5}{El teorema de dualidad}

En este número, $X$ designa una variedad proyectiva no singular de dimensión
$n$.

\result{TEOREMA}{2} La clase fundamental $\varepsilon_X$ de
$H^{n}(X,\Omega_X^{n})$ es una base de este espacio vectorial.

(La demostración se dará más adelante.) En virtud del teorema anterior, se puede
identificar $H^{n}(X,\Omega_X^{n})$ con el cuerpo $k$. Consideremos ahora los
acoplamientos contemplados en el n\textsuperscript{o} 2, que dan en particular un
acoplamiento

\src{182}
\[
\operatorname{Ext}_{\Osh_X}^{p}(X;\Osh_X,F)
\times\operatorname{Ext}_{\Osh_X}^{n-p}(X;F,\Omega_X^{n})
\longrightarrow
\operatorname{Ext}_{\Osh_X}^{n}(X;\Osh_X,\Omega_X^{n}),
\]
es decir,
\begin{equation}\tag{5.1}
  H^{p}(X,F)
  \times\operatorname{Ext}_{\Osh_X}^{n-p}(X;F,\Omega_X^{n})
  \longrightarrow H^{n}(X,\Omega_X^{n}).
\end{equation}
Teniendo en cuenta el teorema 2, este acoplamiento define un homomorfismo
\begin{equation}\tag{5.2}
  \operatorname{Ext}_{\Osh_X}^{n-p}(X;F,\Omega_X^{n})
  \longrightarrow\bigl(H^{p}(X,F)\bigr)'.
\end{equation}
Este homomorfismo es funtorial en $F$ y conmuta con los homomorfismos de conexión
relativos a las sucesiones exactas $0\longrightarrow F'\longrightarrow F
\longrightarrow F''\longrightarrow0$.

\result{TEOREMA}{3} El homomorfismo (5.2) es un isomorfismo.

Se recupera, en particular, el resultado de Serre:

\result{COROLARIO}{} Sean $E$ un fibrado vectorial algebraico sobre $X$ y
$\Osh_X(E)$ el haz de gérmenes de secciones regulares de $E$; entonces se tienen
isomorfismos canónicos:
\begin{equation}\tag{5.3}
  \bigl(H^{p}(X,\Osh_X(E))\bigr)'
  =H^{n-p}\bigl(X,\Omega_X^{n}\otimes\Osh_X(E')\bigr).
\end{equation}
Basta aplicar el teorema 3 y el corolario 1 a la proposición 1.

Los teoremas 2 y 3 resultarán del enunciado siguiente:

\noindent\textbf{(D) El homomorfismo siguiente}
\begin{equation}\tag{5.2 bis}
  \operatorname{Ext}_{\Osh_X}^{n-p}(X;F,\Omega_X^{n})
  \longrightarrow\bigl(H^{p}(X,F)\bigr)'\otimes L,
  \qquad L=H^{n}(X,\Omega_X^{n}),
\end{equation}
deducido del acoplamiento (5.1) es un isomorfismo.

Demostremos, en efecto, que (D) implica el teorema 2. Sea
$k_x=\Osh_{(x)}$ el haz de anillos locales de la variedad reducida al punto
$x\in X$; consideremos el homomorfismo canónico
$\Osh_X\longrightarrow k_x$ y el homomorfismo asociado
\begin{equation}\tag{5.3 bis}
  H^{0}(X,\Osh_X)\longrightarrow H^{0}(X,k_x),
\end{equation}
su transpuesto se identifica con el homomorfismo
\begin{equation}\tag{5.4}
  \operatorname{Ext}_{\Osh_X}^{n}(X;k_x,\Omega_X^{n})\otimes L'
  \longrightarrow
  \operatorname{Ext}_{\Osh_X}^{n}(X;\Osh_X,\Omega_X^{n})\otimes L'
\end{equation}
deducido del homomorfismo sobre los $\operatorname{Ext}^{n}$ asociado a
$\Osh_X\longrightarrow k_x$:
\begin{equation}\tag{5.5}
  \operatorname{Ext}_{\Osh_X}^{n}(X;k_x,\Omega_X^{n})
  \longrightarrow
  \operatorname{Ext}_{\Osh_X}^{n}(X;\Osh_X,\Omega_X^{n})
  =H^{n}(X,\Omega_X^{n}).
\end{equation}

\src{183}
Como (5.3 bis) es un isomorfismo,\footnote{El texto impreso da una
segunda fórmula « (5.3) » en la página 182 y remite aquí a ella; el identificador
« (5.3 bis) » distingue esta fórmula de la primera (5.3).} lo mismo ocurre
con (5.4), y por tanto también con (5.5). Como $s_{(x)}$ es una base de
$\operatorname{Ext}_{\Osh_X}^{n}(X;k_x,\Omega_X^{n})$ en virtud de (4.2), se
sigue que su imagen $\varepsilon_X$ es una base de
$H^{n}(X,\Omega_X^{n})$.

Queda por demostrar el enunciado (D), que resultará de manera puramente formal de los hechos
elementales resumidos en los lemas siguientes. En ellos se supone que $X$ es una
parte cerrada (singular o no) del espacio proyectivo $P$ de dimensión $r$.
Se utiliza la notación $\Osh_P(m)$ para el haz sobre $P$ denotado por $\Osh(m)$ en
[3], y la notación $\Osh_X(m)$ para el haz análogo sobre $X$.

\result{LEMA}{2} El enunciado (D) es verdadero si $X=P$ y $F=\Osh_P(m)$.

Este lema se comprueba mediante un cálculo directo. El cálculo explícito de los
$H^{i}(P,\Osh_P(m))$ se encuentra en [3], pero puede hacerse de manera más
elemental. El cálculo del producto cup
\[
  H^{i}(P,\Osh_P(m))\times H^{j}(P,\Osh_P(m'))
  \longrightarrow H^{i+j}(P,\Osh_P(m+m'))
\]
necesario para explicitar el acoplamiento (5.1) no presenta
dificultades.\footnote{El texto impreso dice $\Osh_P(r+r')$; los dos factores
$\Osh_P(m)$ y $\Osh_P(m')$ exigen $\Osh_P(m+m')$.}

\result{LEMA}{3} Todo haz algebraico coherente $F$ sobre $X$ es isomorfo a
un haz cociente de un haz $\Osh_X(-m)^{k}$, donde se puede suponer que $m$ es
tan grande como se quiera.

Resulta del hecho de que $F\otimes\Osh_X(m)$ está « generado por sus secciones » para
$m$ grande, cf. [3].

\result{LEMA}{4} Sea $i>0$. Entonces
$H^{r-i}(P,\Osh_P(-m))=0$ para $m$ grande; y para todo haz algebraico
coherente $B$ sobre $X$, se tiene
\[
  \operatorname{Ext}_{\Osh_X}^{i}(X;\Osh_X(-m),B)=0
\]
para $m$ grande.

El primer enunciado resulta de los cálculos explícitos mencionados anteriormente; para el
segundo, se observa que se tiene un isomorfismo
\[
  \operatorname{Ext}_{\Osh_X}^{i}(X;\Osh_X(-m),B)
  =H^{i}(X,B\otimes\Osh_X(m))
\]
(proposición 1, corolario 1), de donde se sigue la conclusión en virtud de un resultado bien
conocido de [3]. Combinando los dos lemas anteriores, se obtiene el

\result{COROLARIO}{} Sea $i>0$. Entonces el funtor
$F\longmapsto H^{r-i}(P,F)$ sobre la categoría de los haces algebraicos coherentes
sobre $P$ es borrable, y lo mismo ocurre con el funtor
$\operatorname{Ext}_{\Osh_X}^{i}(X;F,B)$ sobre la categoría de los haces
algebraicos coherentes sobre $X$.

\src{184}
\result{LEMA}{5} Sean $A$ y $B$ dos haces algebraicos coherentes sobre
$X$, y pongamos $A(m)=A\otimes\Osh_X(m)$. Entonces, para $m$ suficientemente grande,
el homomorfismo canónico
\[
  \operatorname{Ext}_{\Osh_X}^{1}(X;A(-m),B)
  \longrightarrow H^{0}\bigl(X,\Extsh_{\Osh_X}^{1}(A(-m),B)\bigr)
  =H^{0}\bigl(X,\Extsh_{\Osh_X}^{1}(A,B)(m)\bigr)
\]
es un isomorfismo.

Esto resulta inmediatamente de la sucesión espectral de la proposición 1 aplicada a
$A(-m)$ y $B$, pues se tendrá
\[
  E_{2}^{p,q}(A(-m),B)
  =H^{p}\bigl(X,\Extsh_{\Osh_X}^{q}(A(-m),B)\bigr)
  =H^{p}\bigl(X,\Extsh_{\Osh_X}^{q}(A,B)(m)\bigr),
\]
que es nulo para $p>0$ y $m$ grande.

Demostremos ahora (D) en el caso en que $X=P$. Probamos primero que (5.2 bis)
es un isomorfismo para $p=n$; como los dos miembros son entonces
funtores exactos por la izquierda (pues $H^{r+1}(P,F)=0$), del lema 3 resulta
que basta demostrar la afirmación para $F=\Osh_P(-m)$, pero esta se encuentra
contenida en el lema 2. Como los homomorfismos (5.2 bis) son funtoriales y
compatibles con los operadores de conexión, y como para $p<n$ los dos miembros de
(5.2 bis) son funtores borrables en $F$ (corolario del lema 4), se
sigue entonces, mediante un razonamiento estándar, que (5.2 bis) es un isomorfismo
para todo $p$. Esto demuestra el teorema de dualidad para el espacio proyectivo.
Supongamos ahora que $X$ es cualquiera, pero no singular. Por el teorema de
dualidad para $P$, se tiene un isomorfismo
\[
  H^{n}(X,F)'=H^{n}(P,F)'
  =\operatorname{Ext}_{\Osh_P}^{r-n}(P;F,\Omega_P^{r}).
\]
En virtud del lema 1 (n\textsuperscript{o} 4), el último miembro se identifica con
\[
  \operatorname{Hom}_{\Osh_P}(P;F,\omega_X^{n})
  =\operatorname{Hom}_{\Osh_X}(X;F,\Omega_X^{n})
  =\operatorname{Ext}_{\Osh_X}^{0}(X;F,\Omega_X^{n}),
\]
de donde se obtiene un isomorfismo
\begin{equation}\tag{5.6}
  H^{n}(X,F)'=\operatorname{Hom}_{\Osh_X}(X;F,\Omega_X^{n})
  =\operatorname{Ext}_{\Osh_X}^{0}(X;F,\Omega_X^{n}).
\end{equation}
Tomando $F=\Omega_X^{n}$, se obtiene un isomorfismo
\begin{equation}\tag{5.7}
  \gamma:H^{n}(X,\Omega_X^{n})\xrightarrow{\sim}k.
\end{equation}

\src{185}
Se comprueba que el isomorfismo (5.6) no es otro que (5.2 bis) para $p=n$, cuando
se toma $L=k$ gracias a (5.7). Por consiguiente, (5.2 bis) es un isomorfismo para
$p=n$. Para demostrar que es un isomorfismo para todo $p$, basta nuevamente
demostrar que, para $p<n$, los dos miembros de (5.2 bis) son funtores
borrables en $F$, y a fortiori (teniendo en cuenta el lema 3) que ambos miembros
son nulos cuando se toma $F=\Osh_X(-m)$ con $m$ grande. Ahora bien, para el primer
miembro, esto es cierto en virtud del lema 4, y para el segundo miembro se escribe,
utilizando el teorema de dualidad sobre $P$:
\[
  H^{p}(X,\Osh_X(-m))'
  =\operatorname{Ext}_{\Osh_P}^{r-p}
    (P;\Osh_X(-m),\Omega_P^{r}).
\]
El segundo miembro es nulo para $p<n$ y $m$ grande, como resulta del lema 5
(donde se toma $X=P$) y del hecho de que $\Osh_X$, como haz algebraico
coherente sobre $P$, tiene dimensión cohomológica $\leq r-n$ (pues $X$ es no
singular), de donde
\[
  \Extsh_{\Osh_P}^{r-p}(\Osh_X,\Omega_P^{r})=0
  \qquad\text{para }p<n.
\]

\fsection{6}{El teorema de dualidad para las variedades singulares}

Sea $X$ una parte cerrada de dimensión $n$ del espacio proyectivo $P$ de
dimensión $r$. La fórmula (5.6) debe escribirse aquí
\begin{equation}\tag{6.1}
  H^{n}(X,F)'
  \simeq\operatorname{Hom}_{\Osh_X}(X;F,\omega_X^{n})
  =\operatorname{Ext}_{\Osh_X}^{0}(X;F,\omega_X^{n}),
\end{equation}
donde se ha puesto
\begin{equation}\tag{6.2}
  \omega_X^{n}=E^{0}(\Osh_X)
  =\Extsh_{\Osh_P}^{r-n}(\Osh_X,\Omega_P^{r}).
\end{equation}
Como se dijo en la proposición 5, el haz así introducido no depende
en realidad de la inmersión elegida de $X$ en una variedad no singular $P$.
Tomando, en (6.1), $F=\omega_X^{n}$, se obtiene
\begin{equation}\tag{6.2 bis}
  H^{n}(X,\omega_X^{n})'
  \simeq\operatorname{Hom}_{\Osh_X}(X;\omega_X^{n},\omega_X^{n}),
\end{equation}
\footnote{El texto impreso numera también esta fórmula « (6.2) »; el sufijo
« bis » distingue las dos apariciones.} de donde resulta la existencia de un elemento
distinguido en $H^{n}(X,\omega_X^{n})'$, correspondiente al homomorfismo
identidad de $\omega_X^{n}$ en sí mismo:
\begin{equation}\tag{6.3}
  \gamma:H^{n}(X,\omega_X^{n})\longrightarrow k.
\end{equation}

\src{186}
Consideremos entonces los acoplamientos definidos por la composición de los $\operatorname{Ext}$:
\begin{equation}\tag{6.4}
  H^{p}(X,F)\times
  \operatorname{Ext}_{\Osh_X}^{n-p}(X;F,\omega_X^{n})
  \longrightarrow H^{n}(X,\omega_X^{n}),
\end{equation}
y compongámoslos con el homomorfismo $\gamma$ de (6.3); se obtienen
homomorfismos funtoriales, compatibles con los operadores de borde:
\begin{equation}\tag{6.5}
  \operatorname{Ext}_{\Osh_X}^{n-p}(X;F,\omega_X^{n})
  \longrightarrow H^{p}(X,F)'
\end{equation}
(que generalizan (5.2)). Se comprueba que para $p=n$ se obtiene así
el isomorfismo (6.1). Sentado esto, se tiene el

\result{TEOREMA}{3 bis} Las cuatro condiciones siguientes sobre $X$ son
equivalentes para un entero dado $k\geq0$:
\begin{enumerate}
  \item[\textit{i.}] El homomorfismo funtorial (6.5) es un isomorfismo para
    $n-k\leq p\leq n$.
  \item[\textit{ii.}] Se tiene $H^{p}(X,\Osh_X(-m))=0$ para $m$ grande y para
    $n-k\leq p<n$.
  \item[\textit{iii.}] El funtor $H^{p}(X,F)$ sobre la categoría de los haces
    algebraicos coherentes sobre $X$ es borrable para $n-k\leq p<n$.
  \item[\textit{iv.}] Se tiene
    $E^{i}(\Osh_X)=\Extsh_{\Osh_P}^{r-n+i}(\Osh_X,\Omega_P^{r})=0$
    para $0<i\leq k$.
\end{enumerate}

\noindent\textbf{DEMOSTRACIÓN.} --- \textit{i.}$\Rightarrow$\textit{ii.} en
virtud del lema 4, \textit{ii.}$\Rightarrow$\textit{iii.} en virtud del lema 3,
\textit{iii.}$\Rightarrow$\textit{i.} mediante un razonamiento estándar bien conocido,
teniendo en cuenta que los dos miembros de (6.5) son entonces funtores borrables
para $n-k\leq p<n$ (el primero lo es en virtud del lema 4). Por último, demostremos
\textit{ii.}$\Leftrightarrow$\textit{iv.} Esto resulta del corolario de la
proposición siguiente:

\result{PROPOSICIÓN}{6} Sea $F$ un haz algebraico coherente sobre $X$, y
sea $i$ un entero. Entonces, para $m$ suficientemente grande, se tiene un isomorfismo:
\begin{equation}\tag{6.6}
  H^{i}(X,F(-m))'
  \simeq H^{0}\bigl(X,E^{n-i}(F)(m)\bigr),
\end{equation}
donde se pone (compárese n\textsuperscript{o} 3, proposición 5):
\begin{equation}\tag{6.7}
  E^{j}(F)=\Extsh_{\Osh_P}^{r-n+j}(F,\Omega_P^{r}).
\end{equation}
En efecto, en virtud del teorema de dualidad sobre $P$, el primer miembro de (6.6)
es isomorfo a
$\operatorname{Ext}_{\Osh_P}^{r-i}(P;F(-m),\Omega_P^{r})$, por lo que (6.6) resulta
del lema 5.

\src{187}
\result{COROLARIO}{} Para que se tenga $H^{i}(X,F(-m))=0$ para $m$ grande, es
necesario y suficiente que $E^{n-i}(F)=0$.

Recordemos que los $E^{j}(F)$ no dependen en realidad de la inmersión proyectiva
considerada. La condición del corolario es de naturaleza puramente local; por consiguiente,
si se verifica para $F$, se verifica para todo haz localmente
isomorfo a un haz $F^{n}$. En particular, si esta condición se verifica
para $\Osh_X$, se verifica para todo haz algebraico coherente localmente
libre. Este es, por ejemplo, el caso para todo $i<n$ si $X$ es no singular, para
$i=0$ si ninguna componente de $X$ se reduce a un solo punto, y para $i=0,1$
si $X$ es normal y todas sus componentes tienen dimensión $>1$ (cf.
[3]).\footnote{El texto impreso dice aquí « $S$ es normal »; el párrafo
solo se refiere a la variedad $X$.} Para que esto ocurra para todo
$i<k$, es necesario y suficiente, por definición, que los anillos locales
$\Osh_x$ ($x\in X$) tengan « codimensión homológica » $\geq k$ (cf. [4]
para esta noción). Si $k=n$, esto significa, en virtud del teorema 3 bis, que el
teorema de dualidad es verdadero para $X$, es decir, que (6.5) es un isomorfismo para
todo $p$ y todo $F$. Se pueden dar numerosas condiciones equivalentes sobre
los anillos locales $\Osh_x$ para que esto ocurra (Nagata), por ejemplo que
satisfagan el teorema de equidimensionalidad de Cohen--Macaulay. Así ocurre,
por ejemplo, si localmente $X$ es una « intersección completa » en una
variedad ambiente no singular.

\fsection{7}{La dualidad de Poincaré}

Sea $X$ una variedad proyectiva no singular de dimensión $n$. Entonces
\[
  H^{**}(X)=H^{*}(X,\Omega_X^{*})
\]
es un álgebra anticonmutativa bigraduada de dimensión finita, que graduamos
por el grado total, de modo que $H^{p,q}(X)=H^{p}(X,\Omega_X^{q})$ tiene grado
$p+q$. Los grados de $H^{*}(X)$ están comprendidos entre $0$ y $2n$. En virtud del
teorema 2 y del corolario del teorema 3, $H^{*}(X)$ es un « álgebra de
Poincaré » de dimensión $2n$, es decir, $H^{2n}(X)$ está provisto de un isomorfismo con
el cuerpo base $k$, y el producto
\[
  H^{m}(X)\times H^{2n-m}(X)\longrightarrow H^{2n}(X)=k
\]
es una dualidad entre $H^{m}(X)$ y $H^{2n-m}(X)$. Además, si $Y$ es otra
variedad proyectiva no singular, la fórmula de Künneth para los
haces algebraicos coherentes da
\begin{equation}\tag{7.1}
  H^{*}(X\times Y)=H^{*}(X)\otimes H^{*}(Y),
\end{equation}
isomorfismo compatible con las estructuras de álgebras de Poincaré. Además,
$H^{*}(X)$ es, como álgebra conmutativa, un funtor contravariante en
$X$, pues un morfismo $f:Y\longrightarrow X$ define de manera evidente un
homomorfismo de álgebras graduadas

\src{188}
\begin{equation}\tag{7.2}
  f^{*}:H^{*}(X)\longrightarrow H^{*}(Y).
\end{equation}
Como se trata de álgebras de Poincaré, se obtiene por transposición un
homomorfismo de espacios vectoriales:
\begin{equation}\tag{7.3}
  f_{*}:H^{*}(Y)\longrightarrow H^{*}(X).
\end{equation}
Se vio en el n\textsuperscript{o} 4 que el efecto de $f^{*}$ sobre las clases de
cohomología correspondientes a ciclos no singulares se interpreta
geométricamente tomando las clases de cohomología que corresponden a sus
imágenes inversas. En el caso actual importa demostrar que (7.3) corresponde
análogamente a la operación « imagen directa » de ciclos. Esto resulta, al menos bajo
condiciones adecuadas de no singularidad, del caso particular siguiente:

\result{TEOREMA}{4} Si $f$ es la inclusión de una subvariedad no
singular $Y^{m}$ de $X^{n}$ en $X^{n}$, entonces, denotando por $1_Y$
el elemento unidad de $H(Y)$, se tiene
\begin{equation}\tag{7.4}
  f_{*}(1_Y)=P_X(Y),
\end{equation}
donde el segundo miembro es la clase de cohomología en $X$ asociada a $Y$.

Esta fórmula equivale a
\begin{equation}\tag{7.4 bis}
  \bigl\langle\xi^{m,m},P_X(Y)\bigr\rangle\varepsilon_Y
  =f^{*}(\xi^{m,m})
  \qquad
  \bigl(\xi^{m,m}\in H^{m}(X,\Omega_X^{m})\bigr),
\end{equation}
\footnote{El texto impreso dice $f_{*}(\xi^{m,m})$. Puesto que
$\xi^{m,m}\in H^{m}(X,\Omega_X^{m})$ y que ambos miembros pertenecen a
$H^{m}(Y,\Omega_Y^{m})$, la aplicación requerida es $f^{*}$.} donde
$\varepsilon_Y$ es el elemento fundamental en $H^{m}(Y,\Omega_Y^{m})$, y
constituye, en el caso de las variedades proyectivas no singulares, una nueva
definición de la clase de cohomología asociada a $Y$. Para demostrar el
teorema 4, se interpreta, gracias al teorema 3, el transpuesto del homomorfismo
\[
  H^{m}(X,\Omega_X^{m})
  \longrightarrow H^{m}(Y,\Omega_Y^{m})
  =H^{m}(X,\Omega_Y^{m})
\]
como el homomorfismo
\begin{equation}\tag{7.5}
\begin{split}
  H^{n-m}(X,\Omega_X^{n-m})
  &\simeq
  \operatorname{Ext}_{\Osh_X}^{n-m}(X;\Omega_X^{m},\Omega_X^{n})\\
  &\longleftarrow
  \operatorname{Ext}_{\Osh_X}^{n-m}(X;\Omega_Y^{m},\Omega_X^{n})
  \simeq
  \operatorname{Hom}_{\Osh_X}(X;\Omega_Y^{m},\Omega_Y^{m}).
\end{split}
\end{equation}
Se comprueba que el elemento $1_Y$ del dual de $H^{m}(Y,\Omega_Y^{m})$ se identifica
con el elemento del miembro de la derecha correspondiente al endomorfismo identidad de
$\Omega_Y^{m}$, y por otra parte que la imagen de este elemento en
$H^{n-m}(X,\Omega_X^{n-m})$ es precisamente $P_X(Y)$. Estas compatibilidades, que podrían
haberse dado ya en el n\textsuperscript{o} 4, pueden enunciarse, y son
en realidad válidas, para variedades no singulares cualesquiera; la segunda

\src{189}
resulta, por ejemplo, de la conmutatividad del diagrama siguiente de endomorfismos
canónicos:
\begin{equation}\tag{7.6}
\resizebox{\textwidth}{!}{$
\left\{
\begin{array}{ccccc}
\operatorname{Ext}_{\Osh_X}^{n-m}(X;\Omega_X^{m},\Omega_X^{n})
&\longleftarrow&
\operatorname{Ext}_{\Osh_X}^{n-m}(X;\Omega_Y^{m},\Omega_X^{n})
&\simeq&
\operatorname{Hom}_{\Osh_X}(X;\Omega_Y^{m},\Omega_Y^{m})\\[3pt]
\Big\downarrow&&&&\Big\downarrow\\[3pt]
\operatorname{Ext}_{\Osh_X}^{n-m}(X;\Osh_X,\Omega_X^{n-m})
&\longleftarrow&
\operatorname{Ext}_{\Osh_X}^{n-m}(X;\Osh_Y,\Omega_X^{n-m})
&\simeq&
\operatorname{Hom}_{\Osh_X}(X;\Omega_X^{m},\Omega_Y^{m})
\end{array}
\right.
$}
\end{equation}

Se obtiene así un equivalente exacto del formalismo de la dualidad de Poincaré
sobre las variedades orientadas compactas. En particular, el teorema 4 permite
determinar la clase de cohomología asociada a la diagonal de $X\times X$. Mediante
un razonamiento bien conocido, se deduce, por ejemplo, una fórmula de Lefschetz:

\result{TEOREMA}{5} Sea $f$ un endomorfismo de una variedad proyectiva no
singular $X$, tal que los puntos fijos de $f$ tengan multiplicidad 1.
Entonces el número de estos puntos fijos es congruente, módulo la característica de
$k$, con la suma alternada de las trazas de los endomorfismos de los $H^{i}(X)$ definidos
por $f$.

La restricción sobre $f$ que ha sido necesario imponer se debe a las dificultades mencionadas en
la observación del n\textsuperscript{o} 4. No obstante, se observará que la fórmula de
Lefschetz sigue siendo verdadera si $f$ es « homotópico » a un endomorfismo cuyos
puntos fijos tienen todos multiplicidad 1.

\fsection{8}{Generalización del teorema de dualidad}

Sea $X$ una variedad algebraica no singular tal que todo haz
algebraico coherente $F$ sobre $X$ sea isomorfo a un cociente de un haz
algebraico coherente localmente libre (lo cual ocurre si $X$ es localmente
cerrado en un espacio proyectivo). Entonces todo haz algebraico coherente $F$ sobre
$X$ admite una resolución finita $L$ mediante haces localmente libres, y para
dos resoluciones de esta clase siempre se puede encontrar una tercera y un
homomorfismo de esta en las dos primeras compatible con las
aumentaciones. Análogamente, si $L$ es una resolución finita localmente libre de
$F$, y si se da un homomorfismo $F'\longrightarrow F$, entonces existe
una resolución finita localmente libre $L'$ de $F'$ y un homomorfismo
$L'\longrightarrow L$ compatible con $F'\longrightarrow F$, que se puede
suponer sobreyectivo si $F'\longrightarrow F$ lo es. Esto permite definir, dados
dos enteros $r,s\geq0$, dos multifuntores cohomológicos en los
argumentos $A_1,\ldots,A_r$; $B_1,\ldots,B_s$,

\src{190}
que varían en la categoría de los haces algebraicos coherentes sobre $X$, y cuyos
valores están en la categoría de los haces algebraicos coherentes sobre $X$,
resp. en la categoría de los módulos sobre $H^{0}(X,\Osh_X)$, mediante las fórmulas:
\begin{equation}\tag{8.1}
\resizebox{\textwidth}{!}{$
\left\{
\begin{aligned}
\underline{T}_{r}^{s*}(A_1,\ldots,A_r;B_1,\ldots,B_s)
  &={}H^{*}\!\left(
    \Homsh_{\Osh_X}\!\left(
      L(A_1)\otimes\cdots\otimes L(A_r),
      L(B_1)\otimes\cdots\otimes L(B_s)
    \right)\right),\\[3pt]
T_{r}^{s*}(A_1,\ldots,A_r;B_1,\ldots,B_s)
  &={}\mathbf R^{*}\Gamma\!\left(
    \Homsh_{\Osh_X}\!\left(
      L(A_1)\otimes\cdots\otimes L(A_r),
      L(B_1)\otimes\cdots\otimes L(B_s)
    \right)\right).
\end{aligned}
\right.
$}
\end{equation}
En esta fórmula, $L(F)$ designa una resolución finita localmente libre del
haz algebraico coherente $F$, y $\mathbf R^{*}\Gamma(K)$ designa
la hipercohomología del espacio $X$ con respecto al complejo de haces $K$. Si
$r$ o $s$ es nulo, se sustituye el producto tensorial de los $L(A_i)$, resp. de los
$L(B_j)$, por $\Osh_X$. En particular, $\underline{T}_{0}^{0}$ y $T_{0}^{0}$
son funtores graduados de 0 argumentos, $\underline{T}_{0}^{0}$ se reduce al
grado 0, donde es el haz $\Osh_X$, mientras que $T_{0}^{0}$ es idéntico a
$H^{*}(X,\Osh_X)$. El hecho de que los segundos miembros de (8.1) no dependan
de las resoluciones elegidas es, en cualquier caso, evidente para la primera línea
(pues la cuestión es entonces local), y para la segunda resulta de las observaciones
generales precedentes, teniendo en cuenta la sucesión espectral de
hipercohomología del complejo de haces
\[
  K=\Homsh_{\Osh_X}\bigl(L(A_1)\otimes\cdots,
      L(B_1)\otimes\cdots\bigr),
\]
que converge a la hipercohomología de $X$ con respecto a $K$ y cuyo término
inicial es $H^{p}(X,H^{q}(K))$, es decir,
\begin{equation}\tag{8.2}
  E_{2}^{p,q}
  =H^{p}\!\left(X,
    (\underline{T}_{r}^{s})^{(q)}
      (A_1,\ldots,A_r;B_1,\ldots,B_s)\right).
\end{equation}
Se ve entonces que esta sucesión espectral tampoco depende de las resoluciones
elegidas; su límite es el primer miembro de (8.1). Se definen
fácilmente los operadores de conexión relativos a los diversos argumentos $A_i$, $B_j$,
observando que toda sucesión exacta
$0\longrightarrow F'\longrightarrow F\longrightarrow F''\longrightarrow0$
puede resolverse mediante una sucesión exacta de complejos finitos localmente libres.

En el sistema de los funtores $\underline{T}_{r}^{s*}$, resp.
$T_{r}^{s*}$, se definen operaciones análogas a las del cálculo tensorial, cuya
definición es inmediata a partir de las fórmulas definitorias (8.1). Así, se tiene
una composición (que generaliza la considerada en el n\textsuperscript{o} 2):
\begin{equation}\tag{8.3}
\resizebox{\textwidth}{!}{$
\begin{aligned}
&T_{r}^{s*}(A_1,\ldots,A_r;B_1,\ldots,B_s)
\times T_{r'}^{s'*}(A'_1,\ldots,A'_{r'};B'_1,\ldots,B'_{s'})\\
&\qquad\longrightarrow
T_{r+r'}^{s+s'*}
(A_1,\ldots,A_r,A'_1,\ldots,A'_{r'};
 B_1,\ldots,B_s,B'_1,\ldots,B'_{s'}),
\end{aligned}
$}
\end{equation}
que satisface las propiedades evidentes de asociatividad y de compatibilidad con
los homomorfismos funtoriales, los homomorfismos de conexión y las sucesiones
espectrales. Asimismo, se tienen operaciones de simetría cuya explicitación se deja al
lector. Se tiene además una operación de contracción cada vez que
uno de los argumentos $A_i$

\src{191}
es igual a uno de los argumentos $B_j$:
\begin{equation}\tag{8.4}
\resizebox{\textwidth}{!}{$
\begin{aligned}
&T_{r}^{s*}(A_1,\ldots,A_{i-1},C,A_{i+1},\ldots,A_r;
 B_1,\ldots,B_{j-1},C,B_{j+1},\ldots,B_s)\\
&\qquad\longrightarrow
T_{r-1}^{s-1*}(A_1,\ldots,\widehat{A_i},\ldots,A_r;
 B_1,\ldots,\widehat{B_j},\ldots,B_s).
\end{aligned}
$}
\end{equation}
Además, si un argumento $A_i$ es un haz localmente libre, puede
suprimirse a condición de sustituir uno de los $A_j$ ($j\neq i$) por
$A_j\otimes A_i$, o uno de los $B_k$ por $B_k\otimes A_i'$ (donde
$A_i'=\Homsh_{\Osh_X}(A_i,\Osh_X)$), y se tiene una regla de cálculo análoga en
el caso en que uno de los argumentos $B_i$ es localmente libre. En particular, siempre
se puede suprimir un argumento que sea idéntico a $\Osh_X$. Si todos los
argumentos son localmente libres, salvo a lo sumo uno de los argumentos $B_i$, la regla
que acaba de enunciarse da un isomorfismo funtorial
\begin{equation}\tag{8.5}
  T_{r}^{s*}(A_1,\ldots,A_r;B_1,\ldots,B_s)
  =H^{*}\bigl(X,A_1'\otimes\cdots\otimes A_r'
    \otimes B_1\otimes\cdots\otimes B_s\bigr)
\end{equation}
(pues se reduce al caso en que $r=0$, $s=1$, donde esto es inmediato; también se
puede utilizar directamente la sucesión espectral de término inicial (8.2)). Para los
$\underline{T}_{r}^{s}$ se definen operaciones correspondientes a las anteriores.
Las relaciones entre las diversas operaciones así
introducidas son las mismas que para las operaciones análogas del cálculo
tensorial.

Sea $n$ la dimensión de $X$. Aplicando sucesivamente una composición
tensorial (8.3) y contracciones (8.4) sobre los argumentos repetidos, se
obtiene un acoplamiento
\begin{equation}\tag{8.6}
\resizebox{\textwidth}{!}{$
  (T_{r}^{s})^{p}(A_1,\ldots;B_1,\ldots)
  \times
  (T_{s}^{r})^{n-p}(B_1,\ldots;A_1,\ldots,A_r\otimes\Omega_X^{n})
  \longrightarrow H^{n}(X,\Omega_X^{n}).
$}
\end{equation}

\result{TEOREMA}{6} Si $X$ es una variedad proyectiva no singular, entonces
los acoplamientos (8.6) son dualidades.

Esto resulta de manera puramente formal del corolario del teorema 3. En efecto,
de este corolario resulta fácilmente que, si $K$ es un complejo de haces
algebraicos coherentes localmente libres, entonces la hipercohomología de $X$ con
respecto a $K$ está en dualidad con la hipercohomología de $X$ con respecto a
$K'\otimes\Omega_X^{n}$ mediante los acoplamientos naturales
\begin{equation}\tag{8.7}
  \mathbf R^{p}\Gamma(K)
  \times\mathbf R^{n-p}\Gamma(K'\otimes\Omega_X^{n})
  \longrightarrow\mathbf R^{n}\Gamma(\Omega_X^{n})
  =H^{n}(X,\Omega_X^{n}).
\end{equation}
Esto se ve utilizando la sucesión espectral de término inicial
$H^{p}(X,H^{q}(K))$\footnote{El texto impreso dice
$H^{p}(H^{q}(X,K))$; el término inicial es el de la sucesión espectral (8.2),
esto es, $H^{p}(X,H^{q}(K))$.} y la sucesión espectral análoga para
$K'\otimes\Omega_X^{n}$. Del resultado anterior se deduce el teorema 6
utilizando la definición (8.1).

\src{192}
\subsection*{OBSERVACIONES}

\noindent 1\textsuperscript{o} Para las definiciones anteriores al teorema 6,
no era necesario que $X$ fuese no singular, pues no era
indispensable trabajar únicamente con resoluciones finitas. Pero si $X$
es singular, ya no se podrá afirmar, a priori, que los
$(\underline{T}_{r}^{s})^{p}(A_1,\ldots;B_1,\ldots)$ sean haces
coherentes, pues en el complejo de haces
\[
  \Homsh_{\Osh_X}\bigl(L(A_1)\otimes\cdots,L(B_1)\otimes\cdots\bigr)
\]
habrá una infinidad de componentes de un grado total dado.

\medskip
\noindent 2\textsuperscript{o} Se comprueba fácilmente que, en las fórmulas
(8.1), se puede sustituir uno de los $L(B_i)$ por $B_i$. Teniendo en cuenta la
proposición 3, esto demuestra que se tiene
\begin{equation}\tag{8.8}
\left\{
\begin{aligned}
  \underline{T}_{1}^{1*}(A;B)&=\Extsh_{\Osh_X}^{*}(A,B),\\
  T_{1}^{1*}(A;B)&=\operatorname{Ext}_{\Osh_X}^{*}(X;A,B).
\end{aligned}
\right.
\end{equation}
En particular, tomando en (8.6) $r=s=1$ y $A_1=\Osh_X$, se recupera el
teorema 3. La fórmula (8.8) implica también
\[
  T_{0}^{1*}(B)=H^{*}(X,B),
  \qquad
  T_{1}^{0*}(A)=\operatorname{Ext}_{\Osh_X}^{*}(X;A,\Osh_X).
\]

\medskip
\noindent 3\textsuperscript{o} Se ve en la fórmula (8.1) que, en general, los
funtores $\underline{T}_{r}^{s*}$ y $T_{r}^{s*}$ tienen componentes de
grados positivos y negativos. Utilizando la observación 2, se ve que si la
dimensión de $X$ es $n$, entonces las componentes no nulas de
$\underline{T}_{r}^{s*}$ están comprendidas entre $-(s-1)n$ y $rn$, si $s>0$,
entre $0$ y $rn$, si $s=0$, y las componentes no nulas de $T_{r}^{s*}$
están comprendidas entre $-(s-1)n$ y $(r+1)n$, si $s>0$, entre $0$ y $(r+1)n$,
si $s=0$ (e incluso, salvo error, si $r>0$, entre $-(s-1)n$ y $rn$, resp. entre
$0$ y $rn$).

\section*{BIBLIOGRAFÍA}

\noindent[1] Cartier (Pierre). --- Les groupes
$\operatorname{Ext}^{s}(A,B)$, Séminaire A. Grothendieck : Algèbre homologique,
t. 1, 1957, n\textsuperscript{o} 3.

\noindent[2] Grothendieck (Alexandre). --- Sur quelques points d'algèbre
homologique, \emph{Tohoku Math. J.}, t. 9, 1957, p. 119--183.

\noindent[3] Serre (Jean-Pierre). --- Faisceaux algébriques cohérents,
\emph{Annals of Math.}, t. 61, 1955, p. 197--278.

\noindent[4] Serre (Jean-Pierre). --- Sur la dimension homologique des anneaux
et des modules noethériens, \emph{Proc. Intern. Symp. on Alg. Number Theory}
[1955, Tokyo et Nikko]. --- Tokyo, Science Council of Japan, 1956 ;
p. 175--189.

\clearpage
\src{193}
\section*{ADENDA}

Las dificultades señaladas en la observación de la página 13 están actualmente
completamente resueltas, gracias a una extensión de los teoremas de dualidad a
variedades cualesquiera (con singularidades arbitrarias). Para esta teoría, que puede
formularse en el marco general de la « teoría de esquemas », remitimos
a los « Éléments de Géométrie algébrique » de J. Dieudonné y A. Grothendieck,
actualmente en preparación. Se encontrarán algunas indicaciones en:

\begin{quote}
Grothendieck (A.). --- Cohomology theory of algebraic varieties,
\emph{Congrès int. des Mathématiciens}, 1958, Edinburgh (de próxima aparición).
\end{quote}

\begin{flushright}
[Abril de 1959]
\end{flushright}

\begin{center}\rule{36mm}{0.5pt}\end{center}

\def\FgaUnitKey{e149}\clearpage
\chapter*{Fe de erratas de la exposición n\textsuperscript{o} 149}
\addcontentsline{toc}{chapter}{Fe de erratas de la exposición n\textsuperscript{o} 149}
\setcounter{footnote}{0}
\src{298}

\noindent\textbf{1. Teorema de dualidad para los haces algebraicos
coherentes.}

\begin{center}
[Seminario Bourbaki, t. 9, 1956/57, n\textsuperscript{o} 149, 25 p.]
\end{center}

\noindent\textbf{Observación sobre la página 14.} --- Como señalé en mi
conferencia en el Congreso Internacional de Matemáticos de 1958,\footnote{
Grothendieck (Alexander). --- The cohomology theory of abstract algebraic
varieties, \emph{Proceedings of the International Congress of Mathematicians}
[1958, Edinburgh], p. 103--118. --- Cambridge, at the University Press, 1960.}
las cuestiones planteadas aquí están ahora completamente resueltas.

El lector encontrará más información sobre la dualidad de los haces
coherentes en \emph{loco citato}\footnotemark[1], p. 112--115, en EGA III,
2\textsuperscript{a} parte (en preparación) y, mientras tanto, en SGA 1962.
Un tratamiento más sistemático se encontrará en un capítulo posterior de EGA
(capítulo IX en el plan previsto).

\def\FgaUnitKey{182}\unit{Exposición 182 — Geometría formal y geometría algebraica}
\sid{182}
\src{193}
\seminar{(Mayo de 1959)}
\paperhead{GEOMETRÍA FORMAL Y GEOMETRÍA ALGEBRAICA}%
  {por Alexander GROTHENDIECK}

\fsection{1}{Esquemas}

Se sabe que un espacio algebraico afín definido sobre un cuerpo $k$ está
esencialmente determinado por su álgebra afín $A$ (anillo de las funciones
regulares definidas sobre $k$), pues los morfismos de espacios algebraicos
$X\longrightarrow Y$ corresponden biunívocamente a los homomorfismos de
$k$-álgebras $A(Y)\longrightarrow A(X)$. El álgebra afín correspondiente a un
espacio algebraico es una $k$-álgebra de tipo finito y, desde el punto de vista
«clásico», no tiene elementos nilpotentes; recíprocamente, toda álgebra de este
tipo se obtiene como álgebra afín de un espacio algebraico definido sobre $k$.
Se dispone así de un conocido diccionario que permite interpretar en términos
de álgebra conmutativa las situaciones relativas a espacios algebraicos afines.
Desde hace mucho se ha observado que de este modo se obtienen enunciados más
generales, pues por lo común ya no es necesario suponer que los anillos en juego
sean del tipo recién considerado, y la hipótesis noetheriana suele bastar. En
particular, haya o no un cuerpo de base dado, no hay razón para excluir el caso
en que estos anillos contengan elementos nilpotentes. Hasta ahora, los geómetras
se habían negado a tener en cuenta estas indicaciones y se han obstinado en
limitarse a considerar álgebras afines sin elementos nilpotentes, es decir,
espacios algebraicos cuyos haces estructurales no tienen elementos nilpotentes
(y, las más de las veces, incluso espacios algebraicos «absolutamente
irreducibles»). El conferenciante considera que esta actitud ha constituido un
serio obstáculo para el desarrollo de métodos verdaderamente naturales en
Geometría algebraica.

Sea $A$ un anillo conmutativo. Es bien sabido que el conjunto
$X=\operatorname{Spec}(A)$ de los ideales primos de $A$ está provisto de una
topología natural, la «topología de Zariski» o topología espectral. Por otra
parte, hay un haz de anillos conmutativos $\Osh_X$ sobre $X$, cuyo tallo en
$\mathfrak p\in X$ es el anillo localizado $A_{\mathfrak p}$ y cuyo anillo de
secciones se identifica con $A$. Así, $X$ se convierte en un \emph{espacio
anillado}, llamado \emph{espectro primo} de $A$. Un homomorfismo de anillos
$f:A\longrightarrow B$ define un morfismo de espacios anillados
$f':\operatorname{Spec}(B)\longrightarrow\operatorname{Spec}(A)$; la
aplicación subyacente de conjuntos no es otra que
$\mathfrak p\longmapsto f^{-1}(\mathfrak p)$. Los morfismos

\src{194}
de espacios anillados de $\operatorname{Spec}(B)$ en
$\operatorname{Spec}(A)$ obtenidos de este modo son exactamente aquellos para
los cuales los homomorfismos $\Osh_x\longrightarrow\Osh_y$ ($x=f'(y)$) son
locales (es decir, la imagen inversa del ideal maximal es el ideal maximal).

Se llama \emph{esquema afín} a un espacio anillado isomorfo a algún
$\operatorname{Spec}(A)$, y \emph{preesquema} a un espacio anillado localmente
afín, es decir, tal que cada punto tiene un entorno abierto que, con la
estructura inducida, es un esquema afín. Los \emph{morfismos} de preesquemas se
definen de manera evidente; localmente corresponden a homomorfismos de anillos.

Cuando se fija un preesquema $S$ y se consideran morfismos de preesquemas
$X\longrightarrow S$, el espacio $S$ desempeña el papel de un cuerpo o de un
anillo de base (o, mejor, de un espacio de base en una fibración). Se dice
entonces que $X$ es un $S$-preesquema; si $S=\operatorname{Spec}(A)$, esto
significa también que $\Osh_X$ es un haz de $A$-álgebras. Así, todo preesquema
puede considerarse de manera única como un $\mathbf Z$-preesquema. Naturalmente,
los $S$-preesquemas forman una categoría; además, se demuestra que en esta
categoría siempre existe el producto de dos objetos $X,Y$, que se denota por
$X\times_S Y$. Esta noción de producto permite definir el cambio de base de un
$S$-preesquema correspondiente a un morfismo $S'\longrightarrow S$: en efecto,
$X\times_S S'$ puede considerarse como un $S'$-preesquema.

Se dice que $X$ es \emph{separado} sobre $S$ si la diagonal de
$X\times_S X$ es cerrada. Se llama \emph{esquema} a un preesquema separado
sobre $\mathbf Z$; es entonces separado sobre cualquier base. Para simplificar,
en adelante solo hablaremos de esquemas y, además, los supondremos
\emph{noetherianos}, es decir, uniones finitas de abiertos afines que son
espectros de anillos noetherianos. Se dice que $X$ es \emph{de tipo finito
sobre} $S$ si, para todo abierto afín $U$ de $S$, su imagen inversa en $X$ es
una unión finita de abiertos afines cuyos anillos son álgebras de tipo finito
sobre el anillo de $U$. Son estos $S$-esquemas los que se prestan a un estudio
propiamente geométrico. En particular, para todo $s\in S$, la fibra
$f^{-1}(s)$ de $X$ sobre $s$ es un esquema algebraico sobre el cuerpo residual
$\kappa(s)$ del anillo local $\Osh_s$ de $s$ en $S$. Así, en cierta medida, $X$
puede considerarse como una familia de «espacios algebraicos» $f^{-1}(s)$, donde
el parámetro $s$ recorre $S$ (es decir, desde el punto de vista local, el
conjunto de los ideales primos de un anillo dado). Naturalmente, los
$\kappa(s)$ pueden tener características distintas. Si
$S=\operatorname{Spec}(k)$, donde $k$ es un cuerpo, se recupera esencialmente la
noción usual de «espacio algebraico», con la única diferencia de que ahora el
haz estructural puede tener elementos nilpotentes.

\src{195}
Inspirándose en nociones bien conocidas, se define la noción de
\emph{morfismo proyectivo} y, más generalmente, la de \emph{morfismo propio}.
Tal morfismo es de tipo finito; además, transforma partes cerradas en partes
cerradas y conserva esta propiedad bajo cualquier cambio de base.

Siendo $X$ un esquema (noetheriano, como siempre), el haz $\Osh_X$ es un
\emph{haz coherente de anillos} en el sentido de [2]. Por tanto, los haces
coherentes de módulos sobre $X$ son también los haces que localmente son
isomorfos al conúcleo de un homomorfismo
$\Osh_X^{m}\longrightarrow\Osh_X^{n}$.

\fsection{2}{Esquemas formales}

Sean $X$ un esquema y $X'$ una parte cerrada de $X$. Entonces existe un
subhaz coherente $J$ de $\Osh_X$ tal que
$X'=\operatorname{Supp}(\Osh_X/J)$ (e incluso existe uno máximo). Provisto de
$\Osh_X/J$, el espacio $X'$ se convierte en un esquema, denotado por $X_0$;
tal esquema se llama «subesquema cerrado de $X$». Para todo $n$ también se
puede considerar $X'$ provisto de $\Osh_X/J^{n+1}$, que se denota por $X_n$;
es un subpreesquema cerrado de $X$ cuyo conjunto subyacente sigue siendo $X'$,
pero que tiene otro haz estructural, a saber,
$\Osh_{X_n}=\Osh_X/J^{n+1}$. Evidentemente, los $\Osh_{X_n}$ forman un sistema
proyectivo de haces de anillos sobre $X'$, cuyo límite proyectivo
$\overline{\Osh}_X$ se llama \emph{completación formal de $\Osh_X$ a lo largo
de $X'$}. Provisto de este haz de anillos, $X'$ se llama \emph{completación
formal de $X$ a lo largo de $X'$}; es, pues, un espacio anillado, pero en
general no un esquema. Para todo haz coherente $F$ sobre $X$ se puede considerar
análogamente la completación formal $\overline F=\varprojlim_n F_n$ de $F$ a lo
largo de $X'$ (donde $F_n=F\otimes_{\Osh_X}\Osh_X/J^{n+1}$); es un haz de
módulos sobre $\overline X$. Sus secciones se llaman «secciones formales de $F$
a lo largo de $X'$» y se identifican con los elementos de
$\varprojlim_n\Gamma(X',F_n)$. Para $F=\Osh_X$ se obtienen las «funciones
holomorfas» de $X$ a lo largo de $X'$ en el sentido de Zariski (terminología que
no seguiremos por sus interferencias con la terminología clásica).

Se llama \emph{esquema formal} (se sobreentiende: noetheriano) a un espacio
topológico $\mathfrak X$, provisto de un haz de anillos topológicos
$\Osh_{\mathfrak X}$ que satisface la condición siguiente: hay un isomorfismo
de haces de anillos topológicos
$\Osh_{\mathfrak X}=\varprojlim_n\Osh_n$, donde los $\Osh_n$ forman un sistema
proyectivo de haces de anillos sobre $\mathfrak X$, cada uno de los cuales hace
de $\mathfrak X$ un esquema $\mathfrak X_n$, y tal que, para $m\geq n$, el
homomorfismo $\Osh_m\longrightarrow\Osh_n$ es sobreyectivo y tiene por núcleo
$J_m^{n+1}$, donde $J_m$ es

\src{196}
el núcleo de $\Osh_m\longrightarrow\Osh_0$. Se demuestra entonces que
$\Osh_{\mathfrak X}$ es un \emph{haz coherente de anillos locales
noetherianos}.

En virtud de las definiciones, una completación formal $\overline X$ como la
anterior es un esquema formal y, recíprocamente, todo esquema formal es
\emph{localmente} de este tipo. De hecho, dar un esquema formal \emph{afín} (es
decir, tal que $\mathfrak X_0$ sea afín, lo cual implica que todos los
$\mathfrak X_n$ lo son) equivale a dar un anillo topológico noetheriano,
separado y completo para la topología $J$-ádica.

Las definiciones habituales de morfismo, morfismo de tipo finito, morfismo
propio, etc., para los esquemas ordinarios se extienden sin dificultad a los
esquemas formales.

\fsection{3}{Los tres teoremas fundamentales}

Sea $f:X\longrightarrow Y$ un morfismo propio de esquemas (noetherianos, como
siempre), sea $Y'$ una parte cerrada de $Y$ y sea $X'$ su imagen inversa en
$X$; consideremos las completaciones formales correspondientes $\overline Y$ y
$\overline X$. Entonces $f$ induce un morfismo
$\overline f:\overline X\longrightarrow\overline Y$ de esquemas formales, que
además es propio. Si $F$ es un haz coherente sobre $X$, entonces $\overline F$
es un haz coherente sobre $\overline X$. En el teorema 1 se olvidan $X,Y,F$ y
solo se consideran el morfismo propio $\overline f$ de esquemas formales y el
haz coherente $\overline F$ sobre $\overline X$. (No obstante, el conferenciante
solo ha escrito una demostración completa para el caso en que se parte de
$X,Y,f,F$.)

\result{TEOREMA}{1 (teorema de finitud)}
\begin{enumerate}
  \item[\textit{i.}] Los $R^{q}\overline f_{*}(\overline F)$ son haces
    coherentes sobre $\overline Y$.
  \item[\textit{ii.}] Los homomorfismos naturales
    \[
      R^{q}\overline f_{*}(\overline F)
      \longrightarrow \varprojlim_n R^{q}f_{n*}(F_n)
    \]
    son isomorfismos.
\end{enumerate}

En este enunciado se supone elegido un subhaz coherente $J$ de $\Osh_Y$ que
define $Y'$; por imagen inversa se obtiene un subhaz coherente de $\Osh_X$ que
define $X'$ y, por consiguiente, se definen $F_n,X_n,Y_n$ y
$f_n:X_n\longrightarrow Y_n$ como en el n\textsuperscript{o} 2. Si se parte de
un morfismo propio cualquiera entre dos esquemas formales, son evidentes los
pequeños cambios que deben hacerse al explicitar estas notaciones.

El teorema 1 se refería únicamente a la «cohomología formal». El teorema siguiente

\src{197}
la relaciona con la «cohomología algebraica» y es análogo a un conocido teorema
de Serre [4] sobre la comparación entre cohomología algebraica y cohomología
analítica.

\result{TEOREMA}{2 (Primer teorema de comparación)} Los
$R^{q}f_{*}(F)$ son haces coherentes sobre $Y$ (lo cual es un caso particular
del teorema 1), y los homomorfismos naturales
\[
  \overline{R^{q}f_{*}(F)}
  \longrightarrow \varprojlim_n R^{q}f_{n*}(F_n)
\]
son isomorfismos.

\result{COROLARIO}{1} Se tienen isomorfismos canónicos:
\[
  \overline{R^{q}f_{*}(F)}=R^{q}\overline f_{*}(\overline F).
\]

Para $q=0$, este corolario es una generalización del «teorema fundamental de
las funciones holomorfas» de Zariski, del que deduciremos una generalización del
«teorema de conexión» de Zariski. Obsérvese además que, mientras el teorema 1
(\textit{ii}) es trivial para $q=0$, en modo alguno ocurre lo mismo con el
teorema 2 ni con su formulación equivalente (corolario 1). En efecto, la
demostración procede por inducción descendente sobre $q$ (es trivial para $q$
grande, pues entonces ambos miembros son nulos), de modo que el caso $q=0$
aparece como el último paso de la inducción y, por así decirlo, como el caso «más
difícil».

\result{COROLARIO}{2} Supongamos $Y=\operatorname{Spec}(A)$ y que $Y'$ está
definido por el ideal $J$ de $A$. Entonces, para todo haz coherente $F$ sobre
$X$, los $H^{q}(X,F)$ son $A$-módulos de tipo finito cuyas completaciones
$J$-ádicas son los $H^{q}(\overline X,\overline F)$.

Aplicando, por último, este corolario a
$H=\Homsh_{\Osh_X}(F,G)$, se obtiene:

\result{COROLARIO}{3} Supongamos $Y=\operatorname{Spec}(A)$ y que $Y'$ está
definido por el ideal $J$ de $A$. Sean $F,G$ dos haces coherentes sobre $X$;
entonces $\operatorname{Hom}(F,G)$ es un módulo de tipo finito sobre $A$, cuya
completación $J$-ádica se identifica con
$\operatorname{Hom}(\overline F,\overline G)$.

Naturalmente, la aplicación natural
$\operatorname{Hom}(F,G)\longrightarrow
\operatorname{Hom}(\overline F,\overline G)$ es la que asocia a un
homomorfismo $u:F\longrightarrow G$ su prolongación «por continuidad»
$\overline u:\overline F\longrightarrow\overline G$ (con lo cual
$\overline F$ se convierte en un funtor de $F$).

Supongamos ahora que $A$ es separado y completo para su topología $J$-ádica;
entonces los corolarios 2 y 3 precedentes dan:
\[
  H^{q}(X,F)=H^{q}(\overline X,\overline F),
  \qquad
  \operatorname{Hom}(F,G)
  =\operatorname{Hom}(\overline F,\overline G).
\]

\src{198}
Esta última identidad muestra que la categoría de los haces coherentes sobre
$X$ se identifica con una \emph{subcategoría} (cuyos morfismos son los
morfismos inducidos) de la categoría de los haces coherentes sobre
$\overline X$. De hecho, se tiene incluso:

\result{TEOREMA}{3} Para que un haz de módulos sobre $\overline X$ sea
coherente, es necesario y suficiente que sea isomorfo a un haz de la forma
$\overline F$, donde $F$ es un haz coherente sobre $X$ (determinado salvo
isomorfismo canónico en virtud del teorema 2, corolario 3). [Recordemos que
ahora $Y=\operatorname{Spec}(A)$, siendo $A$ un anillo topológico noetheriano,
separado y completo para la topología $J$-ádica.]

\result{COROLARIO}{1} Los subesquemas cerrados de $X$ corresponden
biunívocamente a los subesquemas formales cerrados de $\overline X$.

En efecto, corresponden a los subhaces coherentes de $\Osh_X$ y,
respectivamente, de $\overline{\Osh}_X$. Considerando los grafos de morfismos
como subesquemas cerrados, se deduce del corolario 1:

\result{COROLARIO}{2} Sean $X,Z$ dos esquemas propios sobre $A$ (anillo
noetheriano, separado y completo para la topología $J$-ádica). Entonces la
aplicación $g\longmapsto\overline g$ define una correspondencia biunívoca entre
los $Y$-morfismos de $X$ en $Z$ y los $\overline Y$-morfismos de $\overline X$
en $\overline Z$.

En otros términos, los esquemas algebraicos propios sobre $A$ aparecen como una
\emph{subcategoría} (cuyos morfismos son los morfismos inducidos) de la
categoría de los esquemas formales propios sobre $\overline Y$. Sin embargo,
debe advertirse que existen esquemas formales propios sobre $\overline Y$ que
no son «algebraizables», es decir, isomorfos a algún $\overline X$ con $X$
propio sobre $A$ (del mismo modo que existen variedades analíticas complejas
compactas que no provienen de variedades algebraicas definidas sobre el cuerpo
de los complejos). Tales esquemas formales aparecen de manera natural en la
«teoría de módulos». Señalemos, no obstante, un caso particular interesante en
el que un esquema formal es algebraizable:

\result{TEOREMA}{4} Sea $A$ un anillo noetheriano local completo de cuerpo
residual $k$, y sea $\mathfrak X$ un esquema formal propio sobre $A$ (provisto
de su topología $\mathfrak m(A)$-ádica). Se supone:
\begin{enumerate}
  \item[\textit{i.}] que los anillos locales de $\Osh_{\mathfrak X}$ son
    módulos planos sobre $A$ o, lo que equivale a lo mismo, si se dota a
    $\Osh_{\mathfrak X}$ y a $A$ de la filtración definida por las potencias del
    ideal maximal de $A$, para los graduados asociados se tiene la relación
    \[
      \operatorname{gr}(\Osh_{\mathfrak X})
      \simeq
      \operatorname{gr}^{0}(\Osh_{\mathfrak X})
      \otimes_k\operatorname{gr}(A).
    \]
  \item[\textit{ii.}] Considerando
    $\mathfrak X_0=\mathfrak X\otimes_A k$ como un esquema algebraico sobre
    $k$, se tiene
    \[
      H^2(\mathfrak X_0,\Osh_{\mathfrak X_0})=0.
    \]
  \item[\textit{iii.}] $\mathfrak X_0$ es proyectivo.
\end{enumerate}
Bajo estas condiciones, $\mathfrak X$ es algebraizable y, más precisamente, es
isomorfo a $\overline X$, donde $X$ es un $A$-esquema proyectivo.

Las condiciones \textit{(ii)} y \textit{(iii)} se satisfacen, en particular,
si $\mathfrak X_0$ es una curva simple sobre $k$, y el teorema 4 puede aplicarse
en particular en la «teoría de módulos» para las curvas de género dado\ldots{}
Indiquemos cómo se demuestra el teorema 4: se demuestra (cf. la proposición 3
más adelante) que \textit{(i)} y \textit{(ii)} implican que todo haz coherente
sobre $\mathfrak X_0$, localmente isomorfo al haz fundamental, se obtiene por
reducción a partir de un haz de la misma naturaleza sobre $\mathfrak X$. Se
parte entonces de un haz «amplio» sobre $\mathfrak X_0$ (existe uno por
\textit{(iii)}), se lo levanta a un haz invertible sobre $\mathfrak X$ y,
utilizando el teorema 1, se demuestra que un múltiplo de este último define una
inmersión de $\mathfrak X$ en la completación formal de un esquema
$\mathbf P_A^r$ («proyectivo tipo» de dimensión $r$ sobre $A$).

Para la demostración de los teoremas 1 a 3, véase [1].

\src{199}
\fsection{4}{Aplicación al teorema de conexión y al «Main theorem» de
Zariski}

Sea $f:X\longrightarrow Y$ un morfismo propio de esquemas. Por el teorema de
finitud, $f_*(\Osh_X)=\mathcal A$ es un haz coherente sobre $Y$; además, es un
haz de álgebras conmutativas, por lo que corresponde a un $Y$-esquema
$g:Y'\longrightarrow Y$ finito sobre $Y$ (definido por la condición de ser afín
sobre $Y$, es decir, que la imagen inversa de un abierto afín sea afín, y por
$g_*(\Osh_{Y'})=\mathcal A$). Es inmediato que $f$ se factoriza entonces
canónicamente como $f=gf'$, donde $f':X\longrightarrow Y'$ es un morfismo de
$X$ en $Y'$ que esta vez satisface $f'_*(\Osh_X)=\Osh_{Y'}$. Esta factorización
de $f$ se llama \emph{factorización de Stein} de $f$. Aplicando el primer
teorema de comparación, corolario 1, a $f'$ y a la parte de $Y'$ reducida a un
punto $y'$, se obtiene que $f'^{-1}(y')=X'$ es conexo (de otro modo, las
secciones formales de $X$ a lo largo de $X'$ no formarían un anillo local,
¡mientras que la completación de
$f'_*(\Osh_X)_{y'}=\Osh_{Y',y'}$ es local!). Hemos demostrado:

\result{TEOREMA}{5 («Teorema de conexión» de Zariski)} Sea
$f:X\longrightarrow Y$ un morfismo propio. Entonces $f$ se factoriza de manera
única (salvo isomorfismo) como $f=gf'$, donde $g:Y'\longrightarrow Y$ es finito
y $f':X\longrightarrow Y'$ satisface $f'_*(\Osh_X)=\Osh_{Y'}$ (de donde
$g_*(\Osh_{Y'})=f_*(\Osh_X)$). Las fibras de $f'$ son conexas; es decir, el
conjunto de las componentes conexas de una fibra $f^{-1}(y)$ de $f$ está en
correspondencia biunívoca con el conjunto de los puntos de $Y'$ sobre $y$, esto
es, con el conjunto de los ideales maximales de $f_*(\Osh_X)_y$.

De aquí se deducen inmediatamente las variantes habituales del teorema de
conexión. Señalemos aquí únicamente el

\src{200}
\result{COROLARIO}{1} Para que un punto $x$ de $X$ esté aislado en su fibra
$f^{-1}(y)$, es necesario y suficiente que la fibra $f'^{-1}(y')$ (donde
$y'=f'(x)$) se reduzca a $x$, o bien que $f'$ induzca un isomorfismo de un
entorno de $x$ sobre un entorno de $y'$. El conjunto de estos puntos es un
abierto $U$, y $f'$ induce un isomorfismo de $U$ sobre un abierto de $Y'$.

Para demostrar que $f'$ es un isomorfismo local en $x$, se observa que $f'$
induce un isomorfismo $\Osh_{y'}\longrightarrow\Osh_x$, como se ve gracias a
$f'_*(\Osh_X)=\Osh_{Y'}$ y al hecho de que (siendo $f'$ una aplicación cerrada
cuya fibra en $y'$ se reduce a $x$) los $f'^{-1}(V)$ recorren un sistema
fundamental de entornos de $x$ cuando $V$ recorre un sistema fundamental de
entornos de $y'$. De aquí se deduce enseguida el resultado siguiente, debido a
Chevalley en el caso «geométrico».

\result{COROLARIO}{2} Para que $f$ sea un morfismo finito, es necesario y
suficiente que sea propio y que sus fibras sean finitas.

Si así ocurre, $f'$ es efectivamente un isomorfismo por lo anterior.

Sea $f:X\longrightarrow Y$ un morfismo no necesariamente propio, pero
supongamos que $X$ está contenido como abierto en un $Y$-esquema propio
$\widetilde f:\widetilde X\longrightarrow Y$ (lo cual ocurre, en particular,
si $f$ es cuasiproyectivo). Aplicando el corolario 1, se ve que $\widetilde f'$
induce un isomorfismo del conjunto $U$ de los puntos de $X$ aislados en su fibra
sobre un abierto de $Y'$ (y que $U$ es, en efecto, abierto). Se deduce la
siguiente variante global del «Main Theorem» de Zariski.

\result{TEOREMA}{6 («Main Theorem» de Zariski)} Sea
$f:X\longrightarrow Y$ un morfismo de tipo finito. Entonces el conjunto $U$ de
los puntos de $X$ que están aislados en su fibra es abierto y, si $f$ es
separado,\footnote{La fe de erratas de 1962 sustituye aquí la hipótesis impresa
«$f$ cuasiproyectivo» por la hipótesis más débil «$f$ separado», en virtud de
SGA VIII, 6.2: todo morfismo cuasifinito y separado es cuasiproyectivo.} $U$ es
$Y$-isomorfo a una parte abierta de un esquema $Y'$ finito sobre $Y$.

Como un morfismo de tipo finito es localmente afín y, \textit{a fortiori},
localmente cuasiproyectivo, del teorema 6 se deducen enseguida las variantes
locales usuales del Main Theorem.

\fsection{5}{Aplicación al estudio cohomológico de los morfismos propios y
planos}

Sea $f:X\longrightarrow Y$ un morfismo propio, y sea $F$ un haz coherente sobre
$X$; se supone que $F$ es $Y$-plano, es decir, los $F_x$ son módulos planos sobre los
anillos $\Osh_y$ ($y=f(x)$). Esto significa también que, para todo $y\in Y$, si
filtramos $F$ a lo largo de la fibra $f^{-1}(y)$ mediante los
$\mathfrak m_y^nF$ (siendo $\mathfrak m_y$ el ideal maximal de $\Osh_y$), el
graduado asociado es isomorfo a
$(F/\mathfrak m_yF)\otimes_{\kappa(y)}\operatorname{gr}(\Osh_y)$; en otros
términos, se tiene
\[
  \mathfrak m_y^nF/\mathfrak m_y^{n+1}F
  =F_y\otimes_{\kappa(y)}
  (\mathfrak m_y^n/\mathfrak m_y^{n+1})
\]
para todo entero $n$, donde $X_y$ denota la fibra $f^{-1}(y)$ considerada como
esquema propio sobre el cuerpo residual $\kappa(y)$ de $y$, y $F_y$ el haz
$F/\mathfrak m_yF$ inducido por $F$ sobre $X_y$. Teniendo en cuenta este
isomorfismo y el teorema 2, se obtienen cotas y, en ocasiones, cálculos de los
$R^qf_*(F)$ en un entorno de $y$, conociendo la cohomología de $X_y$ con
coeficientes en $F_y$. El teorema 2 adopta aquí la forma
\[
  \widehat{R^qf_*(F)_y}
  =\varprojlim_n H^q(X_y,F/\mathfrak m_y^nF).
\]
Señalaremos únicamente la consecuencia siguiente:

\src{201}
\result{PROPOSICIÓN}{1} Sean $f:X\longrightarrow Y$ un morfismo propio y $F$
un haz sobre $X$, coherente y $Y$-plano. Sea $y\in Y$, sea $q$ un entero, y
supongamos que $H^q(X_y,F_y)=0$. Entonces $R^qf_*(F)$ es nulo en un entorno de
$y$ y, para todo $n$, el homomorfismo natural
\[
  R^{q-1}f_*(F)_y\longrightarrow
  H^{q-1}(X_y,F/\mathfrak m_y^nF)
\]
es sobreyectivo.

En particular, si $f$ es un morfismo plano (es decir, $\Osh_X$ es $Y$-plano),
entonces todo haz coherente localmente libre $F$ sobre $X$ es $Y$-plano. Sean
$F,G$ dos haces de este tipo; aplicando la proposición 1 a
$\Homsh_{\Osh_X}(F,G)$ y a $q=1$, se obtiene:

\result{TEOREMA}{7} Sea $f$ un morfismo propio y plano, sean $F,G$ dos haces
coherentes localmente libres sobre $X$, sea $y\in Y$ y supongamos que
\[
  H^1\bigl(X_y,\Homsh_{\Osh_{X_y}}(F_y,G_y)\bigr)=0,
\]
entonces todo homomorfismo $u_0:F_y\longrightarrow G_y$ está inducido por un
homomorfismo
\[
  u:F|V\longrightarrow G|V,
  \qquad V=f^{-1}(U),
\]
donde $U$ es un entorno de $y$.

\result{COROLARIO}{} Si $u_0$ es un isomorfismo (resp. un monomorfismo, un
epimorfismo), entonces lo mismo ocurre con $u$ para $U$ suficientemente
pequeño.

En particular:

\src{202}
\result{COROLARIO}{2} Sea $E_0$ un haz coherente localmente libre sobre $X_y$
tal que
\[
  H^1\bigl(X_y;\Homsh_{\Osh_{X_y}}(E_0,E_0)\bigr)=0,
\]
entonces dos haces localmente libres cuyas restricciones a $X_y$ sean isomorfas
a $E_0$ son isomorfos en un entorno de $X_y$.

Así:

\result{COROLARIO}{3} Supongamos $H^1(X_y,\Osh_{X_y})=0$. Entonces dos haces
invertibles sobre $X$ (es decir, localmente isomorfos a $\Osh_X$) cuyas
restricciones a $X_y$ sean isomorfas son isomorfos en un entorno de $X_y$.\footnote{El texto impreso no precisa el entorno. Se restituye aquí el alcance local del corolario 2, utilizado también en la demostración de la proposición 2.}

Se deduce:

\result{PROPOSICIÓN}{2} Sean $Y$ un esquema conexo y $E$ un haz coherente localmente libre de rango al menos $2$ sobre $Y$; consideremos el fibrado de espacios proyectivos
$X=\mathbf P(E)$ asociado a $E$, provisto de su conocido haz invertible
$\Osh_X(1)$. Entonces todo haz invertible $L$ sobre $X$ es isomorfo a un haz de
la forma
\[
  f^*(L')\otimes\Osh_X(n),
\]
donde $L'$ es un haz invertible sobre $Y$ y $n$ un entero. Este último está
determinado de manera única, y $L'$ está determinado salvo isomorfismo.

El corolario 3 precedente demuestra que $L$ es isomorfo a algún $\Osh_X(n)$ en
un entorno de cada fibra. El resto es casi formal.

La proposición 2 permite determinar los $Y$-morfismos de $X=\mathbf P(E)$ en
otro fibrado proyectivo. En particular, se obtiene:

\result{COROLARIO}{} Sea $u$ un automorfismo de $X=\mathbf P(E)$. Entonces
existen un haz invertible $L'$ sobre $Y$ y un isomorfismo $v$ de $E$ sobre
$E\otimes L'$ tales que $u$ es el isomorfismo correspondiente
\[
  \mathbf P(E)\xrightarrow{\sim}\mathbf P(E\otimes L')=\mathbf P(E);
\]
el par $(v,L')$ está determinado salvo isomorfismo.

Sea $\Gamma$ el conjunto de las clases de fibrados invertibles $L'$ sobre $Y$
tales que $E\otimes L'$ sea isomorfo a $E$. Sus elementos son de torsión, pues,
si $n$ es el rango de $E$, se ve (tomando las potencias exteriores $n$-ésimas)
que debe tenerse $L'^{\otimes n}\simeq\Osh_Y$. El corolario puede expresarse
entonces diciendo que se tiene una sucesión exacta de grupos
\[
  e\longrightarrow
  \operatorname{Aut}(E)/\Gamma(Y,\Osh_Y^*)
  \longrightarrow \operatorname{Aut}_Y(X)
  \longrightarrow \Gamma\longrightarrow e
\]
(deducida, por lo demás, de la sucesión exacta de cohomología obtenida a partir
de la siguiente sucesión exacta de haces de grupos
\[
  e\longrightarrow \Osh_Y^*
  \longrightarrow \underline{\operatorname{Aut}}(E)
  \longrightarrow \underline{\operatorname{Aut}}_Y(X)
  \longrightarrow e,
\]
donde $\Osh_Y^*$ es el haz de las «unidades» de $\Osh_Y$, identificado con el
centro de $\underline{\operatorname{Aut}}(E)$).

\src{203}
\fsection{6}{Aplicación a teoremas de existencia y unicidad para haces y
esquemas sobre un anillo completo para la topología $J$-ádica}

El teorema 7 daba un resultado de unicidad para haces coherentes localmente
libres utilizando los teoremas 1 y 2. Utilizando el teorema 3, deducimos ahora
teoremas de existencia de haces, de morfismos de esquemas o de esquemas. En lo
sucesivo, $A$ denota un anillo noetheriano local, separado y completo. El método
general consiste siempre en efectuar construcciones formales, lo cual consiste
esencialmente en hacer geometría algebraica sobre un anillo artiniano, y extraer
de ellas conclusiones de naturaleza «algebraica» mediante los tres teoremas
fundamentales.

\result{PROPOSICIÓN}{3} Sea $\mathfrak X$ un esquema formal propio y plano
sobre $A$, y sea $F_0$ un haz localmente libre sobre $X_0$ tal que
\[
  H^2\bigl(X_0,\Homsh_{\Osh_{X_0}}(F_0,F_0)\bigr)=0.
\]
Entonces existe un haz localmente libre $F$ sobre $\mathfrak X$ que induce
sobre $X_0$ un haz isomorfo a $F_0$. (Además, este $F$ es único salvo
isomorfismo si
$H^1(X_0,\Homsh_{\Osh_{X_0}}(F_0,F_0))=0$.)

Se construyen sucesivamente haces localmente libres $F_n$ sobre los $X_n$, cada
uno de los cuales induce el anterior. La construcción de $F_n$ encuentra una
obstrucción en
\[
  H^2\bigl(X_0,\Homsh_{\Osh_{X_0}}(F_0,F_0)\bigr)
  \otimes_{A/J}(J^n/J^{n+1}),
\]
que es, por tanto, nula por hipótesis. Utilizando ahora el teorema 3, se
obtiene:

\result{COROLARIO}{1} Sea $X$ un esquema propio y plano sobre $A$, y sea $F_0$
como arriba. Entonces existe un haz localmente libre $F$ sobre $X$ que induce
sobre $X_0$ un haz isomorfo a $F_0$. Además, este $F$ es único salvo
isomorfismo si
$H^1(X_0,\Homsh_{\Osh_{X_0}}(F_0,F_0))=0$.

Sea $X_0$ un esquema de tipo finito sobre el cuerpo $k$; se supone que $X_0$ es
simple (por lo cual entendemos absolutamente simple) sobre $k$, pero no
necesariamente propio sobre $k$. Sea $A$ un anillo artiniano local de cuerpo
residual $k$. Nos interesa hallar los esquemas $X$ planos sobre $A$ tales que
$X\otimes_A k=X_0$ (este es el punto de partida de la «teoría de módulos» o de
las «variaciones de estructura» de $X_0$). Es equivalente dar un tal $X$, o dar
sobre el espacio topológico $X_0$ un haz $\Osh_X$ provisto de las estructuras
siguientes:
\begin{enumerate}
  \item[\textit{i.}] Es un haz de $A$-álgebras.
  \item[\textit{ii.}] Está provisto de un homomorfismo de aumentación
    $\Osh_X\longrightarrow\Osh_{X_0}$ (compatible con las estructuras de
    $A$-álgebras); estos datos están sujetos a las condiciones siguientes: la
    aumentación induce un isomorfismo
    $\Osh_X\otimes_A k\xrightarrow{\sim}\Osh_{X_0}$; $\Osh_X$ es plano sobre
    $A$, es decir, el graduado asociado a $\Osh_X$, filtrado por las potencias
    del ideal maximal $\mathfrak m$ de $A$, es isomorfo a
    $\operatorname{gr}^0(\Osh_X)\otimes_k\operatorname{gr}(A)$; es decir, se
    tienen isomorfismos
    \[
      \mathfrak m^n\Osh_X/\mathfrak m^{n+1}\Osh_X
      =\Osh_{X_0}\otimes_k(\mathfrak m^n/\mathfrak m^{n+1}).
    \]
\end{enumerate}
El hecho fundamental es el siguiente:

\src{204}
\result{TEOREMA}{8} Sea $X_0$ un esquema de tipo finito y simple sobre el
cuerpo $k$, y supongamos que $X_0$ es afín. Sea $A$ un anillo local artiniano de
cuerpo residual $k$. Entonces existe un $A$-esquema $X$ plano sobre $A$ tal que
$X\otimes_A k=X_0$, y dos esquemas de este tipo son necesariamente isomorfos.

Obsérvese que el isomorfismo en cuestión no es canónico, pues en general $X$
tendrá $A$-automorfismos no triviales que inducen la identidad sobre $X_0$. Por
otra parte, en general no hay una elección «canónica» de un $X$ que satisfaga
las condiciones dadas, salvo cuando $A$ es una $k$-álgebra (caso de
características iguales), en cuyo caso puede tomarse $X=X_0\otimes_k A$, es
decir, $\Osh_X=\Osh_{X_0}\otimes_k A$ (sea o no afín $X_0$). En el caso de
características distintas, ignoro en general, cuando $X_0$ no es afín, si se
puede «levantar» $X_0$ a un $X$ definido sobre $A$. Sin embargo, sea $n$ un
entero $>0$, pongamos $A_{n-1}=A/\mathfrak m^n$, y supongamos que se ha
levantado $X_0$ a un $A_{n-1}$-esquema plano $X_{n-1}$; se pretende levantar
$X_{n-1}$ a un $A_n$-esquema plano $X_n$. Por el teorema 8 ya se sabe que esto
es posible localmente; por otra parte, se comprueba fácilmente que, si $U_n$
levanta un abierto $U_{n-1}$ de $X_{n-1}$, entonces el haz de los grupos de
automorfismos de $U_n$ (que inducen la identidad sobre $U_{n-1}$) es
canónicamente isomorfo a
\[
  \Theta_{X_0/k}\otimes_k(\mathfrak m^n/\mathfrak m^{n+1})
\]
restringido a $U_{n-1}$ y, en particular, es conmutativo
($\Theta_{X_0/k}$ denota el haz de gérmenes de $k$-derivaciones sobre $X_0$).
Se sigue fácilmente que hay una obstrucción para construir $X_n$ que levante
$X_{n-1}$, la cual se encuentra en
\[
  H^2(X_0,\Theta_{X_0/k})\otimes_k
  \mathfrak m^n/\mathfrak m^{n+1}.
\]
Por consiguiente:

\result{COROLARIO}{1} Sea $X_0$ un esquema de tipo finito simple sobre $k$, y
supongamos $H^2(X_0,\Theta_{X_0/k})=0$. Entonces, para todo anillo local
artiniano $A$ de cuerpo residual $k$, existe un $A$-esquema plano $X$ tal que
$X\otimes_A k=X_0$.

Además, si se ha podido hallar un $X$ plano sobre $A$ que levante $X_0$,
entonces, en virtud del teorema 8, el conjunto de las clases (salvo isomorfismo)
de $A$-esquemas planos que levantan $X_0$ se identifica con
\[
  H^1\bigl(X_0,\underline{\operatorname{Aut}}(X)\bigr),
\]
donde, naturalmente, $\underline{\operatorname{Aut}}(X)$ denota el haz de
gérmenes de automorfismos del haz de $A$-álgebras $\Osh_X$ compatibles con la
aumentación. La filtración de $\Osh_X$ define una filtración de
$\underline{\operatorname{Aut}}(X)$, y el cociente de este haz por el $n$-ésimo
subgrupo de la filtración es $\underline{\operatorname{Aut}}(X_n)$; el graduado
asociado a esta filtración es conmutativo y se identifica con
$\Theta_{X_0/k}\otimes_k\operatorname{gr}(A)$. En particular, si
$\mathfrak m^{n+1}$ es la primera potencia de $\mathfrak m$ que es nula,
entonces $F^n(\underline{\operatorname{Aut}}(X))$ (el último término de la
filtración) pertenece al centro de $\underline{\operatorname{Aut}}(X)$ y es
isomorfo a $\Theta_{X_0/k}\otimes_k\mathfrak m^n$; es también el haz de
gérmenes de automorfismos de $X$ que inducen la identidad sobre
$X_{n-1}=X\otimes_A A/\mathfrak m^n$. Utilizando estos resultados se obtienen
enseguida los enunciados siguientes:

\src{205}
\result{COROLARIO}{2} Sea $X_0$ un esquema de tipo finito simple sobre $k$, sea
$A$ un anillo local artiniano de cuerpo residual $k$ e ideal maximal
$\mathfrak m$, supongamos $\mathfrak m^{n+1}=0$, sea
$A_{n-1}=A/\mathfrak m^n$ y, por último, sea $X_{n-1}$ un
$A_{n-1}$-esquema plano tal que $X_{n-1}\otimes_A k=X_0$. Entonces el conjunto
de las clases (salvo un isomorfismo que induzca la identidad sobre $X_{n-1}$)
de $A$-esquemas planos $X_n$ tales que
$X_n\otimes_A A_{n-1}=X_{n-1}$ es vacío o es un espacio homogéneo principal
bajo
\[
  H^1(X_0,\Theta_{X_0/k})\otimes_k\mathfrak m^n.
\]

(Obsérvese que, en general, no hay un elemento origen privilegiado en este
último espacio, pues no hay una manera privilegiada de levantar $X_{n-1}$ a
$X_n$).

\result{COROLARIO}{3} Sea $X_0$ un esquema de tipo finito simple sobre $k$, y
supongamos $H^1(X_0,\Theta_{X_0/k})=0$. Entonces, para todo anillo local
artiniano $A$ de cuerpo residual $k$, existe a lo sumo (salvo isomorfismo) un
$A$-esquema plano $X$ tal que $X\otimes_A k=X_0$.

Los corolarios 1 y 3 implican inmediatamente los enunciados, en apariencia más
generales, que se obtienen suponiendo solamente que $A$ es un anillo local
noetheriano completo de cuerpo residual $k$, con tal de introducir allí $X$
como un esquema formal sobre $A$:

\result{TEOREMA}{9} Sean $k$ un cuerpo y $X_0$ un esquema de tipo finito simple
sobre $k$. Para todo anillo local noetheriano completo $A$ de cuerpo residual
$k$, sea $F(A)$ el conjunto de las clases (salvo un isomorfismo que induzca la
identidad sobre $X_0$) de esquemas formales $\mathfrak X$ sobre $A$, de tipo
finito y planos sobre $A$, tales que $\mathfrak X\otimes_A k=X_0$. Con estas
notaciones:
\begin{enumerate}
  \item[\textit{i.}] Si $H^1(X_0,\Theta_{X_0/k})=0$, entonces, para todo $A$,
    $F(A)$ tiene a lo sumo un elemento.
  \item[\textit{ii.}] Si $H^2(X_0,\Theta_{X_0/k})=0$, entonces, para todo $A$,
    $F(A)$ tiene al menos un elemento.
\end{enumerate}

\src{206}
\result{COROLARIO}{1} Supongamos que $X_0$ es propio sobre $k$. Bajo la
condición \textit{(i)}, para todo $A$ existe a lo sumo (salvo un isomorfismo que
induzca la identidad sobre $X_0$) un esquema $X$ propio y plano sobre $A$ tal
que $X\otimes_A k=X_0$.

Se utiliza el teorema 3, corolario 2. Por ejemplo:

\result{COROLARIO}{2} Si $X$ es un esquema propio y plano sobre $A$ tal que
$X\otimes_A k$ es isomorfo al esquema proyectivo tipo $\mathbf P_k^r$ de
dimensión $r$ sobre $k$, entonces $X$ es isomorfo a $\mathbf P_A^r$.

(También se puede deducir este resultado de la proposición 3, corolario 1).

\result{COROLARIO}{3} Sea $X_0$ un esquema simple y proyectivo sobre $k$, y
supongamos que se tiene
\[
  H^2(X_0,\Osh_{X_0})=H^2(X_0,\Theta_{X_0/k})=0.
\]
Entonces, para todo $A$, existe un esquema proyectivo y plano $X$ sobre $A$ tal
que $X\otimes_A k=X_0$.

Se combinan el teorema 9\textit{(ii)} y el teorema 4. En particular:

\result{COROLARIO}{4} Sea $X_0$ el esquema de una curva algebraica completa y
simple sobre $k$. Entonces, para todo anillo local noetheriano completo $A$ de
cuerpo residual $k$, existe un «esquema de curva simple» $X$ sobre $A$ tal que
$X\otimes_A k=X_0$.

\result{OBSERVACIONES}{}

1) Los corolarios 3 y 4 son especialmente interesantes si $k$ es de
característica $p\ne0$, tomando como $A$ un anillo de valoración discreta de
característica $0$ y cuerpo residual $k$; por ejemplo, el «menor $A$ posible»,
a saber, aquel cuyo ideal maximal está generado por $p$. (De hecho, por los
teoremas de Cohen, basta conocer los corolarios 3 y 4 para un anillo $A$ de este
tipo). Señalemos a este respecto que, según los especialistas, no se sabe si
existen esquemas sobre un cuerpo $k$ que no sean la reducción módulo $p$ de un
esquema plano definido sobre un tal anillo $A$. Al menos, los resultados de este
número proporcionan un medio de investigar sistemáticamente esta cuestión. Se
debería comenzar por examinar si la primera obstrucción que aparece en
$H^2(X_0,\Theta_{X_0/k})$ es necesariamente nula.

\emph{Complemento de 1962.} --- J.-P. Serre ha construido una variedad
proyectiva no singular, de dimensión 3, sobre un cuerpo algebraicamente cerrado
$k$ de característica $p>0$, que no se obtiene por reducción de un esquema
propio sobre un anillo local íntegro de cuerpo residual $k$ y cuerpo de
fracciones de característica cero. Mumford habría encontrado un ejemplo análogo
con una superficie proyectiva no singular.

2) Obsérvese que el teorema 3, y la técnica correspondiente, solo es válido para
un anillo de base completo (local, para fijar ideas).

Para pasar de resultados conocidos para la completación de un anillo local a
los resultados correspondientes para ese mismo anillo local, se necesitaría un
cuarto «teorema fundamental», cuyo enunciado definitivo aún está por hallar.

\src{207}
3) Compárense los resultados de este número (en particular, los corolarios 1 y
2 precedentes) y del siguiente con los resultados de Kodaira--Spencer sobre la
variación de las estructuras complejas. Utilizando el teorema conjetural al que
se acaba de aludir, bajo las condiciones del corolario 1, pero sin suponer ya
que $A$ sea completo, se debería poder concluir que existe un anillo $A'$ que
contiene a $A$, finito y no ramificado sobre $A$, tal que $X\otimes_A A'$ y
$X'\otimes_A A'$ son $A'$-isomorfos (siendo $X,X'$ dos $A$-esquemas propios y
planos dados, tales que $X\otimes_A k=X'\otimes_A k=X_0$). Esto es al menos lo
que puede demostrarse cuando $X_0=\mathbf P_k^r$, utilizando el corolario de la
proposición 2. En todo caso, cuando $H^1(X_0,\Theta_{X_0/k})=0$, puede
demostrarse que las fibras de $X,X'$ sobre todo punto $y$ de
$Y=\operatorname{Spec}(A)$ son isomorfas, al menos después de pasar a la
clausura algebraica del cuerpo residual $\kappa(y)$. (Se tiene un resultado
local, aparentemente más fuerte, sin suponer ya necesariamente que $A$ sea
local). En cuanto a las «variaciones de estructura» del espacio proyectivo,
señalemos además la cuestión siguiente, sugerida por un problema correspondiente
de Kodaira--Spencer. Sea $X$ un esquema propio y plano sobre el anillo local
íntegro $A$ de cuerpo de fracciones $K$ y cuerpo residual $k$, y supongamos que
$X\otimes_A K$ es isomorfo a $\mathbf P_K^r$; ¿es cierto que
$X\otimes_A k=X_0$ es isomorfo a $\mathbf P_k^r$ (al menos sobre la clausura
algebraica de $k$)? En este problema se puede suponer que $A$ es un anillo de
valoración discreta completo. Se plantea un problema análogo cuando $X_0$ es
una variedad abeliana.

\emph{Complemento de 1962.} --- El conferenciante es ahora menos optimista
respecto de los resultados conjeturados arriba. No obstante, Hironaka ha
resuelto afirmativamente el problema relativo a las variaciones de estructura
del espacio proyectivo, y Koizumi el problema análogo para las variedades
abelianas.

\fsection{7}{Aplicación a la «Teoría de módulos»}

Como el propio conferenciante apenas acaba de abordar esta teoría, nos veremos
obligados a limitarnos a algunas indicaciones. Para simplificar, nos situamos
sobre un cuerpo $k$, es decir, trabajamos en igual característica, aunque el
teorema 8 permite tratar también el caso general sin cambio esencial, al
parecer. Actualmente no hemos superado la etapa «formal»; sin embargo, el
conferenciante espera poder construir a partir de ella verdaderos esquemas de
módulos en ciertos casos y, en particular, construir para cada entero $g$ un
esquema sobre los enteros que desempeñe el papel de esquema de módulos universal
para las curvas simples de género $g$.

\emph{Complemento de 1962.} --- Mumford ha construido los esquemas de módulos
para las curvas de género $g$ (cf. SHMI). El teorema 10 demuestra además que
los «esquemas de Jacobi de nivel $n$» de la teoría de Módulos son no singulares,
e incluso simples sobre $\mathbf Z$.

\src{208}
Retomemos la situación y las notaciones del teorema 9, suponiendo ahora que $A$
es un álgebra local de rango finito sobre $k$, que para simplificar se supone
algebraicamente cerrado. Entonces $F(A)$ puede considerarse como un funtor
covariante de $A$, con valores en la categoría de conjuntos: un homomorfismo de
$k$-álgebras $A\longrightarrow B$ define una aplicación
$F(A)\longrightarrow F(B)$, pues todo $A$-esquema plano $X$ tal que
$X\otimes_A k=X_0$ da lugar a un $B$-esquema $X\otimes_A B$ con las mismas
propiedades. Supongamos que se puede hallar una $k$-álgebra local noetheriana
completa $\mathcal O$ y un isomorfismo funtorial
\begin{equation}
  \tag{*}
  \operatorname{Hom}(\mathcal O,A)\xrightarrow{\sim}F(A)
\end{equation}
(donde el primer miembro denota los homomorfismos de $k$-álgebras). Se ve
fácilmente que un tal $\mathcal O$ está determinado salvo isomorfismo canónico;
se llama entonces al espectro formal $\mathfrak Y$ de $\mathcal O$ (es decir,
al espacio topológico reducido a un punto y provisto de un haz de anillos
topológicos reducido a $\mathcal O$) el \emph{esquema formal de módulos} para
$X_0$. (Obsérvese que no existe necesariamente). Sea $\mathfrak r$ el ideal
maximal de $\mathcal O$; para cada $n$, sea
$\mathcal O_n=\mathcal O/\mathfrak r^{n+1}$ (por tanto $\mathcal O_0=k$).
Entonces el homomorfismo canónico $\mathcal O\longrightarrow\mathcal O_n$ es un
elemento de $\operatorname{Hom}(\mathcal O,\mathcal O_n)$, por lo que define un
elemento de $F(\mathcal O_n)$, es decir, un $\mathcal O_n$-esquema plano $X_n$
cuya reducción módulo $\mathfrak r$ es $X_0$. Estos $X_n$ se deducen unos de
otros por extensión de escalares (es decir, aquí por reducciones), de donde
resulta que provienen de un esquema formal $\mathfrak X$ bien determinado sobre
el esquema formal de módulos $\mathfrak Y$; $\mathfrak X$ es plano sobre
$\mathfrak Y$ y $\mathfrak X_0=X_0$. Como se ve enseguida, el isomorfismo (*)
se obtiene entonces asociando a todo homomorfismo de $k$-álgebras
$\mathcal O\longrightarrow A$ la clase del $A$-esquema
$\mathfrak X\otimes_{\mathcal O}A$ (es decir, a todo morfismo de $k$-esquemas
$Y'=\operatorname{Spec}(A)\longrightarrow\mathfrak Y$ se asocia el
$Y'$-esquema $\mathfrak X\otimes_{\mathfrak Y}Y'$ obtenido por cambio de base).
Además, se ve que el isomorfismo (*) y su descripción precedente siguen siendo
válidos si solo se supone que $A$ es una $k$-álgebra local noetheriana y
completa (no necesariamente artiniana). Naturalmente, como de costumbre,
$\mathcal O$ puede tener perfectamente \textit{a priori} elementos nilpotentes,
y parece probable que existan casos en que $\mathcal O$ sea él mismo artiniano
sin ser idéntico a $k$. Esto muestra hasta qué punto el punto de vista de
Kodaira--Spencer (que se limita a tomar como $A$ anillos regulares) es
inadecuado \textit{a priori} en el caso general.

\src{209}
Queda por dar condiciones suficientes para que exista un esquema formal de
módulos para $X_0$, que se supone propio sobre $k$. En general, es fácil dar
condiciones necesarias y suficientes sencillas sobre un funtor
$A\longmapsto F(A)$ (de $k$-álgebras locales de rango finito, con valores en
los conjuntos) para que pueda escribirse en la forma
$\operatorname{Hom}(\mathcal O,A)$ para un $\mathcal O$ adecuado. No las
detallaremos aquí. Señalemos solamente que, en el caso que nos ocupa, estas
condiciones imponen sobre $X_0$ condiciones cohomológicas no triviales, y parece
poco probable que se satisfagan siempre, aunque el conferenciante no ha
construido ningún contraejemplo. En cambio, parece plausible que la condición
$H^0(X_0,\Theta_{X_0/k})=0$ sea una condición suficiente (aunque en modo alguno
necesaria) para la existencia de un esquema formal de módulos. Nos limitamos a
enunciar aquí un teorema en un caso particularmente sencillo (cuyo análogo en
la teoría de espacios analíticos es bien conocido, cf. Kodaira--Spencer), que se
establece sin dificultad mediante los resultados del número anterior:

\result{TEOREMA}{10} Sea $X_0$ un esquema propio y simple sobre el cuerpo $k$
tal que
\[
  H^0(X_0,\Theta_{X_0/k})=H^2(X_0,\Theta_{X_0/k})=0.
\]
Entonces existe un esquema formal de módulos para $X_0$, correspondiente a un
anillo local regular $\mathcal O$ (es decir, un álgebra de series formales sobre
$k$).

Como ya hemos señalado, no es cierto en general que el esquema formal
$\mathfrak X$ sobre $\mathcal O$ sea algebraizable; pero se sabe que sí lo es
cuando $X_0$ es proyectivo y $H^2(X_0,\Osh_{X_0})=0$ (teorema 4), por ejemplo
cuando $X_0$ tiene dimensión 1. Esto da cierto fundamento a la esperanza de
construir un esquema de módulos sobre los enteros para las curvas de género
dado\ldots{}

Obsérvese también que métodos como los expuestos en este número pueden aplicarse
a la construcción y el estudio de las variedades de Picard, así como a muchas
otras construcciones. Volveremos sobre ello próximamente.

\fsection{8}{Aplicación al grupo fundamental}

Las técnicas expuestas permiten abordar el estudio sistemático del grupo
fundamental siguiendo el modelo de la teoría topológica. Los dos primeros
teoremas enunciados en este número son generalizaciones de resultados de un
trabajo reciente de Lang--Serre.

Sea $X$ un esquema. Un $X$-esquema $X'$ se llama revestimiento no ramificado de
$X$ si
\begin{enumerate}
  \item[\textit{i.}] $X'$ es finito sobre $X$, es decir, está definido por un
    haz coherente de álgebras $\mathcal A=\mathcal A(X')$ sobre $X$;
  \item[\textit{ii.}] $\mathcal A$ es un haz localmente libre sobre $X$;
  \item[\textit{iii.}] Para todo $x\in X$, el cociente
    \[
      \mathcal A_x/\mathfrak m_x\mathcal A_x
      =\mathcal A_x\otimes_{\Osh_x}\kappa(x)
    \]
    es un álgebra separable sobre $\kappa(x)$.
\end{enumerate}

\src{210}
Esta noción de revestimiento no ramificado (debida a Serre y al conferenciante)
posee todas las propiedades elementales que cabe razonablemente esperar, cuya
lista no daremos. Limitémonos a decir que da lugar a una teoría de Galois
modelada sobre la teoría de Galois clásica (y que la contiene; las
demostraciones son más bien más sencillas que las comúnmente aceptadas para
esta última) y sobre la teoría de Galois de los revestimientos topológicos. Más
precisamente, llamemos \emph{punto geométrico} de un esquema $X$ a un morfismo
$a$ del espectro $\xi$ de un cuerpo algebraicamente cerrado $\Omega$ en $X$, es
decir, al dato de una extensión algebraicamente cerrada del cuerpo residual
$\kappa(x)$ de un punto $x=|a|$ de $X$ (llamado \emph{localidad} del punto
geométrico $a$). Si $X'$ es entonces un revestimiento no ramificado de $X$, se
le puede asociar el conjunto $E_a(X')$ de los «puntos geométricos de $X'$ sobre
$a$», es decir, el conjunto de los pares formados por un $x'\in X'$ sobre $x$
y un $\kappa(x)$-homomorfismo $\kappa(x')\longrightarrow\Omega$. Se obtiene así
(para $(X,a)$ fijos) un funtor $F(X,a)$ de la categoría $R(X)$ de los
revestimientos no ramificados $X'$ de $X$ en la categoría de los conjuntos
finitos. Si $X$ es conexo, el par formado por $R(X)$ y $F(X,a)$ posee las
propiedades formales necesarias para ser isomorfo al par análogo definido por
un grupo topológico compacto totalmente discontinuo $\pi$ adecuado (es decir,
un límite proyectivo de grupos finitos): se toma la categoría $C(\pi)$ de los
conjuntos finitos $E$ sobre los cuales $\pi$ actúa continuamente, y el funtor
identidad $F(\pi)(E)=E$ de esta categoría en la categoría de los conjuntos
finitos. Además, el grupo $\pi$ está determinado salvo isomorfismo canónico por
la condición de que $(C(\pi),F(\pi))$ sea isomorfo a un par dado. En este caso,
$\pi$ se llama \emph{grupo fundamental} del esquema conexo $X$ en el punto
geométrico $a$, y se denota por $\pi_1(X,a)$. Si $X$ no es conexo, se lo
reemplaza por la componente conexa de $x=|a|$. Si, por el contrario, $X$ es
conexo, los grupos $\pi_1(X,a)$ correspondientes a dos puntos geométricos
$a',a''$ de $X$ son isomorfos (el isomorfismo está determinado salvo
automorfismos interiores), de modo que, como de costumbre, puede escogerse $a$
de la manera más conveniente, por ejemplo en el punto genérico de $X$, supuesto
irreducible. Naturalmente, $\pi_1(X,a)$ es un funtor covariante del esquema
punteado $(X,a)$. Todo enunciado relativo a la clasificación de los
revestimientos no ramificados\footnote{El texto impreso dice «revestimientos
inseparables»; el párrafo trata de la categoría $R(X)$ de los revestimientos no
ramificados y de su equivalencia con los conjuntos finitos dotados de una acción
continua de $\pi_1(X,a)$.} puede traducirse entonces al lenguaje de la teoría de
grupos siguiendo el diccionario bien conocido (salvo que hay que tener en cuenta
que aquí se trata de grupos topológicos).

\src{211}
Nuestro objetivo es desarrollar el análogo de la sucesión exacta de homotopía
de los espacios fibrados para un morfismo propio $f:X\longrightarrow Y$.
Evidentemente, al no saber qué son los grupos de homotopía superiores, solo
obtendremos resultados necesariamente incompletos. Para poder aplicar los
teoremas fundamentales del n\textsuperscript{o} 3, debemos explicitar primero
algunos lemas elementales relativos a los esquemas sobre cuerpos o anillos de
Artin (¡conforme al procedimiento general!).

\result{LEMA}{1} Sea $(X',a')$ un revestimiento no ramificado punteado asociado
a una representación punteada de $\pi_1(X,a)$ en un conjunto finito $E$
(provisto del punto marcado $e$). Entonces el morfismo canónico
$\pi_1(X',a')\longrightarrow\pi_1(X,a)$ identifica el primer grupo con el
estabilizador de $e$ en $\pi_1(X,a)$ (y es, por tanto, inyectivo).

\result{LEMA}{2} Sea $X$ un esquema algebraico sobre el cuerpo $k$, y sea $k'$
una extensión radicial de $k$. Entonces todo revestimiento no ramificado de
$X\otimes_k k'$ proviene por imagen inversa (es decir, por extensión de
escalares) de un revestimiento no ramificado de $X$, determinado salvo
isomorfismo canónico.

De estos dos lemas resulta, en particular, que para toda extensión algebraica
$K$ de $k$, y para todo punto geométrico $a'$ de $X'=X\otimes_k K$ que se
proyecta sobre el punto geométrico $a$ de $X$, el homomorfismo funtorial
$\pi_1(X',a')\longrightarrow\pi_1(X,a)$ es inyectivo.

\result{LEMA}{3} Sea $X$ un esquema completo sobre un anillo local artiniano
$A$, tal que $H^0(X,\Osh_X)=A$; sea $X'$ un revestimiento no ramificado de $X$,
y sea $A'=H^0(X',\Osh_{X'})$, que es, por tanto, un anillo finito sobre $A$
(que \textit{a priori} puede estar ramificado sobre $A$). Sean $X_0,X'_0$ los
subesquemas reducidos asociados a $X,X'$ (obtenidos dividiendo por los haces de
elementos nilpotentes de $\Osh_X,\Osh_{X'}$), y sea $k$ un subcuerpo de
$A/\mathfrak r(A)$ sobre el cual $A/\mathfrak r(A)$ sea finito (así, $X_0$ es
un esquema algebraico completo sobre $k$, y $X'_0$ es un revestimiento no
ramificado suyo). Sea, por último, $\Omega$ una extensión algebraicamente
cerrada de $k$, y consideremos el revestimiento no ramificado
$X'_0\otimes_k\Omega$ de $X_0\otimes_k\Omega$. Entonces las dos condiciones
siguientes son equivalentes:
\begin{enumerate}
  \item[\textit{i.}] $X'_0\otimes_k\Omega$ está completamente descompuesto
    sobre $X_0\otimes_k\Omega$.
  \item[\textit{ii.}] El morfismo natural
    $X'\longrightarrow X\otimes_A A'$ es un isomorfismo.
\end{enumerate}
Bajo estas condiciones, $A'$ es una extensión no ramificada de $A$. Por último,
si $X'$ es conexo, la condición \textit{(i)} equivale a la condición siguiente,
en apariencia más débil:
\begin{enumerate}
  \item[\textit{i bis.}] $X'_0\otimes_k\Omega$ admite una sección regular
    sobre $X_0\otimes_k\Omega$.
\end{enumerate}
Cuando se satisface la condición \textit{(ii)}, diremos que el revestimiento no
ramificado $X'$ de $X$ es \emph{geométricamente trivial}.

\src{212}
\result{LEMA}{4} Sea $f:X\longrightarrow Y$ un morfismo propio tal que
$f_*(\Osh_X)=\Osh_Y$. Sean $a$ un punto geométrico de $X$ y $b$ su proyección
sobre $Y$. Entonces $\pi_1(X,a)\longrightarrow\pi_1(Y,b)$ es sobreyectivo.

En efecto, hay que demostrar lo siguiente: si un revestimiento no ramificado
$Y'$ de $Y$ (correspondiente a un haz de álgebras localmente libre
$\mathcal A$) es tal que $X\otimes_Y Y'$ es disconexo, entonces $Y'$ también lo
es. En efecto, $f^*(\mathcal A)$ será suma directa de dos haces de anillos no
nulos y, por tanto, también lo será su imagen directa; ahora bien, esta última
no es otra que $\mathcal A\otimes f_*(\Osh_X)=\mathcal A$.

\result{LEMA}{5} Sea $X$ un esquema completo sobre un cuerpo $k$; se supone que
$H^0(X,\Osh_X)$ es un anillo local $A$ y que $A/\mathfrak r(A)$ es radicial
sobre $k$. Sea $\Omega$ una clausura algebraica de $k$, y sea
$\overline X=X\otimes_k\Omega$ (es conexo). Elijamos un punto geométrico
$\overline a$ de $\overline X$ que se proyecte sobre el punto geométrico $a$ de
$X$; entonces se tiene una sucesión exacta
\[
  e\longrightarrow\pi_1(\overline X,\overline a)
  \longrightarrow\pi_1(X,a)
  \longrightarrow\pi_1(k,b)\longrightarrow e
\]
(donde $\pi_1(k,b)$ es el grupo de Galois de $\Omega$ sobre $k$).

La inyectividad del primer homomorfismo se ha visto con los lemas 1 y 2, la
exactitud en el medio resulta del lema 3 y, por último, la sobreyectividad del
último homomorfismo (que es la única que utiliza el hecho de que
$A/\mathfrak r(A)$ sea radicial) resulta del lema 4.

\result{PROPOSICIÓN}{4} Sea $f:X\longrightarrow Y$ un morfismo propio y plano
tal que, para todo $y\in Y$, $H^0(f^{-1}(y),\Osh_{f^{-1}(y)})$ sea un álgebra
separable sobre el cuerpo residual $\kappa(y)$ (lo cual ocurre, por ejemplo, si
$f^{-1}(y)$ es un esquema separable sobre $\kappa(y)$, es decir, reducido y tal
que los cuerpos correspondientes a sus componentes irreducibles sean
extensiones separables de $\kappa(y)$). Entonces el revestimiento $Y'$ de $Y$
asociado a $f_*(\Osh_X)$ es no ramificado.

La demostración es fácil gracias al teorema 2.

\src{213}
Esta proposición, junto con el lema 1, reduce prácticamente el estudio
homotópico de los morfismos propios y planos (con fibras separables) al caso en
que $f_*(\Osh_X)=\Osh_Y$ (pues, utilizando la factorización de Stein, se
reemplaza $Y$ por $Y'$).

\result{OBSERVACIÓN}{} Un morfismo plano de tipo finito cuyas fibras son
esquemas separables (resp. simples) se llama \emph{separable} (resp.
\emph{simple}). Se demuestra que, si $f$ es plano y $f^{-1}(y)$ es separable
(resp. simple), existe un entorno de $f^{-1}(y)$ en el que $f$ es separable
(resp. simple). El mismo resultado es válido para «absolutamente normal»
(«teorema de Bertini»).

Sea $f:X\longrightarrow Y$ un morfismo propio tal que
\begin{equation}
  \tag{i}
  f_*(\Osh_X)=\Osh_Y,
\end{equation}
y sea $X'$ un esquema finito sobre $X$. Sea $Y'$ el revestimiento de $Y$
correspondiente a la factorización de Stein de $X'\longrightarrow Y$
(cf. teorema 5). Sea $y\in Y$; entonces el conjunto de las componentes conexas
de la fibra $F'$ de $X'$ sobre $y$ se identifica con el conjunto de los puntos
$y'$ de $Y'$ sobre $y$ (teorema 5). Consideremos el morfismo evidente
\begin{equation}
  \tag{*}
  X'\longrightarrow X\times_Y Y'
\end{equation}
deducido de los morfismos naturales $X'\longrightarrow X$ y
$X'\longrightarrow Y'$\footnote{El texto impreso omite la prima en el segundo
$X'$; el morfismo hacia $Y'$ es el de la factorización de Stein de
$X'\to Y$.}; será un isomorfismo siempre que $X'$ sea de la forma
$X\times_Y Y''$, donde $Y''$ es un revestimiento no ramificado de $Y$, y
entonces $Y'$ no será sino $Y''$ y (*) será la identidad. Queremos dar
precisamente condiciones bajo las cuales $X'$ sea de la forma que acabamos de
indicar, es decir, que $Y'$ sea no ramificado y (*) un isomorfismo. Para ello,
introduzcamos la fibra $F$ de $X$ en $y$; es un esquema propio sobre
$\kappa(y)$ del cual $F'$ es un revestimiento (no ramificado si $X'$ lo es).
Sea $F'_1$ una componente conexa de $F'$ correspondiente a un punto $y'_1$ de
$Y'$ sobre $y$. Supongamos además
\begin{enumerate}
  \item[\textit{(ii)}] $X'$ es no ramificado sobre $X$ en los puntos de $F'_1$
    (por tanto, $F'_1$ es un revestimiento no ramificado de $F$), y
  \item[\textit{(iii)}] $F'_1$ es un revestimiento geométricamente trivial de
    $F$ (cf. lema 3).
\end{enumerate}

\result{TEOREMA}{11} En estas condiciones, existe un entorno abierto $U'$ de
$y'_1$ en $Y'$ tal que (*) es un isomorfismo sobre $U'$. Además, $Y'$ es no
ramificado en $y'_1$ sobre $Y$ (pero puede ser ramificado en los demás puntos
$y'$ de $Y'$ sobre $y$).

Naturalmente, las condiciones \textit{(ii)} y \textit{(iii)} son también
necesarias para que sea válida la conclusión del teorema. La demostración del
teorema es fácil, con ayuda del lema 3 y del teorema 2.

\src{214}
\result{COROLARIO}{1} Supongamos siempre satisfecha \textit{(i)}. Para que un
revestimiento no ramificado sobre $X$ sea isomorfo a la imagen inversa de un
revestimiento no ramificado $Y'$ de $Y$, es necesario y suficiente que $X'$
induzca sobre cada fibra $f^{-1}(y)$ un revestimiento geométricamente trivial.

Como, según el teorema 11, el conjunto de los puntos de $Y$ para los cuales se
cumple esta condición es abierto, basta verificarla en los puntos $y$ que son
cerrados\ldots{} Señalemos el enunciado siguiente, equivalente al corolario 1:
el núcleo del homomorfismo (sobreyectivo según el lema 4)
$\pi_1(X)\longrightarrow\pi_1(Y)$ es el subgrupo invariante cerrado engendrado
por las imágenes en $\pi_1(X)$ de los $\pi_1(\overline{f^{-1}(y)})$, donde
\[
  \overline{f^{-1}(y)}
  =f^{-1}(y)\otimes_{\kappa(y)}\overline{\kappa(y)},
\]
$\overline{\kappa(y)}$ designa una clausura algebraica de $\kappa(y)$. Se
observará que, al no poder elegir el mismo punto base para todas las fibras,
los homomorfismos $\pi_1(\overline{f^{-1}(y)})\longrightarrow\pi_1(X)$ solo
quedan determinados en cualquier caso (una vez elegido un punto base para $X$
y, por tanto, para $Y$) módulo composición por un automorfismo interior en
$\pi_1(X)$.

\result{COROLARIO}{2} En las condiciones generales del teorema 11, supongamos
además que $Y,X,X'$ son íntegros, y sean $K,L,L'$ sus cuerpos. Entonces existe
una subextensión separable $K'$ de $K$ en $L'$, linealmente disjunta de $L$, tal
que $L'=LK'$ (de donde $L'=L\otimes_K K'$).

(Se aplica la última parte del lema 3 a la fibra geométrica de $X$). El caso de
aplicación más interesante del teorema 11 se obtiene cuando $f$ es un morfismo
separable. Entonces $X'$ es también separable sobre $Y$; por tanto, en virtud
de la proposición 4, $Y'$ es no ramificado sobre $Y$, y el segundo miembro
$X\times_Y Y'$ de (*) es no ramificado sobre $X$. Se concluye fácilmente:

\result{COROLARIO}{3} Supongamos, además de \textit{(i)}, que $f$ es separable.
Sea $X'$ un revestimiento no ramificado y conexo de $X$; para que
$X'$\footnote{El texto impreso dice aquí «$X$»; el sujeto del criterio es el
revestimiento $X'$ introducido en la frase anterior.} sea imagen inversa de un
revestimiento no ramificado $Y'$ de $Y$, es necesario y suficiente que el
revestimiento $\overline F'$ inducido sobre una fibra geométrica
$\overline F=f^{-1}(y)$ admita una sección regular.

Se observará que no era necesario suponer que $\overline F'$ fuera
geométricamente trivial sobre $\overline F$ (lo cual será cierto \textit{a
posteriori}, aunque \textit{a priori} esta condición sea mucho más fuerte). El
corolario 3 es, de hecho, equivalente al enunciado siguiente:

\src{215}
\result{COROLARIO}{4} Sea $f:X\longrightarrow Y$ un morfismo propio y separable
tal que $f_*(\Osh_X)=\Osh_Y$, sea $\overline F$ la fibra geométrica de un punto
$y\in Y$, y elijamos un punto geométrico en $\overline F$, del cual se obtienen,
mediante los morfismos $\overline F\longrightarrow X\longrightarrow Y$, puntos
geométricos en $X,Y$, que se tomarán como puntos base para los grupos
fundamentales de $\overline F,X,Y$. En estas condiciones se tiene la sucesión
exacta
\[
  \boxed{
    \pi_1(\overline F)\longrightarrow\pi_1(X)
    \longrightarrow\pi_1(Y)\longrightarrow e
  }.
\]

De ello se deducen fácilmente los dos enunciados siguientes de Serre--Lang,
liberados aquí de toda hipótesis de normalidad:

\result{COROLARIO}{5} Sean $X,Y$ esquemas conexos sobre un cuerpo $k$;
supongamos que el esquema reducido $X_{\mathrm{réd}}$ es separable sobre $k$
(lo cual es automáticamente cierto si $k$ es perfecto), y que $X$ o $Y$ es
propio sobre $k$.\footnote{Condición añadida en la fe de erratas de 1962.}
Elijamos un punto geométrico $a$ (resp. $b$) en $X$ (resp. $Y$), de donde se
obtienen un punto geométrico $c=(a,b)$ en $X\times_kY$ y un morfismo natural
\[
  \pi_1(X\times_kY,c)\longrightarrow
  \pi_1(X,a)\times\pi_1(Y,b)
\]
(deducido de los morfismos funtoriales de $\pi_1(X\times Y,c)$ en $\pi_1(X,a)$
y $\pi_1(Y,b)$). Este morfismo es inyectivo, e incluso biyectivo si $k$ es
algebraicamente cerrado.

(La sobreyectividad en este último caso es casi trivial). De ello se deduce,
con Serre--Lang:

\result{COROLARIO}{6} Sea $X$ un esquema algebraico conexo y propio sobre un cuerpo
algebraicamente cerrado $k$, y sea $K$ una extensión algebraicamente cerrada
de $k$; entonces los grupos fundamentales de $X$ y $X\otimes_kK$ son los
mismos, es decir, todo revestimiento no ramificado de este último esquema
proviene por extensión de escalares de un revestimiento no ramificado (único
salvo isomorfismo) de $X$.

\result{OBSERVACIONES}{}

1) Utilizando la proposición 4, se ve que la hipótesis
$f_*(\Osh_X)=\Osh_Y$ del corolario 4 no es esencial. En el caso general, en vez
de colocar el grupo unidad $e$ después de $\pi_1(Y)$, hay que continuar con
\[
  \pi_0(\overline F)\longrightarrow\pi_0(X)
  \longrightarrow\pi_0(Y)\longrightarrow e.
\]
como en Topología algebraica.

\src{216}
2) En general, por el momento no se puede decir nada sobre el núcleo de
$\pi_1(\overline F)\longrightarrow\pi_1(X)$, en el que debería intervenir un
$\pi_2(Y)$. Parece, sin embargo, que debería poder demostrarse que
$\pi_1(\overline F)\longrightarrow\pi_1(X)$ es inyectivo si $Y$ es el espectro
de un anillo local $A$, apoyándose en el teorema 12 de más abajo (que afirma que
así es si $A$ es completo).

El teorema 11 solo utilizaba los teoremas 1 y 2. Vamos a aplicar ahora el
teorema 3, apoyándonos en el lema elemental siguiente:

\result{LEMA}{6} Sea $X$ un esquema, y $X_0$ el esquema reducido correspondiente
(es decir, en el que se han anulado los elementos nilpotentes). Entonces todo
revestimiento no ramificado $X'_0$ de $X_0$ está inducido por un revestimiento
no ramificado $X'$ de $X$, determinado salvo isomorfismo canónico.

Este lema, de naturaleza puramente local, desempeña aquí un papel análogo al del
teorema 8 en teoría de módulos. Combinándolo con el teorema de existencia
(teorema 3), se deduce:

\result{TEOREMA}{12} Sea $A$ un anillo local noetheriano y completo de cuerpo
residual $k$. Sea $X$ un esquema propio sobre $A$. Entonces todo revestimiento
no ramificado $X'_0$ de $X_0=X\otimes_A k$ está inducido por un revestimiento no
ramificado $X'$ de $X$, único salvo isomorfismo.

En otros términos:

\result{COROLARIO}{1} Elijamos un punto geométrico en $X_0$ como punto base para
los grupos fundamentales de $X_0$ y $X$. Entonces el homomorfismo canónico
$\pi_1(X_0)\longrightarrow\pi_1(X)$ es un isomorfismo.

Apliquemos ahora a $X_0$ el lema 5 (suponiendo, para simplificar, que
$H^0(X_0,\Osh_{X_0})=k$), y observemos que, por ser $A$ completo, sus
extensiones no ramificadas corresponden a las extensiones no ramificadas de su
cuerpo residual, es decir, se tiene $\pi_1(Y)=\pi_1(k)$ (donde
$Y=\operatorname{Spec}(A)$). Se obtiene la sucesión exacta:
\[
  e\longrightarrow\pi_1(\overline X_0)
  \longrightarrow\pi_1(X)\longrightarrow\pi_1(Y)\longrightarrow e.
\]

\result{COROLARIO}{2} Sea $f:X\longrightarrow Y$ un morfismo propio y plano,
sean $y_1$ un punto de $Y$ e $y_0$ una especialización de $y_1$, y consideremos
las fibras «geométricas» correspondientes $\overline X_1$ y $\overline X_0$;
se supone que $\overline X_0$ es separable y conexa (lo que implica que
$\overline X_1$ satisface las mismas condiciones). Entonces se puede hallar un
homomorfismo de grupos
$\pi_1(\overline X_1)\longrightarrow\pi_1(\overline X_0)$, definido salvo un
automorfismo interior, y este homomorfismo es sobreyectivo.

\src{217}
Podría esperarse que este homomorfismo fuera siempre biyectivo; por desgracia,
en general no es así si $\kappa(y_0)$ tiene característica $>0$. Obtendremos,
sin embargo, más abajo una cota para el núcleo de este homomorfismo (al menos si
$\overline X_0$ es simple), de la que se sigue que, si $\kappa(y_0)$ tiene
característica $0$, el homomorfismo anterior es biyectivo (resultado que también
podría demostrarse por vía trascendente). En todo caso, obtenemos ya una cota
para el $\pi_1$ de una fibra especial por medio del de una fibra genérica.
Utilizando, por ejemplo, el hecho de que una curva algebraica en característica
$p$ se puede levantar a una curva en característica $0$ (teorema 9, corolario
4), se obtiene por vía trascendente:

\result{COROLARIO}{3} Sea $X_0$ el esquema de una curva completa simple sobre
un cuerpo algebraicamente cerrado de característica cualquiera, y sea $g$ el
género de $X_0$; entonces $\pi_1(X_0)$ admite $2g$ generadores topológicos,
ligados por la relación bien conocida.

De ello se deduce, mediante una técnica bien conocida de secciones
hiperplanas:

\result{COROLARIO}{4} Sea $X$ un esquema proyectivo simple sobre un cuerpo
algebraicamente cerrado de característica cualquiera; entonces $\pi_1(X)$
admite un número finito de generadores topológicos.

Busquemos el núcleo del homomorfismo
$\pi_1(\overline X_1)\longrightarrow\pi_1(\overline X_0)$. Para ello se puede
suponer que $Y$ es el espectro de un anillo de valoración discreta completo
$V=\Osh_y$ ($y=y_0$). La cuestión equivale a la siguiente: dado un revestimiento
no ramificado $\overline X'_1$ de $\overline X_1$ (que, si se quiere, puede
suponerse galoisiano), determinar en qué condiciones proviene de un
revestimiento no ramificado de $X_0$. \textit{A priori}, el revestimiento dado
proviene, por extensión de escalares, de un revestimiento no ramificado $X'_1$
de $X_1\otimes_K K'$, donde $K'$ es una subextensión finita de la clausura
algebraica $\overline K$ del cuerpo de fracciones $K$ de $V$; si
$\overline X'_1$ fuera galoisiano de grupo $G$, se puede elegir $X'_1$
galoisiano de grupo $G$. Se tiene entonces: para que
$X'_1\otimes_{K'}\overline K=\overline X'_1$ provenga de un revestimiento no
ramificado de $X_0$, es necesario y suficiente que se pueda hallar una extensión
finita $K''$ de $K'$ en $\overline K$ tal que
$X''_1=X'_1\otimes_{K'}K''$ sea de la forma $X''\otimes_{V''}K''$, donde $V''$
es la clausura normal de $V$ en $K''$ y $X''$ es un revestimiento no ramificado
de $X\otimes_VV''$. Supongamos, por ejemplo, que $X_0$ es absolutamente normal;
de ello resulta que $X\otimes_VV''$ es normal (pues es plano sobre $V''$ y tiene
fibra especial normal), y su cuerpo de funciones es idéntico al $K''(X_1)$ de
\[
  X_1\otimes_KK''
  =(X\otimes_VK)\otimes_KK''
  =X\otimes_VK''
  =(X\otimes_VV'')\otimes_{V''}K''.
\]
Sea $L''=K''(X'_1)$ el cuerpo de funciones de $X'_1\otimes_{K'}K''$; es una
extensión finita separable de $K''(X_1)$, y la condición anterior significa
también que $L''$ es una extensión no ramificada del cuerpo de funciones de
$X\otimes_VV''$ (es decir, la normalización de $X\otimes_VV''$ en $L''$ es no
ramificada sobre $X\otimes_VV''$). Basta, por lo demás, verificar que $L''$ es
no ramificado en los puntos de la fibra especial de $X\otimes_VV''$ (pues es no
ramificado sobre la fibra genérica $X_1\otimes_KK''$). Si ahora $X_0$ es simple,
del «teorema de pureza» de Nagata--Zariski resulta incluso que basta verificar
que $L''$ es no ramificado sobre el anillo local $O''$ del punto genérico de la
fibra especial de $X\otimes_VV''$, que es un anillo de valoración discreta,
normalización en $K''(X_1)$ del anillo local $O\subset K(X_1)$ del punto
genérico de la fibra especial de $X$. Por tanto, se ha obtenido:

\src{218}
\result{COROLARIO}{5} En las condiciones y con las notaciones anteriores, para
que el revestimiento no ramificado
$X'_1\otimes_{K'}\overline K$ de
$\overline X_1=X_1\otimes_K\overline K$ provenga de un revestimiento no
ramificado de $\overline X_0$, es necesario y suficiente que exista una
subextensión finita $K''$ de $\overline K/K'$ tal que $K''(X'_1)$ sea no
ramificado sobre el anillo de valoración discreta $O''\subset K''(X_1)$.

Observemos ahora que $O''$ es la normalización en $K''(X_1)$ del anillo de
valoración discreta $O'\subset K'(X_1)$ (normalización de $O$ en $K'(X_1)$), y
que $O'$ contiene la normalización $V'$ de $V$ en $K'$; un uniformizante $u$ de
$V'$ es también un uniformizante de $O'$. Supongamos ahora que $X'_1$ es
galoisiano de grupo de Galois $G$, de orden $n$ primo con la característica $p$
de $\kappa(y_0)$ (que es también la característica del cuerpo residual de
$O'$). Por tanto, $K'(X'_1)$ está «moderadamente ramificado» sobre $O'$, de lo
que resulta fácilmente («lema de Abhyankar») que, si se le adjunta una raíz
$n$-ésima $v$ de un uniformizante de $O'$, se vuelve no ramificado sobre la
normalización de $O'$ en $K'(X_1)(v)$. Ahora bien, se puede tomar como $v$ una
raíz $n$-ésima de un uniformizante de $V'$, lo que prueba que se satisface la
condición del corolario 5. (Esta posibilidad de utilizar el lema de Abhyankar y
el teorema de pureza me la ha vendido Serre). Para expresar el resultado
obtenido, introduzcamos, para todo grupo compacto totalmente discontinuo $\pi$,
el grupo cociente $\widetilde\pi$ de $\pi$ por el subgrupo cerrado engendrado
por sus $p$-subgrupos de Sylow, es decir, el límite proyectivo de los grupos
cocientes discretos de $\pi$ que son de orden primo con $p$. Con esta notación,
se obtiene:

\result{TEOREMA}{13} Sea $f:X\longrightarrow Y$ un morfismo propio y plano,
sean $y_1$ un punto de $Y$ e $y_0$ una especialización de $y_1$, y supongamos
$\overline X_0$ conexo y simple. Entonces el homomorfismo
$\widetilde\pi_1(\overline X_1)\longrightarrow
\widetilde\pi_1(\overline X_0)$ deducido del homomorfismo sobreyectivo del
teorema 12, corolario 2, es un isomorfismo.

En otros términos:

\src{219}
\result{COROLARIO}{1} La clasificación de los revestimientos galoisianos no
ramificados, de grupo de Galois de orden primo con la característica $p$ de
$\kappa(y_0)$, es la misma para $\overline X_0$ y para $\overline X_1$.

En particular, si $\kappa(y_0)$ tiene característica cero, se obtiene por vía
algebraica que $\pi_1(\overline X_1)\longrightarrow\pi_1(\overline X_0)$ es
biyectivo.

Señalemos, por último, que las técnicas utilizadas proporcionan también el
resultado siguiente, más general que el teorema 13:

\result{TEOREMA}{14} Sea $f:X\longrightarrow Y$ un morfismo propio y simple,
sea $D$ un subesquema cerrado de $X$, simple sobre $Y$, de codimensión 1 en
todos sus puntos. Para una fibra $Z=f^{-1}(z)$ de $f$, sea $Z'=Z-Z\cap D$ y sea
$\pi_1^t(Z')$ el cociente del grupo fundamental $\pi_1(Z')$ que clasifica los
revestimientos no ramificados de $Z'$ que están «moderadamente ramificados»
sobre $\overline{Z\cap D}$. Sean $y_0,y_1$ como en el teorema 13. Entonces se
tiene un homomorfismo sobreyectivo (definido salvo un automorfismo interior)
\[
  \pi_1^t(\overline X'_1)\longrightarrow\pi_1^t(\overline X'_0),
\]
y el homomorfismo correspondiente
\[
  \widetilde\pi_1^t(\overline X'_1)
  \longrightarrow\widetilde\pi_1^t(\overline X'_0)
\]
es un isomorfismo.

De ello se deducen las variantes correspondientes de los corolarios del
teorema 13 y del corolario 4 del teorema 12. Del mismo modo, utilizando el
teorema 9, corolario 3, se obtiene por vía trascendente:

\result{COROLARIO}{} Sea $X_0$ el esquema de una curva completa simple sobre un
cuerpo algebraicamente cerrado de característica cualquiera, y sea
$S=(s_i)_{1\le i\le n}$ una parte finita con $n$ elementos de $X_0$. Entonces
$\pi_1^t(X_0-S)$ admite $2g+n$ generadores topológicos
$x_i,y_i$ ($1\le i\le g$) y $\sigma_j$ ($1\le j\le n$), ligados por la relación
\[
  \left(\prod_i x_i y_i x_i^{-1}y_i^{-1}\right)
  \sigma_1\cdots\sigma_n=1,
\]
donde los $\sigma_j$ son generadores de los grupos de inercia correspondientes
a los $s_j$. Para todo grupo finito $G$ de orden primo con la característica,
engendrado por elementos $x_i,y_i,\sigma_j$ que satisfacen la relación anterior,
existe un revestimiento galoisiano no ramificado de $X_0-S$, de grupo $G$, cuyos
grupos de inercia en los puntos $s_j$ están engendrados por los $\sigma_j$.

Cuando $X_0$ es de género 0 y $n=3$, se tiene una solución del «problema de los
tres puntos», al menos para los revestimientos galoisianos de orden primo con
la característica. (Aquí, por lo demás, el teorema 9 es inútil, y parece, por
otra parte, que el corolario anterior puede deducirse del caso particular
considerado en el problema de los tres puntos).

\src{220}
\result{OBSERVACIONES}{}

1) Un estudio más completo, en el que sin duda intervendrían revestimientos
galoisianos generalizados de $X,X_0,X_1$ (con grupo de Galois eventualmente
infinitesimal), debería permitir recuperar el núcleo del teorema 12, corolario
2. En cambio, el estudio de los revestimientos que admiten ramificación no
«moderada» parece tener que ser mucho más difícil.

2) El lema 6, junto con un resultado de Grauert relativo a la completación
formal de un esquema proyectivo no singular a lo largo de una sección
hiperplana (o con el teorema, no demostrado por el momento, mencionado en la
observación 2 después del teorema 11), permite también demostrar en geometría
algebraica «abstracta» el teorema clásico de Lefschetz sobre el grupo
fundamental.

\begin{center}
  \textbf{BIBLIOGRAFÍA}
\end{center}

\begin{description}
  \item[{[1]}] Dieudonné (J.) y Grothendieck (A.). --- \emph{Éléments de
    géométrie algébrique}, Publications mathématiques de l'Institut des Hautes
    Études Scientifiques (de próxima aparición).
  \item[{[2]}] Grothendieck (Alexander). --- \emph{The cohomology theory of
    abstract algebraic varieties}. International Congress of Mathematicians
    [1958. Edinburgh] (de próxima aparición).
  \item[{[3]}] Serre (Jean-Pierre). --- \emph{Faisceaux algébriques
    cohérents}, Annals of Math., Series 2, t. 61, 1955, p. 197--278.
  \item[{[4]}] Serre (Jean-Pierre). --- \emph{Géométrie algébrique et
    géométrie analytique}, Ann. Institut Fourier Grenoble, t. 6, 1955--1956,
    p. 1--42.
\end{description}

\def\FgaUnitKey{e182}\clearpage
\chapter*{Fe de erratas de la exposición n\textsuperscript{o} 182}
\addcontentsline{toc}{chapter}{Fe de erratas de la exposición n\textsuperscript{o} 182}
\setcounter{footnote}{1}
\src{298}

\noindent\textbf{2. Geometría formal y geometría algebraica.}

\begin{center}
[Seminario Bourbaki, t. 11, 1958/59, n\textsuperscript{o} 182, 28 p.]
\end{center}

\noindent\textbf{Página 8, teorema 6.} --- La hipótesis «$f$
cuasiproyectivo» puede sustituirse por la hipótesis más débil «$f$
separado», en virtud del resultado siguiente (cf. SGA VIII, 6.2): todo morfismo
$f:X\longrightarrow Y$ que es cuasifinito y separado es cuasiproyectivo.

\src{299}
\noindent\textbf{Página 12, y página 14, observación 1.} --- Señalemos que J.-P.
Serre ha construido una variedad proyectiva no singular, de dimensión 3, sobre
un cuerpo algebraicamente cerrado $k$, de característica $p>0$, que no proviene
por reducción de un esquema propio sobre un anillo local íntegro de cuerpo
residual $k$ y con cuerpo de fracciones de característica
cero.\footnote{Serre (Jean-Pierre). --- Exemples de variétés projectives en
caractéristique $p$ non relevables en caractéristique zéro, \emph{Proc. Nat.
Acad. Sc. U.S.A.}, t. 47, 1961, p. 108--109.} Mumford habría hallado un ejemplo
análogo, con una superficie proyectiva no singular.

\noindent\textbf{Página 15, observaciones 2 y 3.} --- Actualmente soy menos
optimista respecto de los resultados aquí conjeturados. Sin embargo, el
problema relativo a las variaciones de estructura del espacio proyectivo,
señalado al final de la observación 3, ha sido resuelto afirmativamente por
Hironaka, y el problema análogo para las variedades abelianas ha sido resuelto
por Koizumi.

\noindent\textbf{Página 16, línea 3.} --- Señalemos que Mumford ha construido
recientemente los esquemas de módulos para las curvas de género $g$ (cf. SHMI).
El teorema 10 prueba, por lo demás, que los «esquemas de Jacobi de nivel $n$»
de la teoría de Módulos son no singulares (e incluso simples sobre $\mathbf Z$).

\noindent\textbf{Página 23, línea 12, primera línea del corolario
5.} --- Añadir: «siendo $X$ o $Y$ propio sobre $k$».

La sustancia de los §§ 1 a 5 está contenida en la parte publicada de EGA III, y
la de los §§ 6 y 7 está contenida en SGA III. Para el estudio del grupo
fundamental, véanse SGA V, IX, X, XI, así como SGA 1962 (exposiciones X, XII y
XIII), para los teoremas de tipo Lefschetz y numerosas cuestiones abiertas.
Solo la teoría de los revestimientos moderadamente ramificados (cf. teorema 14)
no ha sido todavía objeto de una redacción formal. El corolario del teorema 14,
que da la determinación completa de los revestimientos galoisianos de orden
primo con la característica de una curva algebraica sobre un cuerpo
algebraicamente cerrado, ha sido utilizado de manera esencial en tres ocasiones:

1\textsuperscript{o} en la demostración por Igusa de la desigualdad de Picard
para las superficies proyectivas no singulares en característica cualquiera;

2\textsuperscript{o} en el estudio (desarrollado independientemente por Ogg y
Šafarevič) del grupo de los fibrados principales homogéneos bajo una variedad
abeliana definida sobre un cuerpo de funciones de una variable, en
característica cualquiera;

3\textsuperscript{o} en la demostración reciente, por Artin, de ciertos
teoremas clave sobre la «cohomología de Weil» de las variedades algebraicas.

\def\FgaUnitKey{190}\unit{Exposé 190 -- Técnica de descenso I}
\sid{190}
\src{299}
\seminar{12\textsuperscript{o} año, 1959/60, n\textsuperscript{o} 190\\Diciembre de 1959}
\paperhead{TÉCNICA DE DESCENSO Y TEOREMAS DE EXISTENCIA EN GEOMETRÍA ALGEBRAICA\\
I. GENERALIDADES. DESCENSO POR MORFISMOS FIELMENTE PLANOS}%
  {por Alexander GROTHENDIECK}

Desde un punto de vista técnico, el presente exposé y los que le seguirán
pueden considerarse variaciones sobre el célebre «teorema 90» de Hilbert.
Parece que A. Weil introdujo el método de descenso en geometría algebraica con
el nombre de «descenso del cuerpo base». Weil se restringe, por lo demás, al
caso de extensiones finitas separables de cuerpos. P. Cartier trató después el
caso de extensiones radiciales de altura 1. Por falta del lenguaje de los
esquemas y, más concretamente, por falta de elementos nilpotentes en los
anillos que se permitía considerar, Cartier no había podido formular la
identidad esencial entre estos dos casos.

En la actualidad, parece que la técnica general de descenso que se expondrá
(unida, cuando proceda, a los teoremas fundamentales de la «geometría formal»,
cf. [3]) está en la base de la mayoría de los teoremas de existencia en
geometría algebraica.\footnote{La fe de erratas de 1962 considera excesiva esta
formulación: es válida en gran medida para las técnicas no proyectivas de los
dos primeros exposés de la serie TDTE I a VI, pero no para las técnicas
proyectivas (TDTE IV, V y VI).} Conviene señalar, además, que esta técnica puede
sin duda trasladarse a la «geometría analítica», y cabe esperar que en un
futuro próximo los especialistas sepan demostrar los análogos «analíticos» de
los teoremas de existencia en geometría formal que se darán en el exposé II.
En cualquier caso, los trabajos recientes de Kodaira--Spencer, cuyos métodos
parecen inadecuados para definir y estudiar «variedades de módulos» en un
entorno de sus puntos singulares, muestran con bastante claridad la necesidad
de métodos más próximos a la teoría de esquemas (que deberán completar, por
supuesto, las técnicas trascendentes).

En el presente exposé I trataremos el caso de descenso más elemental (indicado
en el título). Sin embargo, las aplicaciones de los teoremas 1, 2 y 3 que
siguen son ya muy numerosas. Nos hemos limitado a dar algunas a modo de
ejemplo, sin intentar conferirles aquí la máxima generalidad posible.

Utilizaremos libremente el lenguaje de los esquemas, para el cual remitimos al
exposé ya citado, así como a [2]. Señalemos, no obstante, expresamente que los
preesquemas considerados en el presente exposé no son necesariamente
noetherianos, y los

\src{300}
morfismos no son necesariamente de tipo finito. Así, si $A$ es un anillo local
noetheriano, con completado $\widehat A$, habrá que considerar el anillo no
noetheriano $\widehat A\otimes_A\widehat A$, y los morfismos de esquemas afines
correspondientes a las inclusiones entre los anillos considerados.

\fsection{A}{Preliminares sobre las categorías}

\subsection*{1. Categorías fibradas, datos de descenso, morfismos de
$\mathcal F$-descenso}
\addcontentsline{toc}{subsection}{1. Categorías fibradas y datos de descenso}

\result{a. DEFINICIÓN}{1.1} Una \emph{categoría fibrada} $\mathcal F$ de base
$\mathcal C$ consta de una categoría $\mathcal C$; para cada
$X\in\mathcal C$, de una categoría $\mathcal F_X$; para cada
$\mathcal C$-morfismo $f:X\longrightarrow Y$, de un funtor
$f^*:\mathcal F_Y\longrightarrow\mathcal F_X$, denotado también por
\[
  f^*(\xi)=\xi\times_Y X
\]
para $\xi\in\mathcal F_Y$ (considerando $X$ como un «objeto de
$\mathcal C$ sobre $Y$», es decir, provisto del morfismo $f$); y, por último,
para dos morfismos componibles
$X\xrightarrow{f}Y\xrightarrow{g}Z$, de un isomorfismo de funtores
\[
  c_{f,g}:(gf)^*\longrightarrow f^*g^*,
\]
sujetos estos datos a las condiciones siguientes:

\begin{enumerate}
\item[(i)] $\operatorname{id}_X^*=\operatorname{id}$;
\item[(ii)] $c_{f,g}$ es el isomorfismo identidad si $f$ o $g$ es un
isomorfismo identidad;
\item[(iii)] dados tres morfismos componibles
$X\xrightarrow{f}Y\xrightarrow{g}Z\xrightarrow{h}T$, el siguiente diagrama de
isomorfismos, formado mediante isomorfismos de la forma $c_{u,v}$, es
conmutativo:
\[
\begin{array}{ccc}
  (h(gf))^*=((hg)f)^* & \longrightarrow & f^*(hg)^* \\
  \big\downarrow && \big\downarrow \\
  (gf)^*h^* & \longrightarrow & (f^*g^*)h^*=f^*(g^*h^*).
\end{array}
\]
\end{enumerate}

\result{EJEMPLO}{1} Sea $\mathcal C$ una categoría en la que existen los
productos fibrados. Se define entonces una categoría fibrada $\mathcal F$ de
base $\mathcal C$ tomando $\mathcal F_X=$ categoría de los objetos de
$\mathcal C$ sobre $X$; el funtor
$f^*:\mathcal F_Y\longrightarrow\mathcal F_X$ correspondiente a un morfismo
$f:X\longrightarrow Y$ se define mediante el producto fibrado
$Z\longmapsto Z\times_YX$.

\result{EJEMPLO}{2} Sea $\mathcal C$ la categoría de los preesquemas y, para
$X\in\mathcal C$, sea $\mathcal F_X$ la categoría de los haces
cuasi-coherentes de módulos sobre $X$. Si $f:X\longrightarrow Y$ es un morfismo
de preesquemas, $f^*:\mathcal F_Y\longrightarrow\mathcal F_X$ es el funtor
imagen inversa

\src{301}
de haces de módulos. Se obtiene así una categoría fibrada de base
$\mathcal C$.

\result{b. DEFINICIÓN}{1.2} Un diagrama
\[
  E\xrightarrow{u}E'
  \overset{v_1}{\underset{v_2}{\rightrightarrows}}E''
\]
de aplicaciones de conjuntos se llama \emph{exacto} si $u$ es una biyección de
$E$ sobre la parte de $E'$ formada por los $x'\in E'$ tales que
$v_1(x')=v_2(x')$.

\result{DEFINICIÓN}{1.3} Sea $\mathcal F$ una categoría fibrada de base
$\mathcal C$; consideremos un diagrama
\[
  \tag{+}
  S\xleftarrow{\alpha}S'
  \overset{\beta_1}{\underset{\beta_2}{\rightrightarrows}}S''
\]
de morfismos en $\mathcal C$ tales que
$\alpha\beta_1=\alpha\beta_2$; este diagrama se llama
$\mathcal F$-\emph{exacto} si, para cada par $\xi,\eta$ de elementos de
$\mathcal F_S$, el siguiente diagrama de aplicaciones de conjuntos
\[
  \tag{+'}
  \operatorname{Hom}(\xi,\eta)
  \xrightarrow{\alpha^*}
  \operatorname{Hom}(\alpha^*\xi,\alpha^*\eta)
  \overset{\beta_1^*}{\underset{\beta_2^*}{\rightrightarrows}}
  \operatorname{Hom}(\gamma^*\xi,\gamma^*\eta),
\]
donde $\gamma=\alpha\beta_1=\alpha\beta_2$, es exacto. En este último diagrama
hemos identificado $\beta_i^*\alpha^*$ con
$(\alpha\beta_i)^*=\gamma^*$ mediante $c_{\beta_i,\alpha}$, para simplificar.

\result{DEFINICIÓN}{1.4} Sea $\mathcal F$ una categoría fibrada de base
$\mathcal C$; consideremos dos morfismos $\beta_1,\beta_2$ de $S''$ en $S'$,
en $\mathcal C$. Sea
$\xi'\in\mathcal F_{S'}$. Se llama \emph{dato de pegado} sobre $\xi'$
(relativo al par $(\beta_1,\beta_2)$) a un isomorfismo de
$\beta_1^*(\xi')$ sobre $\beta_2^*(\xi')$. Si $\xi',\eta'\in\mathcal F_{S'}$
están provistos cada uno de un dato de pegado, un morfismo
$u:\xi'\longrightarrow\eta'$ en $\mathcal F_{S'}$ se dice
\emph{compatible con los datos de pegado} si el diagrama siguiente es
conmutativo:
\[
\begin{array}{ccc}
  \beta_1^*(\xi') & \longrightarrow & \beta_2^*(\xi') \\
  \big\downarrow && \big\downarrow \\
  \beta_1^*(\eta') & \longrightarrow & \beta_2^*(\eta').
\end{array}
\]
De este modo, los objetos de $\mathcal F_{S'}$ provistos de un dato de pegado
(relativo a $\beta_1,\beta_2$) forman una categoría. Si
$\alpha:S'\longrightarrow S$ es un morfismo tal que
$\alpha\beta_1=\alpha\beta_2$, entonces, para todo $\xi\in\mathcal F_S$, el
objeto $\xi'=\alpha^*(\xi)$

\src{302}
de $\mathcal F_{S'}$ está provisto de un dato de pegado canónico, puesto que
\[
  \beta_i^*\alpha^*(\xi)\simeq(\alpha\beta_i)^*(\xi)=\gamma^*(\xi),
\]
tomando de nuevo $\gamma=\alpha\beta_1=\alpha\beta_2$; además, si
$u:\xi\longrightarrow\eta$ es un morfismo en $\mathcal F_S$, entonces
\[
  \alpha^*(u):\alpha^*(\xi)\longrightarrow\alpha^*(\eta)
\]
es un morfismo en $\mathcal F_{S'}$ compatible con los datos de pegado
canónicos. Se obtiene así un funtor canónico de la categoría $\mathcal F_S$ a
la categoría de los objetos de $\mathcal F_{S'}$ provistos de un dato de pegado
relativo al par $(\beta_1,\beta_2)$. Dicho esto, la definición 1.3 también
puede expresarse diciendo que el diagrama (+) de dicha definición es
$\mathcal F$-exacto si el funtor anterior es «plenamente fiel», es decir, si
define una equivalencia de la categoría $\mathcal F_S$ con una subcategoría de
la categoría de los objetos de $\mathcal F_{S'}$ provistos de un dato de pegado
relativo a $(\beta_1,\beta_2)$.

\result{DEFINICIÓN}{1.5} Se dice que un dato de pegado sobre
$\xi'\in\mathcal F_{S'}$ es \emph{efectivo} (respecto de $\alpha$) si
$\xi'$, provisto de este dato, es isomorfo a un $\alpha^*(\xi)$, con
$\xi\in\mathcal F_S$.

Cuando el diagrama (+) es $\mathcal F$-exacto, el objeto $\xi$ de la definición
anterior queda, por tanto, determinado salvo isomorfismo único, y la categoría
$\mathcal F_S$ es equivalente a la categoría de los objetos de
$\mathcal F_{S'}$ provistos de un dato de pegado efectivo.

\medskip\noindent
c. El caso más importante es aquel en que
\[
  S''=S'\times_S S',
\]
siendo las $\beta_i$ las dos proyecciones $p_1$ y $p_2$ de
$S'\times_S S'$ sobre sus dos factores (suponemos ahora que en
$\mathcal C$ existen los productos fibrados). Se tienen entonces dos
condiciones necesarias inmediatas para que sea efectivo un dato de pegado
$\varphi:p_1^*(\xi')\longrightarrow p_2^*(\xi')$ sobre un
$\xi'\in\mathcal F_{S'}$. Se tiene
\[
  \tag{i}
  \Delta^*(\varphi)=\operatorname{id}_{\xi'},
\]
donde $\Delta:S'\longrightarrow S'\times_S S'$ designa el morfismo diagonal,
y donde se ha identificado $\Delta^*p_i^*(\xi')$ con
$(p_i\Delta)^*(\xi')=\xi'$. Se tiene
\[
  \tag{ii}
  p_{31}^*(\varphi)=p_{32}^*(\varphi)\,p_{21}^*(\varphi),
\]

\src{303}
donde $p_{ij}$ designa la proyección canónica de
$S'\times_S S'\times_S S'$ sobre el producto parcial de sus factores
$i$-ésimo y $j$-ésimo.

\result{DEFINICIÓN}{1.6} Se llama \emph{dato de descenso} sobre
$\xi'\in\mathcal F_{S'}$, relativo al morfismo
$\alpha:S'\longrightarrow S$, a un dato de pegado sobre $\xi'$ relativo al par
$(p_1,p_2)$ de proyecciones canónicas
$S'\times_S S'\longrightarrow S'$ que satisface las condiciones (i) y (ii)
anteriores.

\result{DEFINICIÓN}{1.7} Un morfismo $\alpha:S'\longrightarrow S$ se llama
\emph{morfismo de $\mathcal F$-descenso} si el diagrama de morfismos
\[
  S\xleftarrow{\alpha}S'
  \overset{p_1}{\underset{p_2}{\rightrightarrows}}S'\times_S S'
\]
es $\mathcal F$-exacto (definición 1.3); se dice que $\alpha$ es un
\emph{morfismo de $\mathcal F$-descenso estricto} si, además, todo dato de
descenso (definición 1.6) sobre un objeto de $\mathcal F_{S'}$ es efectivo.

Esta última condición también puede enunciarse, de modo más gráfico, así:
«Da lo mismo» dar un objeto de $\mathcal F_S$ o un objeto de
$\mathcal F_{S'}$ provisto de un dato de descenso.

Obsérvese que si un morfismo $\alpha:S'\longrightarrow S$ admite una «sección»
$s:S\longrightarrow S'$ (es decir, un morfismo $s$ tal que
$\alpha s=\operatorname{id}_S$), entonces es un morfismo de
$\mathcal F$-descenso estricto.\footnote{La fe de erratas de 1962 precisa que
no es necesario suponer previamente que $\alpha$ sea un morfismo de
$\mathcal F$-descenso.} En efecto, si $\xi'\in\mathcal F_{S'}$ está provisto de
un dato de descenso relativo a $\alpha$, «proviene» de
$\xi=s^*(\xi')$.

\medskip\noindent
d. Las nociones anteriores pueden presentarse de manera más intuitiva
introduciendo, para un objeto parámetro $T$ de $\mathcal C$ sobre $S$, el
conjunto
\[
  \operatorname{Hom}_S(T,S')=S'(T),
\]
cuyos elementos se denotarán por $t,t'$, etc. Dado un objeto
$\xi'\in\mathcal F_{S'}$, a cada $t\in S'(T)$ le corresponde entonces un objeto
$t^*(\xi')$ de $\mathcal F_T$, que también se denotará por $\xi'_t$. Un dato de
pegado sobre $\xi'$ (relativo a $p_1,p_2$) queda definido entonces dando, para
todo $T$ sobre $S$ y todo par de puntos $t,t'\in S'(T)$, un isomorfismo
\[
  \varphi_{t',t}:\xi'_t\longrightarrow\xi'_{t'}
\]
(que satisface condiciones evidentes de funtorialidad al variar $T$). Las
condiciones (i) y (ii) de 1.c se escriben entonces
\[
  \tag{i bis}
  \varphi_{t,t}=\operatorname{id}_{\xi'_t}
  \qquad\text{para todo $T$ y todo $t\in S'(T)$},
\]

\src{304}
\[
  \tag{ii bis}
  \varphi_{t,t''}=\varphi_{t,t'}\varphi_{t',t''}
  \qquad\text{para todo $T$ y todos $t,t',t''\in S'(T)$}.
\]
Se ve, además, que (ii bis), tomando $t=t'=t''$, implica
$\varphi_{t,t}=\varphi_{t,t}^2$; puesto que $\varphi_{t,t}$ es un isomorfismo
por hipótesis, de aquí se obtiene la relación (i bis), que es de hecho una
consecuencia de (ii bis) (por tanto, (i) es consecuencia de (ii)). Pero si ya
no se supone a priori que los $\varphi_{t,t'}$ sean isomorfismos (es decir, que
$\varphi:p_1^*(\xi')\longrightarrow p_2^*(\xi')$ sea un isomorfismo), entonces
(ii bis) no implica necesariamente (i bis); sin embargo, la conjunción de
(ii bis) y (i bis) implica que los $\varphi_{t,t'}$ son isomorfismos (pues se
tendrá
$\varphi_{t,t'}\varphi_{t',t}=\varphi_{t,t}=\operatorname{id}_{\xi'_t}$).

\subsection*{2. Diagramas exactos y epimorfismos estrictos, morfismos de
descenso. Ejemplos}
\addcontentsline{toc}{subsection}{2. Diagramas exactos y epimorfismos estrictos}

\result{a. DEFINICIÓN}{2.1} Sea $\mathcal C$ una categoría. Un diagrama
\[
  T\xrightarrow{\alpha}T'
  \overset{\beta_1}{\underset{\beta_2}{\rightrightarrows}}T''
\]
de morfismos en $\mathcal C$ se llama \emph{exacto} si, para todo
$Z\in\mathcal C$, el diagrama correspondiente de aplicaciones de conjuntos
\[
  \operatorname{Hom}(Z,T)\longrightarrow
  \operatorname{Hom}(Z,T')
  \rightrightarrows\operatorname{Hom}(Z,T'')
\]
es exacto (definición 1.2). Se dice entonces que $(T,\alpha)$ (o, por abuso de
lenguaje, que $T$) es un \emph{núcleo} del par de morfismos
$(\beta_1,\beta_2)$.

Evidentemente, este núcleo está determinado salvo isomorfismo único. Si
$\mathcal C$ es la categoría de los conjuntos, la definición anterior es
compatible con la definición 1.2. Se define de manera dual la exactitud de un
diagrama
\[
  S\xleftarrow{\alpha}S'
  \overset{\beta_1}{\underset{\beta_2}{\rightrightarrows}}S''
\]
de morfismos en $\mathcal C$; se dice entonces que $(S,\alpha)$ es un
\emph{conúcleo} del par de morfismos $(\beta_1,\beta_2)$.

\result{DEFINICIÓN}{2.2} Un morfismo $\alpha:S'\longrightarrow S$ se llama
\emph{epimorfismo estricto} si es un epimorfismo y si, para todo morfismo
$u:S'\longrightarrow Z$, la siguiente condición necesaria es también
suficiente para que $u$ se factorice como
$S'\longrightarrow S\longrightarrow Z$: para todo $S''\in\mathcal C$ y todo
par de morfismos $\beta_1,\beta_2:S''\longrightarrow S'$ tal que
$\alpha\beta_1=\alpha\beta_2$, se tiene también $u\beta_1=u\beta_2$.

Si existe el producto fibrado $S'\times_S S'$, equivale a decir que el diagrama
\[
  S\xleftarrow{\alpha}S'
  \overset{p_1}{\underset{p_2}{\rightrightarrows}}S'\times_S S'
\]

\src{305}
es exacto, es decir, que $S$ es un conúcleo del par $(p_1,p_2)$. En cualquier
caso, un morfismo conúcleo es un epimorfismo estricto. Señalemos también que un
epimorfismo estricto que sea un monomorfismo es un isomorfismo. Dejamos al
lector desarrollar la noción dual de \emph{monomorfismo estricto}.

Para precisar las relaciones entre la noción de morfismo de
$\mathcal F$-descenso y la de epimorfismo estricto, introducimos además las
definiciones siguientes:

\result{DEFINICIÓN}{2.3} Un morfismo $\alpha:S'\longrightarrow S$ se llama
\emph{epimorfismo universal} (resp. \emph{epimorfismo estricto universal}) si,
para todo $T$ sobre $S$, existe el producto fibrado $T'=S'\times_ST$ y la
proyección $T'\longrightarrow T$ es un epimorfismo (resp. un epimorfismo
estricto).

En las categorías muy buenas (como la categoría de los conjuntos, la
categoría de los haces de conjuntos sobre un espacio topológico, las categorías
abelianas, etc.), las cuatro nociones de epiyectividad así introducidas
coinciden; por el contrario, todas ellas son distintas en una categoría como
la de los preesquemas, o la de los preesquemas sobre un preesquema no vacío
dado $S$, incluso si nos limitamos a $S$-esquemas finitos sobre $S$.

\result{DEFINICIÓN}{2.4} Un morfismo $\alpha:S'\longrightarrow S$ se llama
\emph{morfismo de descenso} (resp. \emph{morfismo de descenso estricto}) si es
un morfismo de $\mathcal F$-descenso (resp. un morfismo de
$\mathcal F$-descenso estricto) (cf. definición 1.7), donde $\mathcal F$ denota
la categoría fibrada de base $\mathcal C$ de los objetos de $\mathcal C$ sobre
objetos de $\mathcal C$ (n\textsuperscript{o} 1, ejemplo 1).

\result{PROPOSICIÓN}{2.1} Si en $\mathcal C$ existen los productos finitos y
los productos fibrados (finitos), entonces, en $\mathcal C$, los morfismos de
descenso coinciden con los epimorfismos estrictos universales.

\result{b. EJEMPLOS}{} Sea $\mathcal C$ la categoría de los preesquemas. Sea
$S\in\mathcal C$, y sean $S'$ y $S''$ dos preesquemas finitos sobre $S$, es
decir, correspondientes a haces de álgebras $\mathcal A'$, $\mathcal A''$
sobre $S$ que, como haces de módulos, son cuasi-coherentes y de tipo finito (es
decir, coherentes si $S$ es localmente noetheriano). Sea
$\alpha:S'\longrightarrow S$ el morfismo estructural de $S'$, y sean
$\beta_1,\beta_2$ dos $S$-morfismos de $S''$ en $S'$, definidos por
homomorfismos de álgebras $\mathcal A'\longrightarrow\mathcal A''$, denotados
también por $\beta_1,\beta_2$. Utilizando el hecho de que un morfismo finito es
cerrado (primer teorema de Cohen--Seidenberg), se demuestra fácilmente que el
diagrama
\[
  \tag{+}
  S\xleftarrow{\alpha}S'
  \overset{\beta_1}{\underset{\beta_2}{\rightrightarrows}}S''
\]
en $\mathcal C$ es exacto si y solo si lo es el diagrama de haces

\src{306}
\[
  \mathcal O_S=\mathcal A\longrightarrow\mathcal A'
  \overset{\beta_1}{\underset{\beta_2}{\rightrightarrows}}\mathcal A''
\]
sobre $S$. En particular, si $\alpha:S'\longrightarrow S$ es un morfismo
finito correspondiente a un haz $\mathcal A'$ de álgebras sobre $S$, entonces
$\alpha$ es un epimorfismo estricto si y solo si el diagrama de haces
\[
  \mathcal O_S=\mathcal A\longrightarrow\mathcal A'
  \overset{p_1}{\underset{p_2}{\rightrightarrows}}
  \mathcal A'\otimes_{\mathcal A}\mathcal A'
\]
es exacto (es un epimorfismo si y solo si
$\mathcal A\longrightarrow\mathcal A'$ es inyectivo). Si $S$ es afín de anillo
$A$, y por tanto $S'$ es afín de anillo $A'$ finito sobre $A$, entonces
$S'\longrightarrow S$ es un epimorfismo estricto si y solo si
$A\longrightarrow A'$ es un isomorfismo de $A$ sobre el subanillo de $A'$
formado por los $x'\in A'$ tales que
\[
  1_{A'}\otimes_Ax'-x'\otimes_A1_{A'}=0
\]
(es un epimorfismo si y solo si $A\longrightarrow A'$ es inyectivo). Como ya
hemos señalado, incluso si $S$ es el esquema de un anillo local artiniano, un
morfismo finito $S'\longrightarrow S$ que sea un epimorfismo no es
necesariamente un epimorfismo estricto. Sin embargo, puede demostrarse que, si
$S$ es un preesquema noetheriano, todo morfismo finito
$S'\longrightarrow S$ que sea un epimorfismo es la composición de una sucesión
finita de epimorfismos estrictos (también finitos). Esto muestra, además, que
la composición de dos epimorfismos estrictos no es necesariamente un
epimorfismo estricto.

\result{c}{} Si $(+)$ es un diagrama exacto de morfismos finitos, entonces,
para todo morfismo plano $T\longrightarrow S$ de preesquemas, el diagrama
obtenido de $(+)$ mediante el cambio de base $T\longrightarrow S$ sigue siendo
exacto. De ello se deduce que, si $X,Y$ son dos $S$-preesquemas, siendo $X$
plano sobre $S$, entonces el siguiente diagrama de aplicaciones de conjuntos
(donde $X',Y'$ son las imágenes inversas de $X,Y$ sobre $S'$, y $X'',Y''$ sus
imágenes inversas sobre $S''$)
\[
  \operatorname{Hom}_S(X,Y)\longrightarrow
  \operatorname{Hom}_{S'}(X',Y')
  \rightrightarrows
  \operatorname{Hom}_{S''}(X'',Y'')
\]
es exacto. En particular, si $\mathcal F$ denota la categoría fibrada cuya base
es la categoría $\mathcal C$ de los preesquemas y tal que, para
$X\in\mathcal C$, $\mathcal F_X$ es la categoría de los $X$-preesquemas
planos, entonces el diagrama $(+)$ es $\mathcal F$-exacto. (Este resultado deja
de ser cierto si no se impone la hipótesis de planitud; en particular, un
epimorfismo estricto finito no es necesariamente un morfismo de descenso.) Se
ve asimismo que $(+)$ es $\mathcal F$-exacto si $\mathcal F$ denota la
categoría fibrada para la cual $\mathcal F_X$ es la categoría de los haces
cuasi-coherentes y planos sobre el preesquema $X$ (también aquí la hipótesis

\src{307}
de planitud es esencial). En uno y otro caso, la cuestión de la efectividad de
un dato de pegado (y, más concretamente, de un dato de descenso cuando
$S''=S'\times_SS'$) sobre un objeto plano sobre $S'$ es delicada (y su
respuesta en diversos casos particulares es uno de los objetos principales de
los presentes exposés). El conferenciante ignora si, para todo epimorfismo
estricto finito $S'\longrightarrow S$, todo dato de descenso sobre un haz
cuasi-coherente plano sobre $S'$ es efectivo (incluso suponiendo que $S$ es el
espectro de un anillo local artiniano y limitándose a los haces localmente
libres de rango 1). Más generalmente, sean $A$ un anillo y $A'$ una
$A$-álgebra (todo es conmutativo) tal que el diagrama de aplicaciones
\[
  A\longrightarrow A'\rightrightarrows A'\otimes_AA'
\]
sea exacto, lo que equivale también a que el morfismo
$S'\longrightarrow S$ correspondiente entre los espectros de $A,A'$ sea un
morfismo de $\mathcal F$-descenso, donde $\mathcal F$ es la categoría fibrada
de los haces cuasi-coherentes planos. Sea $M'$ un $A'$-módulo plano provisto de
un dato de descenso a $A$, es decir, de un isomorfismo
\[
  \varphi:M'\otimes_AA'\xrightarrow{\sim}A'\otimes_AM'
\]
de $A'\otimes_AA'$-módulos que satisface las condiciones (i) y (ii) de 1.c
(cuya explicitación en términos de módulos dejamos al lector). ¿Es efectivo
este dato (respecto de la categoría fibrada de los haces cuasi-coherentes
planos)? Sea $M$ el subconjunto de $M'$ formado por los $x'\in M'$ tales que
\[
  \varphi(x'\otimes_A1_{A'})=1_{A'}\otimes_Ax',
\]
que es un sub-$A$-módulo de $M'$. La inyección canónica
$M\longrightarrow M'$ define un homomorfismo de $A'$-módulos
$M\otimes_AA'\longrightarrow M'$. La efectividad de $\varphi$ significa
entonces lo siguiente: $M$ es un $A$-módulo plano y el homomorfismo anterior es
un isomorfismo.

\result{OBSERVACIÓN}{} En las consideraciones anteriores no habíamos impuesto
ninguna hipótesis de planitud a los morfismos del diagrama $(+)$, lo que nos
obligaba, para disponer de una técnica de descenso, a imponer hipótesis de
planitud a los objetos sobre $S,S'$ considerados. En el parágrafo 2 impondremos
una hipótesis de planitud a $\alpha:S'\longrightarrow S$, lo que nos permitirá
disponer de una técnica de descenso para objetos sobre $S,S'$ que ya no están
sometidos a condición alguna de planitud. En todos los casos interviene una
hipótesis de planitud. Esta es una de las principales razones de la importancia
de la noción de planitud en geometría algebraica (¡cuyo papel no podía
manifestarse mientras uno se limitara a cuerpos base, sobre los cuales, en
efecto, todo es plano!).

\src{308}
\subsection*{3. Aplicación a los morfismos étale}
\addcontentsline{toc}{subsection}{3. Aplicación a los morfismos étale}

Sean $A$ un anillo local y $B$ un álgebra local sobre $A$ cuyo ideal maximal
induce el de $A$. Diremos que $B$ es \emph{étale} sobre $A$ (en lugar de
«no ramificada», expresión usada en otros lugares) si satisface las condiciones
siguientes:
\begin{enumerate}
\item[(i)] $B$ es plana sobre $A$;
\item[(ii)] $B/\mathfrak mB$ es una extensión finita separable de
$A/\mathfrak m=k$ (donde $\mathfrak m$ denota el ideal maximal de $A$).
\end{enumerate}
Cuando $A$ y $B$ son noetherianos y $k$ es algebraicamente cerrado, esto
significa que el homomorfismo $\widehat A\longrightarrow\widehat B$ entre los
completados que prolonga $A\longrightarrow B$ es un isomorfismo. Un morfismo de
tipo finito $f:T\longrightarrow S$ se dice \emph{étale en $x\in T$}, o también
que $T$ es \emph{étale sobre $S$ en $x$}, si $\mathcal O_x$ es étale sobre
$\mathcal O_{f(x)}$; se dice que $f$ es \emph{étale}, o se llama a $f$ un
\emph{morfismo étale}, o se dice que $T$ es \emph{étale sobre $S$}, si $f$ es
étale en todo $x\in T$. Señalemos, además, que si $S$ es localmente
noetheriano, el conjunto de puntos de $T$ en los que $f$ es étale es abierto, y
el uso del «main theorem» de Zariski permite precisar la estructura de $T/S$
en un entorno de tal punto (mediante una ecuación de tipo bien conocido).

Cuando $S$ es un esquema de tipo finito sobre el cuerpo de los números
complejos, le corresponde un espacio analítico $\overline S$ en el sentido de
Serre [5], salvo que $\overline S$ puede tener elementos nilpotentes en su haz
estructural, lo que no cambia nada esencial en [5]. Se ve entonces fácilmente
que $f$ es étale si y solo si
$\overline f:\overline T\longrightarrow\overline S$ lo es, es decir, si todo
punto de $\overline T$ admite un entorno en el que $\overline f$ induce un
isomorfismo sobre un abierto de $\overline S$. En particular, a todo
\emph{recubrimiento étale} $T$ de $S$ (es decir, un morfismo finito étale
$f:T\longrightarrow S$) le corresponde un recubrimiento étale $\overline T$ de
$\overline S$, que es conexo si y solo si $T$ lo es [5]. Se ve además
fácilmente que, si $T,T'$ son dos esquemas étale sobre $S$, entonces la
aplicación natural
\[
  \operatorname{Hom}_S(T,T')\longrightarrow
  \operatorname{Hom}_{\overline S}(\overline T,\overline{T'})
\]
es biyectiva, es decir, el funtor $T\longmapsto\overline T$ de la categoría de
los esquemas étale sobre $S$ a la categoría de los espacios analíticos étale
sobre $\overline S$ es «plenamente fiel» y, por tanto, define una equivalencia
de la primera categoría con una subcategoría de la segunda. Un teorema de
Grauert--Remmert [2] implica que, si $S$ es normal, se obtiene así una
equivalencia entre la categoría de los \emph{recubrimientos étale} de $S$ y la
categoría de los recubrimientos étale (finitos) de $\overline S$, es decir, que
todo recubrimiento étale de $\overline S$ es $\overline S$-isomorfo a un
$\overline T$, donde $T$ es un recubrimiento étale de $S$. Demostremos que el
\emph{teorema de Grauert--Remmert sigue siendo válido sin hipótesis de
normalidad sobre $S$}. Sea, en efecto, primero $S'\longrightarrow S$ un
epimorfismo estricto finito; supongamos demostrado el

\src{309}
teorema para $S'$ y demostremos que será cierto para $S$. Sea
$\mathfrak T$ un recubrimiento étale de $\overline S$ y consideremos su imagen
inversa $\mathfrak T'$ sobre $\overline{S'}$, que corresponde a un haz
analítico coherente $\mathcal A'$ de álgebras sobre $\overline{S'}$, imagen
inversa del haz de álgebras $\mathcal A$ sobre $\overline S$ que define
$\mathfrak T$. Por hipótesis, sobre $\overline{S'}$, $\mathfrak T'$ proviene de
un recubrimiento étale $T'$ de $S'$, es decir, $\mathcal A'$ proviene de un haz
coherente de álgebras $A'$ sobre $S'$. Por otra parte, $\mathcal A'$ está
provisto de un dato de descenso canónico relativo a
$\overline{S'}\longrightarrow\overline S$, es decir, de un isomorfismo entre
sus dos imágenes inversas sobre
\[
  \overline{S'}\times_{\overline S}\overline{S'}
  =\overline{(S'\times_SS')},
\]
(que satisface las condiciones (i), (ii)), y este isomorfismo proviene, según
lo dicho, de un isomorfismo entre los haces algebraicos correspondientes, es
decir, de un dato de descenso sobre $A'$ relativo a
$S'\longrightarrow S$. Se comprueba fácilmente que este último es efectivo
(pues el de $\mathcal A'$ lo es, y la efectividad de un dato de descenso, tal
como se explicitó en el número anterior, se reconoce localmente en los
\emph{completados} de los módulos que intervienen). Se obtiene así un haz
coherente de álgebras $A$ sobre $S$, que define un recubrimiento $T$ de $S$,
el recubrimiento buscado. El resultado anterior sigue siendo entonces
manifiestamente válido si $S'\longrightarrow S$ es solo una composición de un
número finito de epimorfismos estrictos finitos, es decir, si es un
epimorfismo finito cualquiera (por el resultado de factorización señalado en
el parágrafo 2). Se sigue que el teorema de Grauert--Remmert continúa siendo
válido si $S$ es un esquema \emph{reducido}, es decir, si $\mathcal O_S$ no
tiene elementos nilpotentes, como se ve introduciendo su normalización $S'$.
De aquí se pasa fácilmente al caso general.

Una demostración enteramente análoga, que utiliza de nuevo el resultado de
factorización para los epimorfismos estrictos finitos y la naturaleza «formal»
de la efectividad de los datos de descenso, permite demostrar el resultado
siguiente: sea $S$ un preesquema localmente noetheriano y sea
$S'\longrightarrow S$ un morfismo \emph{finito, sobreyectivo y radicial} (o,
lo que equivale a lo mismo, un morfismo de tipo finito tal que, para todo $T$
sobre $S$, el morfismo $T'=S'\times_ST\longrightarrow T$ sea un homeomorfismo,
lo cual también se expresa diciendo que $S'\longrightarrow S$ es un
\emph{homeomorfismo universal}). Para todo $T$ étale sobre $S$, consideremos
su imagen inversa $T'=T\times_SS'$, que es étale sobre $S'$. Entonces el
funtor $T\longmapsto T'$ es una \emph{equivalencia de categorías entre la
categoría de los preesquemas $T$ étale sobre $S$ y la categoría de los
preesquemas $T'$ étale sobre $S'$}. (Se utiliza la biyectividad de
\[
  \operatorname{Hom}_S(T_1,T_2)\longrightarrow
  \operatorname{Hom}_{S'}(T'_1,T'_2)
\]
para dos preesquemas $T_1,T_2$ étale sobre $S$, hecho cuya comprobación directa
es fácil, y el hecho de que el teorema enunciado es cierto si
$S'=(S,\mathcal O_S/\mathcal J)$,

\src{310}
donde $\mathcal J$ es un haz coherente nilpotente de ideales de $\mathcal O_S$
([4], lema 6)). Obsérvese, además, que no suponemos aquí que los $T$
considerados sean finitos sobre $S$. Este resultado implica, en particular,
que el morfismo $S'\longrightarrow S$ induce un isomorfismo del grupo
fundamental $[\,]$ de $S'$ sobre el de $S$ («\emph{invariancia topológica del
grupo fundamental de un preesquema}»).

\subsection*{4. Relaciones con la 1-cohomología}
\addcontentsline{toc}{subsection}{4. Relaciones con la 1-cohomología}

\medskip\noindent
a. Sea $\mathcal C$ una categoría en la que siempre existe el producto de dos
objetos, y sea $T\in\mathcal C$. Para todo conjunto finito
$I\ne\varnothing$, se puede considerar $T^I$; al variar $I$, se obtiene así un
funtor covariante de la categoría de los conjuntos finitos no vacíos a
$\mathcal C$, es decir, lo que puede llamarse un \emph{objeto simplicial} de
$\mathcal C$, denotado por $K_T$. Este último depende covariantemente de $T$;
además, si $u,v$ son dos morfismos $T\longrightarrow T'$, los morfismos
correspondientes $K_T\longrightarrow K_{T'}$ son homotópicos. Diremos que $T$
\emph{domina} a $T'$ si $\operatorname{Hom}(T,T')\ne\varnothing$; esta es una
relación de preorden filtrante creciente en $\mathcal C$. De lo anterior se
deduce que, si $T$ domina a $T'$, existe una clase canónica (salvo homotopía)
de homomorfismos de objetos simpliciales $K_T\longrightarrow K_{T'}$; en
particular, si $K_T$ y $K_{T'}$ son tales que cada uno domina al otro, entonces
$K_T$ y $K_{T'}$ son homotópicamente equivalentes. Sea ahora $F$ un funtor
(contravariante, para fijar ideas) de $\mathcal C$ a una categoría abeliana
$\mathcal C'$; entonces
\[
  C^*(T,F)=F(K_T)
\]
es un objeto cosimplicial de $\mathcal C'$ y, por tanto, define de la manera
habitual un complejo (de cocadenas) en $\mathcal C'$, del cual puede tomarse la
cohomología
\[
  H^*(T,F)=H^*(C^*(T,F))=H^*(F(K_T))
\]
(se podrá poner una $\mathcal C$ como subíndice de $H^*$ si existe posibilidad
de confusión). Este es un funtor cohomológico en $F$, cuya varianza al variar
$T$ se deduce de lo dicho sobre los $K_T$; más precisamente, para $F$ fijo y
$T$ variable en $\mathcal C$ (preordenada por la relación de dominación), los
$H^*(T,F)$ forman un sistema inductivo de objetos graduados de $\mathcal C'$;
en particular, si $T$ y $T'$ son tales que cada uno domina al otro, entonces
$H^*(T,F)$ y $H^*(T',F)$ son canónicamente isomorfos.

Supongamos que en $\mathcal C$ existen los productos fibrados. Para
$S\in\mathcal C$ fijo, se puede aplicar lo anterior a la categoría
$\mathcal C_S$ de los objetos de $\mathcal C$ sobre $S$; escribiremos
$C^*(T/S,F)$ y $H^*(T/S,F)$ en lugar de $C^*(T,F)$ y $H^*(T,F)$ cuando
queramos precisar que trabajamos en la categoría $\mathcal C_S$; así,

\src{311}
$C^*(T/S,F)$ es un complejo de cocadenas en $\mathcal C'$ que, en dimensión
$n$, es igual a
\[
  F(T\times_ST\times_S\cdots\times_ST)
\]
(donde el paréntesis contiene $n+1$ factores).

Obsérvese que, como de costumbre, puede definirse $H^0(T/S,F)$ sin suponer
abeliana la categoría $\mathcal C'$: es el núcleo (definición 2.1), si existe,
del par de morfismos $F(p_i)$ ($i=1,2$)
\[
  F(T)\rightrightarrows F(T\times_ST)
\]
correspondientes a las dos proyecciones
$p_1,p_2:T\times_ST\rightrightarrows T$. En particular, se tendrá un morfismo
natural (llamado de \emph{aumento})
\[
  F(S)\longrightarrow H^0(T/S,F)
\]
que será un isomorfismo en los casos favorables (en particular, si
$T\longrightarrow S$ es un epimorfismo estricto y $F$ es «exacto por la
izquierda»). De igual modo, cuando $F$ toma valores en la categoría de los
grupos en una categoría $\mathcal C''$, también puede definirse $H^1(T/S,F)$;
cuando $\mathcal C''$ es la categoría de los conjuntos (es decir, $F$ toma
valores en la categoría de los grupos ordinarios, no necesariamente
conmutativos), $H^1(T/S,F)$ es el cociente del subgrupo $Z^1(T/S,F)$ de
$C^1(T/S,F)=F(T\times_ST)$ formado por los $g$ tales que
\[
  F(p_{31})(g)=F(p_{32})(g)\,F(p_{21})(g)
\]
por el grupo de operadores $F(T)$, que actúa sobre $C^1(T/S,F)$ y, en
particular, sobre el subconjunto $Z^1(T/S,F)$ mediante
\[
  \rho(g')\cdot g=F(p_2)(g')\,g\,F(p_1)(g')^{-1}.
\]

\medskip\noindent
b. Sea, por ejemplo, $\mathcal F$ una categoría fibrada de base $\mathcal C$.
Sean $\xi,\eta\in\mathcal F_S$ y, para todo $S'$ sobre $S$, sea
\[
  F_{\xi,\eta}(S')=
  \operatorname{Hom}(\xi\times_SS',\eta\times_SS').
\]
Así, $F_{\xi,\eta}$ es un funtor contravariante de $\mathcal C_S$ a la
categoría de los conjuntos. Dicho esto, \emph{afirmar que el morfismo de
aumento}
\[
  F_{\xi,\eta}(S)\longrightarrow H^0(S'/S,F_{\xi,\eta})
\]
\emph{es un isomorfismo para todo par de elementos
$\xi,\eta\in\mathcal F_S$ significa que
$\alpha:S'\longrightarrow S$ es un morfismo de $\mathcal F$-descenso}
(definición 1.7).

\medskip\noindent
c. Pongamos asimismo, para $\xi\in\mathcal F_S$ y todo objeto $S'$ de
$\mathcal C$ sobre $S$,
\[
  G_\xi(S')=\operatorname{Aut}(\xi\times_SS') ;
\]
se ha definido así un funtor contravariante $G_\xi$ de $\mathcal C_S$ a la
categoría de los

\src{312}
grupos. Dicho esto, se constata que $Z^1(S'/S,G_\xi)$ \emph{se identifica
canónicamente con el conjunto de los datos de descenso sobre
$\xi'=\xi\times_SS'$ relativos a $S'\longrightarrow S$} (definición 1.6), y
$H^1(S'/S,G_\xi)$ \emph{se identifica con el conjunto de las clases (salvo
isomorfismo) de objetos de $\mathcal F_{S'}$ provistos de un dato de descenso
relativo a $\alpha:S'\longrightarrow S$ que, como objetos de
$\mathcal F_{S'}$, son isomorfos a $\xi'=\xi\times_SS'$}. Por tanto, si
$\alpha:S'\longrightarrow S$ es un morfismo de $\mathcal F$-descenso
(cf. (b)), entonces $H^1(S'/S,G_\xi)$ \emph{contiene como subconjunto el
conjunto de las clases (salvo isomorfismo) de objetos $\eta$ de
$\mathcal F_S$ tales que $\eta\times_SS'$ sea isomorfo en
$\mathcal F_{S'}$ a $\xi\times_SS'$; y esta inclusión es una identidad si y
solo si todo dato de descenso sobre $\xi'=\xi\times_SS'$ relativo a
$\alpha:S'\longrightarrow S$ es efectivo}. (Este será el caso, en particular,
si $\alpha:S'\longrightarrow S$ es un morfismo de $\mathcal F$-descenso
estricto.)

\result{OBSERVACIÓN}{} Los complejos de cocadenas del tipo $C^*(T/S,F)$
contienen como casos particulares la mayoría de los complejos estándar
conocidos (cohomología de Čech, cohomología de grupos, etc.) y desempeñan un
papel importante en geometría algebraica, en particular en la «cohomología de
Weil» de los preesquemas.

\medskip\noindent
d. \result{EJEMPLO}{1} Sea $S'$ un objeto sobre $S\in\mathcal C$, y sea
$\Gamma$ un grupo de automorfismos de $S'$ tal que $S'$ sea «formalmente
principal sobre $S$, de grupo $\Gamma$», es decir, tal que el morfismo natural
\[
  \Gamma\times S'\longrightarrow S'\times_SS',
\]
donde $\Gamma\times S'$ denota la suma directa de $\Gamma$ copias de $S'$, sea
un isomorfismo. (Se supone que en $\mathcal C$ existen las sumas directas que
intervienen aquí.) Sea $F$ un funtor contravariante de $\mathcal C$ a la
categoría de los grupos abelianos. Entonces $C^*(S'/S,F)$ es canónicamente
isomorfo al grupo simplicial estándar de cocadenas homogéneas
$C^*(\Gamma,F(S'))$ y, por tanto, $H^*(S'/S,F)$ es canónicamente isomorfo a
$H^*(\Gamma,F(S'))$.

\medskip\noindent
e. \result{EJEMPLO}{2} Sea $\mathcal C$ la categoría de los preesquemas. Se
denota por $G_a$ («grupo aditivo») el funtor contravariante de $\mathcal C$ a
la categoría de los grupos abelianos definido por
\[
  G_a(X)=H^0(X,\mathcal O_X).
\]
Se define de igual modo el funtor $G_m$ («grupo multiplicativo») mediante
\[
  G_m(X)=H^0(X,\mathcal O_X)^*
\]
(grupo de los elementos invertibles del anillo $H^0(X,\mathcal O_X)$), y, más
generalmente, el funtor $\operatorname{Gl}(n)$ («grupo lineal de orden $n$»)
mediante

\src{313}
\[
  \operatorname{Gl}(n)(X)=
  \operatorname{Gl}\bigl(n,H^0(X,\mathcal O_X)\bigr),
\]
que es un funtor de $\mathcal C$ a la categoría de los grupos (no
necesariamente conmutativos si $n>1$; para $n=1$ se recupera $G_m$). Además,
puede interpretarse $\operatorname{Gl}(n)$ como un funtor de automorfismos
(cf. (c)) considerando la categoría fibrada $\mathcal F$ de base $\mathcal C$
tal que, para $X\in\mathcal C$, $\mathcal F_X$ sea la categoría de los haces
localmente libres sobre $X$: en efecto, se tiene
\[
  \operatorname{Gl}(n)(X)=
  \operatorname{Aut}_{\mathcal F_X}(\mathcal O_X^n).
\]
De (b) se sigue que, si $\alpha:S'\longrightarrow S$ es un morfismo de
$\mathcal F$-descenso (cf. parágrafo 2.c),
$H^1(S'/S,\operatorname{Gl}(n))$ contiene el conjunto de las clases (salvo
isomorfismo) de haces localmente libres sobre $S$ cuya imagen inversa sobre
$S'$ es isomorfa a $\mathcal O_{S'}^n$, y esta inclusión es una igualdad si y
solo si todo dato de descenso sobre $\mathcal O_{S'}^n$ (relativo a
$\alpha:S'\longrightarrow S$) es efectivo. Cuando $S$ es el espectro de un
anillo local, esto significa, por tanto,
$H^1(S'/S,\operatorname{Gl}(n))=(e)$, pues todo haz localmente libre sobre $S$
es entonces trivial.

Señalemos la equivalencia de las condiciones siguientes para un morfismo
$\alpha:S'\longrightarrow S$:
\begin{enumerate}
\item[(i)] el homomorfismo de aumento
\[
  H^0(S,\mathcal O_S)=G_a(S)\longrightarrow H^0(S'/S,G_a)
\]
es un isomorfismo;
\item[(ii)] $S'\longrightarrow S$ es un morfismo de $\mathcal F$-descenso
(siendo $\mathcal F$ la categoría fibrada de base $\mathcal C$ considerada
arriba).
\end{enumerate}
Si $S'\longrightarrow S$ es finito, estas condiciones equivalen también a:
\begin{enumerate}
\item[(iii)] $S'\longrightarrow S$ es un epimorfismo estricto
(cf. parágrafo 2.c).
\end{enumerate}

Supongamos ahora que $S=\operatorname{Spec}(A)$,
$S'=\operatorname{Spec}(A')$; entonces
\[
  C^n(S'/S,G_a)=C^n(A'/A,G_a)=\bigotimes_A^{n+1}A',
\]
siendo el operador coborde
$C^n(A'/A,G_a)\longrightarrow C^{n+1}(A'/A,G_a)$ la suma alternada de los
operadores de cara
\[
  \partial_i(x_0\otimes x_1\otimes\cdots\otimes x_n)
  =x_0\otimes\cdots\otimes x_{i-1}\otimes1_{A'}\otimes
   x_i\otimes\cdots\otimes x_n.
\]
Del mismo modo,
\[
  C^n(S'/S,G_m)=C^n(A'/A,G_m)
  \simeq\left(\bigotimes_A^{n+1}A'\right)^*,
\]
y las operaciones simpliciales en $C^*(A'/A,G_m)$ están inducidas por las de
$C^*(S'/S,G_a)$. Se explicitan análogamente las operaciones simpliciales en
$C^*(A'/A,\operatorname{Gl}(n))$. \emph{En todos los casos que conoce el
conferenciante, se tiene}
\[
  H^i(A'/A,G_a)=0\qquad\text{para $i>0$},
\]
\emph{y si $A$ es local, se tiene}
\[
  H^1(A'/A,G_m)=0
  \quad\text{y, más generalmente,}\quad
  H^1(A'/A,\operatorname{Gl}(n))=(e)
\]
(cuando $S'\longrightarrow S$ es un morfismo de $\mathcal F$-descenso, es
decir, cuando el diagrama
$A\longrightarrow A'\rightrightarrows A'\otimes_AA'$ es exacto; compárese con
el parágrafo 2.c). Obsérvese que \emph{el «teorema 90» de Hilbert no es sino
la relación}
\src{314}
\[
  H^1(S'/S,G_m)=0
\]
\emph{cuando $A$ es un cuerpo y $A'$ una extensión de Galois finita de este
(cf. ejemplo 1), y también puede expresarse diciendo que, en el caso
considerado, $S'\longrightarrow S$ es un morfismo de descenso estricto para la
categoría fibrada de los haces localmente libres de rango 1.} Es bajo esta
última forma como conviene generalizar el teorema de Hilbert, variando tanto
las hipótesis sobre el morfismo $S'\longrightarrow S$ como las relativas a los
haces cuasi-coherentes considerados.

Señalemos, por último, la equivalencia de las propiedades siguientes cuando
$A$ es un anillo local artiniano de ideal maximal $\mathfrak m$ y $A'$ una
$A$-álgebra (denotando, para todo entero $k>0$, por $A_k$ (resp. $A'_k$) los
anillos $A/\mathfrak m^{k+1}$ (resp. $A'/\mathfrak m^{k+1}A'$)):
\begin{enumerate}
\item[(i)] $H^1(A'_k/A_k,G_a)=0$ para todo $k$.
\item[(ii)] $H^1(A'_k/A_k,G_m)=0$ para todo $k$.
\item[(iii)] $H^1(A'_k/A_k,\operatorname{Gl}(n))=(e)$ para todo $k$ y todo
$n$.
\end{enumerate}
Si $S'\longrightarrow S$ es un epimorfismo estricto, las condiciones
anteriores implican incluso que es un morfismo de descenso estricto para los
módulos libres (de tipo finito o no) sobre $A'$.

\result{OBSERVACIÓN}{} La definición de los grupos $H^i(S'/S,G_m)$, cuando
$S,S'$ son esquemas de cuerpos $A,A'$, se debe a Amitsur. El grupo
$H^2(S'/S,G_m)$ es particularmente interesante como variante «global» del
grupo de Brauer; para esta variante se podrá consultar [1], capítulo VII.

\fsection{B}{Descenso por morfismos fielmente planos}

\subsection*{1. Enunciado de los teoremas de descenso}
\addcontentsline{toc}{subsection}{1. Enunciado de los teoremas de descenso}

\result{DEFINICIÓN}{1.1} Un morfismo
$\alpha:S'\longrightarrow S$ de preesquemas se dice \emph{plano} si, para todo
$x'\in S'$, $\mathcal O_{x'}$ es un módulo plano sobre el anillo
$\mathcal O_{\alpha(x')}$ (es decir, si
$\mathcal O_{x'}\otimes_{\mathcal O_{\alpha(x')}}M$ es un funtor exacto en el
$\mathcal O_{\alpha(x')}$-módulo $M$). Un morfismo se dice
\emph{fielmente plano} si es plano y sobreyectivo.

Por ejemplo, si $S=\operatorname{Spec}(A)$ y
$S'=\operatorname{Spec}(A')$, entonces $S'$ es plano sobre $S$ si y solo si
$A'$ es un $A$-módulo plano, y $S'$ es fielmente plano sobre $S$ si y solo si
$A'$ es un $A$-módulo fielmente plano (es decir, el funtor
$A'\otimes_A M$ en el $A$-módulo $M$ es exacto y fiel); esto significa también,
en la terminología de Serre [5], que el par $(A,A')$ es plano. Si $S'$ es
fielmente plano sobre $S$, entonces el funtor imagen inversa de haces
cuasi-coherentes sobre $S$ es exacto y fiel; en otras palabras, para que una
sucesión de homomorfismos de haces

\src{315}
cuasi-coherentes sobre $S$ sea exacta, es necesario y suficiente que lo sea su
imagen inversa sobre $S'$ (en particular, para que un homomorfismo de haces
cuasi-coherentes sobre $S$ sea un monomorfismo, resp. un epimorfismo, resp. un
isomorfismo, es necesario y suficiente que lo sea su imagen inversa sobre
$S'$). Esta propiedad sigue siendo cierta si nos restringimos sobre un abierto
cualquiera de $S$ y, bajo esta forma, caracteriza a los morfismos fielmente
planos.

\result{DEFINICIÓN}{1.2} Un morfismo $\alpha:S'\longrightarrow S$ se dice
\emph{cuasicompacto} si la imagen inversa de toda parte abierta cuasicompacta
$U$ de $S$ es cuasicompacta (es decir, unión finita de abiertos afines).

Evidentemente, basta comprobar esta propiedad para los abiertos afines de $S$.
Por ejemplo, un morfismo afín (es decir, tal que la imagen inversa de un abierto
afín sea afín) es cuasicompacto.

La clase de los morfismos planos, resp. fielmente planos, resp. cuasicompactos,
es estable por composición y por «cambio de base», y contiene, por supuesto,
los isomorfismos.

\result{TEOREMA}{1} Sea $\alpha:S'\longrightarrow S$ un morfismo de
preesquemas, fielmente plano y cuasicompacto. Entonces $\alpha$ es un
\emph{morfismo de descenso estricto} (A, definición 1.7) para la categoría
fibrada $\mathcal F$ de los haces cuasi-coherentes (A, parágrafo 1, ejemplo 2).

Este enunciado significa dos cosas:
\begin{enumerate}
\item[(i)] Si $F$ y $G$ son dos haces cuasi-coherentes sobre $S$, y $F'$ y
$G'$ sus imágenes inversas sobre $S'$, entonces el homomorfismo natural
\[
  \operatorname{Hom}(F,G)\longrightarrow\operatorname{Hom}(F',G')
\]
es una biyección del primer miembro sobre el subgrupo del segundo formado por
los homomorfismos $F'\longrightarrow G'$ compatibles con los datos de descenso
canónicos sobre estos haces, es decir, aquellos cuyas imágenes inversas por
las dos proyecciones de $S''=S'\times_SS'$ sobre $S'$ dan un mismo
homomorfismo $F''\longrightarrow G''$.
\item[(ii)] Todo haz cuasi-coherente $F'$ sobre $S'$, provisto de un dato de
descenso relativo al morfismo $\alpha:S'\longrightarrow S$
(A, definición 1.6), es isomorfo (provisto de este dato) a la imagen inversa de
un haz cuasi-coherente $F$ sobre $S$.
\end{enumerate}
Tomando $F=\mathcal O_S$ en (i), se obtiene:

\result{COROLARIO}{1} Sea $G$ un haz cuasi-coherente sobre $S$; sean
$G'$ y $G''$ sus imágenes inversas sobre $S'$ y sobre
$S''=S'\times_SS'$, y sean $p_1,p_2$ las dos proyecciones de $S''$ sobre
$S'$. Entonces el siguiente diagrama de aplicaciones de conjuntos

\src{316}
\[
  \Gamma(G)\xrightarrow{\alpha^*}\Gamma(G')
  \underset{p_2^*}{\overset{p_1^*}{\rightrightarrows}}\Gamma(G'')
\]
es exacto (A, definición 1.(a)).

Por otra parte, la conjunción de (i), (ii) y de la definición 1.1 da el

\result{COROLARIO}{2} Sea $G$ como en el corolario 1. Entonces existe una
correspondencia biunívoca entre los subhaces cuasi-coherentes de $G$ y los
subhaces cuasi-coherentes de $G'$ cuyas imágenes inversas sobre $S''$, por las
dos proyecciones $p_1,p_2$, dan el mismo subhaz de $G''$.

Por supuesto, se tiene un enunciado equivalente en términos de haces
cociente. Como se ha visto (A, parágrafo 4, (e)), el teorema 1 debe considerarse
una generalización del «teorema 90» de Hilbert, e implica como casos
particulares diversas formulaciones en términos de 1-cohomología. Para la
demostración, se reduce fácilmente al caso en que $S=\operatorname{Spec}(A)$,
$S'=\operatorname{Spec}(A')$, y para (i) se reduce fácilmente a demostrar el
corolario 1, es decir, la exactitud del diagrama
\[
  M=A\otimes_A M\longrightarrow A'\otimes_A M
  \rightrightarrows A'\otimes_AA'\otimes_A M
\]
para todo $A$-módulo $M$, lo cual se deduce del lema más general:

\result{LEMA}{1.1} Sea $A'$ una $A$-álgebra fielmente plana. Entonces, para
todo $A$-módulo $M$, el complejo aumentado por $M$
$C^*(A'/A,G_a)\otimes_A M$ (cf. A, parágrafo 4, (e)) es una
\emph{resolución} de $M$.

Basta demostrar que el complejo aumentado deducido del anterior mediante el
cambio de base de $A$ a $A'$ satisface las mismas conclusiones. Esto conduce a
comprobar el enunciado cuando se reemplazan $A$ por $A'$ y $A'$ por
$A'\otimes_AA'$, por lo que nos reduce a un caso en el que existe un
homomorfismo de $A$-álgebras $A'\longrightarrow A$ (o, en términos
geométricos, al caso en que $S'$ sobre $S$ admite una sección). En este caso,
esto se deduce de las generalidades de A, parágrafo 4, (a). Señalemos de paso
el siguiente corolario, que generaliza un enunciado bien conocido en
cohomología de Galois (compárese con A, parágrafo 4, (e)):

\result{COROLARIO}{} Si $A'$ es fielmente plana sobre $A$, se tiene
$H^0(A'/A,G_a)=A$ y
\[
  H^i(A'/A,G_a)=0\qquad\text{para $i\geq1$}.
\]

Para demostrar la parte (ii) del teorema 1, se procede como para (i),
reduciéndose al caso en que $S'$ sobre $S$ admite una sección, donde se deduce
de (i) (cf. A, parágrafo 1, (c)).

Evidentemente, se puede variar ad libitum el teorema 1 y sus corolarios
introduciendo diversas estructuras adicionales en los haces (o sistemas de
haces) cuasi-coherentes considerados. Por ejemplo, dar sobre $S$ un haz
\src{317}
cuasi-coherente de álgebras conmutativas «equivale» a dar sobre $S'$ un haz de
este tipo provisto de un dato de descenso relativo a
$\alpha:S'\longrightarrow S$. Teniendo en cuenta la correspondencia funtorial
entre tales haces cuasi-coherentes sobre $S$ y los preesquemas afines sobre
$S$, se obtiene la segunda afirmación del teorema siguiente:

\result{TEOREMA}{2} Sea $\alpha:S'\longrightarrow S$ como en el teorema 1.
Entonces $\alpha$ es un \emph{morfismo de descenso} (no estricto en general)
(A, definición 2.4), y es un \emph{morfismo de descenso estricto} para la
categoría fibrada de los esquemas afines sobre preesquemas
(A, definición 1.7).

La primera afirmación del teorema significa lo siguiente: sean $X,Y$ dos
preesquemas sobre $S$; sean $X',Y'$ sus imágenes inversas sobre $S'$, y
$X'',Y''$ sus imágenes inversas sobre $S''=S'\times_SS'$. Entonces el siguiente
diagrama de aplicaciones naturales
\[
  \operatorname{Hom}_S(X,Y)\xrightarrow{\alpha^*}
  \operatorname{Hom}_{S'}(X',Y')
  \underset{p_2^*}{\overset{p_1^*}{\rightrightarrows}}
  \operatorname{Hom}_{S''}(X'',Y'')
\]
es exacto; es decir, $\alpha^*$ es una biyección de
$\operatorname{Hom}_S(X,Y)$ sobre la parte de
$\operatorname{Hom}_{S'}(X',Y')$ formada por los homomorfismos compatibles con
los datos de descenso canónicos sobre $X',Y'$ (es decir, aquellos cuyas
imágenes inversas por las dos proyecciones de $S''$ sobre $S'$ son iguales).
Esto se deduce fácilmente del teorema 1, corolario 1, cuando uno se restringe
a $Y$ afín sobre $S$; en el caso general, hay que combinar el teorema 1 con el
resultado siguiente:

\result{LEMA}{1.2} Sea $\alpha:S'\longrightarrow S$ un morfismo fielmente plano
y cuasicompacto. Entonces $S$ se identifica con un \emph{espacio topológico
cociente} de $S'$; es decir, toda parte $U$ de $S$ tal que
$\alpha^{-1}(U)$ sea abierta es abierta.

Para completar el teorema 2, hay que dar criterios de efectividad para un dato
de descenso sobre un $S'$-preesquema $X'$ (cuando no se supone que $X'$ sea
afín sobre $S'$). Señalemos primero que \emph{tal dato de descenso no es
necesariamente efectivo}, incluso si $S$ es el espectro de un cuerpo $k$, $S'$
el espectro de una extensión cuadrática $k'$ de este y $X'$ un esquema
algebraico propio de dimensión 2 sobre $S'$ (como puede verse, según Serre,
utilizando la superficie no proyectiva de Nagata). \emph{Para que un dato de
descenso sobre $X'/S'$ relativo a $\alpha:S'\longrightarrow S$ (fielmente
plano y cuasicompacto) sea efectivo, es necesario y suficiente que $X'$ sea
unión de abiertos $X'_i$, afines sobre $S'$, que sean «estables» por el dato de
descenso sobre $X'$.} Esto se cumple con certeza (cualesquiera que sean $X'/S'$
y el dato de descenso sobre $X'$) si el morfismo
$\alpha:S'\longrightarrow S$ es \emph{radicial} (es decir, inyectivo y con
extensiones residuales que

\src{318}
son radiciales). Puede demostrarse también que sigue cumpliéndose si
$\alpha:S'\longrightarrow S$ es \emph{finito} y toda parte finita de $X'$,
contenida en una fibra de $X'$ sobre $S$, está contenida en un abierto de $X'$
afín sobre $S$ (este es el \emph{criterio de Weil}). En particular, esto se
cumple si $X'/S'$ es cuasiproyectivo, y en ese caso puede demostrarse que el
preesquema «descendido» $X/S$ también es cuasiproyectivo (y proyectivo si
$X'/S'$ lo es). En resumen:

\result{TEOREMA}{3} Sea $\alpha:S'\longrightarrow S$ un morfismo de
preesquemas, fielmente plano y cuasicompacto. Si $\alpha$ es \emph{radicial},
es un \emph{morfismo de descenso estricto}. Si $\alpha$ es finito, es un
morfismo de descenso estricto respecto de la categoría fibrada de los
preesquemas cuasiproyectivos (o proyectivos) sobre preesquemas.

\result{OBSERVACIONES}{} Ignoro si, en la segunda afirmación anterior, la
hipótesis de que $\alpha$ sea un morfismo \emph{finito} es realmente necesaria;
en cualquier caso, se comprueba de manera «formal» que puede sustituirse por la
siguiente hipótesis, aparentemente más débil: \emph{para todo punto de $S$
existe un entorno abierto $U$, un $U'$ finito y fielmente plano sobre $U$, y un
$S$-morfismo de $U'$ en $S'$.} Un caso típico que no entra en el anterior es
aquel en que $S=\operatorname{Spec}(A)$,
$S'=\operatorname{Spec}(\widehat A)$, donde $A$ es un anillo local noetheriano
y $\widehat A$ su completado; o también aquel en que $S'$ es cuasifinito sobre
$S$ (es decir, localmente isomorfo a un abierto de un $S$-esquema finito) y no
finito. En estos dos casos, el conferenciante también ignora la respuesta a la
pregunta siguiente: sea $X$ un $S$-esquema tal que $X'=X\times_SS'$ sea
proyectivo sobre $S'$; ¿es cierto que $X$ es proyectivo sobre
$S$?\footnote{La fe de erratas de 1962 responde negativamente: un morfismo
$S'\longrightarrow S$ cuasifinito, étale y sobreyectivo, o un morfismo
$\operatorname{Spec}(\widehat A)\longrightarrow\operatorname{Spec}(A)$, no es
siempre un morfismo de descenso estricto, incluso si $A$ es el anillo local de
una curva algebraica sobre un cuerpo algebraicamente cerrado $k$ y
$S=\operatorname{Spec}(A)$. Se puede encontrar así un morfismo propio y liso
$f:X\longrightarrow S$, que hace de $X$ un fibrado principal de base $S$ bajo
una curva elíptica $E$, tal que $f':X'\longrightarrow S'$ sea proyectivo, pero
$f$ no lo sea. Por tanto, también es un ejemplo de fibrado principal homogéneo
no isotrivial bajo un esquema abeliano.}

\subsection*{2. Aplicación al descenso de ciertas propiedades de los morfismos}
\addcontentsline{toc}{subsection}{2. Descenso de propiedades de los morfismos}

Sea $P$ una clase de morfismos de preesquemas. Sea
$\alpha:S'\longrightarrow S$ un morfismo de preesquemas, sea
$f:X\longrightarrow Y$ un morfismo de $S$-preesquemas y sea
$f':X'\longrightarrow Y'$ la imagen inversa de $f$ por $\alpha$. Cabe
preguntarse entonces si la relación «$f'\in P$» implica «$f\in P$». Resulta que
la respuesta es afirmativa en muchos casos importantes cuando se supone que
$\alpha$ es \emph{fielmente plano y cuasicompacto} (siendo a veces superflua
esta última hipótesis). Esto se ve directamente y sin dificultad si $P$ es la
clase de los morfismos sobreyectivos, resp. radiciales (estos casos se deducen
de la sobreyectividad de $\alpha$), resp. planos, resp. fielmente planos, resp.
lisos (estos casos se deducen de la planitud fiel de $\alpha$), resp. de tipo
finito. Utilizando los teoremas 1 y 2 y el lema 1.2, se ve también que ocurre
lo mismo si $P$ es una de las clases siguientes: isomorfismos, inmersiones
abiertas, inmersiones cerradas, inmersiones (si $f$ es de tipo finito e $Y$
localmente noetheriano), morfismos afines, morfismos finitos, morfismos
cuasifinitos, morfismos abiertos, morfismos cerrados, homeomorfismos, morfismos
separados,

\src{319}
morfismos propios. El único caso importante no dilucidado es el de los
morfismos proyectivos y cuasiproyectivos, ya señalado en la observación del
parágrafo 1.

\subsection*{3. Descenso por morfismos finitos fielmente planos}
\addcontentsline{toc}{subsection}{3. Descenso por morfismos finitos fielmente planos}

Sea $\alpha:S'\longrightarrow S$ un morfismo finito correspondiente a un haz de
álgebras $\mathcal A'$ sobre $S$ que, como haz de módulos, es localmente libre
de tipo finito y es $\ne0$ en todas partes; entonces $\alpha$ es un morfismo
fielmente plano y cuasicompacto, al que pueden aplicarse los resultados
anteriores. Dar un haz cuasi-coherente $F'$ sobre $S'$ equivale a dar el haz
cuasi-coherente $\alpha_*(F')$ sobre $S$, provisto de su estructura de
$\mathcal A'$-módulo (observando que
$\mathcal A'=\alpha_*(\mathcal O_{S'})$). Para simplificar, este haz sobre $S$
también se denotará por $F'$. Las dos imágenes inversas $p_i^*(F')$ de $F'$
sobre $S'\times_SS'$ corresponden asimismo a los haces cuasi-coherentes de
$(\mathcal A'\otimes_{\mathcal O_S}\mathcal A')$-módulos
\[
  F'\otimes_{\mathcal O_S}\mathcal A'
  \quad\text{y}\quad
  \mathcal A'\otimes_{\mathcal O_S}F'.
\]
Dar un $(S'\times_SS')$-homomorfismo del primero al segundo equivale a dar un
homomorfismo de $(\mathcal A'\otimes_{\mathcal O_S}\mathcal A')$-módulos y,
teniendo en cuenta que $\mathcal A'$ es localmente libre, esto equivale a dar
un homomorfismo de $(\mathcal A'\otimes\mathcal A')$-módulos:
\[
  \mathcal U=
  \operatorname{Hom}_{\mathcal O_S}(\mathcal A',\mathcal A')
  =\mathcal A'\otimes\check{\mathcal A}'
  \longrightarrow
  \operatorname{Hom}_{\mathcal O_S}(F',F'),
\]
es decir, a dar, para cada sección $\xi$ de $\mathcal U$ sobre un abierto $V$,
un homomorfismo de $\mathcal O_S$-módulos
$T_\xi:F'|_V\longrightarrow F'|_V$ que satisface las condiciones
\begin{equation}
  T_{f\xi}(x)=fT_\xi(x),
  \qquad T_{\xi f}(x)=T_\xi(fx),
  \tag{3.1}
\end{equation}
donde $f,x$ son, respectivamente, secciones de $\mathcal A',F'$ sobre un
abierto de $S$ contenido en $V$. Las condiciones (i) y (ii) de un dato de
descenso (A, parágrafo 1, (c)) se escriben entonces, respectivamente:
\begin{equation}
  T_{1_{\mathcal U}}x=x,
  \qquad\text{es decir,}\qquad
  T_{1_{\mathcal U}}=\operatorname{id}_{F'},
  \tag{3.2}
\end{equation}
\begin{equation}
  T_{\xi\eta}=T_\xi T_\eta.
  \tag{3.3}
\end{equation}
En otras palabras, \emph{un dato de descenso sobre $F'$ equivale a una
representación del haz de $\mathcal O_S$-álgebras
$\mathcal U=\operatorname{Hom}_{\mathcal O_S}(\mathcal A',\mathcal A')$ en el
haz de $\mathcal O_S$-álgebras
$\operatorname{Hom}_{\mathcal O_S}(F',F')$ que satisface las dos condiciones de
linealidad (3.1).} Si se tiene un apareamiento de haces cuasi-coherentes sobre
$S'$:
\[
  F'_1\times F'_2\longrightarrow F'_3
\]
(que puede interpretarse como un apareamiento de haces de
$\mathcal A'$-módulos sobre $S$),

\src{320}
y datos de pegado sobre los $F'_i$, definidos por homomorfismos
\[
  T_i:\mathcal U\longrightarrow
  \operatorname{Hom}_{\mathcal O_S}(F'_i,F'_i)
  \qquad (i=1,2,3),
\]
entonces estos datos son compatibles con el apareamiento dado, en el sentido
evidente del término, si y solo si se verifica la condición siguiente.

Para toda sección $\xi$ de $\mathcal U$ sobre un abierto, denotando por
\[
  \Delta\xi=\sum_i\xi'_i\otimes_{\mathcal A'}\xi''_i
\]
la sección de $\mathcal U\otimes_{\mathcal A'}\mathcal U$ (considerando
$\mathcal U$ como $\mathcal A'$-módulo mediante su estructura izquierda)
definida por la fórmula
\[
  \xi\mathbin{\cdot}(fg)=\sum_i\xi'_i(f)\xi''_i(g)
\]
(donde $f$ y $g$ son dos secciones de $\mathcal A'$ sobre un abierto más
pequeño), se tiene la fórmula
\begin{equation}
  T^{(3)}_\xi(x\mathbin{\cdot}y)
  =\sum_iT^{(1)}_{\xi'_i}x\mathbin{\cdot}T^{(2)}_{\xi''_i}y
  \tag{3.4}
\end{equation}
para todo par de secciones $x,y$ de $F'_1,F'_2$, respectivamente, sobre un
abierto más pequeño.\footnote{El texto impreso dice «de $\mathcal A'$». Los
dominios de $T^{(1)}$ y $T^{(2)}$, así como el apareamiento
$F'_1\times F'_2\longrightarrow F'_3$, obligan a que $x$ e $y$ sean secciones
de $F'_1$ y $F'_2$, respectivamente.} (Esta propiedad podrá expresarse
diciendo que los homomorfismos $T^{(i)}$ son compatibles con la aplicación
diagonal de $\mathcal U$ respecto del apareamiento dado.) En particular, las
fórmulas (3.1) a (3.4) permiten interpretar, en términos de representaciones de
álgebras con aplicaciones diagonales, los datos de descenso sobre un haz
cuasi-coherente de álgebras sobre $S'$ y, por tanto, también (restringiéndose a
las álgebras conmutativas) los datos de descenso sobre un $S'$-esquema afín.

De aquí se pasa a una interpretación análoga de los datos de descenso sobre un
$S'$-preesquema $X'$ cualquiera: dar tal $X'$ equivale a dar un preesquema
$X'$ sobre $S$, provisto de un homomorfismo de $\mathcal O_S$-álgebras
\[
  \mathcal A'\longrightarrow\mathcal O_{X'},
\]
y un dato de descenso sobre $X'$ equivale a dar un homomorfismo de haces
(compatible con el morfismo $h:X'\longrightarrow S'$)
\[
  \mathcal U\longrightarrow
  \operatorname{Hom}_{h^{-1}(\mathcal O_S)}
  (\mathcal O_{X'},\mathcal O_{X'}),
\]
que satisface condiciones análogas a las condiciones (3.1) a (3.4) anteriores.

\textsc{Ejemplo 1 («Weil»).} -- Supongamos que $S'/S$ es un recubrimiento étale
de Galois de grupo de Galois $\Gamma$ (cf. A, parágrafos 3 y 4.d). Entonces, un
dato de descenso sobre un haz cuasi-coherente $F'$ sobre $S'$ (resp. sobre un
$S'$-preesquema $X'$) equivale a dar una representación de $\Gamma$ por
automorfismos de $(S',F')$ (resp. de $(S',X')$) compatible con las operaciones
de $\Gamma$ sobre $S'$. Este resultado

\src{321}
es «formal», es decir, se demuestra en términos de categorías, pero desde el
punto de vista de este número se deduce también de la estructura explícita de
$\mathcal U$ (provisto de su estructura de anillo, del homomorfismo de anillos
$\mathcal A'\longrightarrow\mathcal U$ y de la aplicación diagonal), conocida
por completo gracias al resultado siguiente: $\mathcal U$ admite, como
$\mathcal A'$-módulo izquierdo, una base formada por las secciones de
$\mathcal U$ que corresponden a los elementos de $\Gamma$.

\textsc{Ejemplo 2 («Cartier»).} -- Sea $p$ un número primo; supongamos que
$p\mathcal O_S=0$ (es decir, $\mathcal O_S$ es de característica $p$),
$(\mathcal A')^p\subset\mathcal O_S=\mathcal A$ (es decir, $S'/S$ es radicial
de altura 1) y que el haz de álgebras $\mathcal A'$ sobre $\mathcal A$ admite
localmente una $p$-base, es decir, una familia $(x_i)$ de secciones tal que
$\mathcal A'$ está generada como álgebra por los $x_i$, sometidos a las únicas
condiciones $x_i^p=0$. Suponemos que el conjunto de los $i$ es finito, de
cardinal $n$. Sea $\mathfrak g$ el haz de las $\mathcal A$-derivaciones de
$\mathcal A'$; es un haz de $\mathcal A'$-módulos localmente libre de rango
$n$ y, además, un haz de $p$-álgebras de Lie sobre $\mathcal A$ (pero no sobre
$\mathcal A'$) que satisface la condición
\begin{equation}
  [X,fY]=X(f)Y+f[X,Y].
  \tag{3.5}
\end{equation}

\textsc{Lema.} -- $\mathcal U=
\operatorname{Hom}_{\mathcal O_S}(\mathcal A',\mathcal A')$ está generada,
como $\mathcal O_S$-álgebra provista de un homomorfismo de álgebras
$\mathcal A'\longrightarrow\mathcal U$, por el sub-$\mathcal A'$-módulo
izquierdo $\mathfrak g$, con las relaciones adicionales
\begin{equation}
  \left\{
  \begin{aligned}
    Xf-fX&=X(f),\\
    XY-YX&=[X,Y],\\
    X^p&=X^{(p)}.
  \end{aligned}
  \right.
  \tag{3.6}
\end{equation}

Del lema anterior se deduce que un dato de descenso sobre el haz
cuasi-coherente $F'$ sobre $S'$ equivale a dar, para todo $X\in\mathfrak g$, un
$\mathcal O_S$-endomorfismo $\overline X$ de $F'$ que satisface las condiciones
\begin{align}
  \overline{fX}&=f\overline X, \tag{3.7}\\
  \overline X(fx)&=X(f)x+f\overline X(x), \tag{3.8}\\
  \overline{[X,Y]}&=[\overline X,\overline Y], \tag{3.9}\\
  \overline{X^{(p)}}&=\overline X^p. \tag{3.10}
\end{align}
(Esto es lo que podría llamarse una conexión lineal sobre $F'$, sin curvatura
y compatible con la potencia $p$-ésima.) Se explicita de igual modo la noción
de dato de descenso sobre un $S'$-preesquema $X'$; la relación (3.4) se
sustituye aquí por la condición de que los $\overline X$ sean derivaciones de
$\mathcal O_{X'}$. Como el morfismo $S'\longrightarrow S$ es radicial, el
teorema 3 garantiza que todo dato de descenso de este tipo es

\src{322}
efectivo y, por tanto, define un $S$-preesquema $X$.

Obsérvese que no ha sido necesario imponer ninguna hipótesis de planitud, de no
singularidad ni de finitud de ningún tipo sobre $F'$ resp. $X'$.

\subsection*{4. Aplicación a criterios de racionalidad}
\addcontentsline{toc}{subsection}{4. Aplicación a criterios de racionalidad}

Sea $X$ un $S$-preesquema tal que la imagen directa de $\mathcal O_X$ sobre
$S$ sea $\mathcal O_S$; esta propiedad seguirá siendo cierta tras todo cambio
plano $S'\longrightarrow S$ de la base $S$. Si $F$ es un haz invertible (es
decir, localmente libre de rango 1) sobre $X$, los automorfismos de $F$, que se
identifican con las secciones invertibles de $\mathcal O_X$, corresponden
biunívocamente a las secciones invertibles de $\mathcal O_S$. Sea entonces $s$
una sección de $X$ sobre $S$; llamaremos, de manera gráfica, sección de $F$
sobre $s$ a una sección del haz invertible $s^*(F)$ sobre $S$. De lo anterior
se deduce que, si $F_i$ ($i=1,2$) son dos haces invertibles sobre $X$, cada uno
provisto de una sección sobre $s$, y si $F_1$ y $F_2$ son isomorfos, existe un
único isomorfismo de $F_1$ sobre $F_2$ compatible con las secciones en cuestión
(es decir, que transforma la primera en la segunda). Por otra parte, e
independientemente de la sección $s$, convengamos en considerar equivalentes
dos haces invertibles $F_1$ y $F_2$ sobre $X$ tales que todo punto de $S$ tenga
un entorno abierto $U$ para el cual las restricciones de $F_1$ y $F_2$ a
$X|U$ sean isomorfas. Entonces todo haz invertible $F$ sobre $X$ es equivalente
a un haz invertible $F_1$ provisto de una sección distinguida sobre $s$ (se
toma $F_1=F\,s^*(F)^{-1}$), y $F_1$ está determinado salvo isomorfismo. En
otras palabras, la clasificación de los haces invertibles sobre $X$ salvo
equivalencia es la misma que la clasificación, salvo isomorfismo, de los haces
invertibles provistos de una sección distinguida.

Como estas propiedades siguen siendo ciertas tras un cambio plano
$\alpha:S'\longrightarrow S$ de la base (reemplazando la sección $s$ por su
imagen inversa $s'$ por $\alpha$), se concluye, teniendo en cuenta el teorema 1:

\emph{Sea el preesquema $X/S$ como arriba y supongamos que admite una sección
$s$; sea $\alpha:S'\longrightarrow S$ un morfismo fielmente plano y
cuasicompacto; sea $F'$ un haz invertible sobre $X'=X\times_SS'$. Para que $F'$
sea equivalente a la imagen inversa sobre $X'$ de un haz invertible $F$ sobre
$X$, es necesario y suficiente que sus imágenes inversas $p_1^*(F')$ y
$p_2^*(F')$ sobre $X'\times_XX'=X\times_S(S'\times_SS')$ sean equivalentes. Si
es así, $F$ está determinado salvo equivalencia.} (Se dirá entonces que $F'$
es racional sobre $S$.)

Inspirándose en este principio cuando $\alpha:S'\longrightarrow S$ es como en
el ejemplo 1 o el ejemplo 2 del número anterior, se obtienen los criterios de
racionalidad de Weil o de Cartier. (Obsérvese que estos autores se limitan al
caso en que $S$ y $S'$

\src{323}
son espectros de cuerpos; a fortiori, $S$ es entonces el espectro de un anillo
local, y la relación de equivalencia introducida arriba no es sino la relación
de isomorfía.) En el primer caso, $F'$ es racional sobre $S$ si y solo si sus
transformados por $\Gamma$ son equivalentes a $F'$. Para expresar el criterio
de racionalidad en el segundo caso, se consideran, en general, el morfismo
diagonal
\[
  X'\longrightarrow X''=X'\times_XX'
\]
de $X'/X$, el haz de ideales $\mathcal I$ correspondiente sobre
$X'\times_XX'$ y el haz $\mathcal I/\mathcal I^2$, que se identifica con su
imagen inversa $\Omega^1_{X'/X}$ sobre $X'$. Como las restricciones de los
$F''_i=p_i^*(F')$ ($i=1,2$) a la diagonal son isomorfas (pues son isomorfas a
$F'$), es decir, $F''_1(F''_2)^{-1}=F''$ tiene una restricción a la diagonal
que es trivial, se sigue que la restricción de $F''$ a
$(X'',\mathcal O_{X''}/\mathcal I^2)$ viene dada, salvo isomorfismo, por un
elemento bien determinado $\xi$ de
\[
  H^1(X'',\mathcal I/\mathcal I^2)
  =H^1(X',\Omega^1_{X'/X}).
\]
Además, en este caso se tiene
\[
  \Omega^1_{X'/X}
  =\Omega^1_{S'/S}\otimes_{\mathcal O_S}\mathcal O_X,
\]
y, por tanto, si $\Omega^1_{S'/S}$ es localmente libre sobre $S$ (como en el
caso de Cartier), $\xi$ define una sección de
\[
  R^1f'_*(\mathcal O_{X'})\otimes\Omega^1_{S'/S}
\]
sobre $S'$ (llamada clase de Atiyah--Cartier del haz invertible $F'$ sobre
$X'/S$), cuya anulación es necesaria y suficiente para que sean equivalentes
las imágenes inversas de $F'$ por las dos proyecciones de
\[
  (X'',\mathcal O_{X''}/\mathcal I^2)
  =X\times_S(S'',\mathcal O_{S''}/\mathcal J^2)
\]
sobre $X'$ (donde $\mathcal J$ es el haz de ideales sobre
$S''=S'\times_SS'$ definido por el morfismo diagonal
$S'\longrightarrow S'\times_SS'$). Esta anulación es, pues, trivialmente
necesaria para que sean equivalentes las imágenes inversas de $F'$ sobre el
propio $X''=X\times_SS''$ y, por tanto, también para que $F'$ sea equivalente
a la imagen inversa de un haz invertible $F$ sobre $X$. Además, la clase de
Atiyah--Cartier puede interpretarse también como la obstrucción a la existencia,
localmente sobre $S'$, de una conexión de $F'$ relativa a las derivaciones de
$X'/X$; tal conexión queda, además, determinada cuando se conocen las
derivaciones de $F'$ correspondientes a las prolongaciones naturales a $X'$ de
las derivaciones de $S'/S$. De esto y de los desarrollos del número anterior se
concluye fácilmente que, en el caso del ejemplo 2 de dicho número y cuando
$X/S$ admite una sección, la anulación de la clase de Atiyah--Cartier es también
suficiente para que $F'$ sea racional sobre $S$.

\subsection*{5. Aplicación a la restricción de la base en un esquema abeliano}
\addcontentsline{toc}{subsection}{5. Restricción de la base en un esquema abeliano}

Sea $S$ un preesquema. Se llama esquema abeliano sobre $S$ a un esquema $X$
liso y propio

\src{324}
sobre $S$ cuyas fibras en los puntos $x\in S$ son esquemas de variedades
abelianas sobre los $\kappa(x)$. Supongamos que $S$ es noetheriano y regular
(es decir, que sus anillos locales son regulares); entonces, utilizando el
teorema de conexidad de Murre [4] (al menos en el caso de «características
iguales», en el que el teorema citado está demostrado actualmente), puede
demostrarse que toda sección racional de $X$ sobre $S$ está definida en todas
partes (es decir, es una sección), lo que generaliza un teorema clásico de
Weil. Se deduce, más generalmente, que si $X'$ es un esquema liso sobre $S$,
entonces toda $S$-aplicación racional de $X'$ en $X$ está definida en todas
partes. De ello se obtiene lo siguiente, que generaliza un resultado de
Chow--Iang:\footnote{El texto impreso dice «Igusa--Lang»; la fe de erratas de
1962 prescribe «Chow--Iang».} Sea $S$ noetheriano y regular, y denote $K$ su
anillo de funciones racionales (producto directo de cuerpos); sea $X$ un
esquema abeliano sobre $K$. Si $X$ es isomorfo a un $K$-esquema de la forma
$X_0\times_S\operatorname{Spec}(K)$, donde $X_0$ es un esquema abeliano sobre
$S$, entonces $X_0$ está determinado salvo isomorfismo único.

Utilizando el resultado de unicidad anterior, se ve que la cuestión de la
restricción de la base en $X$ es local sobre $S$ (y, por tanto, que basta saber
efectuar la restricción a los $\operatorname{Spec}(\mathcal O_x)$, con
$x\in S$). Se ve de igual manera que si
$\alpha:S'\longrightarrow S$ es un morfismo liso de tipo finito, si $K'$ es el
anillo de funciones racionales de $S'$, y si $X\otimes_KK'$ es de la forma
$X'_0\times_{S'}\operatorname{Spec}(K')$, entonces $X'_0$ está provisto de un
dato de descenso canónico relativo a $\alpha$. Teniendo en cuenta el teorema 3,
se concluye:

\result{PROPOSICIÓN}{5.1} Sea $S$ un preesquema noetheriano, regular e
irreducible, de cuerpo de funciones racionales $K$; sea $K'$ una extensión
finita de $K$, no ramificada sobre $S$; sea $S'$ la normalización de $S$ en
$K'$ (que es, por tanto, un recubrimiento étale de $S$); y sea $X$ un esquema
abeliano sobre $K$ tal que $X\otimes_KK'$ sea de la forma
$X'_0\times_{S'}\operatorname{Spec}(K')$, donde $X'_0$ es un esquema abeliano
proyectivo sobre $S'$. Entonces $X$ es de la forma
$X_0\times_S\operatorname{Spec}(K)$, donde $X_0$ es un esquema abeliano
proyectivo sobre $S$.

\result{OBSERVACIONES}{} El conferenciante ignora si puede sustituirse la
hipótesis de que $S'\longrightarrow S$ sea un recubrimiento étale sobreyectivo
(que permite utilizar el teorema 3) por la hipótesis de que sea un morfismo
liso, sobreyectivo y de tipo finito (incluso si se supone que es un morfismo
étale), o si la proposición sigue siendo válida sin suponer que $X'_0$ sea
proyectivo sobre $S'$ (condición que tal vez se cumpla automáticamente).

\subsection*{6. Aplicación a criterios de trivialidad local e isotrivialidad}
\addcontentsline{toc}{subsection}{6. Criterios de trivialidad local e isotrivialidad}

Sean $S$ un preesquema, $G$ un «preesquema en grupos» sobre $S$, y $P$ un
preesquema sobre $S$ sobre el que «actúa $G$» (por la derecha). Se dice que $P$
es formalmente principal homogéneo bajo $G$ si el morfismo bien conocido
\[
  G\times_SP\longrightarrow P\times_SP
\]

\src{325}
inducido por la acción de $G$ sobre $P$ es un isomorfismo. En lo sucesivo,
supondremos que $G$ es plano sobre $S$ (y, por tanto, fielmente plano sobre
$S$), y reservaremos el nombre de fibrado principal homogéneo bajo $G$ para un
fibrado formalmente principal homogéneo $P$ que sea fielmente plano y
cuasicompacto sobre $S$. Es inmediato que esto equivale a decir que puede
encontrarse un cambio fielmente plano y cuasicompacto
$S'\longrightarrow S$ de la base $S$, tal que el fibrado formalmente principal
homogéneo
\[
  P'=P\times_SS'
\]
bajo $G'=G\times_SS'$ sea trivial, es decir, isomorfo a $G'$ (es decir, que
admita una sección); en particular, se podrá tomar $S'=P$. Obsérvese además
que, si $S$ es localmente noetheriano, la hipótesis de planitud fiel sobre $P$
equivale a la hipótesis de que
\[
  \widehat P_s=P\times_S\operatorname{Spec}(\widehat{\mathcal O}_s)
\]
sea fielmente plano sobre $\widehat{\mathcal O}_s$ para todo $s\in S$ (donde
$\widehat{\mathcal O}_s$ denota el completado del anillo local
$\mathcal O_s$), como se deduce del hecho de que $\widehat{\mathcal O}_s$ es
fielmente plano sobre $\mathcal O_s$. Además, si $P$ es de tipo finito sobre
$S$ localmente noetheriano, el conjunto de los puntos $s$ que satisfacen la
condición anterior es constructible; por tanto, si $S$ es un «preesquema de
Jacobson» (por ejemplo, un esquema de tipo finito sobre un cuerpo o, más
generalmente, el espectro de un anillo de Jacobson), basta comprobar la
condición en cuestión en los puntos cerrados de $S$. Esto nos reduce al caso en
que la base es el espectro de un anillo local completo $A$. Cuando
$S=\operatorname{Spec}(A)$ ($A$ un anillo local noetheriano completo) y $P$ es
de tipo finito sobre $S$, la planitud fiel de $P/S$ equivale también a la
existencia de un $S'$ finito y plano sobre $S$ tal que $P'$ sea trivial; y, si
además $G$ es liso sobre $S$, puede suponerse que $S'$ es étale sobre $S$. Por
consiguiente, si además el cuerpo residual de $A$ es algebraicamente cerrado
(«caso geométrico»), $P$ es fielmente plano sobre $A$ si y solo si es trivial.
Así pues, si $S$ es un preesquema algebraico sobre un cuerpo algebraicamente
cerrado y $G$ es liso y de tipo finito sobre $S$, se ve que la condición de
planitud fiel sobre $S$ equivale a la condición de trivialidad analítica (SLF)
de Serre ([6], pp. 1--12).

Pueden introducirse otros tipos, más fuertes, de condiciones sobre $P$ que
tengan la naturaleza de una «trivialidad local». Diremos, en particular, que
$P$ es isotrivial (resp. estrictamente isotrivial) si, para todo $s\in S$,
existen un entorno abierto $U$ de $s$ y un morfismo finito y fielmente plano
(resp. un recubrimiento étale sobreyectivo) $U'\longrightarrow U$ tales que
$P'=P\times_SU'$ sea trivial. (Nuestra terminología difiere de la de Serre [6],
quien llama localmente isotrivial a lo que nosotros llamamos estrictamente
isotrivial.) La isotrivialidad estricta es especialmente útil si $G$ es liso
sobre $S$, pero, en cambio, es una noción inadecuada en los demás casos.

Si $G$ es afín sobre $S$, todo fibrado principal homogéneo $P$ bajo $G$ es afín
por el parágrafo 2, de donde, gracias al teorema 2, resulta la posibilidad de
«descender»

\src{326}
tales fibrados por morfismos fielmente planos y cuasicompactos. Tomemos, en
particular, $G=\operatorname{Gl}(n)_S$, definido por la condición de que el
funtor de $S$-preesquemas
\[
  S'\longmapsto\operatorname{Hom}_S(S',G)
\]
(con valores en la categoría de los grupos) se identifique con el funtor
\[
  \operatorname{Gl}(n)(S')
  =\operatorname{Gl}\bigl(n,H^0(S',\mathcal O_{S'})\bigr)
\]
de A, parágrafo 4.e. Utilizando los hechos (i) de que todo fibrado principal
homogéneo bajo $G$ (resp. todo haz localmente libre de rango $n$ sobre $S$) se
vuelve isomorfo al objeto «trivial» $G$ (resp. $\mathcal O_S^n$) tras un cambio
fielmente plano y cuasicompacto adecuado de $S$; (ii) de que pueden descenderse
los objetos del tipo considerado (fibrados principales homogéneos bajo $G$,
resp. haces localmente libres de rango $n$) por tales morfismos; y, por último,
(iii) de que el grupo de automorfismos del fibrado trivial sobre un $S'/S$ es
funtorialmente isomorfo al grupo de automorfismos del haz localmente libre
trivial de rango $n$ sobre $S'$, se concluye «formalmente» que «da lo mismo»
dar sobre $S$ (o sobre un $S'/S$) un fibrado principal homogéneo de grupo $G$
que dar allí un haz localmente libre de rango $n$. (Más precisamente, se tiene
una equivalencia de categorías fibradas.) Se concluye, en particular:

\result{PROPOSICIÓN}{6.1} Todo fibrado principal homogéneo de grupo
$\operatorname{Gl}(n)_S$ es localmente trivial.

Mediante argumentos conocidos se obtiene el mismo resultado para grupos
estructurales tales como $\operatorname{Sl}(n)_S$, $\operatorname{Sp}(n)_S$ y
productos de tales grupos. Se deduce también que, si $F$ es un subgrupo cerrado
de $G=\operatorname{Gl}(n)_S$, plano sobre $S$, tal que existe el cociente
$G/F$, y si $G$ es un fibrado principal homogéneo isotrivial (resp.
estrictamente isotrivial) sobre $G/F$, de grupo $F\times_S(G/F)$, entonces todo
fibrado principal homogéneo de grupo $F$ es isotrivial (resp. estrictamente
isotrivial). Esto se aplica a todos los «grupos lineales» sobre $S$ utilizados
hasta ahora y, en particular, cuando $G=S\times_k\Gamma$, donde $S$ es un
preesquema sobre el cuerpo $k$ y $\Gamma$ un grupo lineal en el sentido clásico
y, en particular, liso, sobre $k$. Esto resuelve, pues, para tales grupos, una
cuestión de Serre (loc. cit.).

Señalemos asimismo que, para la mayoría de los grupos lisos sobre $S$ conocidos,
sean lineales o no, y en cualquier caso para los de la forma
$S\times_k\Gamma$ anterior, puede demostrarse que todo fibrado principal
homogéneo isotrivial es estrictamente isotrivial. Esto resuelve, en particular,
otra cuestión de Serre (loc. cit., pp. 1--14), dado que un fibrado principal
homogéneo obtenido por un descenso a la Cartier (cf. parágrafo 3, ejemplo 2) es
manifiestamente isotrivial.

\result{OBSERVACIÓN}{} Una de las dificultades esenciales en estas cuestiones
(aparte de la cuestión de la existencia de los esquemas cociente $G/F$) es la
falta de criterios de efectividad para un dato de descenso por un morfismo
fielmente plano no finito.

\src{327}
\begin{center}
  \textbf{BIBLIOGRAFÍA}
\end{center}

\begin{description}
  \item[{[1]}] Dieudonné (J.) y Grothendieck (Alexander). --- \emph{Éléments
    de géométrie algébrique}, de próxima aparición en las Publications
    mathématiques de l'Institut des Hautes Études Scientifiques.
  \item[{[2]}] Grauert (H.) und Remmert (R.). --- \emph{Komplexe Räume}, Math.
    Annalen, t. 136, 1958, pp. 245--318.
  \item[{[3]}] Grothendieck (Alexander). --- \emph{Géométrie formelle et
    géométrie algébrique}. Séminaire Bourbaki, t. 11, 1958/59, no 182.
  \item[{[4]}] Murre (J. P.). --- \emph{On a connectedness theorem for a
    birational transformation at a simple point}, Amer. J. Math., t. 80, 1958,
    pp. 3--15.
  \item[{[5]}] Serre (Jean-Pierre). --- \emph{Géométrie algébrique et géométrie
    analytique}, Ann. Inst. Fourier Grenoble, t. 6, 1955--56, pp. 1--42.
  \item[{[6]}] Serre (Jean-Pierre). --- \emph{Espaces fibrés algébriques},
    Séminaire Chevalley, t. 2, 1958: Anneaux de Chow et applications, no 1.
\end{description}

\def\FgaUnitKey{e190}\clearpage
\chapter*{Fe de erratas del Exposé n\textsuperscript{o} 190}
\addcontentsline{toc}{chapter}{Fe de erratas del Exposé n\textsuperscript{o} 190}
\setcounter{footnote}{0}
\src{300}

\noindent\textbf{3. Técnica de descenso y teoremas de existencia en geometría
algebraica [TDTE I]. I: Generalidades. Descenso por morfismos fielmente planos.}

\begin{center}
[Séminaire Bourbaki, t. 12, 1959/60, n\textsuperscript{o} 190, 29 pp.]
\end{center}

\noindent\textbf{Página 1, línea 10.} --- Ahora parece excesivo decir que la
técnica de descenso está «en la base de la mayoría de los teoremas de
existencia en geometría algebraica». Esto es cierto en gran medida para las
técnicas no proyectivas que constituyen el objeto de los dos primeros exposés
de la serie TDTE I a VI, pero no para las técnicas proyectivas (TDTE IV, V,
VI).

\noindent\textbf{Página 5, línea 16.} --- No es necesario suponer que $\alpha$
sea un morfismo de $\mathcal F$-descenso.

\noindent\textbf{Página 20, observación.} --- Un morfismo
$S'\longrightarrow S$ cuasifinito, étale y sobreyectivo, o un morfismo
$\operatorname{Spec}(\widehat A)\longrightarrow\operatorname{Spec}(A)$, no es
siempre un morfismo de descenso estricto, aun si $A$ es un anillo local de una
curva algebraica sobre un cuerpo algebraicamente cerrado $k$ y
$S=\operatorname{Spec}(A)$. Así, se puede encontrar un morfismo propio y liso
$f:X\longrightarrow S$, que hace de $X$ un fibrado principal de base $S$ bajo
una curva elíptica $E$, tal que $f':X'\longrightarrow S'$ sea proyectivo, pero
$f$ no lo sea. Es, por tanto, al mismo tiempo un ejemplo de un fibrado
homogéneo principal no isotrivial bajo un esquema abeliano.

\noindent\textbf{Página 26, línea 9.} --- Léase «Chow--Lang» en lugar de
«Igusa--Lang».

Para diversos detalles acerca de la teoría del descenso, véanse SGA VI, VII y
VIII.

\def\FgaUnitKey{195}\unit{Exposición 195 -- Técnica de descenso II}
\sid{195}
\src{369}
\seminar{12.\textsuperscript{o} año, 1959/60, n.\textsuperscript{o} 195\\Febrero de 1960}
\paperhead{TÉCNICA DE DESCENSO Y TEOREMAS DE EXISTENCIA EN GEOMETRÍA ALGEBRAICA\\
II. EL TEOREMA DE EXISTENCIA EN LA TEORÍA FORMAL DE LOS MÓDULOS}%
  {por Alexander GROTHENDIECK}

\fsection{A}{Funtores representables y pro-representables}

\subsection*{1. Funtores representables}
\addcontentsline{toc}{subsection}{1. Funtores representables}

Sea $\mathcal C$ una categoría. Para todo $X\in\mathcal C$, sea $h_X$ el
funtor contravariante de $\mathcal C$ en la categoría $(\operatorname{Ens})$
de los conjuntos,
\[
  h_X:\mathcal C\longrightarrow(\operatorname{Ens}),
\]
definido por la fórmula
\[
  h_X(Y)=\operatorname{Hom}(Y,X).
\]
Un morfismo $X\longrightarrow X'$ en $\mathcal C$ induce de manera evidente un
homomorfismo de funtores $h_X\longrightarrow h_{X'}$: $h_X$ es un funtor
covariante en $X$; es decir, hemos definido un funtor covariante canónico
\[
  h:\mathcal C\longrightarrow
  \operatorname{Hom}\bigl(\mathcal C^o,(\operatorname{Ens})\bigr)
\]
de $\mathcal C$ en la categoría de los funtores covariantes de la dual
$\mathcal C^o$ de $\mathcal C$ en la categoría de los conjuntos. Recordemos
entonces:

\result{PROPOSICIÓN}{1.1} Este funtor $h$ es \emph{plenamente fiel}; en otros
términos, para todo par $X,X'$ de objetos de $\mathcal C$, la aplicación
natural
\[
  \operatorname{Hom}(X,X')\longrightarrow
  \operatorname{Hom}(h_X,h_{X'})
\]
es biyectiva.

En particular, si un funtor
$F\in\operatorname{Hom}(\mathcal C^o,(\operatorname{Ens}))$ es isomorfo a un
funtor de la forma $h_X$, entonces $X$ queda determinado salvo isomorfismo
único. Se dice en tal caso que el funtor $F$ es \emph{representable}. La
proposición anterior significa que el funtor canónico $h$ define una
equivalencia de la categoría $\mathcal C$ con la subcategoría plena de
$\operatorname{Hom}(\mathcal C^o,(\operatorname{Ens}))$ formada por los
funtores representables. Este hecho fundamenta la noción de «solución de un
problema universal»: tal problema consiste siempre en examinar si un funtor
dado (contravariante como aquí, o covariante en el caso dual) de $\mathcal C$
en $(\operatorname{Ens})$ es representable.

\src{370}
Observemos además que, por la propia definición de producto en una categoría
[1], el funtor $h:X\longmapsto h_X$ conmuta con los productos siempre que estos
existan (y, más generalmente, con los límites proyectivos finitos o infinitos;
en particular, con los productos fibrados, con la formación de «núcleos» [2],
etc., siempre que existan): se tiene un isomorfismo de funtores
\[
  h_{X\times X'}\xrightarrow{\sim}h_X\times h_{X'}
\]
siempre que exista $X\times X'$; es decir, se tienen biyecciones functoriales
en $Y$
\[
  h_{X\times X'}(Y)\xrightarrow{\sim}h_X(Y)\times h_{X'}(Y).
\]
En particular, dar un morfismo
\[
  X\times X'\longrightarrow X''
\]
en $\mathcal C$ (es decir, una «ley de composición» en $\mathcal C$ entre
$X,X',X''$) equivale a dar un morfismo
\[
  h_{X\times X'}=h_X\times h_{X'}\longrightarrow h_{X''},
\]
es decir, a dar, para todo $Y\in\mathcal C$, una ley de composición de
conjuntos
\[
  h_X(Y)\times h_{X'}(Y)\longrightarrow h_{X''}(Y)
\]
de tal modo que, para todo morfismo $Y\longrightarrow Y'$ en $\mathcal C$, el
sistema de aplicaciones entre conjuntos
\[
  h_{X^{(i)}}(Y)\longleftarrow h_{X^{(i)}}(Y')
  \qquad(\text{para }i=0,1,2)
\]
sea un morfismo respecto de las dos leyes de composición correspondientes a
$Y$ y $Y'$. Se ve así que la noción de estructura de «$\mathcal C$-grupo»,
«$\mathcal C$-anillo», etc. sobre un objeto $X$ de $\mathcal C$ se expresa del
modo más cómodo (tanto en la teoría como en la práctica) diciendo que, para
todo $Y\in\mathcal C$, el conjunto $h_X(Y)$ está provisto de una ley de grupo
(resp. de anillo, etc.) en el sentido usual, y que las aplicaciones
$h_X(Y)\longleftarrow h_X(Y')$ correspondientes a morfismos
$Y\longrightarrow Y'$ son homomorfismos de grupos (resp. de anillos, etc.).
Esta es, por ejemplo, la forma más intuitiva y cómoda de definir los diversos
grupos clásicos $G_a$, $G_m$, $\operatorname{Gl}(n)$, etc. sobre un preesquema
base arbitrario $S$ y de escribir las relaciones clásicas entre esos grupos;
de dotar al esquema afín $V(\mathcal F)$ sobre $S$, definido por un haz
cuasicoherente $\mathcal F$, de una estructura de «fibrado vectorial»; y de
definir y estudiar las diversas variedades de banderas asociadas
(grassmannianas, fibrados proyectivos), etc. El yoga general consiste
sencillamente en identificar, mediante el funtor canónico $h$, los objetos de
$\mathcal C$ con ciertos funtores contravariantes: los funtores representables,

\src{371}
de $\mathcal C$ en la categoría de los conjuntos.

El procedimiento habitual de invertir las flechas, necesario por ejemplo en el
caso de los esquemas afines para pasar del lenguaje geométrico al del álgebra
conmutativa, conduce a dualizar las consideraciones anteriores y, en
particular, a introducir también los funtores covariantes representables
$\mathcal C\longrightarrow(\operatorname{Ens})$; es decir, los de la forma
\[
  Y\longmapsto\operatorname{Hom}(X,Y)=h'_X(Y).
\]

\subsection*{2. Funtores pro-representables, proobjetos}
\addcontentsline{toc}{subsection}{2. Funtores pro-representables, proobjetos}

Sea $\widetilde X=(X_i)_{i\in I}$ un sistema proyectivo de objetos de
$\mathcal C$; le corresponde un funtor covariante
\[
  h'_{\widetilde X}=\varinjlim_i h'_{X_i},
\]
más explícitamente,
\[
  h'_{\widetilde X}(Y)=\varinjlim_i h'_{X_i}(Y)
  =\varinjlim_i\operatorname{Hom}(X_i,Y),
\]
de $\mathcal C$ en $(\operatorname{Ens})$. Un funtor de $\mathcal C$ en
$(\operatorname{Ens})$ que sea isomorfo a un funtor de este tipo con $I$
filtrante se llama \emph{pro-representable}. Por el número anterior, estos son
exactamente los funtores isomorfos a límites inductivos filtrantes de funtores
representables. Sea $\widetilde X'=(X'_j)_{j\in J}$ un segundo sistema
proyectivo filtrante en $\mathcal C$ (construido sobre un segundo conjunto de
índices preordenado filtrante $J$); se comprueba fácilmente que hay una
biyección canónica (que generaliza la proposición 1.1)
\[
  \operatorname{Hom}(h'_{\widetilde X'},h'_{\widetilde X})
  =\varprojlim_j\varinjlim_i\operatorname{Hom}(X_i,X'_j).
\]
Esto lleva a introducir la categoría $\operatorname{Pro}(\mathcal C)$ de los
proobjetos de $\mathcal C$; sus objetos son los sistemas proyectivos de objetos
de $\mathcal C$ (sobre conjuntos variables de índices preordenados filtrantes),
y si $\widetilde X=(X_i)_{i\in I}$ y
$\widetilde X'=(X'_j)_{j\in J}$ son dos objetos de esta clase, se pone
\[
  \operatorname{ProHom}(\widetilde X,\widetilde X')
  =\varprojlim_j\varinjlim_i\operatorname{Hom}(X_i,X'_j),
\]
siendo evidente la composición de los prohomomorfismos. Por su propia
construcción, $\widetilde X\longmapsto h'_{\widetilde X}$ puede considerarse
como un funtor contravariante en $\widetilde X$, que establece una equivalencia
entre la categoría dual de la categoría $\operatorname{Pro}(\mathcal C)$ de
los proobjetos de $\mathcal C$ y la categoría de los funtores covariantes
pro-representables de $\mathcal C$ en $(\operatorname{Ens})$. Naturalmente, un
objeto $X$ de $\mathcal C$ define canónicamente un proobjeto, denotado

\src{372}
también por $X$, de modo que $\mathcal C$ es equivalente a una subcategoría
plena de $\operatorname{Pro}(\mathcal C)$. Si entonces
$\widetilde X=(X_i)_{i\in I}$ es un proobjeto cualquiera de $\mathcal C$, se
tendrá (con la identificación anterior)
\[
  \widetilde X=\varprojlim_i X_i,
\]
tomándose el límite proyectivo en $\operatorname{Pro}(\mathcal C)$
(puesto que $h'_{\widetilde X}=\varinjlim_i h'_{X_i}$). Hay que observar que,
aun cuando el límite proyectivo de los $X_i$ exista en $\mathcal C$, en
general no será isomorfo al límite proyectivo $\widetilde X$ en
$\operatorname{Pro}(\mathcal C)$, como ya resulta evidente cuando $\mathcal C$
es la categoría de los conjuntos. Nótese que, por definición,
$\varprojlim_{\mathcal C}X_i=L$ queda definido por la condición de que el
funtor en $Y\in\mathcal C$
\[
  \varprojlim_i\operatorname{Hom}_{\mathcal C}(Y,X_i)
  =\operatorname{Hom}_{\operatorname{Pro}(\mathcal C)}(Y,\widetilde X)
\]
con valores en $(\operatorname{Ens})$ sea representable mediante $L$, es decir,
sea isomorfo a $\operatorname{Hom}_{\mathcal C}(Y,L)$; por consiguiente,
$\varprojlim_{\mathcal C}X_i$ ya está definido en términos del proobjeto
$\widetilde X$ y, más precisamente, depende functorialmente del proobjeto
$\widetilde X$ siempre que esté definido; no hay, pues, inconveniente en
denotarlo por $\varprojlim_{\mathcal C}(\widetilde X)$. Si los límites
proyectivos existen siempre en $\mathcal C$,
$\varprojlim_{\mathcal C}(\widetilde X)$ es un funtor de
$\operatorname{Pro}(\mathcal C)$ en $\mathcal C$, y hay un homomorfismo
canónico de funtores
$\varprojlim_{\mathcal C}(\widetilde X)\longrightarrow\widetilde X$.

Como todo funtor, que supondremos covariante para fijar las ideas,
\[
  F:\mathcal C\longrightarrow\mathcal C'
\]
se prolonga de manera evidente a un funtor
\[
  \operatorname{Pro}(F):\operatorname{Pro}(\mathcal C)
  \longrightarrow\operatorname{Pro}(\mathcal C'),
\]
se sigue que, si en $\mathcal C'$ existen siempre los límites proyectivos, $F$
define también canónicamente un funtor compuesto
\[
  \overline F:\operatorname{Pro}(\mathcal C)\longrightarrow\mathcal C',
  \qquad
  \overline F=\varprojlim_{\mathcal C'}\circ\operatorname{Pro}(F),
\]
que transforma $\widetilde X=(X_i)_{i\in I}$ en
$\varprojlim_{\mathcal C'}F(X_i)$.

Un proobjeto $\widetilde X$ se llama \emph{proobjeto estricto} si es isomorfo a
un proobjeto $(X_i)_{i\in I}$ cuyos morfismos de transición
$X_i\longrightarrow X_j$ son epimorfismos; un funtor definido por tal objeto se
llama \emph{estrictamente pro-representable}. Se puede exigir además que $I$
sea ordenado filtrante y que todo epimorfismo $X_i\longrightarrow X'$ sea
equivalente a un epimorfismo

\src{373}
$X_i\longrightarrow X_j$ para un $j\in I$ adecuado (determinado de manera
única por esta condición). En estas condiciones, el sistema proyectivo
$(X_i)_{i\in I}$ queda determinado salvo isomorfismo único (en el sentido usual
de los isomorfismos de sistemas proyectivos). Se sigue que en
$\operatorname{Pro}(\mathcal C)$ existe siempre el límite proyectivo de un
sistema proyectivo de proobjetos estrictos $\widetilde X^{(\alpha)}$ y que, con
las notaciones anteriores para $F$, $\overline F$, se tendrá
\[
  \overline F\left(\varprojlim_\alpha\widetilde X^{(\alpha)}\right)
  =\varprojlim_{\alpha,\mathcal C'}
  \overline F\left(\widetilde X^{(\alpha)}\right).
\]
En particular, si todo proobjeto de $\mathcal C$ es estricto (cf. el número
siguiente), el funtor prolongado $\overline F$ conmuta con los límites
proyectivos.

\subsection*{3. Caracterización de los funtores pro-representables}
\addcontentsline{toc}{subsection}{3. Caracterización de los funtores pro-representables}

Sean $\mathcal C$, $\mathcal C'$ dos categorías en las que existen los límites
proyectivos finitos (es decir, relativos a conjuntos preordenados finitos, no
necesariamente filtrantes), o, lo que es equivalente, en las que existen los
productos finitos y los productos fibrados finitos (lo cual implica, en
particular, la existencia de un «objeto final» $e$, tal que
$\operatorname{Hom}(X,e)$ se reduzca a un elemento para todo $X$). Sea $F$ un
funtor covariante de $\mathcal C$ en $\mathcal C'$. Las condiciones siguientes
son equivalentes:

\begin{enumerate}
\item[(i)] $F$ conmuta con los límites proyectivos finitos;
\item[(ii)] $F$ conmuta con los productos finitos y los productos fibrados finitos;
\item[(iii)] $F$ conmuta con los productos finitos y, para todo diagrama exacto
\[
  X\longrightarrow X'\rightrightarrows X''
\]
en $\mathcal C$ ([3], A, definición 2.1), el diagrama transformado por $F$
\[
  F(X)\longrightarrow F(X')\rightrightarrows F(X'')
\]
es exacto.
\end{enumerate}

Se dice entonces que $F$ es \emph{exacto por la izquierda}.

En lo sucesivo se supone que en $\mathcal C$ existen los límites proyectivos
finitos. De las definiciones resulta entonces inmediato que un funtor
representable es exacto por la izquierda y, pasando al límite, que un funtor
pro-representable es exacto por la izquierda.

Para obtener un recíproco, sea
\[
  F:\mathcal C\longrightarrow(\operatorname{Ens})
\]
un funtor covariante, y sean $X\in\mathcal C$ y $\xi\in F(X)$. Se dice que
$\xi$ (o el par \src{374}$(X,\xi)$) es \emph{minimal} si, para todo
$X'\in\mathcal C$ y $\xi'\in F(X')$, y todo monomorfismo estricto ([3], A, 2)
$u:X'\longrightarrow X$ tal que $\xi=F(u)(\xi')$, $u$ es un isomorfismo. Por
otra parte, se dice que un par $(X,\xi)$ \emph{domina} a $(X'',\xi'')$
($\xi\in F(X)$, $\xi''\in F(X'')$) si existe un morfismo
$v:X\longrightarrow X''$ tal que $\xi''=F(v)(\xi)$; si $\xi$ es minimal y $F$
es exacto por la izquierda, este morfismo $v$ es entonces único; si $\xi''$ es
minimal, $v$ es sobreyectivo. Se deduce fácilmente la proposición siguiente:

\result{PROPOSICIÓN}{3.1} Para que $F$ sea estrictamente pro-representable, es
necesario y suficiente que satisfaga las dos condiciones:
\begin{enumerate}
\item[(i)] $F$ es exacto por la izquierda;
\item[(ii)] todo par $(X,\xi)$, con $\xi\in F(X)$, está dominado por un par
minimal.
\end{enumerate}

Esta última condición es vacía si todo objeto de $\mathcal C$ es artiniano
(tomando un subobjeto $X'$ de $X$ minimal entre aquellos para los cuales existe
$\xi'\in F(X')$ tal que $\xi$ sea la imagen de $\xi'$). De aquí:

\result{COROLARIO}{} Sea $\mathcal C$ una categoría cuyos objetos son
artinianos y en la que existen los límites proyectivos finitos. Entonces los
funtores pro-representables de $\mathcal C$ en $(\operatorname{Ens})$ son
exactamente los funtores exactos por la izquierda, e incluso son estrictamente
pro-representables.

Este último hecho significa también que \emph{todo proobjeto de $\mathcal C$
es estricto}.

\subsection*{4. Ejemplo: grupos de tipo galoisiano, grupos proalgebraicos}
\addcontentsline{toc}{subsection}{4. Grupos de tipo galoisiano, grupos proalgebraicos}

Si $\mathcal C$ es la categoría de los grupos finitos ordinarios,
$\operatorname{Pro}(\mathcal C)$ es equivalente a la categoría de los grupos
topológicos compactos totalmente disconexos. Son grupos de este tipo y sus
generalizaciones, obtenidas sustituyendo los grupos finitos ordinarios por
esquemas en grupos finitos sobre un preesquema base dado (por ejemplo, los
grupos algebraicos finitos sobre un cuerpo $k$), los que harán las veces de
grupos fundamentales, de homotopía y de homología absolutos y relativos para
los preesquemas. En todos estos ejemplos se aplica el corolario de la
proposición 3.1, y es precisamente mediante el funtor asociado como debe
definirse $\pi_1$ [2]. Lo mismo sucede partiendo de la categoría de los grupos
algebraicos o cuasialgebraicos sobre un cuerpo (o, más generalmente, sobre un
preesquema noetheriano): se obtienen los «grupos proalgebraicos» de Serre [4].

\subsection*{5. Ejemplo: «variedades formales»}
\addcontentsline{toc}{subsection}{5. Ejemplo: variedades formales}

Sean $\Lambda$ un anillo noetheriano y $\mathcal C$ la categoría de las
$\Lambda$-álgebras $A$ que son módulos artinianos de tipo finito sobre
$\Lambda$ (o, más brevemente,

\src{375}
las $\Lambda$-álgebras artinianas). Se cumplen las condiciones del corolario de
la proposición 3.1. La categoría $\operatorname{Pro}(\mathcal C)$ es aquí
equivalente a la categoría de las álgebras topológicas $O$ sobre $\Lambda$
isomorfas a límites proyectivos topológicos
\[
  O=\varprojlim_i O_i
\]
de álgebras $O_i\in\mathcal C$, es decir, cuya topología es lineal, separada y
completa, y tal que para todo ideal abierto $J_i$ de $O$, $O/J_i$ sea un
álgebra artiniana sobre $\Lambda$. El funtor
$\mathcal C\longrightarrow(\operatorname{Ens})$ asociado a tal álgebra no es
otro que
\begin{align*}
  F(A)=h'_O(A)
  &=\{\text{homomorfismos continuos de $\Lambda$-álgebras topológicas }
       O\longrightarrow A\}\\
  &=\varinjlim_i
    \operatorname{Hom}_{\Lambda\text{-álgebras}}(O_i,A).
\end{align*}

Obsérvese, además, que la categoría $\mathcal C$ es esencialmente el producto
de las categorías análogas correspondientes a las completaciones de los
anillos locales $\Lambda_{\mathfrak m}$, para los ideales maximales
$\mathfrak m$ de $\Lambda$; por tanto, si se desea, se puede restringir al caso
en que $\Lambda$ sea un anillo local completo de esta clase. En cualquier caso,
$O$ se descompone canónicamente en el producto topológico de componentes
locales, correspondientes a los «puntos» del esquema formal [2] definido por
$O$. Un punto de esta clase está definido por un elemento $\xi$ de algún
$F(K)$, donde $K\in\mathcal C$ es un cuerpo (por ejemplo, el cuerpo residual
de la componente local considerada); dos pares $(\xi,K)$ y $(\xi',K')$
definen el mismo punto si y solo si están dominados por un mismo
$(\xi'',K'')$, o bien si dominan a un mismo $(\xi''',K''')$. (Si los
$\Lambda/\mathfrak m$ son algebraicamente cerrados, basta tomar la suma
disjunta de los $F(\Lambda/\mathfrak m)$.)

Es importante dar condiciones para que la componente local $O_\xi$ de $O$
correspondiente a un $\xi\in F(K)$ sea un anillo noetheriano. Cuando $\Lambda$
es un anillo local completo (noetheriano, recordémoslo), esto equivale a decir
que $O_\xi$ es isomorfo a un anillo cociente de un anillo de series formales
$\Lambda[[t_1,\ldots,t_n]]$ sobre $\Lambda$. Para dar un criterio de esta clase,
introduzcamos (para todo anillo $A$) la $A$-álgebra $I_A$ de los «números
duales» de $A$ mediante
\[
  I_A=A[t]/t^2A[t].
\]
Sea $\varepsilon:I_A\longrightarrow A$ el homomorfismo de aumento; este define
(si $A\in\mathcal C$) una aplicación
\[
  F(\varepsilon):F(I_A)\longrightarrow F(A),
\]

\src{376}
y, utilizando que $F$ es exacto por la izquierda, se define de manera
intrínseca una estructura de $A$-módulo en la parte
\[
  F(I_A,\xi)=F(\varepsilon)^{-1}(\xi)
\]
de $F(I_A)$ formada por los $\xi'\in F(I_A)$ que «se reducen a $\xi$»;
utilizando la forma explícita de $F$ en términos de $O$, se encuentra que este
$A$-módulo se identifica con
\[
  \operatorname{Hom}_{\Lambda}(\mathfrak m_\xi/\mathfrak m_\xi^2,A),
\]
donde $\mathfrak m_\xi$ es el núcleo del homomorfismo
$\xi:O\longrightarrow A$, es decir, si $A$ es un cuerpo, el ideal maximal de
la componente local $O_\xi$ de $O$. Se deduce inmediatamente la proposición
siguiente:

\result{PROPOSICIÓN}{5.1} Sea $\xi\in F(K)$, donde $K\in\mathcal C$ es un
cuerpo. Para que la componente local correspondiente $O_\xi$ de $O$ sea un
anillo noetheriano, es necesario y suficiente que el conjunto $F(I_K,\xi)$ de
los elementos de $F(I_K)$ que se reducen a $\xi$ sea un espacio vectorial
de dimensión finita sobre $K$. En estas condiciones, se tiene un isomorfismo
canónico
\[
  F(I_K,\xi)=
  \operatorname{Hom}\!\left(
    \frac{\mathfrak m_\xi/\mathfrak m_\xi^2}
         {\operatorname{im}(\mathfrak n_\xi/\mathfrak n_\xi^2)},K
  \right),
\]
donde $\mathfrak n_\xi$ es el ideal maximal de $\Lambda$, núcleo del
homomorfismo $\Lambda\longrightarrow K$.
\footnote{El texto impreso da aquí una fórmula que solo es correcta cuando
$\Lambda$ es un cuerpo. La fe de erratas de 1962 prescribe, en el caso general,
el cociente de $\mathfrak m_\xi/\mathfrak m_\xi^2$ por la imagen de
$\mathfrak n_\xi/\mathfrak n_\xi^2$.}
En particular, la dimensión del $K$-espacio vectorial $F(I_K,\xi)$ es igual a
la dimensión del espacio vectorial así corregido sobre el cuerpo
$O_\xi/\mathfrak m_\xi=K(\xi)$.

Supongamos que $O_\xi$ es noetheriano y, para simplificar la escritura, que
$\Lambda$ es local completo y $O=O_\xi$. Se dice que $O$ es \emph{liso sobre
$\Lambda$} si $O$ es un álgebra finita y étale sobre el álgebra completada del
localizado de $\Lambda[t_1,\ldots,t_n]$ en un ideal maximal de este anillo que
induce el ideal maximal de $\Lambda$; cuando la extensión residual de $O$ sobre
$\Lambda$ es trivial (por ejemplo, si el cuerpo residual de $\Lambda$ es
algebraicamente cerrado), esto equivale a decir que $O$ mismo es isomorfo a un
álgebra de series formales de esta clase.
\footnote{La fe de erratas de 1962 precisa que esta definición de «liso sobre
$\Lambda$» solo es correcta cuando la extensión residual $k'/k$ es separable;
en el caso general, remite a SGA III, 1.1.}
Finalmente, si ya no se supone necesariamente que $O$ sea noetheriano, se dirá
también que $O$ es liso sobre $\Lambda$ si $O$ es isomorfo a un límite
proyectivo topológico de $\Lambda$-álgebras cociente que son noetherianas y
lisas sobre $\Lambda$ en el sentido anterior. Se generaliza inmediatamente al caso
en que $\Lambda$ y $O$ ya no se suponen locales. Dicho esto, se tiene la
proposición siguiente:

\result{PROPOSICIÓN}{5.2} Para que $O$ sea liso sobre $\Lambda$, es necesario
y suficiente que el funtor asociado $F$ transforme epimorfismos en
epimorfismos.

Esto significa que, para todo homomorfismo sobreyectivo
$A\longrightarrow A'$ en $\mathcal C$, $F(A)\longrightarrow F(A')$ es
sobreyectivo. Naturalmente, basta comprobar esta condición cuando $A$ es local
y, procediendo paso a paso, cuando el ideal de $A$ que es núcleo de
$A\longrightarrow A'$ está anulado por el ideal maximal de $A$. En la práctica,
esto reduce la cuestión a comprobar que cierta obstrucción, ligada a invariantes
infinitesimales

\src{377}
de la situación que da origen al funtor $F$, es nula, problema de naturaleza
cohomológica.

Para terminar, digamos unas palabras, en el contexto anterior, sobre los
\emph{anillos de definición}. Sea todavía $F$ un funtor de $\mathcal C$ en
$(\operatorname{Ens})$, pro-representable mediante una $\Lambda$-álgebra
topológica $O$. Entonces, para todo $A\in\mathcal C$ y todo $\xi\in F(A)$,
existe un menor subanillo $A'$ de $A$ tal que $\xi$ sea imagen de un elemento
$\xi'$ de $F(A')$ (que queda entonces determinado de manera única): en efecto,
basta interpretar $\xi$ mediante un homomorfismo de $O$ en $A$ y tomar como
$A'$ la imagen de $O$ por este último. Se dice entonces que $A'$ es el anillo
de definición del objeto $\xi\in F(A)$. Si $u:A\longrightarrow B$ es un
homomorfismo de álgebras y $\eta=F(u)(\xi)$, el anillo de definición de $\eta$
es la imagen por $u$ del anillo de definición de $\xi$. Cuando se parte de un
funtor $F$ de $\mathcal C$ en $(\operatorname{Ens})$, la existencia de los
anillos de definición y las propiedades que acabamos de dar son, con pequeñas
salvedades, equivalentes a que $F$ sea pro-representable; es decir, por lo
general distan mucho de ser triviales.

\fsection{B}{Los dos teoremas de existencia}

Conservemos las notaciones de A, párrafo 5, y partamos de un funtor covariante
\[
  F:\mathcal C\longrightarrow(\operatorname{Ens});
\]
se buscan criterios manejables para que $F$ sea pro-representable, es decir,
expresable mediante una $\Lambda$-álgebra $O$ como arriba. En virtud del
corolario de A, proposición 3.1, para ello es necesario y suficiente que $F$ sea
exacto por la izquierda. En el estado actual de la técnica de descenso (cf. las
cuestiones planteadas en [3], p. 9), este criterio no puede comprobarse
directamente bajo esta forma en los casos más importantes, y se necesitan
criterios aparentemente menos exigentes.

\result{TEOREMA}{1} Para que el funtor $F$ sea pro-representable, es necesario
y suficiente que satisfaga las dos condiciones siguientes:
\begin{enumerate}
\item[(i)] $F$ conmuta con los productos finitos;
\item[(ii)] para toda álgebra $A\in\mathcal C$ y todo homomorfismo
$A\longrightarrow A'$ en $\mathcal C$ tal que el diagrama
\[
  A\longrightarrow A'\rightrightarrows A'\otimes_A A'
\]
sea exacto ([3], A, definición 1.2), el diagrama transformado
\[
  F(A)\longrightarrow F(A')\rightrightarrows F(A'\otimes_A A')
\]

\src{378}
es exacto.
\end{enumerate}
Además, basta comprobar (ii) cuando $A$ es local y se está, además, en uno de
los dos casos siguientes:
\begin{enumerate}
\item[(a)] $A'$ es un módulo libre sobre $A$;
\item[(b)] el módulo cociente $A'/A$ es un $A$-módulo de longitud 1.
\end{enumerate}

La demostración de este teorema es bastante delicada y no puede esbozarse aquí.
Limitémonos a señalar que descansa esencialmente en un estudio de las relaciones
de equivalencia (en el sentido de las categorías) en el espectro de un álgebra
artiniana (estudio que plantea todavía varios problemas cuya solución parece
indispensable para el desarrollo ulterior de la teoría).

En las aplicaciones, la comprobación de la condición (i) es siempre trivial.
La de (ii) se descompone en dos: el caso en que $A'$ es un $A$-módulo libre
corresponde a la técnica de descenso por morfismos planos ([3], teoremas 1, 2 y
3) y no ofrece dificultad. Para tratar el caso (b), se utilizará el resultado
siguiente:

\result{TEOREMA}{2} Sean $A$ un anillo local artiniano de ideal maximal
$\mathfrak m$ y $A'$ una $A$-álgebra que contiene a $A$, tal que
$\mathfrak mA'\subset A$ y que
\[
  A\longrightarrow A'\rightrightarrows A'\otimes_A A'
\]
sea exacto (lo cual ocurre, en particular, si $A'/A$ es un $A$-módulo de
longitud 1). Sea $\mathcal F$ la categoría fibrada ([3], A, definición 1.1) de
los haces cuasicoherentes y planos sobre preesquemas variables. Entonces el
morfismo $\operatorname{Spec}(A')\longrightarrow\operatorname{Spec}(A)$ es un
morfismo de $\mathcal F$-descenso estricto ([3], A, definición 1.7).

En otros términos, el dato de un $A$-módulo plano $M$ es completamente
equivalente al dato de un $A'$-módulo plano $M'$, provisto de un
$A'\otimes_AA'$-isomorfismo de $M'\otimes_AA'$ sobre $A'\otimes_AM'$ que
satisface la condición de transitividad usual para un dato de descenso (loco
citato).

El teorema 2 se demuestra probando primero que
\[
  H^i(A'/A,G_a)=0\qquad\text{para }i\geq 1
\]
([3], A, 4(e)), pues la hipótesis $\mathfrak mA'\subset A$ permite reducirse
fácilmente al caso en que $A$ es un cuerpo $k$ (a saber, $A/\mathfrak m$). Se
aplican entonces las equivalencias señaladas en [3], página 16.

\fsection{C}{Aplicaciones a algunos casos particulares}

\subsection*{1. Observaciones generales sobre los funtores representados por preesquemas}
\addcontentsline{toc}{subsection}{1. Observaciones generales sobre los funtores representados por preesquemas}

Sea $S$ un preesquema localmente noetheriano. Se dice que un preesquema $X$
sobre $S$ es \src{379}localmente de tipo finito sobre $S$ si, para todo
$x\in X$ que se proyecta en $y\in S$,\footnote{El texto impreso dice aquí
« $y\in Y$ » y vuelve a emplear $Y$ al final del párrafo. Como el morfismo
considerado es $X\longrightarrow S$, se lee $S$ en ambas apariciones.} existen
un entorno afín de $y$, de anillo $A$, y un entorno afín de $x$ situado sobre
el anterior, de anillo $B$, tales que $B$ sea un álgebra de tipo finito sobre
$A$. Se encuentran muchos ejemplos importantes de preesquemas localmente de
tipo finito sobre $S$ que no son de tipo finito sobre $S$, como soluciones de
problemas universales clásicos; así, importa poder considerar el esquema de
Picard de una curva como la reunión de una infinidad de componentes conexas
(que no deben confundirse con la componente conexa del elemento neutro, es
decir, la «variedad de Picard»). En ocasiones resulta entonces conveniente
trabajar en la categoría $\mathcal C$ de los preesquemas localmente de tipo
finito sobre $S$, para examinar en ella la cuestión de la representabilidad de
un funtor contravariante $F$. El objetivo principal de estas exposiciones es
desarrollar una técnica general que permita reconocer si tal funtor $F$ es
representable y estudiar las propiedades del $S$-preesquema $X$ correspondiente
a partir de las de $F$. Señalemos de paso que, en este estudio, aparecen
ejemplos no patológicos de preesquemas sobre $S$ no separados sobre $S$, en
particular como «preesquemas de Picard» de excelentes $S$-esquemas; no deben,
por tanto, desterrarse de la geometría algebraica los preesquemas que no son
esquemas.

Sea $X$ un preesquema localmente de tipo finito sobre $S$, y sea
\[
  F(Y)=\operatorname{Hom}_S(Y,X)
\]
el funtor contravariante asociado. Se puede considerar la restricción $F_0$ de
$F$ a la subcategoría $\mathcal C_0$ de $\mathcal C$ formada por los
preesquemas $Y$ sobre $S$ que son artinianos y finitos sobre $S$: cuando
$S=\operatorname{Spec}(\Lambda)$, $\mathcal C_0$ es, pues, la categoría dual
de la categoría de las $\Lambda$-álgebras artinianas considerada en B. Si
$Y=\operatorname{Spec}(A)$, donde $A$ es un anillo local artiniano, $Y$ se
reduce a un solo punto $y$ situado sobre un punto cerrado $s$ de $S$, y un
$S$-homomorfismo de $Y$ en $X$ (es decir, un elemento de $F(Y)$) queda definido
por el dato de un punto $x\in X$ sobre $s$ y de un
$\mathcal O_{s,S}$-homomorfismo de $\mathcal O_{x,X}$ en $A$. Si existe tal
homomorfismo, entonces $x$ es necesariamente un punto cerrado de $X$ (pues su
cuerpo residual es algebraico sobre el de $s$). Esto muestra que la restricción
$F_0$ de $F$ a las «$\Lambda$-álgebras artinianas» es pro-representable, y está
representada por el álgebra topológica sobre $\Lambda$ cuyas componentes
locales son los completados $\widehat{\mathcal O}_{x,X}$ de los anillos locales
de $X$ en los puntos $x$ de $X$ que son cerrados y están situados sobre puntos
cerrados de $S$. Esto muestra que el solo conocimiento de $F_0$ proporciona
información precisa sobre la estructura de $X$ (a saber, la estructura de los
completados de sus anillos locales en los puntos indicados). Obsérvese que,
incluso cuando $S$

\src{380}
es el espectro de un cuerpo algebraicamente cerrado, solo mediante la
consideración sistemática de «variedades» $Y$ tales que $\mathcal O_Y$ pueda
tener elementos nilpotentes, en particular trabajando con los espectros de
anillos artinianos locales, se puede llegar a la «formulación correcta» de los
problemas universales clásicos y captar su aspecto «infinitesimal».

Cuando se parte de un funtor $F$ dado de antemano y se quiere decidir si es
representable, el estudio del funtor $F_0$ (con ayuda de los teoremas 1 y 2)
proporcionará indicaciones casi completas: o bien, como ocurre con frecuencia
(comprobando simplemente, por ejemplo, la naturaleza de los conjuntos
$F(I_K,\xi)$ y su comportamiento funtorial, cf. A), se constata que $F_0$ ya no
es pro-representable (lo que explica el fracaso de los intentos realizados
hasta ahora para definir de manera más o menos natural variedades de módulos
para la clasificación de los fibrados vectoriales de rango $>1$); o bien se
logra comprobar que $F_0$ es representable, pero los espacios vectoriales
$F(I_K,\xi)$ no son de dimensión finita, en cuyo caso hay que conformarse con
la solución «formal»; o bien, por último, $F_0$ está representado por un
producto de anillos locales completos noetherianos, lo que aporta indicios muy
sólidos de que el propio $F$ es representable y, unido a propiedades análogas
pero de naturaleza más global que pensamos desarrollar más adelante, bastará
sin duda para implicar que así sea efectivamente. Finalmente, se encuentran
problemas geométricos interesantes (véanse los párrafos 4 y 5 más abajo) en los
que solo se dispone del funtor $F_0$ (que no procede de ningún funtor «global»
$F$), y en los que habrá motivos para darse por satisfecho si se logra asociarle
un «esquema formal de módulos».

Para terminar estas generalidades, indiquemos cómo la teoría de esquemas
explica anomalías aparentes, como la superficie de Igusa $V$, cuya «variedad
de Picard» $P$ se reduce a un punto y para la cual, sin embargo, se tiene
\[
  H^1(V,\mathcal O_V)\ne 0 ;
\]
en este caso, $P$ es un grupo «puramente infinitesimal» que no se reduce al
grupo unidad, es decir, está definido por un álgebra local $O$ de rango finito
sobre el cuerpo base $k$ y dotado de una aplicación diagonal correspondiente a
la estructura multiplicativa de $P$; si $\mathfrak m$ es el ideal maximal de
$O$, el dual de $\mathfrak m/\mathfrak m^2$ es canónicamente isomorfo a
$H^1(V,\mathcal O_V)$ (cf. el párrafo 3 más abajo). Solo cuando el grupo de
Picard es un grupo algebraico en el sentido clásico, es decir, liso sobre el
cuerpo base $k$, la dimensión de $H^1(V,\mathcal O_V)$ (siempre igual a la de
$\mathfrak m/\mathfrak m^2$) es igual a la del grupo de Picard.

\subsection*{2. Esquemas $\underline{\operatorname{Hom}}_S(X,Y)$,
$\Pi_{X/S}Z$, $\underline{\operatorname{Aut}}_S(X)$, etc.}
\addcontentsline{toc}{subsection}{2. Esquemas Hom, Pi y Aut}

Sean $X,Y$ dos preesquemas sobre $S$; para todo preesquema $T$ sobre $S$,

\src{381}
sean $X_T=X\times_ST$, $Y_T=Y\times_ST$, y consideremos el conjunto
\[
  F(T)=\operatorname{Hom}_T(X_T,Y_T)
      =\operatorname{Hom}_S(X_T,Y)
      =\operatorname{Hom}_S(X\times_ST,Y)
\]
como funtor contravariante en $T$. Si es representable, se denota por
$\underline{\operatorname{Hom}}_S(X,Y)$ el preesquema sobre $S$ que lo
representa; se tiene, por tanto, un isomorfismo funtorial
\[
  \operatorname{Hom}_S
  \bigl(T,\underline{\operatorname{Hom}}_S(X,Y)\bigr)
  \xrightarrow{\sim}
  \operatorname{Hom}_S(T\times_SX,Y).
\]
Hay muchas variantes de este problema universal cuya solución se reduce a él:
el preesquema de los $S$-automorfismos de un $S$-preesquema $X$ (que será un
preesquema en grupos), el preesquema de los $S$-homomorfismos de un
$S$-preesquema en grupos en otro (que será un preesquema en grupos conmutativos
si el segundo esquema en grupos es conmutativo), etc. También se puede
generalizar la definición de $\underline{\operatorname{Hom}}_S(X,Y)$
considerando un preesquema $Z$ sobre el preesquema $X$ sobre $S$, y el funtor
\[
  F(T)=\operatorname{Hom}_{X_T}(X_T,Z_T)
\]
(conjunto de las «secciones» de $Z_T$ sobre $X_T$); si este funtor es
representable, el $S$-preesquema que lo representa se denotará por
$\Pi_{X/S}Z$; se tendrá, pues, por definición, un isomorfismo funtorial
\[
  \operatorname{Hom}_S(T,\Pi_{X/S}Z)
  =\operatorname{Hom}_{X_T}(X_T,Z_T).
\]
Tomando $Z=Y\times_SX$, se recupera
$\underline{\operatorname{Hom}}_S(X,Y)$. De estas definiciones resulta para los
nuevos preesquemas así introducidos un formulario tan trivial como útil, que no
daremos aquí (puesto que es válido en toda categoría en la que existen
productos y productos fibrados). Más seria es la cuestión de la existencia de
esquemas del tipo $\underline{\operatorname{Hom}}_S(X,Y)$. Se constata primero
que, para $X$ fijo, $\underline{\operatorname{Hom}}_S(X,Y)$ solo puede existir
para todo $Y$ sobre $S$ (o $\Pi_{X/S}Z$ para todo $Z$ sobre $X$) si $X$ es plano
sobre $S$. Además, se advierte que solo es razonable esperar la existencia de
una solución, para $Y$ suficientemente general, si $X$ es además propio sobre
$S$. Por otra parte, parece que estas condiciones son suficientes para la
existencia de $\underline{\operatorname{Hom}}_S(X,Y)$ y $\Pi_{X/S}Z$, a
condición de imponer, cuando proceda, una hipótesis del tipo «cuasiproyectivo»
sobre $Y/S$, respectivamente sobre $Z/X$; esto puede comprobarse, en cualquier
caso, en bastantes situaciones (por ejemplo, cuando $Y$ es afín sobre $S$; o,
mediante construcciones directas elementales, cuando $X$ es finito sobre $S$).
He aquí lo que dan los teoremas 1 y 2:

\src{382}
\result{PROPOSICIÓN}{2.1} Sean $\Lambda$ un anillo noetheriano, $X$ e $Y$ dos
preesquemas cualesquiera sobre $\Lambda$, y consideremos el funtor
\[
  F(A)=\operatorname{Hom}_A(X_A,Y_A)
\]
en la categoría $\mathcal C_0$ de las $\Lambda$-álgebras artinianas. Si $X$ es
plano sobre $\Lambda$, este funtor es pro-representable.

Además, se comprueba que para todo $A\in\mathcal C_0$ y $\xi\in F(A)$ se tiene
un isomorfismo canónico
\[
  F(I_A,\xi)=H^0\!\left(
    X_A,
    \underline{\operatorname{Hom}}_{\mathcal O_{X_A}}
    \bigl(\xi^*(\Omega^1_{Y_A/A}),\mathcal O_{X_A}\bigr)
  \right),
\]
donde $\Omega^1_{Y_A/A}$ es el haz de $1$-diferenciales de Kähler de $Y_A$
respecto de $A$. Tomando como $A$ un cuerpo, se obtiene, utilizando A.5.1 y el
teorema de finitud de [2], el corolario siguiente:

\result{COROLARIO}{} Supongamos que $X$ es plano y propio sobre $S$, y que $Y$
es localmente de tipo finito sobre $S$. Entonces $F$ es pro-representable y las
componentes locales del álgebra topológica correspondiente sobre $\Lambda$ son
anillos noetherianos.

\result{OBSERVACIONES}{} Los problemas considerados en este apartado, y muchos
otros, se estudiaban habitualmente, en el marco de la geometría algebraica
clásica, mediante las «coordenadas de Chow» de los ciclos en el espacio
proyectivo, que permitían considerar dichos ciclos como puntos de variedades
proyectivas adecuadas. Este procedimiento, y en general el uso de las
coordenadas de Chow, parece irremediablemente insuficiente desde el punto de
vista de los esquemas, pues destruye los elementos nilpotentes en las
variedades de parámetros y, en particular, no se presta a un estudio
satisfactorio de las variaciones infinitesimales de ciclos (por no hablar de su
naturaleza no intrínseca, ligada al espacio proyectivo). El lenguaje de las
coordenadas de Chow era, por desgracia, el único utilizado por muchos
geómetras algebraicos para estudiar familias de variedades o familias de
ciclos, lo que parece haber constituido un obstáculo serio para aclarar estas
nociones, pese a su indudable interés técnico (probablemente provisional). Si
se quiere obtener el análogo de las variedades de Chow en la teoría de
esquemas, se llega al problema universal siguiente: sea $X$ un preesquema sobre
$S$; para todo preesquema $T$ sobre $S$, se considera el conjunto $F(T)$ de los
subpreesquemas cerrados de $X_T=X\times_ST$ que son planos sobre $T$. Se intenta
representar este funtor en $T$ mediante un preesquema sobre $S$. Más
generalmente, se puede dar un haz cuasicoherente $G$ sobre $X$ y tomar como
$F(T)$ el conjunto de los haces

\src{383}
cuociente de $G_T$ que son planos sobre $T$. Parece que el problema tiene una
solución, mediante un esquema $C$ localmente de tipo finito sobre $S$, cuando
$X$ es propio sobre el preesquema localmente noetheriano $S$ y $G$ es además
coherente. En cualquier caso, suponiendo únicamente que $S$ es localmente
noetheriano, la restricción de $F$ a las «$S$-álgebras artinianas» es
pro-representable; y si, además, $X$ es propio sobre $S$ y $G$ es coherente, las
componentes locales del anillo topológico correspondiente $O$ son
noetherianas. Por supuesto, aun después de demostrar la existencia del
«esquema de Chow» de $X$ sobre $S$, quedará por
encontrar una descomposición en abiertos disjuntos $C_i$ (correspondientes a
la fijación de invariantes continuos, como el grado y la dimensión de los
ciclos que se hacen variar) que sean de tipo finito (y no solo localmente de
tipo finito) sobre $S$, precisar las relaciones de este esquema con las
variedades de Chow clásicas y precisar cuándo un $C_i$ es incluso proyectivo,
o al menos cuasiproyectivo, sobre $S$.\footnote{La fe de erratas de 1962
precisa que el caso proyectivo queda completamente resuelto por los esquemas de
Hilbert (TDTE IV), mientras que ejemplos de Nagata y Hironaka muestran que, sin
una hipótesis de proyectividad, estos funtores pueden no ser representables,
incluso para las subvariedades de dimensión cero de una variedad completa no
singular de dimensión tres; este fenómeno está ligado a la posible
inexistencia del cuadrado simétrico.}

\subsection*{3. Esquemas de Picard}
\addcontentsline{toc}{subsection}{3. Esquemas de Picard}

Sea $f:X\longrightarrow S$ un $S$-preesquema,\footnote{Para un estudio más
completo, la fe de erratas de 1962 remite a TDTE V.} consideremos el haz
multiplicativo $\mathcal O_X^*$ de las unidades del haz estructural de $X$, y
el grupo
\[
  P(X/S)=H^0\bigl(S,R^1f_*(\mathcal O_X^*)\bigr),
\]
llamado grupo de Picard relativo de $X/S$. Un elemento de este grupo queda,
pues, definido dando una cubierta abierta $(U_i)$ de $S$ y un haz invertible
$L_i$ sobre cada $f^{-1}(U_i)$, de modo que, para todo par $(i,j)$,
$L_i|f^{-1}(U_i\cap U_j)$ sea isomorfo a $L_j|f^{-1}(U_i\cap U_j)$, al menos
localmente sobre $U_i\cap U_j$ (es decir, estos dos haces sean «equivalentes»
en el sentido de [3], B, 4). Si $X/S$ admite una sección, entonces $P(X/S)$ no
es sino el conjunto de las clases de haces invertibles sobre $X/S$ módulo
«equivalencia» (loco citato). Pongamos ahora, para todo $T$ sobre $S$,
\[
  F(T)=P(X_T/T) ;
\]
se obtiene un funtor contravariante en $T$, que podrá llamarse funtor de Picard de
$X/S$; si este funtor es representable, el preesquema sobre $S$ que lo
representa se llamará preesquema de Picard de $X/S$ y se denotará por
$\underline P(X/S)$. Se tendrá, pues, un isomorfismo de funtores
\[
  \operatorname{Hom}_S\bigl(T,\underline P(X/S)\bigr)
  \xrightarrow{\sim}P(X_T/T).
\]
La formación de preesquemas de Picard es compatible con la extensión de la
base; en particular, los preesquemas de Picard de las fibras de $X$ sobre $S$
(que son preesquemas sobre los cuerpos residuales $K(s)$ de los $s\in S$) son
las fibras de $\underline P(X/S)$. Por supuesto, como $P(X_T/T)=F(T)$ es un
grupo conmutativo, los preesquemas de Picard son preesquemas en grupos.
Obsérvese también que las jacobianas generalizadas de Rosenlicht no son sino
las componentes conexas de la unidad en los esquemas de Picard de curvas
completas que pueden tener singularidades, lo que debería hacer evidentes la
mayoría de sus propiedades (una vez demostrada la existencia).

\result{OBSERVACIÓN}{} La definición adoptada aquí solo es razonable cuando
todo punto de $S$\footnote{El texto impreso dice «$Y$»; aquí no se ha definido
ninguna base $Y$, y la condición se refiere a la base $S$.} admite un entorno
abierto $U$ sobre el cual $X$ admite una sección. En el caso general, hay que
modificar ligeramente la definición del funtor de Picard para poder obtener
también un teorema de existencia.

Aquí, las condiciones plausibles para la existencia de un preesquema de Picard
son las siguientes: $X$ es propio y plano sobre $S$,
$f_*(\mathcal O_X)=\mathcal O_S$, y $X$ admite localmente una sección sobre
$S$. Esta condición aparece de manera natural al aplicar la técnica de
descenso, para eliminar los automorfismos de un haz invertible $L$ sobre $X$
dotándolo de una sección marcada sobre la sección $s$ ([3], B, 4). Se obtiene,
en particular:

\src{384}
\result{PROPOSICIÓN}{3.1} Supongamos que $X$ es plano sobre
$S=\operatorname{Spec}(\Lambda)$, con $\Lambda$ noetheriano, y que, para todo
$T$ de tipo finito sobre $S$, se tiene
\[
  f_{T*}(\mathcal O_{X_T})=\mathcal O_T
\]
(si $f$ es propio y separable,\footnote{La fe de erratas de 1962 prescribe
explícitamente la hipótesis «$f$ propio y separable».} o si $S$ es el espectro
de un cuerpo, resulta de Künneth que la última condición equivale a
$f_*(\mathcal O_X)=\mathcal O_S$). Entonces el funtor de Picard de $X/S$ en la
categoría de las $\Lambda$-álgebras artinianas es pro-representable.

Además, aquí se tendrá
\[
  F(I_A,\xi)=H^1(X_A,\mathcal O_{X_A}),
\]
en particular:

\result{COROLARIO}{} Si $X$ es propio sobre $S$, entonces las componentes
locales de la $\Lambda$-álgebra topológica $O$ correspondiente al funtor de
Picard son noetherianas.

\result{OBSERVACIONES}{} Se pueden generalizar las definiciones y los resultados
de este apartado a la clasificación de los fibrados principales sobre $X$, con
grupo estructural un esquema en grupos $G$ sobre $S$ que sea afín, plano sobre
$S$ y conmutativo. Cuando $G$ no es conmutativo, y por tanto el fibrado en
grupos adjunto de un fibrado principal (cuyas secciones son los automorfismos
del fibrado principal) deja de ser trivial, la proposición 3.1 ya no es válida
tal cual. No obstante, se puede modificar el problema universal de manera que
siga teniendo solución (al menos, por ahora, en geometría formal). La regla de
oro que debe recordarse, en el contexto de este \src{385}apartado y de los
siguientes, y siempre que se busquen «esquemas de módulos» para clases de
objetos definidos solo salvo isomorfismo, es siempre la siguiente: eliminar
los posibles automorfismos de los objetos que se quieren clasificar mediante
la introducción, si es necesario, de estructuras auxiliares (puntos o
elementos de secciones marcadas, fijación de formas diferenciales, etc.), que
se elegirán lo bastante inocuas como para no modificar sustancialmente el
problema inicial.

\subsection*{4. Módulos formales de una variedad}
\addcontentsline{toc}{subsection}{4. Módulos formales de una variedad}

Sea $\Lambda$ un anillo local noetheriano de cuerpo residual $k$ (con mayor
frecuencia, $\Lambda$ será igual a $k$, o a un $p$-anillo de Cohen), y sea
$X_0$ un preesquema sobre $k$. Para toda $\Lambda$-álgebra artiniana local $A$,
se considera el conjunto $F(A)$ de las clases, salvo isomorfismo, de
$A$-preesquemas $X$ planos sobre $A$, dotados de un isomorfismo
\[
  (*)\qquad
  X\otimes_A\kappa(A)\xleftarrow{\sim}
  X_0\otimes_k\kappa(A),
\]
donde $\kappa(A)$ es el cuerpo residual de $A$; por supuesto, los isomorfismos
entre tales $A$-preesquemas planos deben respetar el isomorfismo precedente
incluido en la estructura. Si $A$ es una $\Lambda$-álgebra artiniana no
necesariamente local, de componentes locales $A_i$, se toma como $F(A)$ el
producto de los $F(A_i)$. De este modo, $F$ se convierte en un funtor
multiplicativo en $A$, que podrá llamarse funtor de módulos para $X_0$ (y
$\Lambda$). Si este funtor es representable, le corresponde una
$\Lambda$-álgebra topológica $O$ local, de cuerpo residual $k$, y el espectro
formal de $O$ se llamará esquema formal de módulos para $X_0$ (y $\Lambda$);
cf. [2] para algunos detalles sobre esta noción.

Aquí, para aplicar la técnica de descenso, los automorfismos «finitos» de $X_0$
son inofensivos, pues no influyen en la existencia de automorfismos (en el
sentido precisado más arriba) de los $A$-preesquemas $X$; cuando $A$ no se
reduce a un cuerpo, la condición necesaria y suficiente para que $X$ no tenga
ningún $A$-automorfismo no trivial es que se tenga
\[
  H^0(X_0,\Theta_{X_0/k})=0,
\]
donde $\Theta_{X_0/k}$ denota el haz de las $k$-derivaciones (= haz tangente)
de $X_0$. Por otra parte, se obtiene fácilmente (al menos si $X_0$ es liso
sobre $k$)
\[
  F(I_A,\xi)=H^1(X_A,\Theta_{X_A/A}).
\]
Se concluye entonces, como de costumbre:

\src{386}
\result{PROPOSICIÓN}{4.1} Supongamos que
$H^0(X_0,\Theta_{X_0/k})=0$. Entonces existe el esquema formal de módulos para
$X_0$. Si, además, $X_0$ es propio sobre $k$, el esquema formal de módulos es
noetheriano.

\result{OBSERVACIONES}{} 1\textsuperscript{o} Cuando no se supone que $X_0$ sea
liso sobre $k$, $F(I_A,\xi)$ se identifica con un sub-$A$-módulo de
\[
  \operatorname{Ext}^1_{\mathcal O_{P_A}}
  \bigl(P_A;\mathcal I_{X_A},\mathcal O_{X_A}\bigr),
\]
donde se pone $P_A=X_A\times_\Lambda X_A$, donde $\mathcal O_{X_A}$ se
considera como un haz coherente sobre $P_A$ mediante el morfismo diagonal
$X_A\longrightarrow P_A$, y donde $\mathcal I_{X_A}$ denota el haz coherente
de ideales sobre $P_A$ definido por el morfismo diagonal. Más precisamente,
una globalización sencilla de la teoría de Hochschild muestra que el
$\operatorname{Ext}^1$ anterior se identifica con el conjunto de las clases,
salvo isomorfismo, de haces de $I_A$-álgebras planas $\mathcal O$ sobre $X_A$,
dotados de un isomorfismo de aumentación
\[
  \mathcal O\otimes_{I_A}A\longrightarrow\mathcal O_{X_A}
\]
(se recuerda que se pone $I_A=A[t]/(t^2)$). El submódulo $F(I_A,\xi)$ es el
que corresponde a los haces de álgebras conmutativas. Las hipótesis de lisura
no son, pues, esenciales en la teoría de módulos, contrariamente a lo que [2]
daba a entender.

2\textsuperscript{o} Recordemos (loco citato) que, en particular, toda curva
algebraica $X_0$ lisa y propia sobre $k$ admite un esquema formal de módulos,
liso sobre $\Lambda$, de dimensión relativa $3g-3$ si el género $g$ es
$\geq2$, y $g$ si $g=0$ o $1$. Estos dos últimos casos ya no entran directamente
en la proposición 4.1. Sin embargo, en el caso de las curvas elípticas
($g=1$), podrían reducirse a ella mediante las consideraciones que siguen.

Por supuesto, se puede variar ad libitum la proposición 4.1 considerando
sistemas de esquemas sobre $k$ dotados de diversas estructuras. Supongamos,
por ejemplo, que $X_0$ sea un esquema abeliano sobre $k$, con punto origen
marcado (es decir, que $X_0$ se considere como esquema en grupos sobre $k$), y
sea $F(A)$ el conjunto de las clases, salvo isomorfismo, de esquemas abelianos
sobre $A$ (es decir, de esquemas en grupos propios y lisos sobre $A$) dotados
de un isomorfismo de esquemas abelianos $(*)$. Se comprueba que la imposición
de una estructura multiplicativa (e incluso de una mera «sección unidad»)
elimina los automorfismos infinitesimales y que, por consiguiente, existe un
esquema formal de módulos, correspondiente a un anillo local completo
noetheriano $O$. Por otra parte, puede demostrarse que, si $X$ es un esquema
propio y liso con fibras «absolutamente conexas» sobre un preesquema localmente
noetheriano $S$, toda estructura multiplicativa sobre $X$ que admita una
sección unidad es necesariamente asociativa y conmutativa (siempre que lo sea,
al menos, sobre una fibra y que $S$ sea conexo), y además queda determinada de
manera única por el conocimiento

\src{387}
de la sección unidad. Además, supongamos que $S$ es el espectro de un anillo
local artiniano $A$ de cuerpo residual $k$, que $X$ es propio sobre $A$ y está
dotado de una sección $s$, y que, por último, $X\otimes_Ak$ está dotado de una
estructura de esquema abeliano sobre $k$ cuyo elemento neutro es el punto de
$X\otimes_Ak$ correspondiente a $s$; un cálculo sencillo de obstrucciones,
unido a un razonamiento debido a Serre, permite demostrar que existe
efectivamente sobre $X$ una estructura multiplicativa que admite la sección
$s$ como sección unidad. (A partir de ahí, utilizando el «teorema de
existencia» de [2] para pasar al caso en que $A$ es noetheriano local completo,
y luego la técnica de descenso de [3] para el caso general, se puede demostrar
el enunciado análogo para todo $S$ localmente noetheriano conexo.) Esto prueba
que el funtor $F(A)$ considerado aquí es isomorfo al funtor análogo definido al
principio de este apartado prescindiendo de la estructura multiplicativa sobre
$X_0$. Se sigue, en particular, que, si $\mathfrak m$ es el ideal maximal de
$O$, $\mathfrak m/\mathfrak m^2$ es canónicamente isomorfo al dual de
$H^1(X_0,\Theta_{X_0/k})$ y tiene, por tanto, dimensión $n^2$, donde
$n=\dim X_0$. Sería muy interesante determinar si $O$ es liso sobre $\Lambda$,
es decir, isomorfo a un álgebra de series formales en $n^2$ indeterminadas
sobre $\Lambda$. El párrafo 1, proposición 5.2, permite dar una formulación
equivalente de este problema como problema de existencia de esquemas abelianos
que se reduzcan según un esquema abeliano dado. En cualquier caso, por métodos
trascendentes se ve que la respuesta es afirmativa si $k$ es de característica
cero. En característica $p\ne0$, basta evidentemente limitarse al caso en que
$\Lambda$ es el anillo de vectores de Witt construido sobre un cuerpo
algebraicamente cerrado $k$. Tal vez haya llegado el momento de que el «funtor
de Greenberg» demuestre su valía\ldots\footnote{La fe de erratas de 1962
anuncia que el esquema formal de módulos de una variedad abeliana sobre un
cuerpo es liso sobre el anillo de Witt: todo esquema abeliano sobre un cociente
artiniano local se levanta al anillo del cual es cociente. La demostración
utiliza la varianza de la clase de obstrucción al levantamiento introducida en
la exposición 182, p. 12. También recuerda las construcciones de Mumford para
las curvas de género $g$ y los esquemas abelianos polarizados (cf. SHMT).}

\subsection*{5. Extensión de los revestimientos}
\addcontentsline{toc}{subsection}{5. Extensión de los revestimientos}

Sean $\mathfrak X$ un preesquema formal noetheriano [2] y $U$ una parte
abierta de $\mathfrak X$, definida localmente por la «no anulación» de una
sección de $\mathcal O_{\mathfrak X}$ que no sea divisor de cero no solo sobre
$\mathfrak X$, sino sobre todo $X_n$,\footnote{Precisión impuesta por la fe de
erratas de 1962.} y que sea, por tanto, lo bastante grande para que toda sección
de $\mathcal O_{\mathfrak X}$ sobre un abierto $V$ que sea nula sobre $U\cap V$
sea nula. Sea $\mathcal J$ un «ideal de definición» para $\mathfrak X$, y sea
\[
  X_n=(\mathfrak X,\mathcal O_{\mathfrak X}/\mathcal J^{n+1}),
\]
que es, por tanto, un preesquema noetheriano ordinario. Si entonces
$\mathfrak X',\mathfrak X''$ son dos revestimientos planos de $\mathfrak X$ (es
decir, dos preesquemas sobre $\mathfrak X$ definidos por haces de álgebras que
son coherentes y localmente libres como haces de módulos), no ramificados sobre
$U$, la aplicación evidente
\[
  \operatorname{Hom}_{\mathfrak X}(\mathfrak X',\mathfrak X'')
  \longrightarrow
  \operatorname{Hom}_{X_0}(X'_0,X''_0)
\]
es inyectiva; en particular, un automorfismo de $\mathfrak X'$ que induce la
identidad sobre $X'_0$ es la identidad. Esto permite aplicar a la situación la
técnica de descenso. Partamos, en particular, de un revestimiento plano $X'_0$
de $X_0$, no ramificado

\src{388}
sobre $U_0$; sea $G(\mathfrak X)$ el conjunto de las clases, salvo isomorfismo
(que induzca la identidad sobre $X'_0$), de revestimientos planos $\mathfrak X'$
de $\mathfrak X$ que inducen $X'_0$ sobre $X_0$ (y son, por tanto,
necesariamente no ramificados sobre $U$). Se define de igual modo $G(V)$ para
toda parte abierta $V$ de $\mathfrak X$, y más generalmente $G(\mathfrak Y)$
para todo preesquema formal $\mathfrak Y$ sobre $\mathfrak X$. Dicho esto, los
resultados de [2] y [3] implican en primer lugar los resultados siguientes:

\noindent a. Cuando $V$ es un abierto variable de $\mathfrak X$, los $G(V)$
forman un haz sobre $\mathfrak X$, sea $G_{\mathfrak X}=G$. La restricción de
este haz a $U$ es el haz constante cuyas fibras se reducen a un elemento.

Más generalmente, la determinación de las fibras de $G_{\mathfrak X}$ es una
cuestión de anillos locales completos; más precisamente:

\noindent b. Para todo $x\in\mathfrak X$, se tiene
\[
  G_x=G\bigl(\operatorname{Spec}(\mathcal O_{\mathfrak X,x})\bigr)
  \subset
  G\bigl(\operatorname{Spec}(\widehat{\mathcal O}_{\mathfrak X,x})\bigr),
\]
(= clases salvo isomorfismo de álgebras finitas y libres $B$ sobre
$\widehat{\mathcal O}_{\mathfrak X,x}$ dotadas de un isomorfismo de
$B\otimes_{\widehat{\mathcal O}_{\mathfrak X,x}}\mathcal O_{X_0,x}$ con
$\widehat{\mathcal O}'_{0x}$, donde $\mathcal O'_0$ es el haz de álgebras sobre
$X_0$ que define $X'_0$).\footnote{La fórmula impresa identifica directamente
$G_x$ con el término completado y omite el sombrero sobre
$\mathcal O'_{0x}$. La inclusión y el completado son los prescritos por la fe
de erratas de 1962.}

\noindent c. Se tiene un isomorfismo canónico
\[
  G_{\mathfrak X}=\varprojlim_nG_{X_n} ;
\]
en otros términos, para todo abierto $V$ de $\mathfrak X$, se tiene
\[
  G(V)=\varprojlim_nG(V_n).
\]

\noindent d. Supongamos que $\mathfrak X$ se obtiene de un esquema ordinario
propio $X$ sobre un anillo local noetheriano completo $\Lambda$ que posee un
ideal de definición $\mathfrak m$, tomando el completado $\mathcal J$-ádico de
$\mathcal O_X$, donde $\mathcal J=\mathfrak m\mathcal O_X$. Entonces
$G(\mathfrak X)$ es canónicamente isomorfo al conjunto de las clases de
revestimientos planos del esquema ordinario $X$ que «se reducen a $X'_0$».

De manera gráfica, puede decirse que (a) y (b) establecen las relaciones
fundamentales entre el aspecto local y el aspecto global del problema, (c) da
las relaciones entre el aspecto «finito» y el aspecto «infinitesimal» y, por
último, (d) recuerda (bajo las condiciones precisadas) la identidad entre el
aspecto «formal» y el aspecto «algebraico».

Supongamos ahora que $\mathfrak X$ está definido sobre un anillo local
noetheriano completo $\Lambda$, con
$\mathcal J=\mathfrak m\mathcal O_{\mathfrak X}$ (por tanto, $X_0$ es un
preesquema sobre $\Lambda/\mathfrak m$). Para toda álgebra $A$ finita sobre
$\Lambda$, pongamos
\[
  F(A)=G(\mathfrak X\times_\Lambda A),
\]
donde el marcado se toma después del cambio de base al álgebra reducida $A_{\mathrm{red}}$, y no literalmente a $A/\mathfrak mA$. Es un funtor covariante en $A$, con valores en la categoría de los conjuntos,
y por (c) este funtor queda completamente determinado si se lo conoce para $A$
artiniana; y decir que este funtor es pro-representable, es decir, de la forma
\[
  F(A)=\operatorname{Hom}_{\Lambda\text{-algèbres top.}}(O,A),
\]

\src{389}
donde $O$ es un álgebra topológica sobre $\Lambda$ del tipo considerado en A,
párrafo 5, equivale a decir que lo mismo ocurre al restringirse a los $A$
artinianos. La conjunción de los teoremas 1 y 2 implica entonces efectivamente:

\result{PROPOSICIÓN}{5.1} Siempre que toda familia artiniana clasificada por el funtor sea plana sobre su álgebra de coeficientes y que el abierto escogido permanezca esquemáticamente denso después de los cambios de base pertinentes, el funtor precedente es pro-representable.

Por supuesto, según (a), si $U=\mathfrak X$, $G(\mathfrak Y)$ se reduce a un
solo elemento para todo $\mathfrak Y$ sobre $\mathfrak X$, y el funtor $F$
tiene entonces escaso interés (se tendrá $O=\Lambda$). Parece que, en casi
cualquier otro caso, el anillo local topológico $O$ no es noetheriano. No
obstante, su existencia pone de manifiesto, de manera sorprendente, la
naturaleza «continua» del conjunto $G(\mathfrak X)$ de las soluciones (que
corresponde intuitivamente al hecho de que hay una elección «continua» en el
modo en que la ramificación se propaga al extender $X'_0$). Compárese este
resultado con el punto de vista de J.-P. Serre [4] en la teoría del cuerpo de
clases local, que pone asimismo de relieve el carácter continuo del grupo de
Galois topológico de la extensión abeliana maximal de un cuerpo local
«geométrico»; allí, el grupo dual (en el sentido de Pontryagin) aparece como un
límite directo de grupos algebraicos (o al menos cuasialgebraicos); también
aquí, la clasificación de las extensiones viene dada por «variedades» de
dimensión infinita. Por otra parte, en lo anterior se puede tomar como
$\mathfrak X$ el espectro formal de un anillo local completo (del cual
$\Lambda$ será, por ejemplo, un subanillo de Cohen), y cabe esperar que los
resultados de este apartado puedan utilizarse en el estudio de las extensiones
de un anillo local completo de dimensión $>1$. Tanto en el caso local como en
el global, quizá permitan formular relaciones precisas entre los fenómenos de
ramificación superior y fenómenos en característica cero (abordables por
métodos trascendentes). En cualquier caso, el análisis preliminar de la
proposición 5.1 es el que permite extender al caso «moderadamente ramificado»
los métodos expuestos en [2] para estudiar el grupo fundamental y resolver por
métodos trascendentes el «problema de los tres puntos».

Para terminar, señalemos que la situación se simplifica cuando $X_0$ tiene dimensión 1: entonces, en virtud de (a), complementado por el pegado formal en los esquemas locales completados perforados,
\[
  G(\mathfrak X)\simeq
  \prod_iG\bigl(\operatorname{Spec}(\widehat{\mathcal O}_{\mathfrak X,x_i})\bigr),
\]
donde los $x_i$ son los puntos de $X_0-U$: pueden darse arbitrariamente las
extensiones «locales» en los puntos de ramificación. Además, la normalidad de $X_0$ por sí sola no implica que el esquema formal de módulos de 5.1 sea liso sobre $\operatorname{Spec}(\Lambda)$; la última afirmación impresa es falsa en general.

\src{390}
\begin{center}
  \textbf{BIBLIOGRAFÍA}
\end{center}

\begin{description}
  \item[{[1]}] Grothendieck (Alexander). --- \emph{Sur quelques points
    d'algèbre homologique}, Tohoku Math. J., t. 9, 1957, p. 119--221.
  \item[{[2]}] Grothendieck (Alexander). --- \emph{Géométrie formelle et
    géométrie algébrique}, Séminaire Bourbaki, t. 11, 1958/59, fasc. 3,
    n\textsuperscript{o} 182.
  \item[{[3]}] Grothendieck (Alexander). --- \emph{Technique de descente et
    théorèmes d'existence en géométrie algébrique, I : Généralités, Descente
    par morphismes fidèlement plats}, Séminaire Bourbaki, t. 12, 1959/60,
    fasc. 1, n\textsuperscript{o} 190.
  \item[{[4]}] Serre (Jean-Pierre). --- \emph{Corps locaux et isogénies},
    Séminaire Bourbaki, t. 11, 1958/59, fasc. 3, n\textsuperscript{o} 185.
\end{description}

\def\FgaUnitKey{e195}\clearpage
\chapter*{Fe de erratas de la exposición n.\textsuperscript{o} 195}
\addcontentsline{toc}{chapter}{Fe de erratas de la exposición n.\textsuperscript{o} 195}
\setcounter{footnote}{0}
\src{300}

\noindent\textbf{4. TDTE II: El teorema de existencia en la teoría formal de
los módulos.}

\begin{center}
[Seminario Bourbaki, t. 12, 1959/60, n.\textsuperscript{o} 195, 22 p.]
\end{center}

\noindent\textbf{Página 8.} --- La fórmula escrita en la proposición 5.1 solo es
correcta cuando $\Lambda$ es un cuerpo; en el caso general, hay que sustituir
$\mathfrak m_\xi/\mathfrak m_\xi^2$ por el cociente de este espacio por la
imagen de $\mathfrak n_\xi/\mathfrak n_\xi^2$, donde $\mathfrak n$ es el ideal
maximal de $\Lambda$. Además, la definición dada para que $O$ sea liso sobre
$\Lambda$ solo es correcta cuando la extensión residual $k'/k$ es separable.
En el caso general, cf. SGA III, 1.1.

\noindent\textbf{Página 14, observación.} --- Los problemas planteados aquí
quedan completamente resueltos en el caso proyectivo mediante los «esquemas de
Hilbert» (cf. TDTE IV). Ejemplos de Nagata y Hironaka muestran, en cambio, que
los funtores considerados no son necesariamente representables si no se impone
una hipótesis de proyectividad, incluso limitándose a la clasificación de las
subvariedades de dimensión cero de una variedad completa no singular de
dimensión 3; esto está ligado al hecho de que el cuadrado simétrico de una
variedad tal puede no existir.

\src{301}
\noindent\textbf{Página 15, n.\textsuperscript{o} 3.} --- Para un estudio más
completo, véase TDTE V.

\noindent\textbf{Página 16, proposición 3.1.} --- En lugar de «si $f$ es
propio», léase «si $f$ es propio y separable».

\noindent\textbf{Páginas 18 y 19, observaciones, 2.\textsuperscript{o}} --- He
demostrado recientemente que el esquema formal de módulos de una variedad
abeliana sobre un cuerpo es efectivamente liso sobre el anillo de Witt; en
otros términos, todo esquema abeliano sobre un anillo artiniano local que sea
cociente de otro proviene, por reducción, de un esquema abeliano sobre este
último. La demostración utiliza únicamente las propiedades de varianza de la
clase de obstrucción al levantamiento introducida en la exposición
n.\textsuperscript{o} 182, página 12. Recordemos asimismo que los esquemas de
módulos de las curvas de género $g$ y de los esquemas abelianos polarizados han
sido construidos por Mumford (cf. SHMT).

\noindent\textbf{Página 19, párrafo 5.} --- Hay que suponer que la sección
considerada para definir $U$ sea un no divisor de cero no solo sobre
$\mathfrak X$, sino sobre todo $X_n$.

\noindent\textbf{Página 20, líneas 12 y 13.} --- En lugar de
\[
  G_x=G\bigl(\operatorname{Spec}(\widehat{\mathcal O}_{\mathfrak X,x})\bigr),
\]
léase
\[
  G_x=G\bigl(\operatorname{Spec}(\mathcal O_{\mathfrak X,x})\bigr)
  \subset
  G\bigl(\operatorname{Spec}(\widehat{\mathcal O}_{\mathfrak X,x})\bigr),
\]
y, en lugar de «$\mathcal O'_{0x}$», léase
«$\widehat{\mathcal O}'_{0x}$».

\def\FgaUnitKey{212}\unit{Exposición 212 -- Preesquemas cociente}
\sid{212}
\src{99}
\seminares{13.er año, 1960/61, n.º 212\\Febrero de 1961}
\paperhead{TÉCNICAS DE CONSTRUCCIÓN Y TEOREMAS DE EXISTENCIA EN GEOMETRÍA ALGEBRAICA\\
III. PREESQUEMAS COCIENTE}%
  {por Alexander GROTHENDIECK}

\subsection*{Introducción}
\addcontentsline{toc}{subsection}{Introducción}

Los problemas tratados en la presente exposición difieren de los considerados en
los dos anteriores en que se intenta representar ciertos funtores covariantes,
y no ya contravariantes, de esquemas variables. Sin embargo, el procedimiento
de paso al cociente es esencial en muchos problemas de construcción en
geometría algebraica, incluidos los de las Exposiciones I y II ([1], [2]). Así, la
cuestión de la \emph{efectividad de un dato de descenso} sobre un
$T$-preesquema $X$, con respecto a un morfismo fielmente plano y cuasicompacto
$T\longrightarrow S$, equivale a la cuestión de la existencia de un cociente
de $X$ (que satisfaga las propiedades razonables examinadas más adelante) por
la relación de equivalencia plana en $X$ definida por el dato de descenso; las
cuestiones planteadas en [1], A, \S~2 c se resolverán sin duda al mismo tiempo
que las cuestiones formuladas en el párrafo 2 de la presente exposición. Del mismo
modo, el \emph{esquema de Picard} de un $S$-esquema $X$ (para la definición,
véase [2], C, \S~3) puede definirse de diversas maneras como cociente de otros
esquemas (de divisores positivos, o de inmersiones en un esquema proyectivo)
por relaciones de equivalencia planas, mientras que la definición y la
construcción de esos esquemas auxiliares son técnicamente más sencillas: son
esquemas de la forma $\underline{\operatorname{Hom}}_S(X,Y)$ y variantes
definidas en [2], C, \S~2, cuya construcción será objeto de la exposición siguiente
(bajo hipótesis adecuadas de proyectividad). Combinando, pues, los resultados
de la presente exposición y de la siguiente se obtiene la construcción de los esquemas
de Picard bajo hipótesis adecuadas.

El problema del paso al cociente en la categoría de los preesquemas presenta
aún varias cuestiones no resueltas. La más importante se menciona en el
párrafo 8. Actualmente es el único obstáculo para la construcción de
\emph{esquemas de módulos sobre los enteros para curvas de género arbitrario,
variedades abelianas polarizadas}, etc. Esto muestra que su solución merece
los esfuerzos de los especialistas en grupos algebraicos.

\src{100}
\subsection*{1. Relaciones de equivalencia y relaciones de equivalencia efectivas}
\addcontentsline{toc}{subsection}{1. Relaciones de equivalencia y relaciones de equivalencia efectivas}

Sean $\mathcal C$ una categoría y $X$ un objeto de $\mathcal C$. Un par de
morfismos
\[
  p_1,p_2:R\rightrightarrows X
\]
se llama \emph{par de equivalencia} en $\mathcal C$, con \emph{destino} $X$
y \emph{fuente} $R$, si, para todo objeto $T$ de $\mathcal C$, las dos
aplicaciones correspondientes
\[
  p_1(T),p_2(T):R(T)\rightrightarrows X(T)
\]
(donde, para todo objeto $Y$ de $\mathcal C$, se pone
$Y(T)=\operatorname{Hom}(T,Y)$) definen una aplicación
\[
  R(T)\longrightarrow X(T)\times X(T)
\]
que induce una biyección de $R(T)$ sobre el grafo de una relación de
equivalencia en el conjunto $X(T)$. Existe una relación de equivalencia
evidente entre los pares de equivalencia con destino $X$; una clase de
equivalencia para esta última relación se llama una \emph{relación de
equivalencia de $\mathcal C$ en $X$}, o simplemente una relación de
equivalencia cuando no pueda producirse confusión.

Si existe $X\times X$, dar una relación de equivalencia en $X$ equivale a dar
un subobjeto $R$ de $X\times X$ tal que, para todo objeto $T$ de $\mathcal C$,
el subconjunto de
\[
  (X\times X)(T)=X(T)\times X(T)
\]
correspondiente a $R(T)$ sea el grafo de una relación de equivalencia en
$X(T)$. Denotando por $p_1$ y $p_2$ los morfismos de $R$ en $X$ inducidos por
las proyecciones $\operatorname{pr}_1$ y $\operatorname{pr}_2$, la condición
anterior afirma que $(p_1,p_2)$ es un par de equivalencia. También pueden
expresarse diagramáticamente en $\mathcal C$ (siempre que existan $X\times X$
y el producto fibrado $(R,p_2)\times_X(R,p_1)$) los axiomas conjuntistas de
una relación de equivalencia para los $R(T)$ en los $X(T)$, conforme al
principio general explicado en [2], A, \S~1. No lo necesitaremos.

Siempre que se tenga un par de morfismos $(p_1,p_2)$ con la misma fuente $R$
y el mismo destino $X$, puede definirse el \emph{conúcleo} del par como el
objeto $Y$ de $\mathcal C$ que representa el funtor covariante en $Z$
\[
  \operatorname{Hom}_{p_1,p_2}(X,Z),
\]
el conjunto de morfismos $u$ de $X$ en $Z$ tales que $up_1=up_2$. Si existe
$Y$, queda determinado salvo isomorfismo único. Lo denotamos por
$X/(p_1,p_2)$ o, por abuso de notación, por $X/R$; esta última notación se usa
principalmente cuando
\src{101}$(p_1,p_2)$ es un par de equivalencia: entonces es habitual
identificar en la notación la relación de equivalencia definida por el par con
$R$. Obsérvese que, si $Y$ se considera como cociente de $X$, en realidad solo
depende de la relación de equivalencia definida por el par $(p_1,p_2)$.

Partamos ahora de un morfismo
\[
  f:X\longrightarrow Y,
\]
que permite considerar $X$ como un «objeto sobre $Y$», y supongamos que existe
el producto fibrado
\[
  R(f)=X\times_Y X
\]
Sean $p_1$ y $p_2$ sus dos proyecciones. Entonces $(p_1,p_2)$ es un par de
equivalencia, llamado par \emph{asociado} al morfismo $f$. Define, por tanto,
una relación de equivalencia, llamada relación asociada a $f$.

Un par de morfismos $(p_1,p_2)$ con destino $X$ y fuente común $R$ se llama
\emph{par de equivalencia efectivo} si
\begin{enumerate}
  \item[(i)] existe el conúcleo $X/(p_1,p_2)=Y$;
  \item[(ii)] existe el producto fibrado $X\times_Y X$;
  \item[(iii)] el morfismo $R\longrightarrow X\times_Y X$ cuyas componentes son
  $p_1$ y $p_2$ es un isomorfismo.
\end{enumerate}
Entonces $(p_1,p_2)$ es efectivamente un par de equivalencia. La relación de
equivalencia que define se llama también \emph{relación de equivalencia
efectiva}.

Un morfismo $f:X\longrightarrow Y$ se llama \emph{epimorfismo efectivo} si
\begin{enumerate}
  \item[(i)] existe el producto fibrado $X\times_Y X=R$;
  \item[(ii)] existe el cociente $X/(p_1,p_2)$, donde $p_1$ y $p_2$ son las dos
  proyecciones de $R$ en $X$;
  \item[(iii)] el morfismo $X/(p_1,p_2)\longrightarrow Y$ inducido por $f$ es un
  isomorfismo.
\end{enumerate}
Entonces $f$ es efectivamente un epimorfismo, e incluso un epimorfismo
estricto (véase [1], A, \S~2.3); el recíproco vale si existe el producto
fibrado $X\times_Y X$. El objeto cociente de $X$ definido por el epimorfismo
$f$ se llama también \emph{cociente efectivo} de $X$.

Las definiciones anteriores implican la siguiente «correspondencia de
Galois»:

\src{102}
\result{PROPOSICIÓN}{1.1} Existe una correspondencia biunívoca que conserva
los órdenes naturales entre el conjunto de relaciones de equivalencia
efectivas $R$ en $X$ y el conjunto de cocientes efectivos $Y$ de $X$: a $R$
le corresponde el cociente efectivo $X/R$, y a $Y$ le corresponde la relación
de equivalencia efectiva definida por la proyección canónica
$X\longrightarrow Y$ (que está definida por el producto fibrado
$X\times_Y X$ provisto de sus dos proyecciones).

En categorías muy buenas (conjuntos, haces de conjuntos, etc.), todo cociente
es efectivo y toda relación de equivalencia es efectiva. Esto deja de ser
cierto en categorías como la categoría de preesquemas sobre un preesquema fijo
$S$, incluso cuando $S$ es el espectro de un cuerpo, y aun cuando uno se
restringe a esquemas finitos sobre $S$. Las cuestiones de efectividad, e
incluso (para los preesquemas no finitos sobre $S$) las cuestiones de
existencia de cocientes, resultan casi siempre delicadas.

\subsection*{2. Ejemplo: preesquemas finitos sobre $S$}
\addcontentsline{toc}{subsection}{2. Ejemplo: preesquemas finitos sobre S}

Sea $\mathcal C$ la categoría de los preesquemas finitos sobre $S$, donde se
supone que $S$ es localmente noetheriano. Es equivalente a la opuesta de la
categoría de haces coherentes de álgebras conmutativas sobre $S$, o, cuando
$S$ es afín de anillo $A$, a la opuesta de la categoría de $A$-álgebras
finitas sobre $A$ (es decir, finitamente generadas como $A$-módulos). Se sigue
inmediatamente que existen en $\mathcal C$ los límites proyectivos finitos y
los límites inductivos finitos. Esto es bien conocido (sin hipótesis alguna
de finitud\ldots) para los primeros. Así, el producto fibrado de preesquemas
$X,Y$ sobre $S$ corresponde al producto tensorial $B\otimes_A C$ de las
álgebras correspondientes; el núcleo de dos morfismos
$X\rightrightarrows Y$, definidos por dos homomorfismos de $A$-álgebras
$u,v:C\rightrightarrows B$, corresponde al cociente de $B$ por el ideal
generado por los elementos $u(c)-v(c)$, etc. Para los límites inductivos
finitos basta considerar, por una parte, las sumas finitas, que corresponden
al producto ordinario de $A$-álgebras, y, por otra, los conúcleos de pares de
morfismos $X\rightrightarrows Y$, que corresponden (como se ve de inmediato)
al subanillo de $C$ formado por los elementos sobre los que coinciden los
homomorfismos $u,v:C\rightrightarrows B$ (este subanillo es finito sobre $A$
por la hipótesis noetheriana). Además, usando la hipótesis noetheriana, puede
demostrarse que los límites inductivos finitos, y en particular los cocientes,
así construidos en la categoría $\mathcal C$ de preesquemas finitos sobre $S$
son en realidad cocientes en la categoría de \emph{todos} los preesquemas.

\src{103}
Como dijimos en [1], existen epimorfismos no efectivos en $\mathcal C$ (o
epimorfismos no estrictos\footnote{El texto impreso omite «no» antes de
«estrictos». En la categoría considerada existen productos fibrados, y el
párrafo anterior establece precisamente la equivalencia entre efectividad y
epimorfía estricta.}; es lo mismo, puesto que existen los productos fibrados).
\emph{No sé si las relaciones de equivalencia son siempre efectivas} cuando
no se impone ninguna hipótesis de planitud. En esta dirección solo he obtenido
resultados positivos muy parciales, indispensables para la demostración del
teorema fundamental de la teoría formal de módulos (véase [2], B, teo. 1).
Señalemos que en el problema planteado es fácil reducirse al caso en que $S$
es el espectro de un anillo local artiniano $A$ con cuerpo residual
algebraicamente cerrado. Pero incluso si $A$ es un cuerpo, no se conoce la
respuesta.

También puede considerarse un preesquema $X$ sobre $S$ que ya no se suponga
finito sobre $S$, pero provisto de una relación de equivalencia $R$ en $X$ tal
que $p_1:R\longrightarrow X$ sea un morfismo finito. Entonces $R$ se llama
una relación de equivalencia finita. Supongamos, para simplificar, que $S$ y
$X$ son afines (y, por tanto, $R$ es afín, de modo que la situación se reduce
a una cuestión de álgebra conmutativa pura). \emph{Incluso en este caso se
ignora si existe un cociente $X/R=Y$ y si el morfismo canónico
$X\longrightarrow Y$ es finito}. (El caso más sencillo es aquel en que $S$ es
el espectro de un cuerpo $k$ y $X$ es el espectro de $k[t]$, es decir, la
recta afín.) Por supuesto, si los dos problemas anteriores tienen soluciones
afirmativas, puede concluirse en la situación presente que $R$ es efectiva.
Obsérvese que el problema de la \emph{existencia} de un cociente $Y$ y de la
finitud de $f:X\longrightarrow Y$ se plantea exactamente en los mismos
términos si, en lugar del grafo de una relación de equivalencia en $X$, solo
se tiene un pregrafo de equivalencia en $X$ en el sentido del párrafo 4.

La cuestión del paso al cociente por una relación de equivalencia finita más
o menos arbitraria se presenta al construir preesquemas «pegando» preesquemas
dados $X_i$ a lo largo de ciertos subpreesquemas cerrados; la ley de pegado se
expresa precisamente mediante una relación de equivalencia finita en el
preesquema suma $X$ de los $X_i$. Cabe esperar también que la solución de los
problemas aquí planteados, y de diversas variantes, sea una condición previa
para desarrollar una técnica general de construcciones no proyectivas en la
dirección iniciada en [2].

El único hecho positivo general conocido por el autor es el siguiente:

\result{PROPOSICIÓN}{2.1} Sea $S$ un preesquema localmente noetheriano, sea
$s$ un punto de $S$ y sea $\Omega$ una extensión algebraicamente cerrada de
$k(s)$. Considérese el
\src{104}«funtor fibra» correspondiente $F$, que asigna a todo $S$-esquema
$X$ finito sobre $S$ el conjunto de puntos con valores en $\Omega$ de $X_s$.
Este funtor (trivialmente exacto por la izquierda) es \emph{exacto por la
derecha}, es decir, conmuta con los límites inductivos finitos y, en
particular, con los conúcleos de pares de morfismos.

Aplicando este resultado a todos los «puntos geométricos» de $S$, se deduce
que la categoría «cociente» $\mathcal C'$ de $\mathcal C$, obtenida razonando
«módulo morfismos radiciales sobreyectivos» (es decir, adjuntando formalmente
inversos a los morfismos en cuestión), es una categoría «geométrica», esto es,
satisface las mismas propiedades «de naturaleza finita» que la categoría de
conjuntos. En particular, toda relación de equivalencia en ella es efectiva.
Esto implica que, si $R$ es una relación de equivalencia en un $X$ finito
sobre $S$, entonces el morfismo canónico
\[
  R\longrightarrow X\times_Y X\qquad (\text{donde }Y=X/R)
\]
es \emph{radicial y sobreyectivo} (de hecho, una inmersión cerrada
sobreyectiva, puesto que es un monomorfismo).

\subsection*{3. El caso de un grupo de operadores}
\addcontentsline{toc}{subsection}{3. El caso de un grupo de operadores}

Supongamos de nuevo que $\mathcal C$ es una categoría arbitraria. Sean $G,X$
objetos de $\mathcal C$, y supongamos que $G$ es un grupo de operadores de
$\mathcal C$ sobre el objeto $X$ de $\mathcal C$. Esto significa (véase [2],
A, \S~1) que, para todo objeto $T$ de $\mathcal C$, se han dado una estructura
de grupo en $G(T)$ y una estructura de $G(T)$-conjunto en $X(T)$, de tal modo
que, al variar $T$, las estructuras en cuestión «varían funtorialmente» con
$T$. Si existen los productos $G\times G$ y $G\times X$ en $\mathcal C$, una
estructura tal puede definirse también mediante un par de morfismos
\[
  G\times G\longrightarrow G,
  \qquad \pi:G\times X\longrightarrow X,
\]
sometidos a la condición de que, para todo objeto $T$ de $\mathcal C$, las
leyes de composición correspondientes en los conjuntos $G(T)$ y $X(T)$ hagan
de $G(T)$ un grupo que actúa en $X(T)$. Traducir este axioma en la
conmutatividad de ciertos diagramas en $\mathcal C$ es fácil pero tedioso y,
de hecho, completamente inútil en todos los casos que conozco.

Supongamos que existe $G\times X$ y consideremos los dos morfismos
\[
  p_1,p_2:G\times X\rightrightarrows X
\]
dados por\footnote{El texto impreso da \texttt{pr\_1}. Puesto que los dos
morfismos tienen codominio $X$, el primero de ellos debe ser
\texttt{pr\_2}; se corrige aquí la errata tipográfica.}
\src{105}
\[
  p_1=\operatorname{pr}_2,\qquad p_2=\pi.
\]
Se comprueba de inmediato que $(p_1,p_2)$ es un par de equivalencia si y solo
si, para todo objeto $T$ de $\mathcal C$, la aplicación
\[
  G(T)\times X(T)\simeq(G\times X)(T)
  \longrightarrow X(T)\times X(T),
\]
definida por este par es inyectiva, es decir, si el grupo $G(T)$ actúa
\emph{libremente} en el conjunto $X(T)$, esto es, si $g\in G(T)$,
$x\in X(T)$ y $g\cdot x=x$ implican que $g$ es el elemento identidad de
$G(T)$. Se dice entonces que $G$ \emph{actúa libremente en $X$} (o que $X$ es
un \emph{espacio principal de $\mathcal C$ bajo $G$}). La relación de
equivalencia asociada al par $(p_1,p_2)$ se llama entonces \emph{relación de
equivalencia definida por el grupo $G$ que actúa libremente en $X$}. Si existe
también $X\times X$ y se considera el morfismo
\[
  p:G\times X\longrightarrow X\times X
\]
definido por el par $(p_1,p_2)$, la condición de que $G$ actúe libremente
significa que $p$ es un \emph{monomorfismo}.

Por supuesto, incluso cuando $G$ no actúa libremente en $X$, se desean
criterios para la existencia de un cociente de $X$ por $G$, es decir, del
conúcleo del par anterior $(p_1,p_2)$.

El conúcleo en cuestión se denotará a menudo por $X/G$, o preferiblemente por
$G\backslash X$ si $G$ actúa por la izquierda (la primera notación se reserva
para el caso en que $G$ actúa por la derecha). Obsérvese que, incluso cuando
existe la «imagen» de $G\times X$ por $p$ (definida, por ejemplo, como el
menor subobjeto de $X\times X$ a través del cual se factoriza $p$), digamos
$R$, esta no es casi nunca una relación de equivalencia en $X$. Si se intenta
entonces pasar directamente al cociente por $R$ (o, más precisamente, por el
par de morfismos de $R$ en $X$ inducidos por las dos proyecciones
$\operatorname{pr}_i$), se pierden irremediablemente las características
especiales del par original $(p_1,p_2)$. Por ello es importante hallar una
generalización de la noción de relación de equivalencia que se aplique
directamente al par definido por un grupo de operadores de $\mathcal C$.

\subsection*{4. Prerrelaciones de equivalencia}
\addcontentsline{toc}{subsection}{4. Prerrelaciones de equivalencia}

Recordemos que una categoría en la que todo morfismo es un isomorfismo se
llama un \emph{grupoide}. Por otra parte, una categoría debe definirse como
constituida por dos conjuntos subyacentes $(X,R)$, el conjunto de
\emph{objetos} y el conjunto de \emph{flechas}, provistos de las estructuras
siguientes:
\src{106}
\begin{enumerate}
\item[(i)] Un par de aplicaciones
\[
  p_1,p_2:R\rightrightarrows X
\]
llamadas \emph{aplicación fuente} y \emph{aplicación destino};

\item[(ii)] Una aplicación
\[
  \pi:(R,p_2)\times_X(R,p_1)\longrightarrow R
\]
llamada \emph{aplicación de composición}.
\end{enumerate}
Estos datos deben satisfacer axiomas conocidos, que no repetiremos aquí y que
podrían expresarse mediante la conmutatividad de ciertos diagramas, junto con
la existencia de una aplicación $D:X\to R$ (necesariamente única) que haga
conmutativos otros dos diagramas. La aplicación $D$ envía un objeto a la
flecha identidad correspondiente y satisface
\[
  p_1D=p_2D=\operatorname{id}_X.
\]
Decir que la categoría es un grupoide significa entonces que existe una
aplicación
\[
  s:R\longrightarrow R
\]
(necesariamente única), llamada la \emph{inversión} de $R$, que transforma
cada flecha en una flecha inversa. Esto podría expresarse mediante la
conmutatividad de otros cuatro diagramas formados a partir de $s$, $\Delta$ y
los datos anteriores; los dos primeros son
\[
  p_1s=p_2,\qquad p_2s=p_1.
\]

Recordados estos puntos, las definiciones generales de ([2], A, \S~1)
muestran, en particular, qué se entiende por una estructura de
$\mathcal C$-\emph{categoría}, respectivamente de
$\mathcal C$-\emph{grupoide}, sobre un par de objetos $X,R$ de una categoría
arbitraria $\mathcal C$: por definición, es el dato, para todo objeto $T$ de
$\mathcal C$, de una estructura de categoría, respectivamente de grupoide, en
el sentido conjuntista, cuyo conjunto de objetos es $X(T)$ y cuyo conjunto de
flechas es $R(T)$, «variando funtorialmente» estas estructuras con la variable
$T$. Esto proporciona, por tanto, dos morfismos
\[
  p_1,p_2:R\rightrightarrows X
\]
llamados \emph{morfismo fuente} y \emph{morfismo destino}, y, cuando existe el
producto fibrado requerido, un morfismo
\src{107}
\[
  \pi:(R,p_2)\times_X(R,p_1)\longrightarrow R
\]
llamado \emph{morfismo de composición}. Estos tres datos bastan entonces para
determinar la estructura de categoría sobre $X,R$, y el axioma que se les
impone es el siguiente: para todo $T$, los tres datos correspondientes para
$X(T),R(T)$ definen una estructura de categoría (respectivamente de grupoide)
sobre este par de conjuntos. Cuando proceda, esto puede expresarse mediante
la conmutatividad de ciertos diagramas en los que interviene un morfismo
determinado de manera única
\[
  D:X\longrightarrow R
\]
y, en el caso de los grupoides, un morfismo determinado de manera única
\[
  s:R\longrightarrow R;
\]
los diagramas se escriben exactamente como en el caso «conjuntista». Esta
laboriosa interpretación de los axiomas es, afortunadamente, inútil en la
práctica. El único interés teórico de poder expresar los datos y axiomas en
términos de morfismos e igualdades de morfismos entre ciertos productos
fibrados es el siguiente: si $F:\mathcal C\to\mathcal C'$ es un funtor exacto
por la izquierda (es decir, que conmuta con productos finitos y productos
fibrados), entonces transforma una $\mathcal C$-categoría (respectivamente un
$\mathcal C$-grupoide) en una $\mathcal C'$-categoría (respectivamente un
$\mathcal C'$-grupoide), siempre que existan productos finitos y productos
fibrados en $\mathcal C$.

En la práctica, es importante saber interpretar los morfismos
$p_1,p_2,\pi,D,s$ como \emph{operaciones simpliciales} en un objeto
semisimplicial o simplicial adecuado de $\mathcal C$, al menos cuando existen
productos fibrados en $\mathcal C$. Para fijar la terminología, introduzcamos
la \emph{categoría $\mathcal S$ de los símplices estándar} como la categoría
cuyos objetos son los conjuntos finitos de la forma
\[
  \Delta_n=[0,n],\qquad\text{donde }n\geq 0
\]
(el intervalo de enteros de $0$ a $n$), y cuyos morfismos o flechas son
\emph{aplicaciones arbitrarias} entre esos conjuntos finitos. Obsérvese que la
categoría $\mathcal S$ es equivalente a la categoría de conjuntos finitos
\emph{no vacíos}, con las aplicaciones de conjuntos finitos como morfismos.
En $\mathcal S$ existe claramente el coproducto de una familia finita
\emph{no vacía} de objetos, así como la suma amalgamada de dos objetos bajo un tercero
(la operación dual del producto fibrado). Denotemos por $\mathcal S'$ la
subcategoría de $\mathcal S$ que tiene los mismos objetos pero cuyos morfismos
son las \emph{aplicaciones crecientes} entre los $\Delta_n$. Esta categoría es
equivalente a la categoría de conjuntos finitos no vacíos totalmente
ordenados. En esta categoría el
\src{108}
coproducto de dos objetos nunca existe, y la suma amalgamada de dos objetos $A,B$ bajo
un tercero $C$ no existe en general (tómese, por ejemplo, $C=\Delta_0$,
$A=B=\Delta_1$, con las dos aplicaciones estructurales $u:C\to A$ y
$v:C\to B$ iguales). Sin embargo, existe en ciertos casos; por ejemplo,
\[
  C=\Delta_0,\qquad A=\Delta_m,\qquad B=\Delta_n,\qquad
  u(0)=m,\qquad v(0)=0,
\]
y en este caso
\[
  \Delta_m\amalg_{\Delta_0}\Delta_n=\Delta_{m+n}.
\]

Un \emph{objeto simplicial} (respectivamente, un \emph{objeto
semisimplicial}) en una categoría $\mathcal C$ es, por definición, un funtor
contravariante $K$ de $\mathcal S$ (respectivamente de $\mathcal S'$) en
$\mathcal C$. Por tanto, un objeto simplicial define por restricción un objeto
semisimplicial, pero la primera noción difiere esencialmente de la segunda por
la presencia de \emph{operaciones de simetría} en los $K_n=K(\Delta_n)$,
correspondientes a las imágenes por el funtor $K$ de los elementos del grupo
simétrico de $n+1$ letras (considerado como el grupo de automorfismos de
$\Delta_n$ en $\mathcal S$).

Para cada $n$, sea $\Delta'_n$ (respectivamente $\Delta''_n$) la categoría
finita cuyo conjunto de objetos es $\Delta_n$, definida por la relación de
orden caótica (respectivamente por el orden total natural) en $\Delta_n$ (el
conjunto de flechas es el grafo de esa relación de orden). Está claro que
$\Delta'_n$ (respectivamente $\Delta''_n$) depende funtorialmente del objeto
$\Delta_n$ de $\mathcal S$ (respectivamente de $\mathcal S'$). Si $Z$ es una
categoría, entonces $\operatorname{Hom}(\Delta'_n,Z)$ (respectivamente
$\operatorname{Hom}(\Delta''_n,Z)$), al variar $\Delta_n$, es un funtor de
$\mathcal S$ (respectivamente de $\mathcal S'$) a la categoría de conjuntos,
es decir, un conjunto simplicial (respectivamente un conjunto
semisimplicial), que se dice \emph{asociado a la categoría $Z$}; denotémoslos
por $Z'$ y $Z''$. Hay, además, un homomorfismo natural evidente del conjunto
semisimplicial asociado a $Z'$ hacia $Z''$, y este último es un isomorfismo si
y solo si $Z$ es un grupoide. Entendido esto:

\result{PROPOSICIÓN}{4.1} El funtor $Z\longmapsto Z''$ de la categoría de
categorías a la categoría de conjuntos semisimpliciales es plenamente fiel y
define una equivalencia entre la categoría de las \emph{categorías} y la
categoría de los conjuntos semisimpliciales, es decir, los funtores
contravariantes $K$ de $\mathcal S'$ en $(\mathrm{Set})$ que transforman las
sumas amalgamadas $A\amalg_C B$ del tipo especificado anteriormente en productos
fibrados de conjuntos. Asimismo, el funtor $Z\longmapsto Z'$ de la categoría
de grupoides a la categoría de conjuntos simpliciales es plenamente fiel y
define una equivalencia entre la categoría de los \emph{grupoides} y la
categoría de los
\src{109}
conjuntos simpliciales, es decir, los funtores contravariantes $K$ de
$\mathcal S$ en $(\mathrm{Set})$ que transforman las sumas amalgamadas en productos
fibrados.

Así, las categorías pueden considerarse como conjuntos semisimpliciales
particulares, y los grupoides como conjuntos simpliciales particulares,
siempre que, por supuesto, se razone «salvo isomorfismo», como exige la
interpretación de ciertas estructuras en términos de otras. El procedimiento
usual de reducción al caso conjuntista implica entonces:

\result{COROLARIO}{4.2} El enunciado anterior sigue siendo válido cuando las
categorías, los grupoides y los conjuntos simpliciales se sustituyen por
$\mathcal C$-categorías, $\mathcal C$-grupoides y objetos simpliciales en
$\mathcal C$, siempre que existan productos fibrados en $\mathcal C$.

El objeto semisimplicial $K$ en $\mathcal C$ asociado a una categoría
$(X,R,\ldots)$ en $\mathcal C$ puede hacerse explícito interpretando la
componente $K_n=K(\Delta_n)$ de $K$ como la potencia fibrada $n$-ésima de
$(R,p_1)$ sobre $X$, o, mejor, mediante la fórmula de recurrencia
\[
  K_0=X,\qquad
  K_n=(K_{n-1},p_{n-1}^{(n-1)})\times_X(R,p_1),
\]
donde los $p_i^{(n-1)}$ ($0\leq i\leq n-1$) son las proyecciones naturales de
$K_{n-1}$ en $X$ (que se definen también recursivamente).\footnote{El texto impreso lleva $(n+1)$, $K_0=R$ y $p_n^{(n-1)}$ en estos tres lugares. La indexación del nervio exige $n$, $K_0=X$ y $p_{n-1}^{(n-1)}$, respectivamente.} De este modo,
$p_1,p_2,\pi,D,s$ se interpretan como las operaciones simpliciales
correspondientes respectivamente a los morfismos siguientes de $\mathcal S$:
la cara $0$ de $\Delta_1$, la cara $1$ de $\Delta_1$, la cara $(0,2)$ de
$\Delta_2$, la degeneración $\Delta_1\to\Delta_0$ y la simetría de
$\Delta_1$. Todas las demás operaciones semisimpliciales (respectivamente,
simpliciales) se obtienen formalmente a partir de las cuatro (respectivamente,
cinco) anteriores mediante composición y productos fibrados.

Una \emph{prerrelación de equivalencia} en un objeto $X$ de una categoría se
define ahora como el dato de un grupoide cuyo objeto de objetos (por así
decirlo) es $X$. Un dato tal comprende, por tanto, entre otras cosas, un
objeto $R$ y dos morfismos
\[
  p_1,p_2:R\rightrightarrows X.
\]
Pero obsérvese que estos datos por sí solos no determinan la estructura en
cuestión, a diferencia de lo que ocurre con los pares de equivalencia. En esta
exposición nos interesa esta noción para obtener criterios de posibilidad del
paso al cociente, es decir, para formar el conúcleo del par $(p_1,p_2)$. Por
ello, el enunciado del problema no hace intervenir los datos adicionales
inherentes a un grupoide. En las demostraciones de los resultados que siguen,
\src{110}
usaremos, no obstante, esos datos adicionales y, en particular, las
operaciones simpliciales (incluidas las operaciones de simetría) hasta la
dimensión $3$ (en las que interviene la cuarta potencia fibrada de $R$ sobre
$X$).

Una relación de equivalencia en un objeto $X$ de $\mathcal C$ define una
prerrelación de equivalencia: basta verlo en el caso conjuntista, donde a una
relación de equivalencia en un conjunto $X$ se asocia el grupoide cuyo
conjunto de objetos es $X$ y cuyo conjunto de flechas es el grafo de la
relación de equivalencia.

Un monoide de $\mathcal C$ $G$ que actúa sobre un objeto $X$ de $\mathcal C$
define una $\mathcal C$-categoría cuyos objetos subyacentes son $R=G\times X$
y $X$ (siempre que exista $G\times X$), y esta es un $\mathcal C$-grupoide si
y solo si $G$ es un grupo. De nuevo, basta comprobarlo en el caso conjuntista.
La composición de las flechas $(g,a)$ y $(g',g\cdot a)$ se define entonces por
\[
  (g',g\cdot a)(g,a)=(g'g,a),
\]
es decir, si $a,b\in X$, entonces $\operatorname{Hom}(a,b)$ es por definición
el transportador de $a$ a $b$, y los morfismos se componen mediante la
composición de los elementos de $G$.

\medskip
\noindent\textbf{OBSERVACIÓN. ---} Las dificultades lógicas planteadas por un
enunciado como (4.1) pueden evitarse entendiendo que todos los objetos
considerados pertenecen a un «Universo» fijo (que es a su vez un conjunto).

\fsection{5}{Cociente por una relación de equivalencia finita y plana}

\result{TEOREMA}{5.1} Sea $X=\operatorname{Spec}(B)$ un esquema afín y sea
$R$ una prerrelación de equivalencia en $X$ cuya componente $R_1$ sea afín,
digamos $R_1=\operatorname{Spec}(C)$. Supongamos que la primera proyección
$p_1:R_1\to X$ es finita y localmente libre; en otras palabras, supongamos
que el homomorfismo de anillos correspondiente $p_1^*:B\to C$ hace de $C$ un
$B$-módulo proyectivo finitamente generado. Sea $A$ el subanillo de $B$ que
es el núcleo del par de homomorfismos
\[
  p_1^*,p_2^*:B\rightrightarrows C
\]
(el conjunto de elementos $b$ tales que $p_1^*(b)=p_2^*(b)$). Sea
$Y=\operatorname{Spec}(A)$ y sea $f:X\to Y$ el morfismo definido por la
inclusión de $A$ en $B$. En estas condiciones:
\begin{enumerate}
\item[(i)] $B$ es entero sobre $A$, es decir, $f$ es un morfismo entero.

\item[(ii)] El morfismo $f$ es sobreyectivo y sus fibras son las clases de
equivalencia
\src{111}
$p_2(p_1^{-1}(x))$ en el conjunto $X$ módulo $R$; además, la topología de $Y$
es el cociente de la de $X$.

\item[(iii)] $Y$ es el cociente de $X$ por $R$ en la categoría de los
preesquemas.

\item[(iv)] Si $R$ proviene de una \emph{relación de equivalencia}, entonces
el morfismo $f:X\to Y$ es finito localmente libre (es decir, $B$ es un
$A$-módulo proyectivo finitamente generado), y la relación de equivalencia es
efectiva, esto es, $R_1\to X\times_Y X$ es un isomorfismo.
\end{enumerate}

Este teorema generaliza el conocido teorema relativo a un grupo finito $G$ que
actúa por automorfismos en el anillo $B$ y al anillo de invariantes $A$, y su
demostración es análoga a la demostración conocida. El enunciado (iii) puede
precisarse como sigue:

\result{COROLARIO}{5.2} El morfismo canónico
$R_1\to X\times_Y X$ es \emph{sobreyectivo}.

Sea todavía $R$ una prerrelación de equivalencia «finita y localmente libre»
en $X$, pero sea ahora $X$ un preesquema arbitrario. Supongamos que puede
hallarse un preesquema $Y$ y un morfismo $f:X\to Y$ tales que $fp_1=fp_2$ y,
además, tales que la sucesión de homomorfismos de haces de anillos sobre $Y$
\[
  \mathcal O_Y\longrightarrow f_*(\mathcal O_X)
  \rightrightarrows g_*(\mathcal O_R)
\]
sea exacta (donde $g=fp_i$). El teorema proporciona entonces conclusiones
análogas a (i)--(iv) anteriores; en particular, por (iii), $Y$ es el cociente
de $X$ por $R$ y, por tanto, queda determinado salvo isomorfismo único. En
estas condiciones se dice que la prerrelación de equivalencia $R$ en $X$ es
\emph{admisible}. Con esta definición:

\result{TEOREMA}{5.3} Sea $X$ un preesquema y sea $R$ una prerrelación de
equivalencia en $X$ tal que $p_1:R_1\to X$ sea finito y localmente libre.
Entonces $R$ es admisible si y solo si toda clase de equivalencia conjuntista
$p_2(p_1^{-1}(x))$ en $X$ módulo $R$ está contenida en un abierto afín
(condición que siempre se satisface si todo subconjunto finito de $X$ está
contenido en un abierto afín, por ejemplo si $X$ es cuasiproyectivo sobre un
esquema afín).

En efecto, se demuestra sin dificultad que toda clase de equivalencia módulo
$R$ en $X$ está entonces contenida en un abierto afín \emph{estable bajo
$R$}, y se construye el cociente $Y$ pegando las piezas obtenidas del Teorema
5.1.
\src{112}

\result{COROLARIO}{5.4} Supongamos que se satisface esta condición y que,
además, $R$ proviene de una relación de equivalencia. Entonces esa relación es
efectiva, es decir, $R_1\to X\times_Y X$ es un isomorfismo, y $f:X\to Y$ es
finito y localmente libre.

Se concluye inmediatamente por «descenso»:

\result{COROLARIO}{5.5} Bajo las hipótesis de (5.4), $X$ tiene rango $n$ en
todo punto sobre $Y$ si y solo si $(R_1,p_1)$ tiene rango $n$ en todo punto
sobre $X$. Si $X,R_1$ son $Z$-preesquemas y $p_1,p_2$ son $Z$-morfismos, de
modo que $Y$ es un $Z$-preesquema, entonces $X$ es plano sobre $Z$ si y solo
si $Y$ es plano sobre $Z$.

En resumen:

\result{ESCOLIO}{} El dato de un morfismo finito, localmente libre y
sobreyectivo $f:X\to Y$ de preesquemas equivale al dato de un preesquema $X$
provisto de una relación de equivalencia $R$ tal que $p_1:R\to X$ sea finito
y localmente libre y toda clase $p_2(p_1^{-1}(x))$ esté contenida en un
abierto afín.

\medskip
\noindent\textbf{OBSERVACIONES 5.6.}

\smallskip
\noindent $1^\circ$ No se ha necesitado ninguna hipótesis noetheriana.

\smallskip
\noindent $2^\circ$ Esta noción de paso al cociente contiene, como caso
particular, el «descenso inseparable» de CARTIER, que consiste en determinar
los morfismos finitos localmente libres $f:X\to Y$ tales que
$f_*(\mathcal O_X)$ admita una $p$-base relativa a $\mathcal O_Y$ (donde $X$
es un preesquema dado cuyo haz $\mathcal O_X$ es anulado por el número primo
$p>0$). Obsérvese que el resultado se expresa fácilmente sin ninguna
hipótesis de regularidad sobre los anillos locales y sin suponer que $X$ sea
un esquema algebraico sobre un cuerpo. La teoría de JACOBSON--BOURBAKI se
obtiene tomando por $X$ el espectro de un cuerpo de característica $p$.

\smallskip
\noindent $3^\circ$ GABRIEL había obtenido anteriormente un caso particular
del Teorema (5.3) en la teoría del paso al cociente para grupos conmutativos
finitos sobre un cuerpo $k$. [Compárese (7.4).]

\fsection{6}{Cociente por una relación de equivalencia propia y plana}

\result{TEOREMA}{6.1} Sea $S$ un preesquema localmente noetheriano, sea $X$
un $S$-esquema cuasiproyectivo y sea $R$ una prerrelación de equivalencia en
$X$ tal que:
\src{113}
\begin{enumerate}
\item[(a)] $p_1:R_1\to X$ es propio y plano;
\item[(b)] $R_1\to X\times_S X$ es finito o, equivalentemente, por (a), tiene
fibras finitas (condición que se satisface automáticamente si $R$ proviene de
una relación de equivalencia).
\end{enumerate}
En estas condiciones:
\begin{enumerate}
\item[(i)] Existe $X/R=Y$ y, si $S$ es noetheriano y $R$ proviene de una
relación de equivalencia, es cuasiproyectivo sobre $S$.\footnote{El texto
impreso no contiene la segunda hipótesis. La fe de erratas vinculada de la Exposición
212 dice que la cuasiproyectividad parece haberse demostrado entonces solo
cuando $R$ proviene de una relación de equivalencia.}

\item[(ii)] El morfismo canónico $f:X\to Y$ es sobreyectivo, propio y abierto;
sus fibras son las clases de equivalencia $p_2(p_1^{-1}(x))$ en $X$ módulo
$R$. Así, $Y$ se identifica con el espacio topológico cociente de $X$ por la
relación de equivalencia conjuntista definida por $R$. Finalmente,
$R_1\to X\times_Y X$ es sobreyectivo.

\item[(iii)] Si $R$ proviene de una relación de equivalencia, entonces esa
relación es efectiva, es decir, $R_1\to X\times_Y X$ es un isomorfismo, y
además $f:X\to Y$ es plano (por tanto, fielmente plano).
\end{enumerate}

Para la demostración, uno se reduce a (5.1) mediante cuasisecciones adecuadas
de $X$ para $R$; la demostración es análoga a la construcción de grupos
algebraicos cociente en el Seminario Chevalley.

En resumen:

\result{ESCOLIO}{} Sea $X$ cuasiproyectivo sobre un $S$ localmente
noetheriano. El dato de un morfismo propio, fielmente plano y sobreyectivo
$f:X\to Y$ de $X$ a un $S$-preesquema $Y$ equivale al dato de una relación de
equivalencia $R$ en $X$ tal que $p_1:R\to X$ sea propio y plano.

El mismo método da el resultado siguiente:

\result{TEOREMA}{6.2} Sea $S$ un preesquema noetheriano, sea $X$ un
preesquema de tipo finito sobre $S$ y sea $R$ una prerrelación de equivalencia
en el $S$-preesquema $X$. Supongamos que:
\begin{enumerate}
\item[(a)] $p_1:R_1\to X$ es plano y de tipo finito;
\item[(b)] el morfismo $R_1\to X\times_S X$ es cuasifinito (es decir, tiene
fibras finitas).
\end{enumerate}
Entonces existe un abierto denso $U$ de $X$ que está \emph{saturado} para $R$
y tal que:
\begin{enumerate}
\item[(i)] Si $R_U$ es la prerrelación de equivalencia inducida por $R$ en
$U$, entonces existe $U/R_U$ y es de tipo finito sobre $S$.

\item[(ii)] El morfismo canónico $U\to U/R_U$ es sobreyectivo y abierto; sus
fibras son las clases de equivalencia conjuntistas para $R_U$ (así,
$U/R_U$ es el espacio topológico cociente de $U$ por la relación de
equivalencia conjuntista definida por $R_U$). Finalmente, el morfismo
\[
  (R_U)_1\longrightarrow U\times_{U/R_U}U
\]
es sobreyectivo.
\src{114}

\item[(iii)] Si $R$ proviene de una relación de equivalencia, puede disponerse
que $U\to U/R_U$ sea fielmente plano y que $R_U$ sea efectiva.
\end{enumerate}

Este es un resultado de naturaleza esencialmente «birracional».

\medskip
\noindent\textbf{OBSERVACIONES 6.3.}

\smallskip
\noindent $1^\circ$ No sé si la hipótesis (b) es innecesaria en (6.1) y
(6.2). Al tomar cocientes por grupos, nos obliga en la práctica a
restringirnos al caso en que todos los estabilizadores son grupos finitos.

\smallskip
\noindent $2^\circ$ Puede preguntarse si existen resultados análogos a (6.1)
y (6.2) sin hipótesis de planitud. No conozco ningún contraejemplo en esta
dirección. Por otra parte, incluso conservando la hipótesis de planitud y
restringiéndose a relaciones de equivalencia para las cuales $p_1:R\to X$ es
plano y cuasifinito (pero no finito), con $X$ afín, puede ocurrir que $R$ no
sea efectiva: tómense las relaciones de equivalencia inducidas en los
subconjuntos abiertos afines que recubren la variedad de Nagata (en la que un
grupo de orden dos actúa de manera «no admisible»).

\fsection{7}{Aplicaciones}

Como dijimos en la introducción, la aplicación más importante de (6.1) es la
construcción de esquemas de Picard, así como la solución de otros diversos
problemas de módulos, sobre los cuales volveremos más adelante.

Se obtiene una demostración sencilla del resultado siguiente de SHIMURA:

\result{PROPOSICIÓN}{7.1} Sea $A$ un esquema abeliano definido sobre un
anillo de valoración discreta $V$ de cuerpo de fracciones $K$. Entonces todo
esquema abeliano $B'$ sobre $K$ que sea isógeno a un cociente de
$A\otimes_V K$ «tiene buena reducción sobre $V$», es decir, es isomorfo a
algún $B\otimes_V K$,\footnote{El texto impreso contiene $B\otimes_V A$.
Como $B$ es un $V$-esquema y la afirmación concierne a su fibra genérica, el
factor de base debe ser $K$.} donde $B$ es un esquema abeliano sobre $V$
(esencialmente único, como recordamos).

Podemos suponer que $B'$ es el cociente de $A_K$ por un subesquema en grupos
$C'$. (N.B.: $C'$ no será en general «reducido», es decir, sus anillos locales
tendrán elementos nilpotentes.) Considérese el subesquema cerrado $C$ de $A$
que es la «clausura» de $C'$, esto es, el menor subesquema cerrado de $A$ tal
que $C_K$ contiene a $C'$. Se tiene $C_K=C'$, y el hecho de que $V$ sea un
anillo de valoración discreta implica fácilmente que $C$ es un subesquema en
grupos de $A$ sobre $V$. Puesto que $A$ es propio sobre
$\operatorname{Spec}(V)=S$, también lo es $C$; además, $A$ es incluso
proyectivo sobre $S$. Podemos, por tanto, aplicar (6.1) para construir
$A/C=B$,
\src{115}
que es el $B$ buscado.

Por último, argumentos esencialmente conocidos permiten deducir de (6.2) el
resultado siguiente:

\result{TEOREMA}{7.2} Sea $S$ el espectro de un anillo artiniano, sean $F$ y
$G$ esquemas en grupos de tipo finito sobre $S$, y sea $u:F\to G$ un
homomorfismo de esquemas en grupos sobre $S$. Supongamos que:
\begin{enumerate}
\item[(i)] $F$ es plano sobre $S$;
\item[(ii)] el núcleo de $u$ es finito.
\end{enumerate}
En estas condiciones existe el esquema cociente $G/F$; el morfismo canónico
$G\to G/F$ es sobreyectivo y abierto, y sus fibras son las clases de
equivalencia conjuntistas definidas por las acciones por la derecha de $F$ en
$G$. Finalmente, si $u$ es un monomorfismo, entonces $G\to G/F$ es plano y
el morfismo
\[
  G\times F\longrightarrow G\times_{G/F}G
\]
es un isomorfismo. En otras palabras, $G$ es un espacio homogéneo principal
sobre $G/F$ con grupo estructural (que actúa por la derecha) $F$, o, más
exactamente, $F\times_S(G/F)$ considerado como esquema en grupos sobre $G/F$
(véase [1], B, \S~6).

\result{COROLARIO}{7.3} En estas condiciones, $G$ es plano sobre $S$ si y
solo si $G/F$ es plano sobre $S$. Si se cumple esta condición, la formación
del cociente por $F$ conmuta con toda extensión de la base $S$; y, si $F$ es
un subgrupo normal de $G$, entonces $G/F$ puede dotarse de la estructura de
\emph{grupo cociente} de $G$ por $F$.

La situación es especialmente sencilla cuando $S$ es el espectro de un
cuerpo, pues entonces todo $S$-preesquema es automáticamente plano sobre $S$.
Se obtiene:

\result{COROLARIO}{7.4} Sean $F,G$ esquemas en grupos de tipo finito sobre un
cuerpo $k$, y sea $u:F\to G$ un homomorfismo de $k$-grupos. Entonces $u$ se
factoriza como $F\to F'\to G$, donde $F\to F'$ es el homomorfismo cociente
por el subgrupo cerrado $\operatorname{Ker}u$ de $F$, y $F'\to G$ es un
homomorfismo de grupos que es una inmersión cerrada. Existe el cociente
$G/F=G/F'$. El formalismo usual (del tipo proporcionado por los teoremas de
Noether) es válido entre grupos algebraicos sobre $k$.

Este resultado permite tratar de manera uniforme tanto el paso al cociente
entre grupos algebraicos en el sentido clásico (es decir, los que son
irreducibles sobre $k$ y lisos sobre $k$) como el paso al cociente por
subgrupos «infinitesimales»
\src{116}
considerados por CARTIER. Resulta ventajoso considerar las «hiperálgebras»
introducidas por CARTIER, a raíz de los trabajos de DIEUDONNÉ sobre los grupos
formales, como grupos en la categoría de los esquemas formales sobre $k$ y,
cuando proceda (si corresponden a hiperálgebras de rango finito sobre $k$),
como grupos algebraicos finitos sobre $k$.

\fsection{8}{Una conjetura}

Responde a la necesidad de formar cocientes por el grupo proyectivo que actúa
en ciertos subesquemas de «esquemas de Hilbert» (que sustituyen a las
variedades de Chow en la teoría de esquemas).

Sea $S$ un preesquema y sea $n$ un entero. A todo preesquema $S'$ sobre $S$
se asocia el grupo $\operatorname{Gl}(n,\Gamma(S',\mathcal O_{S'}))$ de
matrices invertibles de tamaño $n$ por $n$ con entradas en el anillo de
secciones de $\mathcal O_{S'}$. Esto da un funtor contravariante en $S'$, que
se comprueba fácilmente que es representable; corresponde, por tanto, a un
esquema en grupos sobre $S$, de hecho afín sobre $S$, denotado por
$\operatorname{Gl}(n)_S$. Su formación es compatible con el cambio de base,
de modo que, en realidad, todo procede de un esquema en grupos sobre
$\mathbf Z$, denotado por $\operatorname{Gl}(n)$. El grupo
$\operatorname{Gl}(1)$, llamado grupo multiplicativo y denotado a menudo por
$\mathbf G_m$, corresponde al funtor
$S\longmapsto\Gamma(S,\mathcal O_S)^*$, el grupo de las «unidades» sobre $S$.
Hay un homomorfismo evidente
$\operatorname{Gl}(1)\to\operatorname{Gl}(n)$, y se construye fácilmente el
grupo cociente, denotado por $\operatorname{GP}(n-1)$ y llamado \emph{grupo
proyectivo de grado $n-1$ sobre $\mathbf Z$}. Representa el funtor que asocia
a $S$ el grupo de secciones del haz
$\widetilde{\operatorname{Gl}(n)}_S/
\widetilde{\operatorname{Gl}(1)}_S$, donde
$\widetilde{\operatorname{Gl}(n)}_S$ denota el haz de gérmenes de secciones de
$\operatorname{Gl}(n)_S$ sobre $S$. (¡Hay que tener cuidado: las secciones de
$\operatorname{GP}(n-1)_S$ sobre $S$ no proceden, en general, de secciones de
$\operatorname{Gl}(n)_S$ sobre $S$!) Obsérvese que puede demostrarse que
$\operatorname{GP}(n-1)$ representa también el funtor
$S\longmapsto\operatorname{Aut}_S(\mathbf P_S^{n-1})$, donde
$\mathbf P_S^{n-1}$ es el espacio proyectivo estándar de dimensión relativa
$n-1$ sobre $S$, al menos para $S$ noetheriano. Así es como interviene en la
teoría de módulos.

Sea $S$ un esquema noetheriano, que si se desea puede suponerse afín, y sea
$X$ un $S$-preesquema cuasiproyectivo provisto de un haz invertible
$\mathcal L$ muy amplio con respecto a $S$. Supongamos que el grupo
$G=\operatorname{GP}(n)_S$ actúa en $X$ y al mismo tiempo en $\mathcal L$
(compatiblemente con su acción en $X$), y que actúa \emph{libremente} en
$X$.\footnote{El texto impreso dice «en $S$». La acción considerada y la
condición de cociente conciernen a $X$.}
\src{117}

\result{CONJETURA}{8.1}\footnote{La fe de erratas vinculada de la Exposición 212
afirma que esta conjetura es definitivamente falsa. También señala que el
hecho positivo de la Observación 8.1 fue demostrado posteriormente y remite
al teorema de Mumford y a refinamientos posteriores.} En las condiciones
anteriores:

\smallskip
\noindent $1^\circ$ \emph{La relación de equivalencia definida por $G$ es
efectiva, el cociente $Y=X/G$ es de tipo finito sobre $S$, y el morfismo
canónico $f:X\to Y$ es plano y sobreyectivo} (así, $X$ se convierte en un
fibrado principal homogéneo sobre $Y$ con grupo
$G\times_S Y=\operatorname{GP}(n)_Y$).

\smallskip
\noindent $2^\circ$ Sea $\mathcal L'$ el haz invertible sobre $Y$ deducido de
$\mathcal L$ por «descenso fielmente plano» a lo largo de $f$ (véase [1], B,
\S~teo. 1). Entonces $\mathcal L'$ es \emph{preamplio sobre $Y$ con respecto
a $S$}, es decir, existen un entero $m$ y un morfismo cuasifinito de $Y$ a un
esquema proyectivo estándar adecuado $\mathbf P_S^N$ tales que
$(\mathcal L')^{\otimes m}$ es isomorfo a la imagen inversa de
$\mathcal O_{\mathbf P_S^N}(1)$.

Obsérvese que, incluso si $X$ es separado sobre $S$, puede ocurrir que $Y$ no
sea separado sobre $S$ (situación que se presenta en problemas de módulos
nada patológicos). Si se cumple $1^\circ$, entonces $Y$ es separado si y solo
si la relación de equivalencia definida por $G$ tiene grafo cerrado, es decir,
si $G\times X\to X\times X$ tiene imagen cerrada (es entonces una inmersión
cerrada). Si $Y$ es separado, entonces $\mathcal L'$ es preamplio sobre $Y$
con respecto a $S$ si y solo si es amplio, es decir, si una potencia tensorial
adecuada define una inmersión proyectiva. En los problemas de módulos
mencionados en la introducción puede demostrarse que la relación de
equivalencia obtenida tiene, en efecto, grafo cerrado.

\medskip
\noindent\textbf{OBSERVACIONES 8.1. ---} Hemos supuesto
$G=\operatorname{GP}(n)_S$ para fijar las ideas y porque este es actualmente
el caso más importante en la práctica. Una hipótesis razonable sobre $G$
parecería ser, en cambio, que $G$ fuese una de las «formas» sobre $S$ de uno
de los grupos de Tohokû (cuya construcción sobre los enteros fue realizada
también por CHEVALLEY). El único hecho positivo que conozco en la dirección de
la conjetura anterior es el siguiente: sea $X$ \emph{un esquema afín sobre un
cuerpo $k$ de característica $0$, en el que el grupo
$\operatorname{Gl}(n)_k$ o $\operatorname{GP}(n-1)_k$ actúa libremente.
Entonces la relación de equivalencia definida por $G$ es efectiva, el
cociente $X/G$ es también afín y el morfismo $X\to X/G$ es plano y
sobreyectivo}. La demostración utiliza el hecho siguiente (que actualmente
solo se ha demostrado en característica $0$): cuando se hace actuar a $G$ en
el anillo de coordenadas afín $A$ de $G$, considerado como espacio vectorial
sobre $k$, la representación trivial de $G$ aparece una sola vez (en una
serie de composición de un subespacio vectorial estable bajo $G$ y de
dimensión finita sobre $k$). Parece posible que un uso sistemático de
\src{118}
la teoría de representaciones lineales de $G$ conduzca finalmente a una
demostración de la conjetura, al menos sobre un cuerpo base. Cuando la base ya
no es un cuerpo, el conferenciante no sabe nada.

\section*{BIBLIOGRAFÍA}
\addcontentsline{toc}{section}{Bibliografía}

\begin{description}
\item[{[1]}] GROTHENDIECK (Alexander). --- Techniques de descente et théorèmes
d'existence en géométrie algébrique, I: Généralités, Séminaire Bourbaki,
1959/60, no. 190, 29 pp.

\item[{[2]}] GROTHENDIECK (Alexander). --- Techniques de descente et théorèmes
d'existence en géométrie algébrique, II: Le théorème d'existence en théorie
formelle des modules, Séminaire Bourbaki, 1959/60, no. 195, 22 pp.

\item[{[3]}] GROTHENDIECK (A.) and DIEUDONNÉ (J.). --- Éléments de géométrie
algébrique, I: Le langage des schémas. --- Paris, Presses universitaires de
France, 1960 (Institut des hautes Études scientifiques, Publications
mathématiques, 4).
\end{description}

\def\FgaUnitKey{e212}\clearpage
\chapter*{Fe de erratas de la exposición n.º 212}
\addcontentsline{toc}{chapter}{Fe de erratas de la exposición n.º 212}
\setcounter{footnote}{2}
\src{301}

\noindent\textbf{5. TDTE III: Preesquemas cociente.}

\begin{center}
[Seminario Bourbaki, t. 13, 1960/61, n.º 212, 20 pp.]
\end{center}

\noindent\textbf{Página 15, afirmación (i).} --- El hecho de que $Y=X/R$ sea
cuasiproyectivo sobre $S$ parece, por el momento, haberse demostrado solo
cuando $R$ proviene de una relación de equivalencia.

\noindent\textbf{Páginas 18 y 19.} --- Como señalamos al final de la exposición
siguiente, la conjetura considerada es definitivamente falsa. El «hecho
positivo» señalado en la Observación 8.1 parece haber sido demostrado
simultáneamente por diversos autores (NAGATA, ROSENLICHT, GROTHENDIECK,
\ldots).

Señalemos que la aplicación hecha en TDTE V de la teoría del paso al cociente
desarrollada aquí a la construcción de esquemas de Picard puede sustituirse
también por un uso adecuado de los esquemas de Hilbert (cf. SHMT). Como se
indicó en el § 8, la laguna más importante de la teoría presentada
\src{302}
aquí es la falta de un criterio de existencia de cocientes por una relación de
equivalencia no propia, como las relaciones de equivalencia procedentes de
ciertas acciones del grupo proyectivo. Un teorema importante en esta dirección
ha sido obtenido por MUMFORD\footnote{MUMFORD (David). --- An elementary
theorem in geometric invariant theory, \emph{Bull. Amer. Math. Soc.}, t. 67,
1961, pp. 483--486.}. Para un refinamiento de su resultado y diversas
aplicaciones a la teoría, véase SHMT.

\def\FgaUnitKey{221}\unit{Exposición 221 -- Los esquemas de Hilbert}
\sid{221}
\src{249}
\seminares{13.er año, 1960/61, n.º 221\\Mayo de 1961}
\paperhead{TÉCNICAS DE CONSTRUCCIÓN Y TEOREMAS DE EXISTENCIA EN GEOMETRÍA ALGEBRAICA\\
IV. LOS ESQUEMAS DE HILBERT}%
  {por Alexander GROTHENDIECK}

\subsection*{Introducción}
\addcontentsline{toc}{subsection}{Introducción}

Las técnicas expuestas en [2], I y II, eran en lo esencial independientes de
toda hipótesis de proyectividad sobre los esquemas considerados.
Desafortunadamente, en la actualidad no permiten resolver los problemas de
existencia planteados en la Exposición II de [2]. En la presente exposición y
en la siguiente resolveremos estos problemas bajo hipótesis de proyectividad.
Las técnicas empleadas son típicamente proyectivas y prácticamente no hacen
uso de los resultados de I y II de [2]. Aquí construiremos los «esquemas de
Hilbert», destinados a sustituir el uso de las coordenadas de Chow, como se
indicó en [2], II, n.º 2. En la exposición siguiente, la teoría del paso al
cociente en esquemas desarrollada en III, unida a la teoría de los esquemas de
Hilbert, nos permitirá, por ejemplo, construir los esquemas de Picard
(definidos en [2], II, n.º 3) bajo condiciones bastante generales.

En resumen, puede decirse que ahora se dispone de una técnica de construcciones
proyectivas bastante satisfactoria, salvo que todavía
falta\footnote{Véase el aditivo al final de la exposición.} un teorema de paso
al cociente por grupos tales como el grupo proyectivo, que actúan «sin puntos
fijos» (cf. [2], III, n.º 8). La situación parece incluso algo mejor en
geometría analítica (cuando uno se limita al estudio de espacios analíticos
«proyectivos» sobre un espacio analítico dado), pues para los espacios
analíticos desaparece la dificultad del paso al cociente por un grupo que
actúa convenientemente. Por otra parte, tanto en geometría algebraica como en
geometría analítica quedaría por desarrollar una técnica de construcción
válida sin hipótesis de proyectividad.

\fsection{1}{Familias acotadas de haces: propiedades de permanencia}

Sean $k$ un cuerpo y $X$ un $k$-preesquema, de tipo finito para simplificar.
Para toda extensión $K/k$ se obtiene, pues, un $K$-preesquema
$X_K=X\otimes_k K$. Si $\mathcal F$ es un haz coherente sobre $X_K$ y $K'$ es
una extensión de $K$, entonces $\mathcal F\otimes_K K'=\mathcal F_{K'}$ es un
haz cuasicoherente sobre $X_K\otimes_K K'=X_{K'}$. Dicho esto, si $K$ y $K'$
son dos extensiones cualesquiera de $k$, $\mathcal F$ es un haz cuasicoherente
sobre $X_K$ y $\mathcal F'$
\src{250}
es un haz cuasicoherente sobre $X_{K'}$, diremos que $\mathcal F$ y
$\mathcal F'$ son \emph{equivalentes} si existen $k$-homomorfismos de $K$ y
$K'$ en una misma extensión $K''$ de $k$ tales que $\mathcal F_{K''}$ y
$\mathcal F'_{K''}$ sean isomorfos sobre $X_{K''}$. Esta es una relación de
equivalencia, y nos interesaremos por las clases de equivalencia de haces para
esta relación y por los conjuntos de clases de equivalencia. Obsérvese que, si
$X$ es de tipo finito sobre $k$,\footnote{El texto impreso lleva aquí $X_0$;
solo $X$ está definido en este pasaje.} toda clase de haces coherentes puede
definirse por un haz coherente sobre algún $X_K$, donde $K$ es una extensión
de tipo finito de $k$. Por consiguiente, en la definición de las clases de
haces coherentes puede uno limitarse a las extensiones algebraicamente cerradas
de $k$; también puede limitarse a una extensión algebraicamente cerrada fija
$\Omega$ de $k$, de grado de trascendencia infinito. Dos haces coherentes
$\mathcal F$ y $\mathcal F'$ sobre $X_\Omega$ son entonces equivalentes si y
solo si existe un $k$-automorfismo $\sigma$ de $\Omega$ tal que
$\mathcal F\otimes_\Omega(\Omega,\sigma)$ sea isomorfo a $\mathcal F'$.
Se comprueba así que hay una correspondencia biyectiva entre las clases de
haces coherentes para una u otra definición.

Sean $E$ y $E'$ dos conjuntos de clases de haces coherentes sobre $X$.
Consideremos las clases de todos los haces de la forma
$\mathcal F\otimes\mathcal F'$, donde $\mathcal F$ y $\mathcal F'$ son haces
coherentes sobre un mismo $X_K$, la clase de $\mathcal F$ pertenece a $E$ y la
de $\mathcal F'$ pertenece a $E'$. Se obtiene así un conjunto de clases de
haces coherentes, que denotaremos por $E\otimes E'$. Se definen del mismo modo
los $\operatorname{Tor}_i(E,E')$, etc. En general, a toda función $U$ que
asocia a cada sucesión $\mathcal F_1,\ldots,\mathcal F_n$ de $n$ haces
coherentes sobre un mismo $X_K$ un conjunto
$U(\mathcal F_1,\ldots,\mathcal F_n)$ de haces coherentes sobre $X_K$, y que
posee la evidente propiedad de compatibilidad con los isomorfismos de haces y
las imágenes inversas por cambio de base, se le asocia una función, denotada
por el mismo símbolo $U$, que asigna a cada sucesión de $n$ conjuntos
$E_1,\ldots,E_n$ de clases de haces coherentes un conjunto
$U(E_1,\ldots,E_n)$ de clases de haces coherentes.

Nuestro propósito en este apartado es dar una definición de ciertos conjuntos
de clases de haces, llamados \emph{acotados}, y demostrar que las operaciones
$U$ más usuales, aplicadas a conjuntos acotados, vuelven a dar conjuntos
acotados.

Sea $X$ un preesquema de tipo finito sobre un $S$ noetheriano. Para todo
$s\in S$, la fibra $X_s$ es un preesquema de tipo finito sobre $k(s)$, y
consideraremos clases de haces coherentes sobre $X_s$ en el sentido anterior.
Esto da sentido a la expresión «clase de haces coherentes sobre una fibra de
$X/S$» y a las expresiones análogas. Del mismo modo, procediendo por separado
sobre cada fibra, pueden seguir considerándose operaciones como
$E\otimes E'$, etc., que hacen corresponder a sistemas de conjuntos de
\src{251}
clases de haces coherentes sobre las fibras de $X/S$ otro conjunto de clases
de haces coherentes sobre las fibras de $X/S$.

\result{DEFINICIÓN}{1.1} Sea $E$ un conjunto de clases de haces coherentes
sobre las fibras de $X/S$. Se dice que $E$ está \emph{acotado} si existen un
preesquema $S'$ de tipo finito sobre $S$ y un haz coherente $\mathcal F'$ sobre
$X'=X\times_S S'$ tales que $E$ esté contenido en el conjunto de clases de
haces sobre las fibras de $X/S$ definido por $\mathcal F'$.

Por definición, este último hace corresponder a un $s\in S$ las clases de los
haces $\mathcal F'\otimes_{S'}k(s')$, donde $s'$ recorre los puntos de $S'$
situados sobre $s$ (de modo que $k(s')$ es una extensión de $k(s)$ y
$X_s\otimes_{k(s)}k(s')$ se identifica con la fibra
$X'\otimes_{S'}k(s')=X'_{s'}$ de $X'$ en $s'$). Puede decirse que las familias
acotadas son aquellas contenidas en una \emph{familia algebraica} de haces
coherentes parametrizada por un $S'$ de tipo finito sobre $S$.

Una unión finita de familias acotadas es acotada (tómese el preesquema suma de
los preesquemas de parámetros $S_i$ que definen las familias algebraicas
mayorantes). Un cambio de base $T\longrightarrow S$ transforma una familia
acotada relativa a $X/S$ en una familia acotada relativa a $X_T/T$, y el
recíproco es cierto si $T\longrightarrow S$ es sobreyectivo (o, más
generalmente, si su imagen contiene los $s$ que intervienen efectivamente en la
familia dada $E$ para $X/S$). Esto reduce, en principio, la determinación de
las familias acotadas al caso en que $S$ es el espectro de un álgebra de tipo
finito sobre el anillo de los enteros $\mathbf Z$.

Si $E$ y $E'$ son familias acotadas de clases de haces relativas a $X/S$,
entonces $E\otimes E'$ también está acotada. En efecto, supongamos que $E$ y
$E'$ están mayoradas, respectivamente, por las familias algebraicas definidas
por $T\longrightarrow S$ y $\mathcal F$ sobre $X_T$, y por
$T'\longrightarrow S$ y $\mathcal F'$ sobre $X_{T'}$. Se ve entonces que
$E\otimes E'$ está mayorada por la familia algebraica definida por
$T''=T\times_S T'\longrightarrow S$ y el haz $\mathcal F''$ sobre
$X_{T''}$, producto tensorial de las imágenes inversas de $\mathcal F$ y
$\mathcal F'$ sobre $X_{T''}$. Este razonamiento solo es válido porque el
funtor $\mathcal F\otimes\mathcal F'$ en $\mathcal F,\mathcal F'$ es exacto a
derecha y, por tanto, conmuta con la extensión de la base (y, en particular,
con el paso a las fibras). No es aplicable directamente a operaciones locales
tales como $\operatorname{Tor}_i(E,E')$, $\operatorname{Hom}(E,E')$ y
$\operatorname{Ext}^i(E,E')$. No obstante, puede demostrarse que estas
operaciones siguen transformando conjuntos acotados en conjuntos acotados,
procediendo como para $E\otimes E'$ pero utilizando además resultados del tipo
siguiente (todos ellos contenidos en [3], IV, 6.11): una familia acotada $E$
está siempre mayorada por una familia algebraica definida por un haz coherente
$\mathcal F$ sobre un $X_T$ ($T$ de tipo finito sobre $S$) que es plano
relativamente a $T$. (Se «corta en trozos» el espacio de parámetros inicial.)
Tales
\src{252}
propiedades de planitud de haces adecuados aseguran, en efecto, que operaciones
como $\operatorname{Tor}_i(\mathcal F,\mathcal F')$ conmuten con cualquier
cambio de base. El mismo método se aplica a operaciones de naturaleza global:
imágenes directas e imágenes directas derivadas de haces coherentes por
morfismos propios, $\operatorname{Ext}$ globales relativos a morfismos propios
(cf. [5], III, \S~6), etc. Todas estas operaciones transforman familias
acotadas de haces en familias acotadas de haces. (N. B. -- Aquí los
preesquemas sobre los que se toman los diversos haces pueden cambiar por las
operaciones del tipo considerado.)

Los dos enunciados siguientes se demuestran esencialmente mediante la misma
técnica de planitud; para la descomposición primaria sobre las fibras de un
morfismo de tipo finito, véase en particular [5], IV.

\result{PROPOSICIÓN}{1.2} Sean $E$ y $E'$ conjuntos acotados de clases de
haces sobre las fibras de $X/S$, donde se supone que $X$ es propio sobre $S$.
Entonces:

\medskip
\noindent (i) La familia de núcleos, conúcleos e imágenes de homomorfismos
$\mathcal F\longrightarrow\mathcal F'$, donde la clase de $\mathcal F$
pertenece a $E$ y la de $\mathcal F'$ pertenece a $E'$, está acotada.

\smallskip
\noindent (ii) La familia de haces $\mathcal F''$ que son extensiones de un
$\mathcal F$ por un $\mathcal F'$, donde la clase de $\mathcal F$ pertenece a
$E$ y la de $\mathcal F'$ pertenece a $E'$, está acotada.

Mediante un cambio de base conveniente puede suponerse que $E$ y $E'$ están
definidos, respectivamente, por haces coherentes $\mathcal G$ y $\mathcal G'$
sobre un $X_T/T$, donde $T$ es de tipo finito sobre $S$. Además, puede
suponerse que se satisfacen ciertas hipótesis de planitud que implican que la
formación de los haces
$f_{T*}(\operatorname{Hom}_{X_T}(\mathcal G,\mathcal G'))$ y
$\operatorname{Ext}^{1}_{f_T}(\mathcal G,\mathcal G')$ conmuta con el cambio
de base para un morfismo arbitrario $T'\longrightarrow T$. Asimismo, puede
suponerse que estos haces coherentes sobre $T$ son localmente libres. Sean
$T_0$ y $T_1$ los fibrados vectoriales sobre $T$ cuyos haces de gérmenes de
secciones son, respectivamente, los haces precedentes. Se define entonces
canónicamente un homomorfismo
$\mathcal G_{T_0}\longrightarrow\mathcal G'_{T_0}$ de haces coherentes sobre
$X_{T_0}$ y una extensión
\[
0\longrightarrow\mathcal G'_{T_1}\longrightarrow\mathcal G''
\longrightarrow\mathcal G_{T_1}\longrightarrow0
\]
de haces coherentes sobre $X_{T_1}$ que posee una propiedad universal
evidente.\footnote{El texto impreso invierte los términos extremos de esta sucesión. Se trata de extensiones de $\mathcal G$ por $\mathcal G'$, como indica $\operatorname{Ext}^1(\mathcal G,\mathcal G')$.} El segundo haz define una familia algebraica que mayora la familia
considerada en (ii). Lo mismo vale para el núcleo, el conúcleo y la imagen del
homomorfismo
\src{253}
precedente, cuya consideración demuestra (i) (siempre que se suponga que el
conúcleo es plano relativamente a $T_0$, caso al que se reduce de nuevo
cortando $T_0$ en trozos\ldots).

\result{PROPOSICIÓN}{1.3} Sea $E$ una familia acotada de clases de haces sobre
las fibras de $X/S$. Entonces las clases de los haces estructurales de los
$(\operatorname{supp}\mathcal F)_{\mathrm{red}}$, donde $\mathcal F$ es un
haz coherente sobre un $X_K$, con $K$ algebraicamente cerrado, cuya clase
pertenece a $E$, forman una familia acotada.

Aquí $(\operatorname{supp}\mathcal F)_{\mathrm{red}}$ denota el soporte de
$\mathcal F$ provisto de la estructura reducida inducida; es decir, su haz
estructural es el cociente de $\mathcal O_{X_K}$ por el mayor haz de ideales
que define $\operatorname{supp}\mathcal F$. Puede demostrarse el resultado
análogo a (1.3) para los haces deducidos canónicamente de $\mathcal F$ por la
teoría de la descomposición primaria; por ejemplo, los
$\mathcal F/\mathcal F_i$, donde los $\mathcal F_i$ son los subhaces de
$\mathcal F$ primarios para las componentes del soporte de $\mathcal F$ y
mínimos con esta propiedad; o los $\mathcal O_{X_K}/\mathfrak p$, donde
$\mathfrak p$ es un haz de ideales primos asociado a $\mathcal F$; o los
$\mathcal O_{X_K}/\mathfrak q$, donde $\mathfrak q$ es un haz de ideales
primarios asociado a una componente del soporte de $\mathcal F$ (siendo el
cuerpo de referencia algebraicamente cerrado).

\fsection{2}{Familias acotadas y polinomio de Hilbert}

En lo que sigue suponemos que $X$ es proyectivo sobre $S$ y está provisto de
un haz muy amplio, denotado $\mathcal O_X(1)$. Para toda extensión $K$ de una
extensión residual $k(s)$ de un punto $s$ de $S$, se considera sobre $X_K$ el
haz correspondiente $\mathcal O_{X_K}(1)$, que sigue siendo muy amplio.

A todo haz coherente $\mathcal F$ sobre $X_K$ se le asocia la función
\[
P_{\mathcal F}(n)=\text{característica de Euler--Poincaré de $\mathcal F(n)$
sobre $X_K$},
\]
que es un polinomio en el entero $n$, llamado \emph{polinomio de Hilbert} de
$\mathcal F$. Para valores grandes de $n$, $P(n)$ no es sino la dimensión de
$H^0(X_K,\mathcal F(n))$ sobre $K$, puesto que los
$H^i(X_K,\mathcal F(n))$ se anulan para $i>0$ y $n$ grande.

Supongamos ahora que $\mathcal F$ es un haz coherente sobre $X$, plano sobre
$S$. Entonces los polinomios de Hilbert de los haces $\mathcal F_s$ inducidos
sobre fibras $X_s$ pertenecientes a una misma componente conexa de $S$ son
todos iguales [5], III, \S~7. De ello se sigue (sin hipótesis de planitud) que
el conjunto de polinomios de Hilbert de los haces $\mathcal F_s$, $s\in S$,
es finito para todo haz coherente $\mathcal F$ sobre $X$.
\src{254}
Recordemos además que, si $\mathcal F$ es un haz coherente sobre $X$, es
isomorfo a un haz cociente de un haz $\mathcal O_X(-n)^N$, para $n$ y $N$
suficientemente grandes. Por tanto, los haces $\mathcal F_s$ inducidos sobre
las fibras también son cocientes del haz $\mathcal O(-n)^N$ sobre la fibra.\footnote{El texto impreso omite aquí el exponente $N$; la restricción a la fibra conserva la potencia directa de rango $N$ del haz precedente.}

De estas dos observaciones se deduce la parte «solo si» del teorema siguiente:

\result{TEOREMA}{2.1} Sea $X$ proyectivo sobre un $S$ noetheriano y sea
$\mathcal O_X(1)$ muy amplio sobre $X$ relativamente a $S$. Sea $E$ un
conjunto de clases de haces sobre las fibras de $X/S$. Para que $E$ esté
acotado, es necesario y suficiente\footnote{La exposición impresa dice
solamente «es necesario»; la fe de erratas de 1962 de la Exposición 221 prescribe «es
necesario y suficiente».} que satisfaga las condiciones siguientes:

\medskip
\noindent (a) Existe un haz coherente $\mathcal L$ sobre $X$ (que puede
suponerse de la forma $\mathcal O_X(-n)^N$) tal que $E$ esté contenido en la
familia de clases de haces coherentes cocientes de haces de la forma
$\mathcal L_K$.

\smallskip
\noindent (b) Los polinomios de Hilbert $P_{\mathcal F}$ de los haces
$\mathcal F$ cuya clase pertenece a $E$ son elementos de un mismo conjunto
finito de polinomios.

Queda por demostrar la parte «si», que será un caso particular de un resultado
más preciso. Para todo Módulo coherente $\mathcal F$ sobre un preesquema de
tipo finito sobre un cuerpo $K$, y todo entero $r$, sea $N_r$ el sub-Módulo de
$\mathcal F$ cuyas secciones sobre un abierto son las secciones de $\mathcal F$
sobre dicho abierto cuyo soporte tiene dimensión $<r$. Se tiene, pues,
$N_r=\mathcal F$ para $r>\dim\operatorname{supp}\mathcal F$, mientras que
$N_r=0$ para $r\leq0$, y se obtiene una filtración creciente finita de
$\mathcal F$ cuyos factores $N_{r+1}/N_r$\footnote{El texto impreso invierte
el cociente, escribiendo $N_r/N_{r+1}$, lo que es incompatible con la
filtración creciente que acaba de definirse.} tienen como ciclos primos
asociados exactamente los ciclos primos asociados a $\mathcal F$ que tienen
dimensión $r$. Pondremos
\[
\mathcal F_{(r)}=\mathcal F/N_r.
\]
Entonces los ciclos primos asociados a $\mathcal F_{(r)}$ son exactamente los
ciclos primos asociados a $\mathcal F$ que tienen dimensión $\geq r$; en
particular, es igual a $\mathcal F$ si y solo si los ciclos primos asociados a
$\mathcal F$ tienen dimensión $\geq r$. Con esta notación:

\result{TEOREMA}{2.2} Bajo las hipótesis de (2.1), sea $s$ un entero y
supongamos que $E$ satisface la condición (a) y la siguiente forma debilitada
de (b):

\medskip
\noindent $(b_s)$ Los polinomios de Poincaré $P_{\mathcal F}$ de los haces
$\mathcal F$ cuya clase pertenece a $E$ tienen acotados sus coeficientes en
grados $\geq s-1$.

En estas condiciones, los haces $\mathcal F_{(s)}$ (cuando la clase de
$\mathcal F$ recorre $E$) forman una familia acotada. Además, los coeficientes
de grado $s-2$ de los $P_{\mathcal F}$ están minorados.
\src{255}
Así pues:

\result{COROLARIO}{2.3} Supongamos que los haces $\mathcal F$ cuya clase
pertenece a $E$ son tales que todos sus ciclos primos asociados tienen
dimensión $d$ con $s\leq d\leq r$. Entonces, en (2.1) (b), basta considerar los
coeficientes de los $P_{\mathcal F}$ en grados comprendidos entre $s-1$ y $r$.

El resto de este apartado está dedicado a esbozar la demostración de (2.2).
Los lemas clave son los dos lemas siguientes, el primero de los cuales es bien
conocido (y resume el contenido matemático útil de las coordenadas de Chow):

\result{LEMA 2.4 (CHOW)}{} Consideremos los haces estructurales de subesquemas
$Y$ de fibras $X_K$ ($K$ recorre las extensiones algebraicamente cerradas de
cuerpos residuales de $S$), donde $Y$ es reducido y todas sus componentes
tienen la misma dimensión $r$ (considerándose $\mathcal O_Y$ como un haz
cociente de $\mathcal O_X$). Si los grados de los $Y$ están acotados, los $Y$
forman una familia acotada.

Aquí, el grado $a$ de $Y$ puede definirse del modo más cómodo mediante el
coeficiente del término principal de
\[
P_{\mathcal O_Y}=a\frac{n^r}{r!}+\cdots.
\]

\result{LEMA}{2.5} Sea $\mathcal L$ un haz coherente sobre $X$, y sea $E$ un
conjunto de clases de haces cocientes $\mathcal F$ de haces $\mathcal L_K$
($K$ una extensión residual de $S$). Supongamos que las fibras de $X$ sobre
$S$ tienen dimensión $\leq r$, y escribamos
\[
P_{\mathcal F}(n)=a_{\mathcal F}\frac{n^r}{r!}
+b_{\mathcal F}\frac{n^{r-1}}{(r-1)!}
+\text{términos de grado $<r-1$}.
\]
Entonces el coeficiente $a_{\mathcal F}$ está acotado, mientras que
$b_{\mathcal F}$ está minorado. Si $b_{\mathcal F}$ está acotado, la familia
de los $\mathcal F_{(r)}$ está acotada.

\emph{Demostración.} -- Tras sustituir $S$ por una unión de subesquemas de
$S$ que recubren $S$, puede suponerse que existe un morfismo finito
$f:X\longrightarrow\mathbf P^r_S$ tal que $\mathcal O_X(1)$ es isomorfo a la
imagen inversa de $\mathcal O_{\mathbf P^r_S}(1)$. Por tanto, para todo haz
coherente $\mathcal F$ sobre $X$ se tiene
$P_{\mathcal F}=P_{f_*(\mathcal F)}$. Por otra parte, se comprueba fácilmente
(mediante la técnica del apartado anterior) que un conjunto de haces
$\mathcal F$ sobre $X$ está acotado si y solo si lo está el conjunto de los
$f_*(\mathcal F)$. Finalmente,
\[
f_*(\mathcal F)_{(r)}=f_*(\mathcal F_{(r)}).
\]
\src{256}
Esto nos reduce al caso $X=\mathbf P^r_S$. Además, puede suponerse que
$\mathcal L=\mathcal O_{\mathbf P^r_S}(k)^s$ para ciertos $k$ y $s$. El
coeficiente $a_{\mathcal F}$ satisface
\[
0\leq a_{\mathcal F}\leq s
\]
y, por tanto, está acotado. Es equivalente decir que el coeficiente de
$n^{r-1}$ en los $P_{\mathcal F}(n)$ está minorado (respectivamente, acotado),
o decir lo mismo de
$P_{\mathcal F}(n-k)=P_{\mathcal F(-k)}(n)$. Nos reducimos así al caso
\[
\mathcal L=\mathcal O_{\mathbf P^r_S}^{s}.
\]

Consideremos la sucesión exacta
\[
0\longrightarrow N_r\longrightarrow\mathcal F
\longrightarrow\mathcal F_{(r)}\longrightarrow0,
\]
de donde
\[
P_{\mathcal F}=P_{\mathcal F_{(r)}}+P_{N_r}.
\]
Como el coeficiente de $n^{r-1}$ en $P_{N_r}$ es positivo (pues
$\dim\operatorname{supp}N_r\leq r-1$), se tiene
\[
b_{\mathcal F_{(r)}}\leq b_{\mathcal F}.
\]
Así pues, para demostrar las afirmaciones del lema podemos sustituir
$\mathcal F$ por $\mathcal F_{(r)}$; es decir, podemos suponer que los haces
cocientes $\mathcal F$ de $\mathcal L$ considerados son libres de torsión.

Como $\mathbf P^r_K$ es normal, se sigue que $\mathcal F$ es localmente libre
de rango $a=a_{\mathcal F}$ sobre un abierto $U=\mathbf P^r_K-Y$, donde $Y$
tiene codimensión $\geq2$. Por tanto, $\bigwedge^a\mathcal F$ es un haz sobre
$\mathbf P^r_K$ cuya restricción a $U$ es invertible; luego (como
$\mathbf P^r_K$ es regular y $Y$ tiene codimensión $\geq2$) es isomorfo a la
restricción de un haz invertible sobre $\mathbf P^r_K$, determinado salvo
isomorfismo. Este último es de la forma $\mathcal O_{\mathbf P^r_K}(d)$ para
un entero $d$ bien determinado. Como $\bigwedge^a\mathcal F$ es un cociente
de $\bigwedge^a\mathcal O_{\mathbf P^r_K}^{s}\simeq
\mathcal O_{\mathbf P^r_K}^{N}$\footnote{El texto impreso sustituye aquí el
exponente $s$ por un $n$ no definido; $s$ es el rango del haz libre que se
proyecta sobre $\mathcal F$.}, donde $N=\binom{s}{a}$, posee $N$ secciones
canónicas, y estas definen, por tanto, secciones de
\src{257}
$\mathcal O_{\mathbf P^r_K}(d)$ sobre $U$ que son restricciones de secciones
$s_i$ ($1\leq i\leq N$) de $\mathcal O_{\mathbf P^r_K}(d)$ (pues
$\mathbf P^r_K$ es normal y $Y$ tiene codimensión $\geq2$). Estas secciones
generan $\mathcal O_{\mathbf P^r_K}(d)$ en los puntos de $U$, de modo que no
son todas nulas; esto implica $d\geq0$. Por otra parte, un cálculo sencillo da
\[
b_{\mathcal F}=a_{\mathcal F}(r+1)/2+d.
\]
En particular, esto demuestra que $b_{\mathcal F}>0$, por lo que
$b_{\mathcal F}$ está minorado. Está acotado si y solo si $d$ lo está;
demostremos que, en tal caso, $\mathcal F$ pertenece a una familia acotada.
Podemos suponer que $a_{\mathcal F}$ y $b_{\mathcal F}$ son fijos, digamos
iguales a $a$ y $b$, de modo que $d$ es fijo. El conocimiento de las $N$
secciones $s_i$ de $\mathcal O_{\mathbf P^r_K}(d)$, es decir, de un
homomorfismo
\[
s:\bigwedge^a\mathcal L_K\longrightarrow\mathcal O_{\mathbf P^r_K}(d),
\]
permite recuperar $\mathcal F$ como la coimagen del homomorfismo compuesto
correspondiente
\[
\mathcal L_K\longrightarrow
\operatorname{Hom}\!\left(\bigwedge^{a-1}\mathcal L_K,
\bigwedge^a\mathcal L_K\right)\longrightarrow
\operatorname{Hom}\!\left(\bigwedge^{a-1}\mathcal L_K,
\mathcal O_{\mathbf P^r_K}(d)\right),
\]
donde el primero es el homomorfismo canónico procedente del producto exterior
y el segundo se deduce de $s$. La conclusión se sigue ahora de (1.2) (i).

Combinando los dos lemas precedentes se obtiene:

\result{LEMA}{2.6} Bajo las hipótesis preliminares de (2.1), supongamos que,
para todo $\mathcal F$, se tiene
\[
P_{\mathcal F}(n)=a_{\mathcal F}\frac{n^r}{r!}
+b_{\mathcal F}\frac{n^{r-1}}{(r-1)!}
+\text{términos de grado $<r-1$},
\]
y que los coeficientes $a_{\mathcal F}$ están acotados. Entonces los
coeficientes $b_{\mathcal F}$ están minorados. Si los $b_{\mathcal F}$ están
acotados, los $\mathcal F_{(r)}$ forman una familia acotada.

Podemos suponer que los cuerpos base $K$ de los haces $\mathcal F$ son
algebraicamente cerrados. Dótese cada
$\operatorname{supp}\mathcal F_{(r)}$, unión de las componentes de
$\operatorname{supp}\mathcal F$ de dimensión $r$, de la estructura reducida
inducida. Entonces los grados de los $\operatorname{supp}\mathcal F_{(r)}$
están acotados superiormente por $a$, por lo que, en virtud de (2.4), los
$\operatorname{supp}\mathcal F_{(r)}$ forman un conjunto acotado. Además, para
cada componente de $\operatorname{supp}\mathcal F_{(r)}$, la longitud de
$\mathcal F_{(r)}$ a lo largo de dicha componente es $\leq a$. Por tanto, si
$\mathcal J_{\mathcal F}$ es el ideal que define
$\operatorname{supp}\mathcal F_{(r)}$, entonces $\mathcal F_{(r)}$
\src{258}
puede considerarse como un Módulo sobre el subesquema $Y_{\mathcal F}$ de $X$
definido por $\mathcal J_{\mathcal F}^{a}$. Como en el lema anterior, también
podemos reducirnos al caso $\mathcal F=\mathcal F_{(r)}$, de modo que
$\mathcal F$ proviene de un Módulo sobre $Y_{\mathcal F}$. Los
$Y_{\mathcal F}$ corresponden a una familia acotada de Módulos cocientes de
los $\mathcal O_{X_K}$ y, por consiguiente, provienen de un subesquema cerrado
$Y$ de un esquema $X\times_S T$. Podemos entonces aplicar (2.5) a $Y/T$ y
$\mathcal L\otimes_XY$, lo que da la conclusión.

Ya podemos demostrar (2.2) por inducción sobre la cota superior $r$ de
$\dim\operatorname{supp}\mathcal F$. El enunciado es trivial para $r<0$;
supongamos, pues, que $r\geq0$ y que el enunciado ha sido demostrado para
$r'<r$. Por (2.6), los $\mathcal F_{(r)}$ forman un conjunto acotado; por
tanto, en virtud de (1.2) (i), también lo forman los núcleos de los
homomorfismos $\mathcal L_K\longrightarrow\mathcal F_{(r)}$. Existe, pues, un
Módulo coherente $\mathcal L'$ sobre $X$ tal que los núcleos en cuestión y,
por consiguiente, también los
\[
N_r(\mathcal F)=\operatorname{Ker}(\mathcal F\longrightarrow\mathcal F_{(r)}),
\]
son cocientes de los Módulos $\mathcal L'_K$. Como los
$\mathcal F_{(r)}$ están acotados, los $P_{\mathcal F_{(r)}}$ permanecen
acotados, y la fórmula
\[
P_{\mathcal F}=P_{\mathcal F_{(r)}}+P_{N_r}
\]
demuestra que los $P_{N_r}$ satisfacen la misma condición $(b_s)$ que los
$P_{\mathcal F}$. Por tanto, los $N_r$ satisfacen las condiciones (a) y
$(b_s)$, y por la hipótesis de inducción los $(N_r)_{(s)}$ permanecen acotados.
Ahora bien, $\mathcal F_{(s)}$ es una extensión de $\mathcal F_{(r)}$ por
$(N_r)_{(s)}$; luego, por (1.2) (ii), los $\mathcal F_{(s)}$ permanecen
acotados. Para la última afirmación de (2.2), llevamos simultáneamente la inducción sobre el acotamiento y sobre la cota inferior del coeficiente siguiente. Los núcleos acotados de los morfismos de fuente fija hacia $\mathcal F_{(s)}$ proporcionan un mismo haz coherente que se aplica sobreyectivamente sobre cada $N_s$, pero (1.2) (i) por sí solo no hace acotada la familia de estos cocientes. Reaplicamos por inducción la afirmación sobre los coeficientes a estos cocientes; la aditividad del polinomio de Hilbert y el Lema 2.6 dan entonces la cota inferior requerida para el coeficiente siguiente. Esto termina la demostración.

\fsection{3}{Esquemas de Hilbert: definición y teorema de existencia}

Sea $X$ un preesquema sobre otro preesquema $S$, y sea $\mathcal F$ un Módulo
cuasicoherente sobre $X$. Denotemos por
\[
\operatorname{Quot}(\mathcal F/X/S)
\]
el conjunto de Módulos cuasicoherentes cocientes de $\mathcal F$ que son
planos sobre $S$. Sea ahora $S'\longrightarrow S$ un morfismo de cambio de
base, y pongamos $X'=X\times_S S'$ y
$\mathcal F'=\mathcal F\otimes_{\mathcal O_S}\mathcal O_{S'}$. Así, $X'$ es
un preesquema sobre $S'$ provisto de un Módulo
\src{259}
$\mathcal F'$, y podemos considerar
$\operatorname{Quot}(\mathcal F'/X'/S')$. Pongamos
\[
\underline{\operatorname{Quot}}_{\mathcal F/X/S}(S')
=\operatorname{Quot}(\mathcal F'/X'/S')
\qquad\text{(donde $X'=X\times_SS'$)}.
\]
Si ahora $S''\longrightarrow S'$ es un $S$-morfismo, entonces
$X''=X\times_SS''$ es isomorfo a $X'\times_{S'}S''$, y $\mathcal F''$ es
isomorfo a $\mathcal F'\otimes_{\mathcal O_{S'}}\mathcal O_{S''}$. Como el
funtor de imagen inversa
\[
\mathcal G'\longmapsto
\mathcal G'\otimes_{\mathcal O_{S'}}\mathcal O_{S''}
\]
de la categoría de Módulos cuasicoherentes sobre $X'$ a la categoría de
Módulos cuasicoherentes sobre $X''$ es exacto a derecha y transforma Módulos
planos sobre $S'$ en Módulos planos sobre $S''$, hay una aplicación natural
\[
\operatorname{Quot}(\mathcal F'/X'/S')\longrightarrow
\operatorname{Quot}(\mathcal F''/X''/S''),
\]
por lo que $\underline{\operatorname{Quot}}_{\mathcal F/X/S}(S')$ es un
funtor contravariante de $S'$ (preesquema sobre $S$), con valores en la
categoría de conjuntos. En adelante supondremos que $X$ es proyectivo sobre un
$S$ noetheriano y que $\mathcal F$ es coherente; para simplificar, nos
limitaremos a los $S'$ sobre $S$ que sean localmente noetherianos.

\result{TEOREMA}{3.1} Bajo estas hipótesis, el funtor contravariante
$\underline{\operatorname{Quot}}_{\mathcal F/X/S}$ sobre la categoría de los
$S$-preesquemas localmente noetherianos es representable por un $S$-preesquema
$\operatorname{Quot}_{\mathcal F/X/S}$ que es unión disjunta de una sucesión
de $S$-esquemas proyectivos (en particular,
$\operatorname{Quot}_{\mathcal F/X/S}$ es localmente de tipo finito sobre
$S$).

Obtendremos tal descomposición como sigue. Sea $\mathcal O_X(1)$ un haz
invertible sobre $X$, muy amplio relativamente a $S$. Para todo polinomio
$P(n)$ con coeficientes racionales, sea
$\operatorname{Quot}^P(\mathcal F/X/S)$ el subconjunto de
$\operatorname{Quot}(\mathcal F/X/S)$ formado por los cocientes coherentes
$\mathcal G$ de $\mathcal F$ que son planos sobre $S$ y cuyo polinomio de
Hilbert en todo $s\in S$ es igual a $P$. Pongamos
\[
\underline{\operatorname{Quot}}^P_{\mathcal F/X/S}(S')
=\operatorname{Quot}^P(\mathcal F'/X'/S').
\]
Esto da un subfuntor de
$\underline{\operatorname{Quot}}_{\mathcal F/X/S}$. La propiedad de
invariancia de los polinomios de Hilbert recordada en el Apartado 2 implica lo
siguiente: el funtor $\underline{\operatorname{Quot}}_{\mathcal F/X/S}$ es
representable si y solo si lo son los
$\underline{\operatorname{Quot}}^P_{\mathcal F/X/S}$; en tal caso, el
$S$-preesquema representante $\operatorname{Quot}_{\mathcal F/X/S}$ es
isomorfo a la unión disjunta de los
$\operatorname{Quot}^P_{\mathcal F/X/S}$ que representan los funtores
$\underline{\operatorname{Quot}}^P_{\mathcal F/X/S}$. Así, (3.1) se deduce
del teorema siguiente:
\src{260}
\result{TEOREMA}{3.2} Con las notaciones anteriores, el funtor
$\underline{\operatorname{Quot}}^P_{\mathcal F/X/S}$ está representado por un
$S$-preesquema proyectivo $\operatorname{Quot}^P_{\mathcal F/X/S}$.

El resto de este apartado está dedicado a la demostración de (3.2).

Sea $\nu$ un entero. Para todo $S'$ sobre $S$, denotemos por $A_\nu(S')$ el
conjunto de cocientes $\mathcal G=\mathcal F'/\mathcal K$ de
$\mathcal F'=\mathcal F\otimes_{\mathcal O_S}\mathcal O_{S'}$ que son
coherentes y planos sobre $S'$, y que satisfacen las condiciones siguientes:
\begin{align*}
\text{a.}\quad &R^if'_*(\mathcal G(n))=0 &&\text{para $i>0$ y $n\geq\nu$},\\
\text{b.}\quad &R^if'_*(\mathcal K(n))=0 &&\text{para $i>0$ y $n\geq\nu$},\\
\text{c.}\quad &f'_*(\mathcal K(\nu+k))
=\mathcal S'_k f'_*(\mathcal K(\nu)) &&\text{para $k\geq0$}.
\end{align*}
Para la última relación, se supone que $X$ está escrito como el espectro primo
homogéneo de un álgebra graduada cuasicoherente $\mathcal S_*$ sobre $S$,
graduada positivamente y generada por $\mathcal S_1$, y se pone
$\mathcal S'_* =\mathcal S_*\otimes_{\mathcal O_S}\mathcal O_{S'}$, de modo
que $X'$ es el espectro primo homogéneo de $\mathcal S'_*$. Para demostrar
(3.2), se reduce fácilmente al caso $X=\mathbf P^r_S$ (pues $S$ es una unión
de abiertos $U$ tales que $X|U$ es un subesquema cerrado de $\mathbf P^r_U$ y
$\mathcal O_X(1)$ está inducido por $\mathcal O_{\mathbf P^r_U}(1)$), y al
caso en que $\mathcal F$ es de la forma
$\mathcal O_{\mathbf P^r_S}^{N}$ y es, por tanto, plano sobre $S$. En la
discusión precedente, los haces $\mathcal K$ son entonces también planos sobre
$S'$. Las relaciones de Künneth ([5], III, \S~7) demuestran que las condiciones
(a) y (b) son estables por cambio de base e implican que, para $n\geq\nu$, la
formación de $f'_*(\mathcal G(n))$ y $f'_*(\mathcal K(n))$ conmuta con la
extensión de la base (loc. cit.); por tanto, (c) también es estable por
extensión de la base.

En otras palabras, $A_\nu(S')$ es un funtor contravariante de $S'$, más
precisamente un subfuntor de
$A(S')=\underline{\operatorname{Quot}}^P_{\mathcal F/X/S}(S')$. Al variar
$\nu$, se obtiene una sucesión creciente de subconjuntos $A_\nu(S')$ de
$A(S')$ cuya unión es $A(S')$, por un teorema bien conocido de Serre [5], III,
\S~2. Supongamos ahora que $\mathcal G$ es un cociente coherente de
$\mathcal F'$, plano sobre $S'$, y que $s$ es un punto de $S'$ tal que el
cambio de base $\operatorname{Spec}(k(s))\longrightarrow S'$ da un cociente
$\mathcal G_s$ de $\mathcal F_s$ que satisface las condiciones (a), (b) y (c),
es decir, que pertenece a $A_\nu(\operatorname{Spec}(k(s)))$. Existe entonces
un entorno abierto $U$ de $s$ sobre el cual se satisfacen las mismas
condiciones para $\mathcal G|f'^{-1}(U)$; es decir, este cociente pertenece a
$A_\nu(U)$. Para (a) y (b), esto resulta del «teorema de las funciones
holomorfas» [5], III, \S~7; la condición (c) resulta del lema de Nakayama y del
hecho de que, en cualquier caso, se sabe que
$f'_*(\mathcal K(n+k))=\mathcal S'_k f'_*(\mathcal K(n))$
\src{261}
para $n$ suficientemente grande y $k\geq0$ (cf. [5], III, \S~2). De estas
observaciones se deduce lo siguiente (compárese [4], IV): el funtor $A$ es
representable si y solo si los funtores $A_\nu$ son representables; en tal
caso, el $S$-preesquema $Q$ que representa $A$ es unión creciente de
subesquemas abiertos $Q_\nu$ que representan los $A_\nu$.

Sea
\[
M_* =\bigoplus_{n\geq0}f_*(\mathcal F(n))=\mathcal S_*^{N},
\]
de modo que
\[
M'_* =M_*\otimes_{\mathcal O_S}\mathcal O_{S'}
=\bigoplus_{n\geq0}f'_*(\mathcal F'(n))=\mathcal S'_*{}^{N}.
\]
De (a) se sigue que

\noindent (a') $f'_*(\mathcal G(n))$ es localmente libre de rango $P(n)$ para
$n\geq\nu$, [5], III, \S~7; y de (b), con $i=1$, que

\noindent (a'') $f'_*(\mathcal G(n))$ es un Módulo cociente de $M'_n$.

Además, el conocimiento de este Módulo cociente para $n=\nu$ determina, por
(c), los sub-Módulos $f'_*(\mathcal K(n))$ de $M'_n$ para $n\geq\nu$ y, por
tanto, determina $\mathcal K$ y luego $\mathcal G$. Se obtiene así una
aplicación inyectiva
\[
A_\nu(S')\longrightarrow\operatorname{Grass}_{P(\nu)}(M'_\nu)
\]
de $A_\nu(S')$ en el conjunto de Módulos cocientes localmente libres de
$M'_\nu$ de rango $P(\nu)$ y, por tanto, un homomorfismo funtorial
\[
i_\nu:A_\nu(S')\longrightarrow
\underline{\operatorname{Grass}}_{P(\nu)}(M_\nu)(S'),
\]
donde el funtor de la derecha está representado por el esquema de Grassmann
$\operatorname{Grass}_{P(\nu)}(M_\nu)$ (compárese [4], V), que es proyectivo
sobre $S$. Afirmo:

\result{LEMA}{3.3} El funtor $A_\nu(S')$ es representable, y el morfismo
$Q_\nu\longrightarrow\operatorname{Grass}_{P(\nu)}(M_\nu)$ que representa el
homomorfismo $i_\nu$ es una inmersión (lo que implica que $Q_\nu$ es
cuasiproyectivo sobre $S$).

Esta afirmación equivale a la siguiente (compárese [4], IV). Supongamos dado
un Módulo cociente $N$ de $M'_\nu$, localmente libre de rango $P(\nu)$. Existe
entonces un subpreesquema $Z$ de $S'$ tal que, para todo preesquema localmente
noetheriano $T'$
\src{262}
sobre $S'$, la imagen inversa de $N$ sobre $T'$ pertenece a
$\operatorname{Im}A_\nu(T')$ si y solo si $T'\longrightarrow S'$ se factoriza
a través del subpreesquema $Z$. Cambiando las notaciones, podemos suponer que
$S'=S$; es decir, se nos da un cociente $N_\nu$ de $M_\nu$ por un sub-Módulo
$R_\nu$. Este proviene de un elemento de $A_\nu(S)$ si y solo si satisface las
dos condiciones siguientes:

\medskip
\noindent (i) $M_{\nu+k}/\mathcal S_kR_\nu$ es localmente libre de rango
$P(\nu+k)$ para $k\geq0$.

\smallskip
\noindent (ii) El subhaz $\mathcal K$ de $\mathcal F$ definido por el
sub-Módulo graduado
\[
R_* =\sum_{k\geq0}\mathcal S_kR_\nu
\]
de $M_*$ ([5], II, \S~3), y el cociente
$\mathcal G=\mathcal F/\mathcal K$, satisfacen las condiciones (a) y (b)
anteriores.

Estas condiciones son manifiestamente necesarias. Recíprocamente, si se
satisfacen, el haz $\mathcal G$ definido en (ii), al ser isomorfo al haz
asociado al $\mathcal S_*$-Módulo graduado $N_*$ que es la suma de los
$M_{\nu+k}/\mathcal S_kR_\nu$, es plano sobre $S$ (pues sus fibras son
sumandos directos de localizaciones de $N_*$\footnote{El texto impreso lleva $N$; aquí se localiza el Módulo graduado $N_*$ que acaba de definirse.} en ideales primos homogéneos de
$\mathcal S_*$), y tiene polinomio de Hilbert $P$ por (i). En virtud de (ii),
se satisfacen (a') y (a''), de modo que, para $n\geq\nu$, el homomorfismo
natural $N_n\longrightarrow f_*(\mathcal G(n))$ es un homomorfismo
sobreyectivo entre Módulos localmente libres del mismo rango y, por tanto, un
isomorfismo. Así,
\[
f_*(\mathcal K(\nu+k))=\mathcal S_kR_\nu
\quad\text{para todo $k\geq0$}.
\]
Esto demuestra que $\mathcal G\in A_\nu(S)$ y que $R_\nu$ es el elemento de
$\operatorname{Grass}_{P(\nu)}(M_\nu)$ definido por $\mathcal G$.

El criterio (i)--(ii) también se aplica después de un cambio de base
$S'\longrightarrow S$. Demostraremos primero que la condición (i) después del
cambio de base $S'\longrightarrow S$ puede expresarse exigiendo que
$S'\longrightarrow S$ se factorice a través de cierto subpreesquema $Z$ de
$S$. Una vez hecho esto, sustituir $S$ por $Z$ nos reduce al caso en que (i) ya
se satisface sobre $S$; puesto que es estable por cambio de base, queda por
expresar la condición (ii). Sea $U$ el conjunto de los $s\in S$ tales que la
cohomología en dimensión $>0$ de los haces inducidos sobre la fibra $X_s$ por
$\mathcal G(n)$ y $\mathcal K(n)$ se anula para $n\geq\nu$. Ya hemos observado
que $U$ es abierto, y la condición (ii) se satisface tras un cambio de base
$S'\longrightarrow S$ si y solo si $S'\longrightarrow S$ se factoriza a
través de $U$. Esto demuestra el Lema 3.3. Queda por demostrar el lema
siguiente:

\result{LEMA}{3.4} Sea $S$ un preesquema localmente noetheriano provisto de un
álgebra cuasicoherente graduada positivamente $\mathcal S_*$, generada por
$\mathcal S_1$; sea $M_*$ un $\mathcal S_*$-Módulo graduado cuasicoherente de
tipo finito, $P$ un polinomio con coeficientes racionales y $\nu$ un entero.
Existe entonces un subpreesquema manifiestamente único $Z$ de $S$
\src{263}
con la propiedad siguiente: para todo preesquema $S'$ sobre $S$, el Módulo
$M_n\otimes_{\mathcal O_S}\mathcal O_{S'}$ es localmente libre de rango $P(n)$
para todo $n\geq\nu$ si y solo si $S'\longrightarrow S$ se factoriza a través
de $Z$.

Podemos suponer claramente que $S$ es afín y, por tanto, noetheriano. Entonces:

\result{LEMA}{3.5} Para todo entero $N\geq\nu$, sea $U_N$ el subconjunto
abierto de $S$ formado por los $s\in S$ tales que
\[
\operatorname{rank}_{k(s)}
M_{ns}\otimes_{\mathcal O_{S,s}}k(s)\leq P(n)
\quad\text{para todo $\nu\leq n\leq N$}.
\]
Entonces la sucesión decreciente de abiertos $U_N$ es estacionaria.

Se sabe [3], IV, que $S$ posee una partición finita en subesquemas reducidos
$S_i$ tales que cada $M\otimes_{\mathcal O_S}\mathcal O_{S_i}$ es plano sobre
$S_i$. Podemos, pues, suponer que $M$ es plano y, por tanto, que los $M_n$ son
planos. Finalmente, podemos suponer claramente que $S$ es conexo. Existen
entonces un entero $n_0$ y un polinomio $Q$ tales que
\[
\operatorname{rank}_{k(s)}
M_{ns}\otimes_{\mathcal O_{S,s}}k(s)=Q(n)
\quad\text{para $n\geq n_0$},
\]
[5], III, \S~7. Supongamos primero que $P<Q$\footnote{El texto impreso lleva
$P\neq Q$; la fe de erratas de 1962 prescribe $P<Q$, sustituye el caso «$P=Q$» por
el caso complementario $P(n)\geq Q(n)$ para $n$ grande, y sustituye las
desigualdades estrictas de (*) por desigualdades no estrictas.}, de modo que
$P(n)<Q(n)$ para $n$ grande. Entonces, evidentemente, $U_N=\varnothing$ para
$N$ grande, y la sucesión $U_N$ es estacionaria. En el caso complementario,
$P(n)\geq Q(n)$ para $n$ grande, de modo que $U_N=U_{n_0}$ para $N\geq n_0$,
y la sucesión $U_N$ vuelve a ser estacionaria.

En particular, el conjunto $U_\infty$ de los $s\in S$ tales que
\[
\tag{*}
\operatorname{rank}_{k(s)}
M_{ns}\otimes_{\mathcal O_{S,s}}k(s)\leq P(n)
\quad\text{para todo $n\geq\nu$}
\]
es abierto, pues es la intersección de los $U_N$. Para demostrar (3.4),
podemos, por tanto, sustituir $S$ por el abierto $U_\infty$\footnote{El texto impreso lleva $U$; el abierto considerado aquí es $U_\infty$, la intersección de los $U_N$.}, reduciéndonos al caso en
que la desigualdad (*) se satisface en todo $s\in S$. Además:

\result{LEMA}{3.6} Sea $M$ un Módulo sobre un preesquema localmente
noetheriano $S$, y sea $r$ un entero. Existe entonces un subpreesquema
manifiestamente único $Z$ de $S$ con la propiedad siguiente: para todo $S'$
sobre $S$, el Módulo $M\otimes_{\mathcal O_S}\mathcal O_{S'}$ es localmente
libre de rango $r$ si y solo si $S'\longrightarrow S$ se factoriza a través de
$Z$. Si
\[
\operatorname{rank}_{k(s)}
M_s\otimes_{\mathcal O_{S,s}}k(s)\leq r
\quad\text{para todo $s$},
\]
entonces $Z$ es un subpreesquema cerrado de $S$ (suponiéndose $M$ coherente).
\src{264}
En efecto, el razonamiento anterior nos reduce al caso en que la desigualdad
(*) se satisface para todo $s\in S$ (sustituyendo $S$, si es necesario, por el
subconjunto abierto de los puntos donde se satisface). Podemos entonces suponer
que $M$ interviene en una sucesión exacta
\[
\mathcal O_S^q\longrightarrow\mathcal O_S^r\longrightarrow M\longrightarrow0.
\]
La condición considerada para $S'$ sobre $S$ significa también que, en la
sucesión exacta correspondiente
$\mathcal O_{S'}^q\longrightarrow\mathcal O_{S'}^r
\longrightarrow M'\longrightarrow0$, la segunda flecha es un isomorfismo; es
decir, que la primera es nula. Por consiguiente, el subpreesquema cerrado $Z$
de $S$ definido por el ideal generado por las entradas de la matriz que define
$\mathcal O_S^q\longrightarrow\mathcal O_S^r$ posee la propiedad requerida.

Volvamos ahora a la demostración de (3.4). Sea $Z_n$ el subesquema cerrado de
$S$ asociado por (3.6) al Módulo $M_n$ y al entero $r=P(n)$, y sea $Z'_N$ el
ínfimo de los $Z_n$ para $\nu\leq n\leq N$. Los $Z'_N$ forman una sucesión
decreciente de subesquemas cerrados de $S$\footnote{El texto impreso lleva $Z$; todos los $Z_n$ y sus ínfimos $Z'_N$ son subesquemas cerrados de $S$.} y, por tanto, es necesariamente
estacionaria. Sea $Z$ su valor constante para $N$ grande. Este es el $Z$
buscado en (3.4). Esto termina la demostración de (3.4) y, por consiguiente,
también la de (3.3).

Hemos demostrado, pues, que
$\underline{\operatorname{Quot}}^P_{\mathcal F/X/S}$ está representado por un
$S$-preesquema $Q$ que es unión creciente de subesquemas abiertos $Q_\nu$
cuasiproyectivos sobre $S$. Para ir más lejos, invoquemos el Teorema 2.1, del
cual se deduce fácilmente que $Q$ es cuasicompacto (pues es la imagen del
preesquema $S'$ de tipo finito sobre $S$ que parametriza la familia de haces
cocientes de los $\mathcal F_K$ con polinomio de Hilbert $P$). Así, $Q$ es
igual a uno de los $Q_\nu$ y, por tanto, es cuasiproyectivo sobre $S$. Para
demostrar que es proyectivo sobre $S$, queda por demostrar que es propio sobre
$S$. Para ello basta invocar el criterio valuativo de la propiedad de ser
propio en la forma
[5], II, 7.3.8. Solo hay que verificar:

\result{LEMA}{3.7} Sea $S$ el espectro de un anillo de valoración discreta,
sea $s$ su punto genérico, sea $X$ un preesquema sobre $S$, sea $\mathcal F$
un Módulo cuasicoherente sobre $X$, y sea $\mathcal G_s$ un Módulo
cuasicoherente cociente de
$\mathcal F_s=\mathcal F\otimes_{\mathcal O_S}k(s)$ sobre $X_s$. Existe
entonces un único Módulo cuasicoherente cociente $\mathcal G$ de $\mathcal F$,
plano sobre $S$, cuya restricción a $X_s$ es $\mathcal G_s$.
\src{265}
En efecto, si $\mathcal G_s=\mathcal F_s/\mathcal K_s$, basta tomar el mayor
subhaz $\mathcal K$ de $\mathcal F$ que induce $\mathcal K_s$, [5], I, 9.4.2,
y poner $\mathcal G=\mathcal F/\mathcal K$. Se comprueba fácilmente que este
haz posee las propiedades requeridas.

El Teorema 3.2 y, por consiguiente, (3.1), quedan ahora completamente
demostrados.

La demostración establece al mismo tiempo:

\result{PROPOSICIÓN}{3.8} Bajo las hipótesis de (3.2), sean
$Q=\operatorname{Quot}^P_{\mathcal F/X/S}$, $X_Q=X\times_SQ$ y
$\mathcal F_Q=\mathcal F\otimes_{\mathcal O_S}\mathcal O_Q$; y sea
$\mathcal G$ el cociente coherente de $\mathcal F_Q$, plano sobre $Q$, con
polinomio de Hilbert relativo $P$, tal que $(Q,\mathcal G)$ representa el
funtor $\underline{\operatorname{Quot}}^P_{\mathcal F/X/S}$. Existe un entero
$\nu$ tal que, para $n\geq\nu$, $(f_Q)_*(\mathcal G(n))$ es un Módulo sobre
$Q$, localmente libre de rango $P(n)$ y muy amplio relativamente a $S$; es
decir, define una inmersión de $Q$ en un esquema de Grassmann
$\operatorname{Grass}_{P(n)}(M)$ sobre $S$. En particular, para $n\geq\nu$,
el haz
\[
\bigwedge^{P(n)}(f_Q)_*(\mathcal G(n))
\]
sobre $Q$ es invertible y muy amplio relativamente a $S$.

Como en (3.2), podemos reducirnos al caso en que $\mathcal F$ es plano sobre
$S$; entonces basta elegir un entero $\nu$ tal que $A_\nu=A$ (con las
notaciones anteriores).

La aplicación más importante de (3.2) es el caso
$\mathcal F=\mathcal O_X$. Escribimos entonces
\[
\operatorname{Quot}_{\mathcal O_X/X/S}=\operatorname{Hilb}_{X/S},
\qquad
\operatorname{Quot}^P_{\mathcal O_X/X/S}=\operatorname{Hilb}^P_{X/S},
\]
y obtenemos una descomposición
\[
\operatorname{Hilb}_{X/S}=\coprod_P\operatorname{Hilb}^P_{X/S}.
\]
Por definición, $\operatorname{Hilb}_{X/S}$ representa el funtor
$\underline{\operatorname{Hilb}}_{X/S}(S')$, conjunto de los subesquemas
cerrados de $X'=X\times_SS'$ que son planos sobre $S'$; y
$\operatorname{Hilb}^P_{X/S}$ representa el subfuntor correspondiente de
subesquemas cerrados con polinomio de Hilbert dado $P$. Estos preesquemas se
llaman también, respectivamente, \emph{preesquema de Hilbert} de $X$ sobre
$S$ y \emph{preesquema de Hilbert de índice $P$}. La terminología se justifica
por el papel desempeñado en la teoría por los polinomios de Hilbert. Su
diferencia de naturaleza respecto de las variedades de Chow clásicas (que
parametrizan ciclos y no variedades) es del mismo orden que la existente entre
\src{266}
el anillo de Chow de las clases de ciclos sobre una variedad y el anillo de
clases de haces sobre la variedad (tal como aparece en el teorema de
Riemann--Roch [1]). En efecto, cuando $X=\mathbf P^r_S$ y $S$ es el espectro
de un cuerpo, conocer el polinomio de Hilbert de un Módulo coherente
$\mathcal F$ sobre $X$ equivale también a conocer las clases de Chern de
$\mathcal F$, o la clase de $\mathcal F$ en el anillo de clases de haces
coherentes sobre $X$.

\result{OBSERVACIONES}{3.9} Obsérvese asimismo que la construcción de
$\operatorname{Quot}_{\mathcal F/X/S}$ y
$\operatorname{Quot}^P_{\mathcal F/X/S}$ se redujo al caso
$X=\mathbf P^r_S$, $\mathcal F=\mathcal O_X^N$, con $\mathcal O_X(1)$ el haz
muy amplio usual; más precisamente, los
$\operatorname{Quot}_{\mathcal F/X/S}$ generales aparecen como subesquemas
cerrados de estos. Como la formación de los
$\operatorname{Quot}^P_{\mathcal F/X/S}$ es manifiestamente compatible con
los cambios de base $S'\longrightarrow S$, nos reducimos además a
$S=\operatorname{Spec}(\mathbf Z)$. En definitiva, pues, se reduce al estudio
de los esquemas proyectivos siguientes sobre $\mathbf Z$:
\[
\mathfrak Q^P_{r,N}=
\operatorname{Quot}^{P}_{(\mathcal O_{\mathbf P^r_{\mathbf Z}})^N/
\mathbf P^r_{\mathbf Z}/\operatorname{Spec}(\mathbf Z)},
\]
y, más particularmente, de los esquemas de Hilbert absolutos
\[
\operatorname{Hilb}^P_r=\mathfrak Q^P_{r,1}.
\]
Sería muy interesante un estudio más detallado de estos esquemas, comenzando
por la determinación de sus componentes conexas (¿son conexos?),\footnote{La
fe de erratas de 1962 señala que Hartshorne demostró la conexidad de los
$\operatorname{Hilb}^P_r$ sobre un cuerpo algebraicamente cerrado y determinó
los pares $(r,P)$ para los cuales
$\operatorname{Hilb}^P_r\neq\varnothing$.} y de sus componentes irreducibles
(por Serre [7], puede haber componentes irreducibles situadas enteramente sobre
un primo no nulo $p$). Recordemos la pregunta de Weil de si las componentes
irreducibles de las fibras de $\operatorname{Hilb}^P_r$ sobre
$s\in\operatorname{Spec}(\mathbf Z)$ corresponden a extensiones «regulares»
del cuerpo primo, es decir, si son «relativamente conexas». Estas cuestiones
quizá sean más abordables para los esquemas de Hilbert que para las
«variedades de Chow».

\fsection{4}{Variantes}

\noindent a. Bajo las hipótesis de (3.1), sea $U$ un subconjunto abierto de
$X$, y sea $A'$ el subfuntor de
$A=\underline{\operatorname{Quot}}_{\mathcal F/X/S}$ tal que $A'(S')$ es el
conjunto de Módulos cocientes $\mathcal G$ de $\mathcal F'$, planos sobre
$S'$, cuyo soporte está contenido en $U'$. Manifiestamente, $A'$ está
representado por un subconjunto abierto del preesquema
$\operatorname{Quot}_{\mathcal F/X/S}$ que representa $A$. De ello resulta
que los Teoremas (3.1) y (3.2)
\src{267}
siguen siendo válidos cuando se supone que $X$ es cuasiproyectivo sobre $S$ en
lugar de proyectivo sobre $S$, siempre que en las conclusiones también se
sustituya «proyectivo» por «cuasiproyectivo» y que
$\underline{\operatorname{Quot}}_{\mathcal F/X/S}(S')$ denote ahora el
conjunto de cocientes coherentes $\mathcal G$ de $\mathcal F'$, planos sobre
$S'$, cuyo soporte es propio sobre $S'$.

\medskip
\noindent b. Más generalmente, pueden imponerse a los cocientes $\mathcal G$
de $\mathcal F'=\mathcal F\otimes_{\mathcal O_S}\mathcal O_{S'}$, planos
sobre $S'$, toda clase de condiciones naturales adicionales estables por
cambio de base, obteniéndose así subfuntores de
$\underline{\operatorname{Quot}}_{\mathcal F/X/S}$ cuya representabilidad
puede buscarse. En muchos casos, el criterio usual demuestra que estos
funtores vuelven a estar representados por subconjuntos abiertos de
$\operatorname{Quot}_{\mathcal F/X/S}$. Esto ocurre, en particular, cuando se
impone una de las propiedades adicionales siguientes:

\smallskip
\noindent $1^\circ$ Las dimensiones de los ciclos primos asociados a los
Módulos $\mathcal G_{s'}$ ($s'\in S'$) inducidos sobre las fibras $X'_{s'}$
pertenecen a un conjunto dado de enteros.

\smallskip
\noindent $2^\circ$ (Cuando $\mathcal F=\mathcal O_X$, de modo que
$\mathcal G$ corresponde a un subpreesquema cerrado $Y$ de $X'$): $Y$ es un
preesquema liso ([3], IV) sobre $S'$, respectivamente normal sobre $S'$ (es
decir, las fibras $Y_{s'}$ son normales «sobre $k(s')$», lo que significa
normales después de toda extensión del cuerpo base), respectivamente (cuando
$X$ es plano sobre $S$) una intersección completa local relativa en $X'$ sobre
$S'$ (es decir, las fibras $Y_{s'}$ son intersecciones completas locales en los
$X'_{s'}$).\footnote{El texto impreso omite las primas en estas propiedades de $Y$ sobre la base variable $S'$ y dentro de $X'$. La hipótesis relativa al $X$ original plano sobre $S$ permanece sin primas.}

Otras condiciones harían intervenir propiedades cohomológicas de los Módulos
$\mathcal G_{s'}$ inducidos sobre los $X'_{s'}$, etc. Por supuesto, una
conjunción de condiciones, cada una de ellas representada por un subconjunto
abierto $U_i$ de $\operatorname{Quot}_{\mathcal F/X/S}$, está representada por
la intersección de estos subconjuntos abiertos. Por ejemplo, asignar a todo
$S'$ sobre $S$ el conjunto de subesquemas cerrados $Y$ de
$X'=X\times_SS'$ que son recubrimientos étale ([3], I) de $S'$ de rango dado
$r$ da un funtor contravariante representable de $S'$.

\medskip
\noindent c. Los preesquemas $\operatorname{Hom}_S(X,Y)$,
$\Pi_{X/S}(Z/X)$ y $\operatorname{Isom}_S(X,Y)$ definidos en [2], II, C,
n.º 2 existen bajo hipótesis de proyectividad convenientes y aparecen como
subconjuntos abiertos de preesquemas de Hilbert adecuados. Puesto que
\[
\operatorname{Hom}_S(X,Y)=\Pi_{X/S}((X\times_SY)/X),
\]
el caso de $\operatorname{Hom}_S(X,Y)$ se reduce al de
$\Pi_{X/S}(Z/X)$. Para todo $S'$ sobre $S$, el conjunto de secciones de
$Z'=Z\times_SS'$ sobre $X'=X\times_SS'$ está en correspondencia biyectiva con
el conjunto de subpreesquemas $\Gamma$ de $Z'$ (necesariamente cerrados si $Z$
es separado sobre $X$) tales que el morfismo $\Gamma\longrightarrow X'$
inducido por $Z'\longrightarrow X'$ sea un isomorfismo. De este modo,
\src{268}
cuando $X$ es plano y propio sobre $S$ y $Z$ es cuasiproyectivo sobre $S$,
$\Pi_{X/S}(Z/X)$ existe y aparece como subpreesquema abierto de
$\operatorname{Hilb}_{Z/S}$. Por consiguiente, cuando $X$ es proyectivo y
plano sobre $S$ e $Y$ es cuasiproyectivo sobre $S$,
$\operatorname{Hom}_S(X,Y)$ existe y aparece como subpreesquema abierto de
$\operatorname{Hilb}_{(X\times_SY)/S}$. Cuando $X$ e $Y$ son ambos
proyectivos sobre $S$, se sigue inmediatamente que
$\operatorname{Isom}_S(X,Y)$ también existe y está representado por un
subconjunto abierto de $\operatorname{Hom}_S(X,Y)$. Análogamente, cuando $X$
es plano y proyectivo sobre $S$ e $Y$ es cuasiproyectivo sobre $S$, el
$S$-preesquema $\operatorname{Imm}_S(X/Y)$ correspondiente al subfuntor de
$\operatorname{Hom}_S(X,Y)$ cuyos elementos son los $S'$-homomorfismos
$X'\longrightarrow Y'$ que son inmersiones también está representado por un
subconjunto abierto de $\operatorname{Hom}_S(X,Y)$.

Sea $\mathcal L$ (respectivamente, $\mathcal M$) un haz invertible sobre $X$
(respectivamente, $Y$), muy amplio relativamente a $S$. Entonces
$\mathcal L\otimes_{\mathcal O_S}\mathcal M$ es un haz muy amplio sobre
$X\times_SY$ relativamente a $S$. Por consiguiente, para todo polinomio $P$
con coeficientes racionales, $\operatorname{Hilb}^P_{(X\times_SY)/S}$ está
definido y es un preesquema cuasiproyectivo sobre $S$. Induce, por tanto,
sobre $\operatorname{Hom}_S(X,Y)$ un subconjunto abierto y cerrado,
cuasiproyectivo sobre $S$, que denotaremos por
$\operatorname{Hom}_S(X,Y)^P$. Así, las secciones de
$\operatorname{Hom}_S(X,Y)^P$ sobre $S$ son los $S$-morfismos
$g:X\longrightarrow Y$ tales que, para todo entero $n$,
\[
\chi\!\left(\left(\mathcal L\otimes_{\mathcal O_X}g^*(\mathcal M)\right)^{
\otimes n}\right)=P(n).
\]

De este modo se obtienen generalizaciones del teorema de Matsusaka que afirma
que los automorfismos de una variedad proyectiva «polarizada» forman un grupo
algebraico. Aquí la afirmación tiene un sentido manifiestamente más preciso,
pues este grupo se define como solución de un problema universal. Obsérvese
también que, sobre un cuerpo algebraicamente cerrado, el grupo de
automorfismos considerado anteriormente se obtiene del «verdadero» grupo
definido aquí pasando al cociente por los elementos nilpotentes. Esto explica
por qué las construcciones antiguas tienen pocas posibilidades de funcionar
sobre un cuerpo base no perfecto: el ideal de elementos nilpotentes que
aparece después de una extensión del cuerpo base no tiene por qué estar
«definido sobre $k$». La misma observación se aplica a la mayoría de las
construcciones de estilo antiguo.

\fsection{5}{Estudio infinitesimal de los esquemas de Hilbert}

Se deduce del resultado siguiente:

\src{269}
\result{PROPOSICIÓN}{5.1} Sea $S$ un preesquema, sea $S_0$ un subpreesquema
definido por un ideal cuasicoherente $\mathcal J$ de cuadrado nulo, sea $X$ un
$S$-preesquema y sea $\mathcal F$ un Módulo cuasicoherente sobre $X$. Pongamos
$X_0=X\times_SS_0$ y
$\mathcal F_0=\mathcal F\otimes_{\mathcal O_S}\mathcal O_{S_0}$, y sea
$\mathcal G_0=\mathcal F_0/\mathcal K_0$ un Módulo cuasicoherente cociente de
$\mathcal F_0$, plano sobre $S_0$. Para todo subconjunto abierto $U$ de $X$,
sea $\mathcal E(U)$ el conjunto de Módulos cuasicoherentes cocientes
$\mathcal G$ de $\mathcal F|U$, planos sobre $S$, tales que
$\mathcal G\otimes_{\mathcal O_S}\mathcal O_{S_0}=\mathcal G_0$. Al variar
$U$, los $\mathcal E(U)$ son, pues, las secciones de un haz $\mathcal E$ sobre
$X$.\footnote{El texto impreso lleva $U$; al variar los abiertos $U$ de $X$, los conjuntos $\mathcal E(U)$ definen un haz sobre $X$.} El haz de grupos
\[
\mathcal A=\operatorname{Hom}_{\mathcal O_{X_0}}
(\mathcal K_0,\mathcal G_0\otimes_{\mathcal O_{S_0}}\mathcal J)
\]
actúa naturalmente sobre $\mathcal E$, que se convierte así en un haz
«formalmente homogéneo principal bajo $\mathcal A$» (es decir, para todo
abierto $U$ de $X$, $\mathcal E(U)$ es vacío o es un conjunto homogéneo
principal bajo $\mathcal A(U)$).

De ello se sigue que:

\result{COROLARIO}{5.2} Supongamos que localmente sobre $X$ existe una
extensión $\mathcal G$ de $\mathcal G_0$ a un cociente de $\mathcal F$ plano
sobre $S$ (es decir, que las fibras del haz $\mathcal E$ no son vacías).
Existe entonces una clase de obstrucción canónica
\[
c(\mathcal G_0)\in H^1(X,\mathcal A)
\]
cuya anulación es necesaria y suficiente para que exista una extensión global
$\mathcal G$ de $\mathcal G_0$ a un cociente de $\mathcal F$ plano sobre $S$.
Si esta clase se anula, el conjunto $\mathcal E(X)$ de todas las extensiones
posibles es un conjunto homogéneo principal bajo
\[
\mathcal A(X)=\operatorname{Hom}_{\mathcal O_{X_0}}
(\mathcal K_0,\mathcal G_0\otimes_{\mathcal O_{S_0}}\mathcal J).
\]
\footnote{En esta última fórmula, el texto impreso lleva $\mathcal O_X$ en
lugar de $\mathcal O_{X_0}$, aunque la definición inmediatamente anterior de
$\mathcal A$ fija $\mathcal O_{X_0}$.}
En particular, existe una extensión global si $H^1(X,\mathcal A)=0$.

\result{COROLARIO}{5.3} Supongamos que
$Q=\operatorname{Quot}_{\mathcal F/X/S}$ existe (cf. 4 (a)); por ejemplo, que
$X$ es cuasiproyectivo sobre un $S$ localmente noetheriano y que $\mathcal F$
es coherente. Sea $x\in Q$ correspondiente a una extensión residual
$K=k(x)$ de $k(s)$ ($s\in S$):
\src{270}
así, $x$ está definido por un Módulo coherente cociente
$\mathcal G_0=\mathcal F_0/\mathcal K_0$ del Módulo
$\mathcal F_0=\mathcal F_K$ sobre el $K$-preesquema $X_K$. Sea $\mathcal A$
el haz coherente sobre $X_K$ definido por
\[
\mathcal A=\operatorname{Hom}_{\mathcal O_{X_0}}(\mathcal K_0,\mathcal G_0).
\]
Entonces el espacio tangente de Zariski a la fibra $Q_s$ en $x$ (el dual
sobre $K$ de $\mathfrak m/\mathfrak m^2$, donde $\mathfrak m$ es el ideal
maximal de $\mathcal O_{Q_K,x}$) es canónicamente isomorfo a
$H^0(X_K,\mathcal A)$.

La descripción del espacio tangente de Zariski admite la generalización
siguiente. Dado un $S$-morfismo $g:S'\longrightarrow Q$, es decir, una sección
$g'$ de $Q'=Q\times_SS'$ sobre $S'$, el Módulo
\[
\Omega=g^*(\Omega^1_{Q/S})=g'^*(\mathcal J/\mathcal J^2)
\]
(donde $\mathcal J$ es el ideal sobre $Q'$ definido por la sección $g'$ de
$Q'$ sobre $S'$) está caracterizado por la fórmula siguiente, funtorial en el
Módulo coherente $\mathcal M$ sobre $S'$:
\[
\operatorname{Hom}_{\mathcal O_{S'}}(\Omega,\mathcal M)
\simeq H^0(X',\mathcal A),
\]
donde $\mathcal A$ es el Módulo sobre $X'=X\times_SS'$ definido por
\[
\mathcal A=\operatorname{Hom}_{\mathcal O_{X'}}
(\mathcal K,\mathcal G\otimes_{\mathcal O_{S'}}\mathcal M).
\]
\footnote{El texto impreso pone $\mathcal O_S$ en este producto tensorial;
aquí $\mathcal G$ y $\mathcal M$ son Módulos sobre $S'$, y la construcción de
deformación de cuadrado nulo exige el producto tensorial sobre
$\mathcal O_{S'}$.}
Aquí $\mathcal G=\mathcal F'/\mathcal K$ es el Módulo cociente de
$\mathcal F'=\mathcal F\otimes_{\mathcal O_S}\mathcal O_{S'}$ que corresponde
a $g$. En efecto, aplíquese (5.1) sustituyendo $S_0$ por $S'$ y $S$ por el
preesquema $D(\mathcal M)=(S',\mathcal O_{S'}+\mathcal M)$, donde
$\mathcal M$ se considera como un ideal de cuadrado nulo.

Cuando $\mathcal F=\mathcal O_X$ en (5.1), el dato de $\mathcal G_0$
corresponde a un subpreesquema cerrado $Y_0$ de $X_0$, plano sobre $S_0$,
definido por el ideal $\mathcal J_0=\mathcal K_0$. La fórmula (*) da entonces
\[
\mathcal A=\operatorname{Hom}_{\mathcal O_{X_0}}
(\mathcal J_0/\mathcal J_0^2,
\mathcal O_{Y_0}\otimes_{\mathcal O_{S_0}}\mathcal J),
\]
donde $\mathcal J_0/\mathcal J_0^2$\footnote{El texto impreso omite el subíndice $0$ en las dos menciones del haz conormal y en $\mathcal J_0^{-1}$ más abajo. El Ideal de $Y_0$ en $X_0$ es $\mathcal J_0$; el Ideal $\mathcal J$ de la base es distinto.} se interpreta como el haz conormal de $Y_0$
en $X_0$, denotado también por
\src{271}
$\check{\mathcal N}_{Y_0/X_0}$. Conviene entonces considerar $\mathcal A$
como un Módulo sobre $Y_0$ y calcular $H^0$ y $H^1$ sobre $Y_0$. Además, si
$Y_0$ es localmente una intersección completa en $X_0$, con $X$ plano sobre
$S$, queda garantizada la existencia de extensiones locales en (5.1); asimismo,
$\mathcal J_0/\mathcal J_0^2$ es localmente libre sobre $Y_0$, y puede escribirse
\[
\mathcal A=\mathcal N_{Y_0/X_0}\otimes_{\mathcal O_{S_0}}\mathcal J,
\]
donde el primer factor de la derecha es el haz normal de $Y_0$ en $X_0$.
Utilizando el criterio fundamental de lisura ([3], III, 3.1), se obtiene, por
ejemplo:

\result{COROLARIO}{5.4} Bajo las hipótesis de (5.3), supongamos que
$\mathcal F=\mathcal O_X$, que $X$ es plano sobre $S$ y que el subpreesquema
cerrado $Y_0$ de $X_0$ correspondiente a $\mathcal G_0$ es localmente una
intersección completa. Entonces el espacio tangente de Zariski a $Q_s$ en $x$
es canónicamente isomorfo a $H^0(Y_0,\mathcal N_{Y_0/X_0})$. Si
$H^1(Y_0,\mathcal N_{Y_0/X_0})=0$, entonces el preesquema de Hilbert
$\operatorname{Hilb}_{X/S}$ es liso sobre $S$ en $x$.
($\mathcal N_{Y_0/X_0}$ es el haz normal de $Y_0$ en $X_0$.)

\result{OBSERVACIÓN}{5.5} Este enunciado se aplica, en particular, cuando
$Y_0$ es una intersección completa en $X_0$ definida por una ecuación, es
decir, un «divisor de Cartier» positivo. Entonces $\mathcal N_{Y_0/X_0}$ es
isomorfo al haz sobre $Y_0$ inducido por el haz invertible
$\mathcal J_0^{-1}$ sobre $X_0$ definido por el divisor $Y_0$. Esta situación
se presenta especialmente en el estudio de familias de divisores positivos
sobre una variedad proyectiva no singular $X_0$. El isomorfismo entre el
espacio tangente de Zariski en el punto $x$ de $Q$ (o, equivalentemente, del
subconjunto abierto $D$ de $Q$ correspondiente a los divisores) y
$H^0(Y_0,\mathcal N_{Y_0/X_0})$ era conocido en geometría algebraica clásica
como el «homomorfismo característico» del primero al segundo. Solo estaba
definido cuando $x$ era un punto liso de la variedad de parámetros $T$ de una
«familia continua completa» de divisores, es decir, desde nuestro punto de
vista, de una componente irreducible del esquema $D$ provisto de su estructura
reducida inducida. El
espacio tangente a $T$ en $x$ es entonces un subespacio del espacio tangente a
$D$ en $x$, por lo que el homomorfismo característico clásico es, en efecto,
inyectivo, pero solo es sobreyectivo bajo condiciones adicionales; por ejemplo,
cuando $D$ es íntegro en $x$. De hecho, Zappa [8] construyó un ejemplo
\src{272}
(con $X$ una superficie proyectiva no singular sobre los números complejos) en
el que el homomorfismo característico no es sobreyectivo ni siquiera en el
punto genérico de $T$. Así, $D$ no es íntegro ni siquiera en el punto genérico
de la componente irreducible considerada. Esto demuestra de manera
especialmente llamativa hasta qué punto las variedades con elementos
nilpotentes son necesarias para comprender fenómenos de la teoría más clásica
de superficies.\footnote{La fe de erratas de 1962 añade que Mumford acababa de
construir una componente irreducible del esquema de Hilbert de
$\mathbf P^3_{\mathbf C}$ cuyos puntos generales representan curvas no
singulares de grado 14 y género 24, que es no reducida en su punto genérico;
al hacer explotar una de estas curvas se obtiene también una variedad triple
proyectiva regular cuyo esquema formal de módulos es genéricamente no
reducido.}

\subsection*{5.6}

En (5.4) dimos un criterio de lisura aplicable, en particular, a esquemas de
divisores. Kodaira dio en [6] un criterio diferente, a saber, la anulación de
$H^1(X_0,\mathcal L)$, donde $\mathcal L=\mathcal J_0^{-1}$ es el haz
invertible sobre $X_0$ definido por el divisor $Y_0$. Este criterio es válido
cuando $S$ es el espectro de un cuerpo de característica cero, y fue demostrado
en [6] por métodos trascendentes cuando el cuerpo base es $\mathbf C$. Más
generalmente, siendo de nuevo $S$ arbitrario, la condición de Kodaira basta
para que el morfismo canónico
\[
D\longrightarrow\operatorname{Pic}_{X/S}
\]
del preesquema de divisores al preesquema de Picard de $X/S$ sea liso en el
punto $x$ considerado (como se deduce fácilmente del criterio usual de lisura
una vez conocida la existencia de $\operatorname{Pic}_{X/S}$). Si, además,
$\operatorname{Pic}_{X/S}$ es liso sobre $S$ en la imagen de $x$ (por ejemplo,
si $\operatorname{Pic}_{X/S}$ es liso sobre $S$), entonces $D$ es liso sobre
$S$ en $x$. Por otra parte, Cartier demostró que todo preesquema en grupos
localmente de tipo finito sobre un cuerpo $k$ de característica cero es liso
sobre $k$. Combinando estos dos resultados se recupera el resultado de
Kodaira. Estas observaciones demuestran también que, sobre un cuerpo $K$ de
característica $p>0$, cuando $\operatorname{Pic}_{X/S}$ no es liso sobre $K$\footnote{El texto impreso lleva $k$; el cuerpo base de esta frase en característica positiva es $K$.}
(como cuando $X$ es la superficie de Igusa), la condición
$H^1(X_0,\mathcal L)=0$ implica, por el contrario, que $D$ no es liso en $x$,
y que incluso es no reducido en $x$ si $K$ es algebraicamente cerrado.

Finalmente, registremos el resultado siguiente, que desempeña un papel
importante en el estudio infinitesimal de los espacios fibrados:

\result{PROPOSICIÓN}{5.7} Sea $X$ un preesquema finito y plano sobre un $S$
localmente noetheriano, y sea $Z$ un preesquema sobre $X$ tal que
$\Pi_{X/S}(Z/X)$ exista (como ocurre si $Z$ es cuasiproyectivo sobre $X$). Si
$Z$ es liso sobre $X$, entonces $\Pi_{X/S}(Z/X)$ es liso sobre $S$.

Esto resulta inmediatamente de la definición y del criterio usual de lisura
([3], III, 3.1). Obsérvese que, cuando $X$ es finito y plano sobre $S$, la
existencia de $\Pi_{X/S}(Z/X)$\footnote{El texto impreso lleva aquí $Z/K$; no
se ha definido ningún $K$, y el funtor de secciones considerado desde (5.7) es
$\Pi_{X/S}(Z/X)$.} puede tratarse de manera muy elemental, sin
\src{273}
utilizar la teoría de los esquemas de Hilbert. Por ejemplo, si $X$ es radicial
sobre $S$, entonces $\Pi_{X/S}(Z/X)$ existe sin ninguna restricción sobre $Z$.
Como otro ejemplo, sea $T$ un $S$-preesquema y sea $T_n$ el «entorno
infinitesimal de orden $n$» de la diagonal de $T\times_ST$ en $T\times_ST$,
provisto de los morfismos $p_1,p_2:T_n\longrightarrow T$ inducidos por las dos
proyecciones. Consideremos $T_n$ como un preesquema finito sobre $T$ por medio
de $p_1$, y supongamos además que es plano sobre $T$ (como ocurre cuando $T$
es liso sobre $S$). Para todo preesquema $X$ sobre $T$, pongamos
\[
(X/T/S)^{(n)}=
\Pi_{T_n/T}\bigl(p_2^*(X/T)/T_n\bigr).
\]
\footnote{El texto impreso coloca $T_n/S$ bajo el símbolo $\Pi$. La
Proposición 5.7 se aplica aquí al morfismo finito y plano
$p_1:T_n\to T$, y el resultado es explícitamente un preesquema sobre $T$; por
tanto, el denominador es $T_n/T$.}
Este es un preesquema sobre $T$ llamado \emph{fibrado de gérmenes de secciones
de orden $n$ de $X$ sobre $T$} (relativo a $S$). Depende funtorialmente de $X$
y es liso sobre $T$ si $X$ lo es.

\fsection{6}{Relaciones con la norma y los productos simétricos}

Sea $S$ un preesquema, sean $X$ e $Y$ $S$-preesquemas, y sea
\[
u:(X/S)^n\longrightarrow Y
\]
un $S$-morfismo simétrico de la $n$-ésima potencia cartesiana de $X/S$ en
$Y$. Para simplificar, supongamos que $S$ es localmente noetheriano y que $X$
e $Y$ son de tipo finito sobre $S$. A todo Módulo coherente $\mathcal F$ sobre
$X$ cuyo soporte es finito sobre $S$, que es plano sobre $S$ y tiene rango
relativo $n$ sobre $S$ (es decir, $f_*(\mathcal F)$ es un Módulo localmente
libre de rango $n$ sobre $S$), puede asociársele entonces de manera natural una
sección de $Y$ sobre $S$:
\[
\pi^u_{X/S}(\mathcal F)\in\Gamma(Y/S).
\]
No daremos aquí la definición formal, sino que nos limitaremos a observar que
el formalismo resultante es una generalización natural del formalismo usual de
normas y trazas. Para todo $X/S$ separado, la ley de norma tiene como blanco canónico el espacio de potencias divididas $\Gamma^n_S(X)$. Si $X$ es plano sobre $S$ y existe la $n$-ésima potencia simétrica (por ejemplo, cuando las órbitas de $\mathfrak S_n$ en $(X/S)^n$ están contenidas en abiertos afines), el morfismo canónico $\operatorname{Symm}^n_S(X)\to\Gamma^n_S(X)$ es un isomorfismo, y se obtiene un elemento canónico
\[
\pi_{X/S}(\mathcal F)\in\Gamma(\operatorname{Symm}^n_S(X)/S),
\]
\src{274}
del cual pueden recuperarse los $\pi^u_{X/S}(\mathcal F)$. Otro caso
importante es aquel en que $X$ es un monoide conmutativo sobre $S$, $X=Y$, y
$u$ procede de la ley de composición de $X$. Escribimos entonces simplemente
$\pi(\mathcal F)$ para la sección de $X$ sobre $S$ asociada al Módulo
$\mathcal F$ sobre $X$.

Supongamos ahora que $\mathcal F$ es un Módulo coherente sobre $X$ tal que
$\operatorname{Quot}_{\mathcal F/X/S}$ existe, o al menos tal que el funtor
$\underline{\operatorname{Quot}}^n_{\mathcal F/X/S}$ que asigna a todo $S'$
sobre $S$ el conjunto de haces cocientes coherentes $\mathcal M$ de
$\mathcal F'=\mathcal F\otimes_{\mathcal O_S}\mathcal O_{S'}$ que son planos
sobre $S'$ y tienen rango relativo $n$ está representado por un $S$-preesquema
$\operatorname{Quot}^n_{\mathcal F/X/S}$. (Cuando $X$ es cuasiproyectivo
sobre $S$, $\operatorname{Quot}^n_{\mathcal F/X/S}$ existe y, con las
notaciones del Apartado 3, es $\operatorname{Quot}^P_{\mathcal F/X/S}$, donde
$P$ es el polinomio constante $n$.) Como la formación de
$\pi^u_{X/S}(\mathcal M)$ es compatible con el cambio de base, se obtiene un
morfismo canónico
\[
\pi^u_{X/S}:\operatorname{Quot}^n_{\mathcal F/X/S}\longrightarrow Y,
\]
y, en general,
\[
\pi_{X/S}:\operatorname{Quot}^n_{\mathcal F/X/S}
\longrightarrow\Gamma^n_S(X).
\]
Cuando $X$ es plano sobre $S$, identificamos este blanco con $\operatorname{Symm}^n_S(X)$.
Independientemente de la identificación precedente, cuando $X$ es separado sobre $S$, la construcción mediante la familia finita y plana para $\mathcal F=\mathcal O_X$ da un morfismo
\[
\pi_{X/S}:\operatorname{Hilb}^n_{X/S}
\longrightarrow\operatorname{Symm}^n_S(X).
\]
Este es manifiestamente un isomorfismo para $n=0$ o $n=1$. Pero para
$n>1$\footnote{El texto impreso repite $n\geq1$, contradiciendo el
isomorfismo que acaba de afirmarse para $n=1$.}, incluso cuando $S$ es el
espectro de un cuerpo $k$ y $X$ es liso sobre $S$, no es en general un
isomorfismo, ni siquiera un morfismo inyectivo: un subesquema de dimensión cero
de $X$ (correspondiente, por ejemplo, a un ideal $\mathfrak i$ primario para el
ideal maximal en un anillo local $\mathcal O_{X,x}$, para un punto cerrado $x$
de $X$) no está determinado por el ciclo que define (en este caso, por la
codimensión sobre $k$ de $\mathfrak i$ en $\mathcal O_{X,x}$). Solo puede
afirmarse lo siguiente, siendo de nuevo $S$ arbitrario:

\medskip
\noindent a. Si $X$ es liso sobre $S$, el morfismo norma induce un isomorfismo
del subconjunto abierto de $\operatorname{Hilb}^n_{X/S}$ que clasifica los
recubrimientos étale de rango $n$ contenidos en $X$ (cf. n.º 4, (b)) sobre el
subconjunto abierto de $\operatorname{Symm}^n_S(X)$ correspondiente a los
$n$-ciclos sin componentes múltiples.
\src{275}

\medskip
\noindent b. Si, además, $X$ tiene dimensión relativa 1 sobre $S$, entonces el
morfismo norma es un isomorfismo de $\operatorname{Hilb}^n_{X/S}$ sobre
$\operatorname{Symm}^n_{X/S}$.

Esto se debe a que un subesquema de dimensión cero de una curva algebraica no
singular está determinado por su ciclo correspondiente. La misma observación
se aplica, más generalmente, a los divisores de Cartier positivos sobre un
esquema algebraico no singular (y es posible que, en este caso tan especial,
la variedad de Chow dé el mismo objeto que el esquema de Hilbert).

\fsection{7}{Complementos y cuestiones}

Como observó J.-P. Serre, un ejemplo bien conocido de Nagata demuestra que
pueden encontrarse un esquema $S$, espectro de un cuerpo $k$; un $S$-esquema
$S'$, espectro de una extensión cuadrática $k'$ de $k$; y un $S'$-esquema
propio y liso (pero no proyectivo) $X$ de dimensión 3 tal que
$\Pi_{S'/S}(X/S')$ no existe. Se sigue, a fortiori, que el esquema de
Hilbert $\operatorname{Hilb}^2_{X/S}$ no existe (ni siquiera el $k$-esquema
que representaría los recubrimientos étale de $S$ de rango 2 contenidos en
$X$, y a fortiori tampoco el cuadrado simétrico de $X$; cf. el apartado
anterior). Esto impone serias limitaciones a la posibilidad de realizar
construcciones no proyectivas en geometría algebraica. (Sin embargo, es
plausible que tales limitaciones no aparezcan en geometría analítica, como
tampoco aparecen en geometría formal (cf. [2], II).) Por otra parte, si $X$ es
un esquema propio sobre $S=\operatorname{Spec}k$ y $Z$ es cuasiproyectivo
sobre $X$, entonces $\Pi_{X/S}(Z/X)$ existe y es un esquema que es unión
disjunta de una sucesión de esquemas cuasiproyectivos sobre $S$ (como en el
caso proyectivo (3.1)). En efecto, se reduce al caso en que el propio $X$ es
proyectivo dominando $X$ por un $S$-esquema proyectivo $X'$. No daremos los
detalles de la demostración, que utiliza también el resultado de factorización
para un morfismo finito recogido en [2], I, A, n.º 2 (b). El método funciona
porque, siendo $S$ el espectro de un cuerpo, el $X'$ que interviene en el lema
de Chow es automáticamente plano sobre $S$.

No sé si el resultado sigue siendo cierto sin hipótesis sobre $S$, suponiendo
únicamente que $X$ es propio y plano sobre $S$ y que $Z$ es cuasiproyectivo
sobre $X$. Un caso importante en las aplicaciones es aquel en que $Z$ es un
subesquema cerrado de $X$; si $\Pi_{X/S}(Z/X)$ existe, es necesariamente un
subesquema cerrado de $S$. Puede
\src{276}
construirse directa y muy sencillamente cuando $X$ es proyectivo sobre $S$,
sin utilizar la teoría de los esquemas de Hilbert. Más generalmente, el mismo
método demuestra que, si $Z$ es afín sobre $X$, entonces $\Pi_{X/S}(Z/X)$
existe y es afín sobre $S$. También demuestra que, si $X$ es propio y plano
sobre $S$ (no necesariamente proyectivo sobre $S$), entonces, para todo
fibrado vectorial localmente trivial $Z$ sobre $X$, $\Pi_{X/S}(Z/X)$ existe y
es un fibrado vectorial sobre $S$. Sería deseable retomar y unificar estos
resultados.

\section*{BIBLIOGRAFÍA}
\addcontentsline{toc}{section}{Bibliografía}

\begin{description}
\item[{[1]}] BOREL (A.) y SERRE (J.-P.). -- Le théorème de Riemann--Roch,
\emph{Bull. Soc. math. France}, vol. 86, 1958, pp. 97--136.

\item[{[2]}] GROTHENDIECK (Alexander). -- Technique de construction et
théorèmes d'existence en géométrie algébrique, I, II, III, \emph{Séminaire
Bourbaki}, vol. 13, 1959/60, n.º 190, 29 pp., n.º 195, 22 pp.; y vol. 14,
1960/61, n.º 212, 20 pp.

\item[{[3]}] GROTHENDIECK (Alexander). -- Séminaire de Géométrie algébrique,
I, II, III, IV. -- París, Institut des hautes Études scientifiques, 1960/61
(mimeografiado).

\item[{[4]}] GROTHENDIECK (Alexander). -- Technique de construction en
géométrie analytique, IV, V, \emph{Séminaire Cartan}, vol. 13, 1960/61, n.os 11
y 12.

\item[{[5]}] GROTHENDIECK (A.) y DIEUDONNÉ (J.). -- \emph{Éléments de
géométrie algébrique}, I: Le langage des schémas. -- París, Presses
universitaires de France, 1960 (Institut des hautes Études scientifiques,
Publications mathématiques, 4); II, III, IV (de próxima aparición).

\item[{[6]}] KODAIRA (K.). -- Characteristic linear systems of complete
continuous systems, \emph{Amer. J. of Math.}, vol. 78, 1956, pp. 716--744.

\item[{[7]}] SERRE (Jean-Pierre). -- Exemples de variétés projectives en
caractéristique $p$ non relevables en caractéristique 0, \emph{Proc. Nat. Acad.
Sc. U.S.A.}, vol. 47, 1961, pp. 108--109.

\item[{[8]}] ZAPPA (G.). -- Sull'esistenza, sopra le superficie algebriche, di
sistemi continui completi infiniti, la cui curva è a serie caratteristica
incompleta, \emph{Pont. Acad. Sc. Acta}, vol. 9, 1945, pp. 91--93.
\end{description}

\section*{ADITIVO}
\addcontentsline{toc}{section}{Aditivo}
\begin{center}
{\small [añadido durante la corrección de pruebas]}
\end{center}

Parece ahora que las conjeturas de [2], III, n.º 8 son falsas, incluso para
variedades no singulares sobre un cuerpo de característica cero, tanto en lo
que respecta a la existencia como a la cuasiproyectividad del cociente, e
incluso cuando $G$ actúa con grafo cerrado.

\def\FgaUnitKey{e221}\clearpage
\chapter*{Fe de erratas de la exposición n.º 221}
\addcontentsline{toc}{chapter}{Fe de erratas de la exposición n.º 221}
\setcounter{footnote}{3}
\src{302}

\noindent\textbf{6. TDTE IV: Los esquemas de Hilbert.}

\begin{center}
[Seminario Bourbaki, t. 13, 1960/61, n.º 221, 28 pp.]
\end{center}

\noindent\textbf{Página 6, Teorema 2.1, 3.ª línea.} --- En lugar de
«es necesario», léase «es necesario y suficiente».

\noindent\textbf{Página 15, línea 13.} --- En lugar de «$P\neq Q$», léase
«$P<Q$».

\noindent\textbf{Línea 15.} --- En lugar de «Si, por el contrario, $P=Q$»,
léase «En el caso contrario, se tendrá $P(n)\geq Q(n)$ para $n$ grande».

\noindent\textbf{Línea 18, fórmula (*).} --- En lugar de «$<$» y «$>$»,
léase «$\leq$» y «$\geq$».

\noindent\textbf{5.ª línea desde abajo.} --- En lugar de «$r$», léase «$Z$».

\noindent\textbf{Página 18, Observación 3.9.} --- Señalemos que Hartshorne ha
estudiado las componentes conexas de los esquemas de Hilbert sobre un cuerpo
algebraicamente cerrado. Este autor demuestra que los
$\operatorname{Hilb}^{P}_{r}$ son conexos y determina los pares $(r,P)$ para
los cuales $\operatorname{Hilb}^{P}_{r}\neq\varnothing$\footnote{HARTSHORNE
(R.). --- Connectedness of the Hilbert scheme (Tesis, Harvard, 1962).}.

\noindent\textbf{Página 23, 5.ª línea desde abajo.} --- En lugar de «una
componente», léase «de una componente».

\noindent\textbf{Páginas 23 y 24, Observación 5.5.} --- A propósito del ejemplo
de Zappa, señalemos que Mumford acaba de construir una componente irreducible
del esquema de Hilbert de $\mathbf P^3_{\mathbf C}$ (cuyos puntos generales
representan curvas no singulares de grado 14 y género 24) que es no reducida
en su punto genérico. Haciendo explotar una de las curvas obtenidas, obtiene
también un esquema proyectivo regular de dimensión 3 sobre $\mathbf C$, cuyo
esquema formal de módulos es no reducido en su punto genérico o, lo que
equivale a lo mismo, cuya variedad local de módulos, en el sentido de la
geometría analítica sobre $\mathbf C$ (cf. Seminario Cartan, t. 13, 1960/61),
es no reducida en todos sus puntos.

\def\FgaUnitKey{232}\unit{Exposición 232 -- Esquemas de Picard: teoremas de existencia}
\sid{232}
\src{143}
\seminares{14.\textsuperscript{o} año, 1961/62, n.\textsuperscript{o} 232\\Febrero de 1962}
\paperhead{TÉCNICA DE DESCENSO Y TEOREMAS DE EXISTENCIA EN GEOMETRÍA ALGEBRAICA\\
V. ESQUEMAS DE PICARD: TEOREMAS DE EXISTENCIA}%
  {por Alexander GROTHENDIECK}

\fsection{1}{Grupos y funtores de Picard relativos}

Para todo preesquema (más generalmente, todo espacio anillado) $X$, llamamos al
grupo de clases, consideradas salvo isomorfismo, de Módulos invertibles (es
decir, Módulos localmente isomorfos a $\mathcal O_X$) sobre $X$ el
\emph{grupo de Picard} (absoluto) de $X$, y lo denotamos por
$\operatorname{Pic}(X)$. Se tiene, por tanto, un isomorfismo canónico
\begin{equation}
\operatorname{Pic}(X)\xrightarrow{\sim}H^1(X,\mathcal O_X^*),
\tag{1.1}
\end{equation}
donde $\mathcal O_X^*$ denota el haz de unidades de $\mathcal O_X$ (que se
identifica, en efecto, con el haz de automorfismos del Módulo invertible
estándar $\mathcal O_X$). Es claro que
$X\longmapsto\operatorname{Pic}(X)$ es un funtor contravariante en $X$, y que
el isomorfismo (1.1) es funtorial.

Si $X$ es un preesquema sobre otro preesquema $S$, al variar $S'$ en la
categoría $(\mathrm{Sch})/S$ de preesquemas sobre $S$, la construcción anterior
proporciona un funtor contravariante
$S'\longmapsto\operatorname{Pic}(X\times_SS')$. Este funtor no tiene ninguna
posibilidad de ser «representable» ([4], II, A), pues los automorfismos no
triviales de los Módulos invertibles que se pretende clasificar impiden que sea
de «naturaleza local» ([5], IV, 5.4). Hay, pues, que «hacerlo local». De modo
general, para todo preesquema relativo $X/S$, introduzcamos el grupo relativo
\begin{equation}
\operatorname{Pic}'(X/S)=H^0\bigl(S,R^1f_*(\mathcal O_X^*)\bigr),
\tag{1.2}
\end{equation}
donde $f:X\longrightarrow S$ es el morfismo estructural (compárese [4], II,
C 3). En el pasaje citado se llama a este grupo grupo de Picard relativo; por
razones que aparecerán enseguida, conviene llamarlo aquí \emph{grupo de Picard
relativo restringido} de $X/S$. Al variar $S'$ en $(\mathrm{Sch})/S$, la
asignación $S'\longmapsto\operatorname{Pic}'(X\times_SS'/S')$ es un funtor
contravariante en $S'$, denotado asimismo por $\operatorname{Pic}'_{X/S}$ y
definido, por tanto, esencialmente por
\begin{equation}
\operatorname{Pic}'_{X/S}(S')=
\operatorname{Pic}'(X\times_SS'/S').
\tag{1.3}
\end{equation}

\src{144}
Este funtor es ahora de «naturaleza local», pues se ha efectuado el ajuste
necesario. Intuitivamente, el segundo miembro de (1.3) se interpreta como el
conjunto de «familias algebraicas» de clases de haces invertibles sobre (las
fibras de) $X/S$, indexadas por el preesquema de parámetros $S'/S$. Cuando el
funtor $\operatorname{Pic}'$ es representable, el preesquema sobre $S$ que lo
representa se denota por $\operatorname{Pic}'_{X/S}$ y se llama
\emph{preesquema de Picard} de $X$ sobre $S$; se tiene entonces
\begin{equation}
\operatorname{Hom}_S(S',\operatorname{Pic}'_{X/S})
\simeq\operatorname{Pic}'_{X/S}(S')
=\operatorname{Pic}'(X\times_SS'/S').
\tag{1.4}
\end{equation}

Hay, sin embargo, casos importantes en los que $\operatorname{Pic}'_{X/S}$ no
es representable (por ejemplo, una variedad de «Brauer--Severi» sobre un
cuerpo $k$ sin punto racional sobre $k$), aunque sigue habiendo una definición
natural de preesquema de Picard relativo. Ello se debe a que, al definir el
funtor $\operatorname{Pic}'$ a partir de los grupos de Picard absolutos
$\operatorname{Pic}(X\times_SS'/S')$, aún no se ha localizado suficientemente;
más precisamente, $\operatorname{Pic}'$ no es, en general, «compatible con el
descenso fielmente plano». Hagámoslo explícito.

Sea $\mathfrak M$ la clase de los morfismos de preesquemas que son
\emph{fielmente planos y cuasicompactos}; esta clase es estable por cambio de
base y por composición. Sea $P$ un funtor contravariante de $(\mathrm{Sch})/S$
a la categoría de conjuntos y, para todo $S$-morfismo
$u:T'\longrightarrow T$ con $u\in\mathfrak M$, consideremos el diagrama
\begin{equation}
P(T)\longrightarrow P(T')\rightrightarrows P(T'\times_TT')
\tag{1.5}
\end{equation}
obtenido al aplicar $P$ a
\[
T\longleftarrow T'\mathrel{\substack{\xleftarrow{\,\mathrm{pr}_1\,}\\[-0.35ex]
\xleftarrow[\,\mathrm{pr}_2\,]{}}}T'\times_TT'.
\]
Si $P$ es representable, la teoría del descenso ([4], I, B, teorema 2) implica
que (1.5) es exacto para todo $u\in\mathfrak M$. Esto se expresa diciendo que
$P$ es compatible con $\mathfrak M$; en este caso, que $P$ es «compatible con
el descenso fielmente plano», o también que el «prehaz» $P$ sobre
$(\mathrm{Sch})/S$ es un «haz» para la noción de localización proporcionada
por $\mathfrak M$. Para un $P$ arbitrario, un procedimiento estándar, familiar
por la localización topológica ordinaria, permite asociarle un «haz»
$\mathcal P$ y un homomorfismo de funtores $P\longrightarrow\mathcal P$
universal en el sentido evidente. El cálculo de $\mathcal P$ puede hacerse
explícito como sigue. Para definir $\mathcal P(T)$, para cada $T'$ sobre $T$
tal que $u:T'\longrightarrow T$ pertenezca a $\mathfrak M$, sea

\src{145}
$\overline H^0(T'/T,P)$ el subconjunto de $P(T')$ formado por los elementos
$\xi$ cuyas dos imágenes $\xi_1,\xi_2$ en $P(T'\times_TT')$ tienen la misma
imagen en $P(T'')$ para algún morfismo
$v:T''\longrightarrow T'\times_TT'$ con $v\in\mathfrak M$. [N.B. -- El
conjunto $\overline H^0$ así definido es, por tanto, mayor que el conjunto
$H^0(T'/T,P)$ introducido en [4], I, A, 4(a).] Al variar $T'$ sobre un $T$ fijo
(siempre con $u\in\mathfrak M$), los $\overline H^0(T'/T,P)$ forman un sistema
inductivo cuando se dota al conjunto de los $T'$ del preorden definido por la
dominación, y se pone
\begin{equation}
\mathcal P(T)=\varinjlim_{T'}\overline H^0(T'/T,P).
\tag{1.6}
\end{equation}
La ley funtorial en $T$ de esta expresión es inmediata.

Cuando
\[
P(T)=\operatorname{Pic}(X\times_ST),
\]
el funtor contravariante sobre $(\mathrm{Sch})/S$ definido por (1.6) se llama
\emph{funtor de Picard relativo} de $X$ sobre $S$ y se denota por
$\operatorname{Pic}_{X/S}$; el grupo $\operatorname{Pic}_{X/S}(S)$ se llama
\emph{grupo de Picard relativo} de $X$ sobre $S$ y se denota por
$\operatorname{Pic}(X/S)$. Se obtiene entonces la biyección evidente
\begin{equation}
\operatorname{Pic}_{X/S}(T)\simeq\operatorname{Pic}(X\times_ST/T).
\tag{1.7}
\end{equation}

Así, un elemento de $\operatorname{Pic}(X/S)$ viene definido por un elemento
$\xi'$ de un grupo $\operatorname{Pic}(X\times_SS')$, donde
$S'\longrightarrow S$ es fielmente plano y cuasicompacto, tal que existe un
morfismo fielmente plano y cuasicompacto
$S''\longrightarrow S'\times_SS'$ para el cual coinciden las dos imágenes
inversas de $\xi'$ en $\operatorname{Pic}(X\times_SS'')$. Un elemento $\xi'$ de
$\operatorname{Pic}(X\times_SS')$ y un elemento $\xi_1$ de
$\operatorname{Pic}(X\times_SS_1)$ que satisfagan las condiciones anteriores
definen el mismo elemento de $\operatorname{Pic}(X/S)$ si y solo si existe un
morfismo fielmente plano y cuasicompacto
$S'_1\longrightarrow S'\times_SS_1$ tal que las imágenes de ambos elementos en
$\operatorname{Pic}(X\times_SS'_1)$ sean iguales. A menudo sigue siendo
conveniente trabajar con el funtor $P'=\operatorname{Pic}'_{X/S}$ introducido
antes; se comprueba inmediatamente que el morfismo canónico
$P\longrightarrow P'$ induce un isomorfismo
\begin{equation}
\mathcal P\xrightarrow{\sim}\mathcal P',
\tag{1.8}
\end{equation}
lo que proporciona una descripción, con frecuencia más cómoda, de
$\operatorname{Pic}_{X/S}$ en términos de
$\operatorname{Pic}'_{X/S}=P'$. En virtud de 2.3, al sustituir $P$ por $P'$ en
la descripción anterior de $\operatorname{Pic}(X/S)$, se puede de hecho
\src{146}
tomar $S''=S'\times_SS'$ y $S'_1=S'\times_SS_1$, al menos bajo las condiciones
especificadas en el pasaje citado.

Cuando el funtor $\operatorname{Pic}_{X/S}$ es representable, se dice que
$X/S$ admite un preesquema de Picard, y el preesquema sobre $S$ que representa
el funtor se llama \emph{preesquema de Picard} de $X$ sobre $S$ y sigue
denotándose por $\operatorname{Pic}_{X/S}$. Basta evidentemente que
$P'=\operatorname{Pic}'_{X/S}$ sea representable, pues entonces $P'$ ya es un
«haz», y (1.8) muestra que $P'\longrightarrow\mathcal P'$ es el morfismo
canónico
\begin{equation}
\operatorname{Pic}'_{X/S}\longrightarrow\operatorname{Pic}_{X/S},
\tag{1.9}
\end{equation}
que es entonces un isomorfismo. Nuestra terminología concuerda así con la
introducida anteriormente en (1.4). En general, cuando
$\operatorname{Pic}_{X/S}$ existe, queda definido por el isomorfismo funtorial
\begin{equation}
\operatorname{Hom}_S(S',\operatorname{Pic}_{X/S})\xrightarrow{\sim}
\operatorname{Pic}(X\times_SS'/S').
\tag{1.10}
\end{equation}

\fsection{2}{Relaciones entre los diversos grupos de Picard relativos y absolutos}

\result{PROPOSICIÓN}{2.1} Sea $f:X\longrightarrow S$ un morfismo tal que
$\mathcal O_S\xrightarrow{\sim}f_*(\mathcal O_X)$. Se tiene entonces una
sucesión exacta
\[
0\longrightarrow\operatorname{Pic}(S)\longrightarrow\operatorname{Pic}(X)
\longrightarrow\operatorname{Pic}'(X/S).
\]
Cuando $X$ admite una sección sobre $S$, el último morfismo es sobreyectivo;
es decir, se tiene un isomorfismo
\[
\operatorname{Pic}'(X/S)\xrightarrow{\sim}
\operatorname{Pic}(X)/\operatorname{Pic}(S).
\]

La sucesión exacta puede considerarse como la sucesión exacta en grados bajos
asociada a la sucesión espectral de Leray para $f$ y $\mathcal O_X^*$. La
segunda afirmación es igualmente formal.

\result{PROPOSICIÓN}{2.2} Sea $f:X\longrightarrow S$ un morfismo cuasicompacto
y separado tal que $\mathcal O_S\xrightarrow{\sim}f_*(\mathcal O_X)$, y sea
$S'\longrightarrow S$ fielmente plano y cuasicompacto. Entonces:

\noindent (i) $\operatorname{Pic}'(X/S)\longrightarrow
\operatorname{Pic}'(X\times_SS'/S')$ es inyectivo;

\noindent (ii) si $X$ admite localmente una sección sobre $S$ (es decir, todo
$s\in S$ posee un entorno abierto $U$ tal que $X|U$ admite una sección sobre
$U$), entonces el diagrama
\[
\operatorname{Pic}'(X/S)\longrightarrow
\operatorname{Pic}'(X\times_SS'/S')\rightrightarrows
\operatorname{Pic}'(X\times_SS''/S'')
\]

\src{147}
donde $S''=S'\times_SS'$, es exacto.

La primera afirmación se deduce de las propiedades elementales del descenso
fielmente plano y de la observación general siguiente. Si
$f:X\longrightarrow S$ satisface
$\mathcal O_S\xrightarrow{\sim}f_*(\mathcal O_X)$, entonces el funtor
$\mathcal F\longmapsto f^*(\mathcal F)$ de los Módulos localmente libres de
tipo finito sobre $S$ a los Módulos localmente libres de tipo finito sobre $X$
es plenamente fiel, y su imagen esencial está formada por los Módulos
$\mathcal G$ sobre $X$ tales que $f_*(\mathcal G)$ es localmente libre y el
homomorfismo canónico
\[
f^*\bigl(f_*(\mathcal G)\bigr)\longrightarrow\mathcal G
\]
es un isomorfismo. La segunda afirmación fue demostrada mediante la teoría del
descenso en [4], I, B, 4.

Los resultados de 2.2 también pueden enunciarse como sigue.

\result{COROLARIO}{2.3} Bajo las condiciones de 2.2, el homomorfismo canónico
(1.9) $\operatorname{Pic}'_{X/S}\longrightarrow\operatorname{Pic}_{X/S}$ es
inyectivo, y es biyectivo si $X$ admite localmente una sección sobre $S$.
(Así, en este último caso, el grupo de Picard relativo
$\operatorname{Pic}(X/S)$ y el grupo de Picard relativo restringido
$\operatorname{Pic}'(X/S)$ son idénticos.)

Combinando esto con 2.1 se obtiene:

\result{COROLARIO}{2.4} Bajo las condiciones de 2.2, se tiene una sucesión
exacta
\[
0\longrightarrow\operatorname{Pic}(S)\longrightarrow\operatorname{Pic}(X)
\longrightarrow\operatorname{Pic}(X/S).
\]
Cuando $X$ admite una sección sobre $S$, el último homomorfismo es
sobreyectivo; es decir, se tiene entonces un isomorfismo
\[
\operatorname{Pic}(X/S)\xrightarrow{\sim}
\operatorname{Pic}(X)/\operatorname{Pic}(S).
\]

\result{OBSERVACIÓN}{2.5} Sea $f:X\longrightarrow S$ un morfismo tal que
$\mathcal O_S\xrightarrow{\sim}f_*(\mathcal O_X)$, y sea $g$ una sección de
$X$ sobre $S$. Para un Módulo invertible $\mathcal L$ sobre $X$, llamemos al
isomorfismo $\mathcal O_S\xrightarrow{\sim}g^*(\mathcal L)$ una
$g$-rigidificación de $\mathcal L$, y llamemos a un Módulo invertible
$\mathcal L$ sobre $X$ provisto de una $g$-rigidificación \emph{Módulo
invertible $g$-rigidificado}. Todo automorfismo de una estructura de este
tipo es trivial, y $\operatorname{Pic}'(X/S)$ se identifica con el conjunto de
clases de isomorfismo de Módulos invertibles $g$-rigidificados sobre
$X$\footnote{El texto impreso dice «sobre $S$»; la rigidificación mediante una
sección $g:S\to X$ concierne aquí al Módulo $\mathcal L$ sobre $X$.}. (Este es
el hecho que permitió emplear la teoría del descenso en la demostración de
2.2(ii).) Esto da una nueva interpretación de $\operatorname{Pic}(X/S)$, al
menos cuando $f$ es además cuasicompacto y separado, de modo que
$\operatorname{Pic}(X/S)\simeq\operatorname{Pic}'(X/S)$ por 2.3.

\src{148}
\result{OBSERVACIÓN}{2.6} Sea $f:X\longrightarrow S$ como en 2.2, y sea
$S'\longrightarrow S$ fielmente plano y cuasicompacto, de modo que exista un
$S$-morfismo $S'\longrightarrow X$, o, equivalentemente, que
$X'=X\times_SS'$ admita una sección sobre $S'$. Pongamos
$S''=S'\times_SS'$ y $X''=X\times_SS''$, y consideremos la sucesión exacta
\[
\operatorname{Pic}(X/S)\longrightarrow\operatorname{Pic}(X'/S')
\rightrightarrows\operatorname{Pic}(X''/S'').
\]
Aplicando 2.4 se obtiene una sucesión exacta
\[
\operatorname{Pic}(X/S)\longrightarrow
\operatorname{Pic}(X')/\operatorname{Pic}(S')
\rightrightarrows
\operatorname{Pic}(X'')/\operatorname{Pic}(S''),
\]
y, en particular, todo elemento del grupo de Picard relativo ya «proviene»
de un elemento de $\operatorname{Pic}(X')$. Esto simplifica sustancialmente la
descripción del grupo de Picard relativo dada en la sección anterior, y también
la del funtor de Picard de $X$ sobre $S$, pues para todo $T$ sobre $S$ el
argumento precedente se aplica a $X\times_ST/T$ y al morfismo
$T'=S'\times_ST\longrightarrow T$. Si, por ejemplo, $f$ mismo es fielmente
plano, se puede tomar $S'=X$. Cuando $f$ es además de tipo finito
(respectivamente liso, etc.), esto permite limitarse, en la descripción del
funtor de Picard relativo
$T\longmapsto\operatorname{Pic}(X\times_ST/T)$, a cambios de base
$T'\longrightarrow T$ que sean de tipo finito (respectivamente lisos, etc.).
Cuando $f$ es proyectivo y plano y $S$ es localmente noetheriano, se demuestra
que el morfismo $S'\longrightarrow S$ anterior puede elegirse de modo que $S'$
sea la unión disjunta de recubrimientos planos $S'_i$ de abiertos $S_i$ que
cubren $S$; si $f$ es además separable, los $S'_i$ pueden tomarse étale sobre
$S_i$.

\fsection{3}{El teorema principal de existencia: enunciado}

No se dispone, ni siquiera a título de conjetura, de un enunciado de existencia
de preesquemas de Picard que englobe todos los casos conocidos. Una condición
«prácticamente necesaria», si puede decirse así, es que
$f:X\longrightarrow S$ sea propio, asegurando las propiedades esenciales de
finitud, y plano. Estas condiciones no son suficientes, incluso si $S$ es el
espectro del álgebra de los números duales $k[t]/(t^2)$ sobre un cuerpo $k$
(por ejemplo, el cuerpo $\mathbf C$ de los números complejos) y $X$ tiene
dimensión 1. En el momento de redactar esta exposición, los teoremas de
existencia más importantes para el preesquema de Picard se deducen del teorema
siguiente.

\result{TEOREMA}{3.1} Sea $f:X\longrightarrow S$ un morfismo de preesquemas
localmente noetherianos. Supongamos:

\noindent (i) $f$ es proyectivo;

\src{149}
\noindent (ii) $f$ es plano;

\noindent (iii) las fibras geométricas de $f$ son íntegras.

Bajo estas condiciones, $\operatorname{Pic}_{X/S}$ existe.

La demostración, que se esbozará en las dos secciones siguientes, mostrará al
mismo tiempo lo que sigue. Sea $\xi$ la sección de
$\operatorname{Pic}_{X/S}$ correspondiente a un haz muy amplio
$\mathcal O_X(1)$ sobre $X/S$, es decir, inducido por una inmersión proyectiva
$X\longrightarrow\mathbf P(\mathcal E)$. Existe entonces un abierto $U$ de
$\operatorname{Pic}_{X/S}$, unión disjunta de abiertos cuasiproyectivos sobre
$S$, tal que $U$ es estable por traslación por $\xi$, y
$\operatorname{Pic}_{X/S}$ es la unión creciente de los abiertos $U-n\xi$,
todos isomorfos a $U$. En particular, bajo las condiciones de 3.1,
$\operatorname{Pic}_{X/S}$ es separado sobre $S$.

\result{OBSERVACIÓN}{3.2} Algunos ejemplos muestran (por ejemplo, con $S$ el
espectro de un anillo de valoración discreta y $X$ de dimensión relativa 1
sobre $S$) que, si se omite la hipótesis (iii) de 3.1 y se sustituye por la
hipótesis más débil de que, para todo $s\in S$, el homomorfismo
\[
k(s)\longrightarrow H^0(X_s,\mathcal O_{X_s})
\]
es un isomorfismo, entonces $\operatorname{Pic}_{X/S}$ no es necesariamente
separado sobre $S$. Esto ocurre tanto cuando las fibras geométricas de $f$ son
reducidas pero una fibra geométrica genérica íntegra «estalla» por
especialización en dos componentes irreducibles, como cuando las fibras
geométricas de $f$ son irreducibles pero una fibra geométrica genérica íntegra
se especializa en una «fibra múltiple». El primer caso se presenta, por ejemplo,
con una cónica que degenera en dos rectas secantes; D. Mumford proporcionó un
ejemplo del segundo, con una curva elíptica que degenera en una curva elíptica
doble. Estos ejemplos son válidos en toda característica.

\result{OBSERVACIÓN}{3.3} Bajo las condiciones de 3.1, no sé si
$\operatorname{Pic}_{X/S}$ es una unión disjunta de abiertos de tipo finito, y
por tanto cuasiproyectivos, sobre $S$.\footnote{La fe de erratas de 1962 informa de
que D. Mumford respondió afirmativamente a esta cuestión mediante la teoría de
los esquemas de Picard desarrollada aquí.} Los polinomios de Hilbert
$Q\in\mathbf Q[t]$, como en el caso de los esquemas de Hilbert ([4], IV), dan
una descomposición de $\operatorname{Pic}_{X/S}$ en una unión disjunta de
abiertos $\operatorname{Pic}^{Q}_{X/S}$, y parece plausible que estos sean
cuasiproyectivos sobre $S$; esto se verá al menos en la exposición siguiente
cuando $f$ sea liso. Obsérvese que, si se sustituye la hipótesis (i) por la de
que $X$ sea localmente proyectivo sobre $S$ (lo cual basta para 3.1, pues la
existencia de $\operatorname{Pic}_{X/S}$ es manifiestamente local sobre $S$),
resulta, por el contrario, fácil

\src{150}
dar ejemplos en los que $\operatorname{Pic}_{X/S}$ contiene componentes conexas
que no son de tipo finito sobre $S$. Sea, por ejemplo, $X_0$ una variedad
algebraica no singular y proyectiva sobre un cuerpo algebraicamente cerrado
$k$, provista de un automorfismo $u$ y de un elemento $\xi$ del grupo de
Néron--Severi de $X_0$ tales que los $u^n(\xi)$ sean dos a dos distintos. Se
puede tomar para $X_0$ el producto de una curva elíptica $E$ por sí misma, y
para $u$ el automorfismo
\[
(x,y)\longmapsto(x,y+x)
\]
de $E\times E$. Sea $S$ la unión de dos curvas irreducibles no singulares que
se cortan en dos puntos $a$ y $b$. Existe un recubrimiento principal conexo de
grupo $\mathbf Z$, $P$, sobre $S$; usando las acciones de $\mathbf Z$ sobre $X_0$
definidas por $u$, se obtiene un fibrado asociado sobre $S$ de fibra $X_0$,
trivial sobre $S-a$ y $S-b$ y, en el caso particular considerado, de hecho un
esquema abeliano sobre $S$. Se comprueba fácilmente que
$\operatorname{Pic}_{X/S}$, que también es el fibrado asociado a $P$ y a las
acciones de $\mathbf Z$ sobre $\operatorname{Pic}_{X_0/k}$ por medio de $u$,
contiene una componente conexa isomorfa a
\[
P\times\operatorname{Pic}^{0}_{X_0/k},
\]
donde $\operatorname{Pic}^{0}$ denota la componente conexa de la identidad en
$\operatorname{Pic}$; esta componente no es de tipo finito sobre $S$.
[Fenómenos análogos se producen también en diversos casos de preesquemas de
Picard no separados sobre $S$, como los considerados en 3.2.]

\fsection{4}{Divisores de Cartier relativos y fibrados proyectivos}

Solo necesitaremos divisores positivos y omitiremos este calificativo adicional
en toda la sección.

Sea $X$ un preesquema. Un \emph{divisor de Cartier} sobre $X$, o simplemente un
divisor sobre $X$, es un subpreesquema cerrado $D$ de $X$ definido por un Ideal
$\mathcal J$ que es un Módulo invertible, es decir, localmente generado por una
sección que no es divisor de cero en $\mathcal O_X$. A $D$ asociamos el Módulo
invertible
\[
\mathcal L(D)=\mathcal J^{-1}.
\]
La inyección canónica $\mathcal J\longrightarrow\mathcal O_X$ proporciona un
homomorfismo canónico
\[
s_D:\mathcal O_X\longrightarrow\mathcal J^{-1}=\mathcal L(D),
\qquad\text{es decir,}\quad s_D\in\Gamma(X,\mathcal L(D)).
\]
Además, dar un divisor equivale esencialmente a dar un Módulo invertible
$\mathcal L$ sobre $X$ junto con una sección $s$ que no sea divisor de cero en
ningún punto, asociando recíprocamente a tal par $(\mathcal L,s)$ el
«divisor» de $s$, denotado por $\operatorname{div}(s)$. Para un $\mathcal L$
invertible fijo sobre $X$, el conjunto de

\src{151}
divisores $D$ que definen $\mathcal L$ está en biyección con el conjunto
cociente
\[
\Gamma(X,\mathcal L)^*/\Gamma(X,\mathcal O_X^*),
\]
donde $\Gamma(X,\mathcal L)^*$ es el subconjunto de
$\Gamma(X,\mathcal L)$ formado por las secciones que no son divisores de cero
en ningún punto.

Supongamos ahora dado un morfismo localmente de tipo finito
$f:X\longrightarrow S$, con $S$ localmente noetheriano para simplificar. Sean
$\mathcal J$ un Ideal coherente sobre $X$, $D$ el subesquema de $X$ que define,
$x\in X$, y pongamos $s=f(x)$. Las condiciones siguientes son equivalentes:

\noindent (i) $\mathcal J$ es invertible en $x$ (es decir, $\mathcal J_x$ está
generado por un elemento regular de $\mathcal O_{X,x}$), y $D$ es plano sobre
$S$ en $x$;

\noindent (ii) $X$ y $D$ son planos sobre $S$ en $x$, y $D_s$ es un divisor de
Cartier sobre la fibra $X_s$ en $x$;

\noindent (iii) $X$ es plano sobre $S$ en $x$, y $\mathcal J_x$ está generado
por un elemento $f_x$ que induce sobre $X_s$ un germen que no es divisor de
cero.

Se dice entonces que $D$ es un \emph{divisor de Cartier relativo} sobre $X/S$
(o simplemente un divisor relativo sobre $X/S$) en el punto considerado. La
condición (i) muestra que $D$ es también un divisor relativo en los puntos
vecinos de $x$. Por tanto, si $X$ y $D$ son planos sobre $S$ y $D$ es propio
sobre $S$, el conjunto de los $s\in S$ tales que $D_s$ es un divisor de Cartier
en $X_s$ (es decir, tales que $D$ es un divisor de Cartier relativo en los
puntos de $X_s$) es abierto en $S$. Además, la definición anterior se formuló
precisamente para que los divisores de Cartier relativos sean estables por todo
cambio de base $S'\longrightarrow S$. Sea $\operatorname{Div}(X/S)$ el
conjunto de divisores relativos sobre $X/S$, y definamos el funtor
contravariante en $S'$ sobre $S$ mediante
\[
\operatorname{Div}_{X/S}(S')=
\operatorname{Div}(X\times_SS'/S').
\]

Supongamos que $X$ es plano y propio sobre $S$. Por la caracterización (ii) de
los divisores de Cartier relativos, $\operatorname{Div}_{X/S}$ puede
considerarse como un subfuntor del funtor $\operatorname{Hilb}_{X/S}$ definido
en [4], IV, y la inclusión
\[
\operatorname{Div}_{X/S}\longrightarrow\operatorname{Hilb}_{X/S}
\]
es «representable por inmersiones abiertas» (cf. [5], IV, 3.13), por las
observaciones precedentes. El teorema principal de existencia de [4], IV, da
entonces:

\src{152}
\result{PROPOSICIÓN}{4.1} Supongamos que $f:X\longrightarrow S$ es proyectivo
y plano. Entonces el funtor $\operatorname{Div}_{X/S}$ es representable; más
precisamente, está representado por un abierto de
$\operatorname{Hilb}_{X/S}$.

Utilizando, para un haz muy amplio fijo $\mathcal O_X(1)$ sobre $X/S$, la
descomposición canónica de $\operatorname{Hilb}_{X/S}$ en los abiertos
$\operatorname{Hilb}^{Q}_{X/S}$ correspondientes a los polinomios de Hilbert
$Q\in\mathbf Q[t]$, se obtiene asimismo la descomposición análoga
\[
\operatorname{Div}_{X/S}=
\coprod_{Q\in\mathbf Q[t]}\operatorname{Div}^{Q}_{X/S}
\]
como unión disjunta de abiertos cuasiproyectivos sobre $S$.

La asignación $D\longmapsto\mathcal L(D)$ proporciona además un homomorfismo
funtorial
\begin{equation}
\operatorname{Div}_{X/S}\longrightarrow\operatorname{Pic}_{X/S}
\tag{+}
\end{equation}
que estudiaremos ahora; bajo condiciones bastante generales resultará
relativamente representable ([5], IV, 3). Partamos de un elemento $\xi$ de
$\operatorname{Pic}_{X/S}(S')$, tomando $S'=S$ para simplificar la notación, y
mostremos que el subfuntor correspondiente de $\operatorname{Div}_{X/S}$ es
representable. Consideremos primero el caso en que $\xi$ viene definido por un
Módulo invertible $\mathcal L$ sobre $X$. Supongamos que $X$ es propio y plano
sobre $S$ y que sus fibras geométricas sobre $S$ son íntegras. Esto implica
también ([2], III, párr. 7) que
$\mathcal O_S\xrightarrow{\sim}f_*(\mathcal O_X)$ y que esta relación sigue
siendo válida tras todo cambio de base $S'\longrightarrow S$. Los divisores de
Cartier relativos $D$ sobre $X/S$ para los que $\mathcal L(D)$ y $\mathcal L$
definen el mismo elemento de
$\operatorname{Pic}(X/S)=\operatorname{Pic}_{X/S}(S)$, o, equivalentemente por
2.4, para los que $\mathcal L(D)$ y $\mathcal L$ son localmente isomorfos sobre
$S$, están en biyección con las secciones del haz cociente
$f_*(\mathcal L)^*/\mathcal O_S^*$. Esta correspondencia es compatible con el
cambio de base. Consideraciones generales de tipo «Künneth» del pasaje citado
(véase también [6]) muestran que la propiedad de ser propio de $X/S$ y la
planitud de $\mathcal L$ sobre $S$ implican la existencia de un Módulo coherente
$\mathcal Q$ sobre $S$, definido salvo isomorfismo único, y de un isomorfismo
de haces
\[
f_*(\mathcal L)\xrightarrow{\sim}
\operatorname{Hom}_{\mathcal O_S}(\mathcal Q,\mathcal O_S).
\]
\footnote{El texto impreso coloca $\mathcal O_X$ como subíndice de
$\operatorname{Hom}$; puesto que $\mathcal Q$ es un Módulo sobre $S$, su dual
se toma sobre $\mathcal O_S$.} La formación de $\mathcal Q$ es compatible con
el cambio de base. Aquí $f_*(\mathcal L)^*$ denota el subhaz de conjuntos de
$f_*(\mathcal L)$ cuyas secciones sobre $U$ son las secciones de $\mathcal L$
sobre $f^{-1}(U)$ que definen divisores de Cartier relativos sobre

\src{153}
$f^{-1}(U)/U$, es decir, que inducen secciones que no son divisores de cero
sobre los $X_s$ para $s\in U$. Como las fibras $X_s$ son íntegras, esto
significa simplemente que las secciones inducidas sobre las fibras $X_s$ no
son idénticamente nulas o, equivalentemente, en términos de homomorfismos
locales $\mathcal Q\longrightarrow\mathcal O_S$, que estos homomorfismos son
sobreyectivos, por el lema de Nakayama. Así, el conjunto de secciones de
$f_*(\mathcal L)^*/\mathcal O_S^*$ está en biyección con el conjunto de
Módulos cociente invertibles de $\mathcal Q$, o, por la definición del fibrado
proyectivo $\mathbf P(\mathcal Q)$ asociado al Módulo coherente $\mathcal Q$
(cf. [5], V, 2), con el conjunto de secciones de $\mathbf P(\mathcal Q)$ sobre
$S$. Esta descripción es compatible con las imágenes inversas y proporciona,
por tanto, el teorema siguiente.

\result{TEOREMA}{4.3} Sea $f:X\longrightarrow S$ un morfismo propio y plano de
fibras geométricas íntegras, donde $S$ es localmente noetheriano, y sea
$\mathcal L$ un Módulo invertible sobre $X$. Para todo $S'$ sobre $S$, sea
$T(S')$ el conjunto de divisores relativos $D$ sobre $X\times_SS'/S'$ tales
que $\mathcal L(D)$ sea isomorfo a
$\mathcal L\otimes_{\mathcal O_S}\mathcal O_{S'}$ localmente sobre $S'$, es decir, tales que
$\mathcal L(D)$ y $\mathcal L\otimes_{\mathcal O_S}\mathcal O_{S'}$ definan
el mismo elemento de $\operatorname{Pic}(X\times_SS'/S')$. Existe entonces un
Módulo coherente $\mathcal Q$ sobre $S$, determinado salvo isomorfismo único,
tal que el funtor $T$ está representado por el fibrado proyectivo
$\mathbf P(\mathcal Q)$.

\result{COROLARIO}{4.4} Si, además, $f$ es proyectivo, el homomorfismo funtorial
$\operatorname{Div}_{X/S}\longrightarrow\operatorname{Pic}_{X/S}$ es
representable por morfismos proyectivos.

En efecto, cuando $X$ admite una sección (respectivamente, localmente una
sección) sobre $S$, el homomorfismo anterior está representado por fibrados
proyectivos (respectivamente, fibrados proyectivos locales), en virtud de 4.3 y
2.1. Cuando $f$ es cuasiproyectivo, se reduce fácilmente al caso anterior
mediante un argumento de descenso que utiliza cuasisecciones finitas y planas
de $X$ localmente sobre $S$.

\result{OBSERVACIÓN}{4.5} Bajo las condiciones de 4.3, el Módulo $\mathcal Q$
sobre $S$ no es, en general, localmente libre, como expresa el hecho de que la
dimensión de las fibras reducidas de $\mathcal Q$, es decir, la dimensión de
$H^0(X_s,\mathcal L_s)$ para $s\in S$ variable, pueda dar saltos. Dado un
Módulo coherente $\mathcal Q$ sobre un preesquema localmente noetheriano $S$,
se comprueba fácilmente que, para un $s\in S$ fijo, $\mathcal Q$ es libre en
$s$ si y solo si $\mathbf P(\mathcal Q)$ es plano sobre $S$ en los puntos
situados sobre $s$; en tal caso es incluso liso sobre $S$ en esos puntos. Aquí,
con $\mathcal Q$ definido en términos de $\mathcal L$ como antes, esto significa
también que la formación de la imagen directa $f_*(\mathcal L)$ «conmuta con el
cambio de base
\src{154}
en un entorno de $s$», o, equivalentemente, que
$f_*(\mathcal L)_s\longrightarrow H^0(X_s,\mathcal L_s)$ es sobreyectivo. Esto
ocurre, por ejemplo, si $H^1(X_s,\mathcal L_s)=0$. Suponiendo que existan los
preesquemas en cuestión, estos criterios se aplican en particular a la situación
universal $\operatorname{Div}_{X/S}\longrightarrow\operatorname{Pic}_{X/S}$ y
dan, respectivamente, una condición necesaria y suficiente y una condición
suficiente para que este morfismo sea liso en un punto dado de
$\operatorname{Div}_{X/S}$.

\fsection{5}{Demostración del teorema principal de existencia}

Bajo las condiciones de 3.1, elijamos un Módulo muy amplio $\mathcal O_X(1)$
sobre $X/S$, y sea $\xi$ el elemento correspondiente de
$\operatorname{Pic}(X/S)$. Pongamos, para abreviar,
\[
P(S')=\operatorname{Pic}(X\times_SS'/S'),
\]
y supongamos primero, para simplificar, que $X/S$ admite una sección. Sea
$P^+(S')$ el subconjunto de $P(S')$ formado por las clases de Módulos
invertibles $\mathcal L$ sobre $X\times_SS'$ tales que
\[
R^if'_*(\mathcal L(n))=0\quad\text{para $i>0$ y todo $n\geq0$},
\qquad
f'_*(\mathcal L(n))\neq0\quad\text{para todo $n\geq0$}.
\]
Estas condiciones son estables por cambio de base y definen, por tanto, un
subfuntor $P^+$ de $P$, manifiestamente estable por traslación por $\xi$.
Utilizando los «teoremas A y B» de Serre ([2], III, 2) y las generalidades de
[5], IV, 5, se ve fácilmente que $P$ es representable si y solo si lo es $P^+$;
en ese caso, $P^+$ está representado por un abierto $U$ del preesquema
$\operatorname{Pic}_{X/S}$ que representa a $P$, y este es la unión creciente
de los abiertos $U-n\xi$.

Pongamos, para abreviar, $D=\operatorname{Div}_{X/S}$, y sea $D^+$ la imagen
inversa de $P^+$ por el morfismo canónico $D\longrightarrow P$. Se tiene así un
morfismo
\begin{equation}
D^+\longrightarrow P^+,
\tag{+}
\end{equation}
y ya se sabe que $D^+$ está representado por un abierto $D^+$ del preesquema
$D=\operatorname{Div}_{X/S}$, cuya existencia garantiza 4.1. Se deduce
fácilmente de 4.3 que (+) está representado por morfismos lisos, proyectivos y
sobreyectivos; más precisamente, por fibrados proyectivos asociados a Módulos
localmente libres y no nulos en ningún punto. En efecto, si $\mathcal L$ sobre
$X\times_SS'$ es como al comienzo de esta sección, $f'_*(\mathcal L)$ es un
Módulo localmente libre no nulo cuya formación conmuta con el cambio de base;
en la notación de 4.3, $\mathcal Q$ es entonces isomorfo al dual de
$f'_*(\mathcal L)$. Deduciremos la representabilidad de $P^+$ del criterio
general [5], IV, 4.7. En el pasaje citado, tomemos para $\mathfrak C$ la clase
de los morfismos de preesquemas fielmente

\src{155}
planos y cuasicompactos, que son epimorfismos efectivos por [4], I B. Acabamos
de verificar la condición (a) del pasaje citado, a saber, que (+) es
representable por morfismos pertenecientes a $\mathfrak C$. También se ha
verificado la condición (b), que dice que el funtor $P^+$ es compatible con el
descenso fielmente plano y cuasicompacto; esto es inmediato. Queda por probar
la condición (c) de dicho pasaje: la relación de equivalencia $R$ en el
preesquema $D^+$ deducida del morfismo (+), representable por $\mathfrak C$, es
$\mathfrak C$-efectiva, es decir, es efectiva y
$D^+\longrightarrow D^+/R$ pertenece a $\mathfrak C$. Obsérvese primero que
los abiertos $D^{+Q}$ de $D^+$ correspondientes a los distintos polinomios de
Hilbert virtuales $Q\in\mathbf Q[t]$ son estables por $R$, pues las fibras de
$R$ son conexas. Basta, por tanto, mostrar, para cada $Q$, que la relación de
equivalencia $R^Q$ inducida en $D^{+Q}$ es $\mathfrak C$-efectiva. Ahora bien,
$D^{+Q}$ es cuasiproyectivo, mientras que la relación de equivalencia $R^Q$ es
proyectiva y plana. Se cumplen, pues, las hipótesis de [4], III, 6.1, que dan el
resultado requerido.

En el caso general en que $X/S$ no admite necesariamente una sección, se reduce
fácilmente al caso anterior por descenso, o se repite la demostración precedente
con la modificación necesaria en la definición de $P^+$.

\result{OBSERVACIONES}{5.1} El método empleado es esencialmente el método de
Matsusaka para la construcción proyectiva de las variedades de Picard. El
resultado de [4], III, invocado para pasar al cociente efectivo, también puede
deducirse fácilmente del teorema de existencia de los esquemas de Hilbert
(véase, por ejemplo, [6]); clásicamente, estos cocientes se construían con
coordenadas de Chow. Obsérvese que la formación del abierto
$\operatorname{Pic}^{+}_{X/S}$ de $\operatorname{Pic}_{X/S}$ y su
descomposición en abiertos $\operatorname{Pic}^{+Q}_{X/S}$ cuasiproyectivos
sobre $S$, según los polinomios de Hilbert de los divisores que definen los
Módulos invertibles considerados, es compatible con el cambio de base, lo que
permite aplicar el argumento de descenso.

\result{OBSERVACIONES}{5.2} Es posible que $\operatorname{Pic}_{X/S}$ exista
siempre que $f:X\longrightarrow S$ sea propio y plano y los homomorfismos
\[
k(s)\longrightarrow H^0(X_s,\mathcal O_{X_s})\qquad(s\in S)
\]
sean isomorfismos; esta última condición significa entonces también que
$\mathcal O_S\xrightarrow{\sim}f_*(\mathcal O_X)$ y que esta relación sigue
siendo válida tras todo cambio de base $S'\longrightarrow S$. Esto es, al
menos, lo que se demuestra para espacios analíticos cuando $f$ es además
proyectivo, mediante un método distinto (el de Chow, salvo error), expuesto en
[5], IX, 3.1.\footnote{La fe de erratas de 1962 indica que la conjetura de existencia
formulada al comienzo de esta observación es falsa, aunque los métodos de
Mumford demuestran un teorema ligeramente más débil.} En ese método, el paso al

\src{156}
cociente por una relación de equivalencia propia y plana en un abierto del
esquema de divisores se sustituye por el paso al cociente por el grupo
proyectivo en el esquema de inmersiones de $X$ en $\mathbf P^r_S$. Este método
puede probablemente adaptarse a los esquemas mediante los resultados de
Mumford sobre el paso al cociente por el grupo proyectivo [8]. Por el momento,
solo existe una demostración escrita cuando $X$ tiene «muchas secciones
locales» sobre $S$, por ejemplo, cuando $X$ es separable sobre un anillo local
completo con cuerpo residual algebraicamente
cerrado.\footnote{El texto impreso omite «con cuerpo residual algebraicamente
cerrado»; la fe de erratas dice que el contraejemplo de Mumford hace indispensable
esta restricción.} En principio, este método tendría mayor alcance, pues también
da la existencia de preesquemas de Picard en casos en los que no son separados
y en los que, por tanto, el primer método necesariamente falla. (Técnicamente,
la dificultad consiste en que, cuando las fibras geométricas de $f$ dejan de
ser íntegras, el funtor considerado en 4.3 no está representado por el propio
fibrado proyectivo $\mathbf P(\mathcal Q)$, sino por un abierto suyo, lo que
conduce al delicado problema de tomar el cociente por una relación de
equivalencia plana pero no propia.)

\result{OBSERVACIÓN}{5.3} Obsérvese que la demostración dada aquí no utiliza ni
una construcción previa de las jacobianas de curvas o familias de curvas, ni la
teoría de variedades abelianas o esquemas abelianos. Difiere, por ello,
esencialmente de exposiciones tradicionales como el libro de Lang [7] o la
exposición de Chevalley [1], que siguen el camino esbozado por A. Weil. Incluso
para las jacobianas de curvas no singulares sobre un cuerpo algebraicamente
cerrado, digamos el cuerpo de los números complejos, la construcción dada aquí
es la única conocida que proporciona las propiedades muy fuertes adoptadas
como definición en la sección 1: esencialmente las de Chevalley, pero permitiendo
«variedades de parámetros» con elementos nilpotentes. Que la construcción de
los esquemas de Picard debe preceder, y no seguir, a la teoría de variedades
abelianas es claro a priori, pues los esquemas de Picard no son, en general,
esquemas abelianos ni se reducen a ellos. Esto ya se ve para curvas singulares
sobre un cuerpo algebraicamente cerrado, donde aparecen las «jacobianas
generalizadas» de Rosenlicht, que no son variedades abelianas. Además, la
teoría de variedades abelianas, y más generalmente de esquemas abelianos, se
simplifica mucho cuando se dispone de una teoría general de esquemas de Picard.
En particular, la teoría de dualidad para esquemas abelianos, especialmente los
resultados de tipo Cartier, se vuelve entonces casi formal (véase, por ejemplo,
[10]).

\src{157}
\result{OBSERVACIONES}{5.4} El «principio de compatibilidad» de Igusa para la
jacobiana de una curva que degenera en una curva singular solo puede
comprenderse adecuadamente como un teorema de existencia del esquema de Picard
de un esquema relativo en curvas $X/S$, no necesariamente liso sobre $S$. Es,
por tanto, un caso particular del teorema principal de existencia 3.1 cuando la
curva especializada es íntegra, es decir, en lenguaje clásico, irreducible de
multiplicidad 1. En la actualidad, el caso de una curva especial reducible,
incluso cuando sus componentes tienen multiplicidad 1, es decir, cuando la
curva especial es separable sobre el cuerpo residual, queda fuera de los
teoremas de existencia conocidos, salvo sobre un anillo de valoración discreta
completo con cuerpo residual algebraicamente cerrado; compárese 5.2. Esta
cuestión de existencia aparecerá ciertamente en una construcción geométrica, en
teoría de esquemas, de las «compactificaciones» de Baily--Satake de los esquemas
de módulos de curvas de género $g$. (Por el momento, esta compactificación solo
se conoce para $g=1$, gracias a los trabajos de Igusa.)

\fsection{6}{Teoremas de existencia relativos}

Esbozaremos aquí algunos casos útiles en los que la existencia de ciertos
esquemas de Picard implica la de otros, lo que permite deducir diversos
teoremas de existencia adicionales del teorema principal 3.1.

\result{PROPOSICIÓN}{6.1} Sea $f:X\longrightarrow S$ un morfismo plano y
proyectivo tal que, en la factorización de Stein $f=f''f'$, el morfismo
$f':X\longrightarrow S'$ es plano, con fibras geométricas íntegras, y satisface
por tanto las hipótesis de 3.1, mientras que el morfismo finito
$f'':S'\longrightarrow S$ es plano. Entonces $\operatorname{Pic}_{X/S}$ existe
y, con la notación introducida en [4], II, C 2, se tiene un isomorfismo canónico
\[
\operatorname{Pic}_{X/S}\xrightarrow{\sim}
\Pi_{S'/S}\operatorname{Pic}_{X/S'}.
\]

Para demostrarlo, se establece primero un isomorfismo de funtores
\[
\operatorname{Pic}_{X/S}\xrightarrow{\sim}
\Pi_{S'/S}\operatorname{Pic}_{X/S'},
\]
se aplica 3.1, que implica que $\operatorname{Pic}_{X/S'}$ es representable, y
se usa después la estructura de $\operatorname{Pic}_{X/S'}$ descrita en la
sección 3, que implica que todo subconjunto finito de una fibra de
$\operatorname{Pic}_{X/S'}$ sobre $S$ está contenido en un abierto afín, para
obtener la existencia de $\Pi_{S'/S}\operatorname{Pic}_{X/S'}$.

\src{158}
Por ejemplo, cuando $X$ es un esquema que es unión disjunta de esquemas $X_i$
sobre $S$ que satisfacen las condiciones de 3.1, el enunciado 6.1 se reduce al
enunciado trivial
\[
\operatorname{Pic}_{X/S}\xrightarrow{\sim}
\prod_i\operatorname{Pic}_{X_i/S}.
\]

\result{COROLARIO}{6.2} Sea $f:X\longrightarrow S$ un morfismo proyectivo y
plano con fibras geométricas localmente íntegras (por ejemplo, un morfismo
proyectivo normal). Entonces $\operatorname{Pic}_{X/S}$ existe.

En este caso, además, $S'$ es un recubrimiento étale de $S$; esto ocurre
siempre que $f$ sea separable, es decir, plano con fibras geométricamente
reducidas. Se comprueba que el teorema de estructura enunciado en la sección 3
para $\operatorname{Pic}_{X/S}$ sigue siendo válido, en virtud de la estructura
análoga de $\operatorname{Pic}_{X/S'}$.

Por otra parte, un procedimiento de descenso proporciona un teorema de
existencia relativo cuyo alcance depende de la solución de las cuestiones de
descenso no plano planteadas en [4], I, A (c). Solo explicitaremos el caso
particular siguiente:

\result{PROPOSICIÓN}{6.3} Sea $f:X\longrightarrow S$ un morfismo propio, y sean
$X_1$ y $X_2$ dos subesquemas de $X$, planos sobre $S$, definidos por dos
Ideales coherentes $\mathcal I_1$ y $\mathcal I_2$ tales que
$\mathcal I_1\cap\mathcal I_2=(0)$ y
$\mathcal O_X/(\mathcal I_1+\mathcal I_2)$ es plano sobre $S$; es decir, el
subesquema de $X$ que es el supremo de $X_1$ y $X_2$ es $X$, mientras que su
ínfimo $Z$ es plano sobre $S$.\footnote{La letra final aparece impresa como
«$X$», pero la condición inmediatamente anterior concierne a la planitud de
$\mathcal O_X/(\mathcal I_1+\mathcal I_2)$ sobre $S$ y, por tanto, define $Z$
como plano sobre $S$.} Supongamos además que, para todo $s\in S$, los
homomorfismos
\[
k(s)\longrightarrow H^0(X_{i,s},\mathcal O_{X_{i,s}}),\qquad i=1,2,
\]
son biyectivos. Entonces el homomorfismo natural de funtores
\[
\operatorname{Pic}_{X/S}\longrightarrow
\operatorname{Pic}_{X_1/S}\times\operatorname{Pic}_{X_2/S}
\]
es representable por morfismos afines. Por consiguiente, si
$\operatorname{Pic}_{X_1/S}$ y $\operatorname{Pic}_{X_2/S}$ existen, también
existe $\operatorname{Pic}_{X/S}$, y el morfismo canónico
\[
\operatorname{Pic}_{X/S}\longrightarrow
\operatorname{Pic}_{X_1/S}\times\operatorname{Pic}_{X_2/S}
\]
es afín.

Por descenso fielmente plano se reduce al caso en que $Z$ admite una sección
sobre $S$, que define así secciones de $X$, $X_1$ y $X_2$ sobre $S$ y permite
eliminar los automorfismos de las estructuras consideradas como se explicó en
2.5. La demostración consiste entonces en observar que los datos de un Módulo
invertible

\src{159}
«rigidificado» $\mathcal L$ sobre $X$ equivalen a los datos de una terna
$(\mathcal L_1,\mathcal L_2,u)$, donde $\mathcal L_i$ es un Módulo invertible
«rigidificado» sobre $X_i$ y $u$ es un isomorfismo de $\mathcal L_1|Z$ con
$\mathcal L_2|Z$ compatible con las rigidificaciones. Solo hay que comprobar
que, para $\mathcal L_1$ y $\mathcal L_2$ fijos, el dato de $u$ se expresa por
una sección de un esquema adecuado sobre $S$ que es afín sobre $S$; esto es
fácil. De 6.3 se deduce inmediatamente:

\result{COROLARIO}{6.4} Sea $X$ un esquema propio y separable sobre un cuerpo
$k$, y sean $X_i$ las componentes irreducibles de $X$. Si existen los
$\operatorname{Pic}_{X_i/k}$, también existe $\operatorname{Pic}_{X/k}$, y el
morfismo canónico
\[
\operatorname{Pic}_{X/k}\longrightarrow\prod_i\operatorname{Pic}_{X_i/k}
\]
es afín.

Combinado con 6.2, esto demuestra, por ejemplo, la existencia de
$\operatorname{Pic}_{X/k}$ siempre que $X$ sea un esquema proyectivo y
separable sobre un cuerpo $k$. Cuando $X$ deja de ser separable sobre $k$,
existe asimismo un resultado de reducción que utiliza un argumento debido a
Oort [11]. El método se aplica también sobre un esquema de base arbitrario, caso
útil, por ejemplo, para demostrar en la exposición siguiente el resultado de
finitud anunciado en 3.3. Para evitar un enunciado excesivamente técnico,
limitémonos a un cuerpo de base:

\result{PROPOSICIÓN}{6.5} Sea $X$ un esquema propio sobre un cuerpo $k$, y sea
$X_0$ un subesquema con el mismo conjunto subyacente, definido por tanto por un
Ideal nilpotente sobre $X$. Entonces el morfismo funtorial
\[
\operatorname{Pic}_{X/k}\longrightarrow\operatorname{Pic}_{X_0/k}
\]
es representable por morfismos afines. En particular, si existe
$\operatorname{Pic}_{X_0/k}$, también existe $\operatorname{Pic}_{X/k}$, y el
morfismo $\operatorname{Pic}_{X/k}\longrightarrow\operatorname{Pic}_{X_0/k}$
es afín.

Combinando esto con 6.4 se concluye inmediatamente:

\result{COROLARIO}{6.6} Sea $X$ un preesquema proyectivo sobre un cuerpo $k$.
Entonces $\operatorname{Pic}_{X/k}$ existe.

\result{OBSERVACIÓN}{6.6} Es sumamente plausible que
$\operatorname{Pic}_{X/k}$ exista para todo esquema propio $X$ sobre un cuerpo
$k$.\footnote{La fe de erratas de 1962 informa de que Murre acababa de responder
afirmativamente a esta cuestión.} Los resultados anteriores reducen esta
cuestión al caso en que $k$ es algebraicamente cerrado y $X$ es íntegro. Por el
lema de Chow, existe entonces un esquema íntegro $X'$, proyectivo sobre $k$, y
un morfismo dominante $g:X'\longrightarrow X$. Bastaría, por tanto, mostrar que
el

\src{160}
morfismo funtorial correspondiente
\[
\operatorname{Pic}_{X/k}\longrightarrow\operatorname{Pic}_{X'/k}
\]
es representable, y sin duda representable por morfismos afines, puesto que ya
se sabe que $\operatorname{Pic}_{X'/k}$ es representable. Esto plantea
cuestiones de descenso no plano que por el momento siguen sin resolverse.
Obsérvese que, si se restringe el funtor $\operatorname{Pic}_{X/k}$ a los
preesquemas reducidos, con $X$ propio e íntegro sobre un cuerpo algebraicamente
cerrado $k$, sí se obtiene un funtor representable, como demostró Chevalley [1]
cuando $X$ es normal y como demostró Seshadri [12] en general mediante un
método de descenso. Sin embargo, con nuestra notación, el esquema construido
por estos autores no es $\operatorname{Pic}_{X/k}$, sino
$(\operatorname{Pic}_{X/k})_{\mathrm{red}}$, el esquema reducido correspondiente
a $\operatorname{Pic}_{X/k}$.

\result{OBSERVACIÓN}{6.7} Más generalmente, sea $f:X'\longrightarrow X$ un
morfismo sobreyectivo de preesquemas propios sobre $k$. Consideraciones de
descenso no plano llevan a conjeturar que
\[
\operatorname{Pic}_{X/k}\longrightarrow\operatorname{Pic}_{X'/k}
\]
es un morfismo afín.\footnote{La fe de erratas de 1962 informa de que esta cuestión
tiene respuesta afirmativa y remite al último párrafo de los Comentarios a la
exposición siguiente.} Esto implicaría, en particular, tras dividir por las
componentes conexas de los elementos neutros, que el homomorfismo
correspondiente de grupos de Néron--Severi es inyectivo módulo torsión. Esto se
puede verificar mediante la teoría de intersecciones cuando $X$ y $X'$ son no
singulares. La respuesta no parece conocida en ningún otro caso.

\result{OBSERVACIÓN}{6.8} Contra lo que podría pensarse, los esquemas de Picard
de esquemas algebraicos con elementos nilpotentes son útiles, e incluso
indispensables, en diversas cuestiones. Así, cuando $X$ es un esquema
proyectivo, digamos liso sobre $k$,\footnote{El texto impreso dice «regular»; la
fe de erratas prescribe «liso sobre $k$».} y $Y$ es una sección hiperplana, hay que
considerar los «entornos infinitesimales» $X_n$ de $Y$ de todos los órdenes y
los esquemas de Picard $\operatorname{Pic}_{X_n/k}$. Cuando $X$ es irreducible
de dimensión $\geq4$ (respectivamente $\geq3$), el morfismo canónico
\[
\operatorname{Pic}_{X/k}\longrightarrow\operatorname{Pic}_{X_n/k}
\qquad(n\text{ grande})
\]
es un isomorfismo (respectivamente, induce un isomorfismo sobre las imágenes
inversas de los subgrupos de torsión de Néron--Severi), resultado que será útil
en el estudio cualitativo de los esquemas de Picard en la exposición siguiente.
Asimismo, considerar los esquemas de Picard de ciertas curvas con elementos
nilpotentes, junto con los teoremas fundamentales de geometría formal [3],
permite verificar, en el caso equicaracterístico, una conjetura de Mumford:
para todo anillo local noetheriano normal completo $A$ de dimensión 2, el grupo
de clases de divisores de $A$ puede considerarse como el conjunto de puntos
racionales sobre $k$ de un

\src{161}
grupo algebraico $G$ sobre el cuerpo residual $k$; además, $G$ queda
determinado canónicamente una vez elegido un cuerpo de coeficientes en $A$.
Cuando $A$ tiene dimensión arbitraria, es plausible que exista un grupo
proalgebraico sobre $k$ que desempeñe el mismo papel que $G$. Cuando
$\operatorname{Spec}(A)$ puede «desingularizarse», dicho grupo se construye como
límite proyectivo de esquemas de Picard de esquemas proyectivos adecuados sobre
$k$ con elementos nilpotentes.

\section*{BIBLIOGRAFÍA}
\addcontentsline{toc}{section}{Bibliografía}

\begin{description}
\item[{[1]}] CHEVALLEY (Claude). -- Sur la théorie de Picard,
\emph{Amer. J. of Math.}, vol. 82, 1960, pp. 435--490.

\item[{[2]}] GROTHENDIECK (A.) y DIEUDONNÉ (J.). -- Éléments de géométrie
algébrique, Chapitre I et suivants. -- París, Presses universitaires de France,
1960--1961 (Institut des Hautes Études Scientifiques, Publications
mathématiques 4, 8, \ldots).

\item[{[3]}] GROTHENDIECK (Alexander). -- Géométrie formelle et géométrie
algébrique, \emph{Séminaire Bourbaki}, vol. 11, 1958/59,
n.\textsuperscript{o} 182, 28 pp.

\item[{[4]}] GROTHENDIECK (Alexander). -- Techniques de descente et théorèmes
d'existence en géométrie algébrique, I--IV, \emph{Séminaire Bourbaki}, vol. 12,
1959/60, n.\textsuperscript{o} 190, 29 pp. y n.\textsuperscript{o} 195, 22 pp.;
vol. 13, 1960/61, n.\textsuperscript{o} 212, 20 pp. y
n.\textsuperscript{o} 221, 28 pp.

\item[{[5]}] GROTHENDIECK (Alexander). -- Techniques de construction en
géométrie analytique, I--X, \emph{Séminaire Cartan}, vol. 13, 1960/61:
Déformation des structures analytiques complexes,
n.\textsuperscript{o} 7--17.

\item[{[6]}] GROTHENDIECK (Alexander). -- Séminaire de géométrie algébrique,
editado por Lichtenbaum, Harvard University, 1961 (de próxima aparición).

\item[{[7]}] LANG (Serge). -- \emph{Abelian varieties}. -- Nueva York,
Interscience Publishers, 1959 (Interscience Tracts in pure and applied
Mathematics, 7).

\item[{[8]}] MUMFORD (D.). -- An elementary theorem in geometric invariant
theory, \emph{Bull. Amer. math. Soc.}, vol. 67, 1961, pp. 483--486.

\item[{[9]}] MUMFORD (D.). -- The topology of normal singularities of an
algebraic surface and a criterion for simplicity. -- París, Presses
universitaires de France, 1961 (Institut des Hautes Études Scientifiques,
Publications mathématiques, 9).

\item[{[10]}] MUMFORD (D.) y TATE (J.). -- Séminaire de géométrie algébrique,
Harvard University, Spring Term 1962 (de próxima aparición).

\item[{[11]}] OORT (F.). -- Sur le schéma de Picard,
\emph{Bull. Soc. math. France}, vol. 90, 1962 (de próxima aparición).

\item[{[12]}] SESHADRI (C. S.). -- Tesis, de próxima aparición en
\emph{Annali di Matematica}.
\end{description}

\def\FgaUnitKey{e232}\clearpage
\chapter*{Fe de erratas de la exposición n.\textsuperscript{o} 232}
\addcontentsline{toc}{chapter}{Fe de erratas de la exposición n.\textsuperscript{o} 232}
\src{303}

\noindent\textbf{7. TDTE V: Esquemas de Picard. Teoremas de existencia.}

\begin{center}
[Seminario Bourbaki, vol. 14, 1961/62, n.\textsuperscript{o} 232, 19 pp.]
\end{center}

\noindent\textbf{Página 7, observación 3.3.} --- La cuestión planteada aquí fue
resuelta afirmativamente por MUMFORD (véanse los Comentarios a la exposición
siguiente).

\noindent\textbf{Página 13, observación 5.2.} --- Como señalo al comienzo de la
exposición siguiente, la conjetura de existencia formulada aquí es falsa; no
obstante, MUMFORD consigue demostrar por sus métodos un teorema ligeramente más
débil.

\noindent\textbf{Página 14, línea 7.} --- Después de «completo», añádase «con
cuerpo residual algebraicamente cerrado». El contraejemplo de MUMFORD recién
mencionado muestra, además, que esta restricción es indispensable.

\noindent\textbf{Página 17, observación 6.6.} --- Como señalamos al comienzo de
la exposición siguiente, la cuestión planteada aquí acaba de ser resuelta
afirmativamente por MURRE.

\noindent\textbf{Página 18, observación 6.7.} --- La cuestión planteada aquí
tiene respuesta afirmativa (cf. el último párrafo de los Comentarios a la
exposición siguiente).

\noindent\textbf{Página 18, observación 6.8, cuarta línea.} --- En lugar de
«regular», léase «liso sobre $k$».

\def\FgaUnitKey{236}\unit{Exposición 236 -- Esquemas de Picard: propiedades generales}
\sid{236}
\src{221}
\seminares{14.º año, 1961/62, n.º 236\\Mayo de 1962}
\paperhead{TÉCNICA DE DESCENSO Y TEOREMAS DE EXISTENCIA EN GEOMETRÍA ALGEBRAICA\\
VI. ESQUEMAS DE PICARD: PROPIEDADES GENERALES}%
  {por Alexander GROTHENDIECK}

\fsection{0}{Complementos a la exposición anterior ([3], V)\textsuperscript{(1)}}

Se han producido algunos avances en las cuestiones de existencia de los
preesquemas de Picard planteadas en [3], V:

\medskip
\noindent a. (Mumford). No es cierto en general que, cuando
$f:X\longrightarrow S$ es un morfismo proyectivo separable (= plano con fibras
separables), exista el preesquema de Picard $\operatorname{Pic}_{X/S}$, ni
siquiera si las fibras de $f$ tienen dimensión 1 y $S$ es el espectro de un
anillo de valoración discreta completo. Se obtiene un ejemplo tomando
$S=\operatorname{Spec}\mathbf R[[t]]$ y tomando como $X$ el subesquema de
$\mathbf P^2_S$ definido por la ecuación (donde $x,y,z$ son variables
homogéneas)
\[
x^2+y^2=tz^2.
\]
Se trata de una cónica que degenera geométricamente en dos rectas que se
cortan, aunque la fibra especial sobre $\mathbf R$ es irreducible (está dada
por $x^2+y^2=0$ sobre $\mathbf R$). Se ve fácilmente que, tras la extensión
étale $S'\longrightarrow S$, con $S'=\operatorname{Spec}\mathbf C[[t]]$, existe
el preesquema de Picard de $X'/S'$ y admite una descripción explícita como
unión disjunta de copias de $\widetilde S$, donde $\widetilde S$ se obtiene de
$S$ duplicando el origen un número infinito de veces. Se comprueba fácilmente
que el dato de descenso sobre $\operatorname{Pic}_{X'/S'}$ para
$S'\longrightarrow S$ (dado aquí por la acción del grupo de Galois
$\mathbf Z/2\mathbf Z$ de $S'$ sobre $S$) no es efectivo: el grupo permuta
ciertos puntos duplicados, de modo que algunas órbitas no están contenidas en
ningún abierto afín. Mumford puede demostrar, sin embargo, que si
$f:X\longrightarrow S$ es un morfismo proyectivo separable tal que, para todo
$s\in S$, las componentes irreducibles de $X_s$ son geométricamente
irreducibles sobre $k(s)$, entonces existe $\operatorname{Pic}_{X/S}$. Su
demostración se apoya en un refinamiento de su teorema del cociente; cf. [7].
Obsérvese, por otra parte, que sin ninguna hipótesis sobre las componentes
irreducibles de las fibras $X_s$, sigue siendo posible que exista, no obstante,
el esquema $\operatorname{Pic}^{\tau}_{X/S}$ (que se introducirá más adelante).

\src{222}
\noindent b. (Murre). Sea $X$ un esquema propio sobre un cuerpo. Entonces
existe $\operatorname{Pic}_{X/k}$. La demostración sigue en parte la de
Chevalley [1] y utiliza de manera esencial la estructura de grupo del funtor de
Picard.

Para algunos complementos relativos a la teoría de los esquemas de Picard, en
particular en relación con los esquemas abelianos, será útil consultar [7]. Por
último, una laguna notable de la presente exposición es la ausencia de
«criterios de equivalencia» que permitan comparar el esquema de Picard de un
esquema proyectivo con los de sus secciones hiperplanas. Los teoremas clave
para desarrollar tales criterios se encuentran en [5], a los cuales hay que
añadir los teoremas de existencia para los esquemas de Picard.

\fsection{1}{Propiedades topológicas de los preesquemas conmutativos en grupos}

Sea primero $k$ un cuerpo y $G$ un preesquema en grupos sobre $k$. Como el
elemento neutro $e$, al ser racional sobre $k$, es necesariamente cerrado, se
deduce inmediatamente que la diagonal de $G\times_kG$ es cerrada y, por tanto,
que $G$ es separado: todo preesquema en grupos sobre un cuerpo es separado.
Denotemos por $G^\circ$ la componente conexa del elemento neutro $e$. Como $e$
es racional sobre $k$, $G^\circ$ es de hecho geométricamente conexa; es decir,
la formación de $G^\circ$ es compatible con la extensión del cuerpo de base.
Se deduce asimismo que $G^\circ$ es estable por multiplicación como conjunto y,
cuando $G$ es localmente noetheriano, de modo que $G^\circ$ es abierto, puede
considerarse un subgrupo abierto de $G$. Supongamos en lo sucesivo que $G$ es
localmente de tipo finito sobre $k$. Entonces $G^\circ$ es geométricamente
irreducible y de tipo finito sobre $k$. En efecto, se puede suponer que $k$ es
algebraicamente cerrado y después que $G$ es reducido (pues
$G_{\mathrm{red}}$ es entonces un subgrupo de $G$, ya que
$G_{\mathrm{red}}\times_kG_{\mathrm{red}}$ vuelve a ser reducido); por tanto,
es liso sobre $k$ en un abierto no vacío y, mediante traslación de ese
abierto, en todas partes. Así, $G$ es localmente irreducible, por lo que sus
componentes irreducibles coinciden con sus componentes conexas y $G^\circ$ es
irreducible. Sea $U$ un entorno afín de $e$ en $G^\circ$. Utilizando la
irreducibilidad de $G^\circ$, se ve inmediatamente que $U\cdot U=G^\circ$, lo
que demuestra que $G^\circ$ es cuasicompacto y, por tanto, de tipo finito sobre
$k$.

Para simplificar, supongamos que $G$ es conmutativo. Para todo entero $n>0$,
sea $G^{(n)}$ la imagen inversa de $G^\circ$ por el homomorfismo de potencia
$n$-ésima $\varphi_n$ de $G$; así, $G^{(n)}$ es un subgrupo abierto de $G$.
Pongamos
\[
G^\tau=\bigcup_nG^{(n)}.
\]

\src{223}
\[
G^\sigma=\bigcup_{(n,p)=1}G^{(n)},
\qquad
G^\rho=\bigcup_hG^{(p^h)},
\]
donde $p$ es el exponente característico del cuerpo $k$. Estos son subgrupos
abiertos de $G$ que satisfacen
\[
G^\sigma\cap G^\rho=G^\circ,
\qquad
G^\sigma G^\rho=G^\tau.
\]

\result{OBSERVACIÓN}{} Se puede construir el esquema en grupos cociente
$G/G^\circ=N$ (cf. [3], IV), y definir entonces $G^\tau$, $G^\sigma$ y
$G^\rho$ como las imágenes inversas en $G$, respectivamente, del subgrupo de
torsión de $N$, de la suma de sus componentes $q$-primarias para los primos
$q\neq p$ y de su componente $p$-primaria.\footnote{El texto impreso intercambia
las dos últimas correspondencias: asigna a $G^\sigma$ la componente
$p$-primaria y a $G^\rho$ su complemento primo con $p$. Las definiciones
precedentes de $G^\sigma$ y $G^\rho$ obligan al orden corregido que se da aquí.}
Obsérvese a este respecto que $N$ es un esquema en
grupos discreto separable sobre $k$. Así, tras elegir una clausura algebraica
$\overline k$ de $k$, con grupo de Galois $\pi$, se identifica con un grupo
discreto ordinario sobre el cual $\pi$ actúa por automorfismos. Esto da la
interpretación evidente de la construcción del subgrupo de torsión y de su
descomposición en componentes $q$-primarias. Cuando $G$ es el esquema de
Picard de un esquema propio $X$ sobre $k$, se puede llamar a $N$ el
\emph{esquema de Néron--Severi} (reducido) de $X$ sobre $k$. Cuando
$G^\circ_{\mathrm{red}}$ es un subesquema en grupos de $G^\circ$, lo que ocurre
en particular siempre que $k$ sea perfecto o $G^\circ$ sea propio sobre $k$
(por ejemplo, cuando $X$ sea geométricamente normal), conviene introducir
también el cociente $N'=G/G^\circ_{\mathrm{red}}$, que tiende a comportarse
mejor que $N$ por especialización, es decir, cuando $X$ varía en una familia
de esquemas algebraicos.

Sea ahora $S$ un preesquema localmente noetheriano y sea $G$ un preesquema en
grupos sobre $S$, localmente de tipo finito sobre $S$. No suponemos que $G$ sea
de tipo finito sobre $S$ ni separado sobre $S$. Pongamos
\[
G^\circ=\bigcup_{s\in S}(G_s)^\circ,
\]
y, cuando $G$ sea conmutativo,
\[
G^\tau=\bigcup_{s\in S}(G_s)^\tau,
\qquad
G^\sigma=\bigcup_{s\in S}(G_s)^\sigma,
\qquad
G^\rho=\bigcup_{s\in S}(G_s)^\rho.
\]

\src{224}
Estos son subconjuntos de $G$ estables por la multiplicación en $G$, lo que
evidentemente no implica que puedan ser definidos por subpreesquemas en grupos
de $G$. En particular, parece que en general no existe un subpreesquema en
grupos de $G$ cuyo conjunto subyacente sea $G^\circ$. Por supuesto, si uno de
estos conjuntos es abierto, entonces, con su estructura inducida, es un
subpreesquema abierto en grupos de $G$. Veremos que esto siempre ocurre para
$G^\tau$. Así, desde el punto de vista de los funtores representables, y en
particular desde el punto de vista de la especialización, la equivalencia
numérica se comporta de manera más satisfactoria que la equivalencia
algebraica. He aquí las principales propiedades generales de los conjuntos que
acabamos de definir:

\result{TEOREMA}{1.1} Los conjuntos $G^\circ$, $G^\tau$, $G^\sigma$ y $G^\rho$
son localmente constructibles. Además:

\noindent (i) $G^\circ$ es cuasicompacto sobre $S$. Cuando los $G_s^\circ$ son
propios y $G$ es separado sobre $S$, $G^\circ$ es propio sobre $S$ y, por
tanto, cerrado en $G$.

\noindent (ii) $G^\tau$ es abierto. Si $G^\circ$ es cerrado, también lo es
$G^\tau$.

\noindent (iii) Si $G^\circ$ es cerrado, también lo es $G^\sigma$, siempre que
se esté en característica igual, es decir, que todos los cuerpos residuales de
$S$ tengan la misma característica. Si $G^\circ$ es cerrado y
$G\longrightarrow S$ es universalmente abierto en los puntos de $G^\sigma$
(cf. Corolario 1.5 más adelante), entonces $G^\sigma\longrightarrow S$ es
universalmente abierto.

\noindent (iv) Si $G^\circ$ es cerrado, también lo es $G^\rho$. Supongamos que
se está en característica igual y que, para todo entero $n>0$ con $(n,p)=1$,
el homomorfismo de potencia $n$-ésima de $G$ es abierto. Entonces $G^\rho$ es
abierto.

Damos indicaciones para la demostración. La constructibilidad local de
$G^\circ$ y su cuasicompacidad sobre $S$ están contenidas en el lema siguiente:

\result{LEMA}{1.2} Supongamos $S\neq\varnothing$. Entonces existen un abierto
no vacío $U$ de $S$, un esquema en grupos $H$ de tipo finito sobre $U$ con
fibras conexas y un homomorfismo de esquemas en grupos
$H\longrightarrow G|U$ que es una inmersión abierta de imagen $G^\circ|U$.

Para demostrar el lema, se puede suponer que $S$ es irreducible y denotar por
$\eta$ su punto genérico. Tras el cambio de base
$S'=\operatorname{Spec}(\mathcal O_{S,\eta})\longrightarrow S$, se obtiene un
esquema en grupos $G'$ sobre el anillo local artiniano
$\mathcal O_{S,\eta}=A$, con un subesquema abierto en grupos $G'{}^\circ$ de
tipo finito sobre $A$, como se ha indicado antes. Por tanto, este procede de un
esquema en grupos $H$ de tipo finito sobre un entorno abierto $U$ de $\eta$, y
la inmersión canónica $G'{}^\circ\longrightarrow G'$ procede de una

\src{225}
inmersión abierta $H\longrightarrow G|U$, que es un homomorfismo de esquemas en
grupos después de hacer $U$ suficientemente pequeño. Si $U$ es suficientemente
pequeño, las fibras de $H$ son conexas y tienen todas la misma dimensión,
igual a la de las fibras de $G$. De ello se sigue que, para todo $s\in U$, la
imagen de $H_s$ en $G_s$ es exactamente $G_s^\circ$ (tras reducir de nuevo
$U$), lo que demuestra 1.2.

La segunda afirmación de (i) está contenida en el lema siguiente, aplicado a
un entorno abierto cuasicompacto de $G^\circ$ en $G$:

\result{LEMA}{1.3} Sea $X$ un preesquema separado de tipo finito sobre un
preesquema localmente noetheriano $S$, sea $g$ una sección de $X$ sobre $S$ y
sea $X^\circ$ la unión, en los $X_s$, de las componentes conexas que contienen
$g(s)$. Si $s\in S$ y $X_s^\circ$ es propio sobre $k(s)$, existe un entorno
abierto $U$ de $s$ tal que $X^\circ|U$ es propio sobre $U$ y, por tanto,
cerrado en $X|U$.

Por descenso fielmente plano sobre la base, se reduce al caso en que $S$ es el
espectro de un anillo local completo y $s$ es su punto cerrado. Aplicando [2],
III, 5.5.1, se ve que $X$ es la unión disjunta de dos abiertos $X'$ y $X''$, el
primero propio sobre $S$ y tal que $X'_s=X_s^\circ$. Esto reduce al caso
$X=X'$, es decir, a $X$ propio sobre $S$. Se utiliza entonces un argumento
estándar basado en el criterio valuativo de ser propio para un subconjunto
(omitido en [2], capítulo II).

Demostremos ahora que $G^\tau$ es abierto. Equivalentemente, puesto que la
formación de $G^\tau$ (así como la de $G^\circ$, $G^\sigma$ y $G^\rho$) conmuta
con la extensión de la base, para toda sección $g$ de $G$ sobre $S$, el
conjunto $g^{-1}(G^\tau)$ es abierto. Esto significa dos cosas:

\medskip
\noindent a. Si $y\in S$ y $g(y)\in G_y^\tau$, entonces
$g(y')\in G_{y'}^\tau$ para todo $y'\in\overline{\{y\}}$ suficientemente
próximo a $y$.

\medskip
\noindent b. Si $y'\in S$ y $g(y')\in G_{y'}^\tau$, entonces
$g(y)\in G_y^\tau$ para toda generalización $y$ de $y'$.

Para (a), existe un entero $n>0$ tal que $g^n(y)\in G_y^\circ$. Como
$G^\circ$ es constructible, $(g^n)^{-1}(G^\circ)$ es constructible y, por
tanto, $g^n(y')\in G_{y'}^\circ$ para
$y'\in\overline{\{y\}}$ suficientemente próximo a $y$. Para (b), las
potencias $g(y')^n=g^n(y')$ permanecen en un abierto cuasicompacto de
$G_{y'}$, pues están contenidas en un número finito de clases módulo
$G_{y'}^\circ$. Por consiguiente, existe un abierto cuasicompacto $U$ de $G$
que contiene los $g^n(y')$ y, por tanto, también sus generalizaciones
$g^n(y)$. Así, las potencias de $g(y)$ permanecen en un abierto cuasicompacto
de $G_y$, lo que implica fácilmente que $g(y)\in G_y^\tau$.

\src{226}
Supongamos que $G^\circ$ es cerrado; demostremos lo mismo para $G^\tau$. Como
ya se sabe que $G^\tau$ es abierto y, por tanto, localmente constructible,
basta demostrar que es estable por especialización. Esto se sigue de que es
una unión de conjuntos cerrados, a saber, las imágenes inversas del conjunto
cerrado $G^\circ$ por los homomorfismos de potencia $n$-ésima $\varphi_n$.

El mismo argumento demuestra que $G^\sigma$ y $G^\rho$ son cerrados si
$G^\circ$ es cerrado (en el primer caso, bajo la condición de característica
igual), una vez demostrado que $G^\sigma$ y $G^\rho$ son localmente
constructibles. Para $x\in G^\tau$, sea $\nu(x)$ el menor entero $n>0$ tal que
el homomorfismo de potencia $n$-ésima $\varphi_n$ envía $x$ a $G^\circ$.
Entonces $G^\sigma$ (respectivamente, $G^\rho$) está formado por los
$x\in G$\footnote{El texto impreso dice «$X$», que no está definido en este
pasaje; los cuatro loci considerados son subconjuntos de $G$.} para los cuales
$\nu(x)$ es primo con $p$ (respectivamente, una potencia de $p$). La afirmación
de constructibilidad se deduce ahora del enunciado más preciso:

\result{LEMA}{1.4} La función $\nu$ sobre $G^\tau$ es localmente
constructible.

En efecto, para todo entero $n>0$, el conjunto de los $x\in G^\tau$ tales que
$\nu(x)=n$ es localmente constructible: es la diferencia entre
$\varphi_n^{-1}(G^\circ)$ y la unión de los
$\varphi_d^{-1}(G^\circ)$, donde $d$ recorre los divisores propios de $n$.
Como $G^\circ$ es localmente constructible, también lo son los
$\varphi_d^{-1}(G^\circ)$ y, por consiguiente, también lo es la diferencia
anterior.

Supongamos ahora que $G^\circ$ es cerrado y que $G\longrightarrow S$ es
universalmente abierto en los puntos de $G^\sigma$. Demostremos que
$G^\sigma\longrightarrow S$ es abierto, es decir, que envía un entorno en
$G^\sigma$ de $x\in G^\sigma$ sobre un entorno de $y=f(x)$. Como $G^\sigma$ es
localmente constructible, basta demostrar que para toda generalización $y'$ de
$y$ existe una generalización $x'$ de $x$ en $G^\sigma$ situada sobre $y'$.
Por cambio de base, se reduce al caso en que $S$ es el espectro de un anillo de
valoración discreta y $y,y'$ son, respectivamente, sus puntos cerrado y
genérico. Como $G\longrightarrow S$ es abierto en $x$, existe una
generalización $x_1$ de $x$ en $G$ sobre $y'$. Tras otro cambio de base, se
puede suponer además que existe una sección $g$ de $G$ sobre $S$ con
$x=g(y)$. Si $k(y')$ tiene característica cero, tómese cualquier
generalización $x'$ de $x$ en $G$ sobre $y'$. Pertenece a $G^\tau$ porque
$G^\tau$ es abierto y, por tanto, a $G^\sigma$ porque
$G_{y'}^\sigma=G_{y'}^\tau$. Si $k(y')$ tiene característica $p>0$, escribamos
\[
\nu(g(y'))=p^hm\qquad\text{con }(m,p)=1;
\]

\src{227}
elijamos enteros $a$ y $b$ tales que $ap^h+bm=1$, y pongamos
\[
g_1=g^{ap^h},\qquad g_2=g^{bm},\qquad\text{de modo que}\qquad g=g_1g_2.
\]
Por construcción, $g_1(y')\in G^\sigma$ y $g_2(y')\in G^\rho$. Como
$G^\circ$ es cerrado, se sigue que $g_1(y)\in G^\sigma$ y
$g_2(y)\in G^\rho$. Pero
\[
g(y)=g_1(y)g_2(y)\in G^\sigma,
\]
por lo que $g_2(y)\in G^\sigma$ y, por consiguiente,
\[
g_2(y)\in G_y^\sigma\cap G_y^\rho=G_y^\circ.
\]
La hipótesis, junto con el hecho de que $S$ es el espectro de un anillo de
valoración discreta, implica que $G-(G_y-G_y^\circ)$ es un abierto de $G$ en
el que $G\longrightarrow S$ induce un morfismo abierto. Así, por todo punto de
$G_y^\circ$ pasa una «cuasisección». Después de efectuar una extensión
adicional de la base $S$, se puede suponer, por consiguiente, que existe una
sección $g'_2$ de $G^\circ$ sobre $S$ tal que $g'_2(y)=g_2(y)$. Pongamos
$g'=g_1g'_2$. Entonces, por construcción, $g'(y)=g(y)=x$, y
\[
g'(y')=g_1(y')g'_2(y')\in G^\sigma.
\]
Por tanto, $g'(y')=x'$\footnote{El texto impreso dice «$g(y')=x'$». La sección
construida en la frase anterior, cuya fibra genérica se calcula, es $g'$.} es
una generalización de $x$ en $G^\sigma$ situada sobre $y'$. Esto demuestra
(iii).

Queda por demostrar la última afirmación, (iv). Debemos probar que, si
$x\in G^\rho$, toda generalización $x'$ de $x$ pertenece a $G^\rho$. Tomando
imágenes por $\varphi_{p^h}$ para un $h$ adecuado\footnote{El texto impreso
dice $\varphi_{n^h}$. Para $x\in G^\rho$, la definición de la p.~223 da una
potencia del exponente característico $p$ que envía $x$ a $G^\circ$.}, se
puede suponer que, de hecho, $x\in G^\circ$. Para todo entero $n$ primo con la
característica, $x$ pertenece entonces a la imagen de $\varphi_n$ (pues la
aplicación de potencia $n$-ésima en un grupo conexo de tipo finito sobre un
cuerpo de característica prima con $n$ es sobreyectiva). Como $\varphi_n$ es
abierto, $x'$ pertenece también a la imagen de $\varphi_n$. Más precisamente,
sea $U$ un abierto cuasicompacto de $G^\tau$ que contenga $G_y^\circ$.
Entonces $x'\in\varphi_n(U)$ para todo $n$ primo con $p$. Tomando por $n$ un
múltiplo común de las partes primas con $p$ de los $\nu(z)$, para $z\in U$, se
obtiene $x'\in G^\rho$.

La afirmación (1.1), (iii), se completa mediante el corolario siguiente:

\result{COROLARIO}{1.5} Sea $n>1$ un entero tal que el homomorfismo de potencia
$n$-ésima $\varphi_n:G\longrightarrow G$ sea universalmente abierto (por
ejemplo, étale), pongamos
\[
G^{(n)}=\varphi_n^{-1}(G^\circ),
\]
y supongamos que, en cada fibra geométrica $G_{\bar s}$, el subconjunto
$\bigcup_{h\geq0}|{}_{n^h}G_{\bar s}|$ es denso en
$\bigcup_{h\geq0}G_{\bar s}^{(n^h)}$. Entonces

\src{228}
$G\longrightarrow S$ es universalmente abierto en los puntos de los
$G^{(n^h)}$. En particular, si los $G_s^\circ$ no contienen ninguna componente
aditiva y, para todo entero $n>1$, el homomorfismo $\varphi_n$ es
universalmente abierto en los puntos de característica residual prima con
$n$\footnote{El texto impreso dice «prima con $p$». En este enunciado
cuantificado por $n$, el criterio para que un morfismo sea étale y el argumento que utiliza los
núcleos $ {}_nG$ requieren característica residual prima con $n$.}, entonces
$G\longrightarrow S$ es universalmente abierto en los puntos de $G^\sigma$.
Por tanto, según (1.1), (iii), si $G^\circ$ es cerrado,
$G^\sigma\longrightarrow S$ es universalmente abierto. En estos casos,
$s\longmapsto\dim G_s$ es localmente constante sobre $S$.

En efecto, la hipótesis implica que el núcleo $ {}_nG$ de $\varphi_n$ es
universalmente abierto sobre $S$. Por tanto, $G\longrightarrow S$ es
universalmente abierto en los puntos de $ {}_nG$ y, reemplazando $n$ por
$n^h$, en los puntos de $ {}_{n^h}G$. La hipótesis de densidad dice exactamente que, fibra a fibra, los puntos cuyo
orden es una potencia de $n$ son densos en la unión de los $G^{(n^h)}$. Para
la afirmación particular, la ausencia de componentes aditivas implica esta
densidad para los enteros $n$ invertibles en el campo residual. Se deduce fácilmente que
$G\longrightarrow S$ es universalmente abierto en cada punto de esta unión,
en particular a lo largo de la sección neutra, y, por tanto, que
$s\longmapsto\dim G_s$ es localmente constante.

\result{OBSERVACIÓN}{1.6} Recuérdese que un morfismo $X\longrightarrow S$ se
llama universalmente abierto si envía todo abierto a un abierto y conserva
esta admirable propiedad después de cualquier cambio de base
$S'\longrightarrow S$. Esto también significa (cuando $S$ es localmente
noetheriano y $X\longrightarrow S$ es localmente de tipo finito) que toda
componente irreducible de $X$ domina $S$ y que esta propiedad se conserva
después de cualquier cambio de base $S'\longrightarrow S$. En ambas
afirmaciones basta (bajo las hipótesis de finitud anteriores) comprobar los
cambios de base $S'\longrightarrow S$ en los que $S'$ es el espectro de un
anillo de valoración discreta (completo y con cuerpo residual algebraicamente
cerrado, si se desea~...). La definición se extiende de manera evidente a un
subconjunto $Z$ de $X$ (como el subconjunto $G^\sigma$ de $G$). Es un fenómeno
perfectamente normal, incluso para un morfismo proyectivo liso
$X\longrightarrow S$ con fibras geométricamente conexas (por ejemplo, el
producto fibrado de la familia modular de curvas elípticas consigo misma sobre
$S=\operatorname{Spec}\mathbf C[j]$), que $\operatorname{Pic}_{X/S}$ no sea
universalmente abierto sobre $S$: una componente irreducible de
$\operatorname{Pic}_{X/S}$ puede estar situada enteramente sobre un único
punto de $S$. Esto está relacionado con el hecho de que el rango del grupo de
Néron--Severi de las fibras de $X/S$ puede aumentar bruscamente (fenómenos de
«multiplicación compleja»). Por otra parte, (1.5) nos asegura que, en casos
favorables, $\operatorname{Pic}_{X/S}^\sigma$ (y parece que muy a menudo incluso
$\operatorname{Pic}_{X/S}^\tau$) es universalmente abierto sobre $S$.

\src{229}
Por último, he aquí un caso útil en el que, excepcionalmente, $G^\circ$
consiente en ser abierto:

\result{COROLARIO}{1.7} Supongamos que $G\longrightarrow S$ es universalmente
abierto en los puntos de $G^\circ$ (cf. (1.5)) y que las fibras $G_s$ son
separables y, por tanto, lisas sobre $k(s)$ (esta última condición se satisface
automáticamente en característica residual cero, por un teorema de Cartier).
Entonces $G^\circ$ es abierto en $G$. Si, además, $S$ es reducido, entonces
$G^\circ$ es liso (en particular, plano) sobre $S$.

La primera afirmación puede precisarse del modo siguiente. Si, para un
$s\in S$ dado, $G\longrightarrow S$ es universalmente abierto en los puntos de
$G_s^\circ$ y $G_s^\circ$ es separable sobre $k(s)$, entonces $G^\circ$ es un
entorno de $G_s^\circ$ (y la hipótesis hecha en $s$ seguirá siendo válida en
los puntos próximos). Este enunciado, independiente de toda estructura de
grupo, se demostrará en [2], IV, §7. La última afirmación de (1.7), también
independiente de toda estructura de grupo y de toda cuestión relativa a una
componente conexa, es un caso particular de un criterio de planitud dado en
[2], IV, §5, que implica de manera más general:

\result{COROLARIO}{1.8} Sea $U$ un abierto de una fibra $G_s$ tal que
$G\longrightarrow S$ sea universalmente abierto en los puntos de $U$ (cf.
(1.5)). Si $G_s$ es separable sobre $k(s)$ y $S$ es reducido en $s$, entonces
$G$ es plano sobre $S$ en los puntos de $U$ (y, por tanto, en este caso, liso
sobre $S$ en los puntos de $U$).

\result{OBSERVACIONES}{1.9} No sé si, cuando $G$ es separado sobre $S$,
$G^\circ$ o $G^\tau$ es siempre cerrado en $G$; parece poco probable. En todo
caso, hay contraejemplos evidentes si se suprime la hipótesis de separación.
Por ejemplo, tómese la recta afín con el origen duplicado una infinidad
numerable de veces. Esto da un esquema en grupos $G$ sobre la recta afín cuyas
fibras se reducen todas al grupo neutro salvo una, cuya fibra es $\mathbf Z$.
Este esquema en grupos es un subesquema abierto en grupos, a saber, la clausura
de la sección neutra, del preesquema de Picard del $S$-esquema $X$
correspondiente a una familia de cónicas que degenera en dos rectas que se
cortan.

Además, incluso partiendo de un esquema en grupos finito y plano $G$ sobre el
espectro $S$ de un anillo de valoración discreta $V$, de modo que $G^\circ$ se
reduce a la sección neutra y es, por tanto, cerrado, varias conclusiones pasan
a ser falsas si se suprimen ciertas hipótesis. Supongamos que $V$ tiene
característica igual $p>0$ y sea $G$ el núcleo del homomorfismo
$\mathbf G_a\longrightarrow\mathbf G_a$ definido por el homomorfismo de
funtores
\[
f\longmapsto f^p-tf,
\]
donde $t$ es un uniformizante de $V$.

\src{230}
(Recuérdese que, por definición, el «grupo aditivo» $\mathbf G_a$ sobre $S$
representa el funtor $S'\longmapsto\Gamma(S',\mathcal O_{S'})$.) Así, como
$S$-esquema,
\[
G=\operatorname{Spec}V[x]/(x^p-tx)
\]
es la unión de $p$ «rectas» que se cortan, y representa un grupo
$\mathbf Z/p\mathbf Z$ que degenera en un grupo infinitesimal de tipo aditivo.
En este ejemplo, $G^\circ=G^\sigma$ (y es una de las $p$ rectas), por lo que no
es abierto en $G$. En característica mixta de característica residual $p>0$,
partamos del esquema en grupos $H=\boldsymbol\mu_p$, núcleo de la aplicación de
potencia $p$-ésima en $\mathbf G_m$. Así,
\[
H=\operatorname{Spec}V[x]/(x^p-1);
\]
Este es de nuevo la unión de $p$ rectas que se cortan, y representa un grupo
$\mathbf Z/p\mathbf Z$ (en característica cero) que degenera en un grupo
infinitesimal de tipo multiplicativo en característica $p$. Sea $H'$ el
«esquema en grupos constante» definido por el grupo finito ordinario
$\mathbf Z/p\mathbf Z$, es decir, la unión disjunta de $\mathbf Z/p\mathbf Z$
copias de $S$, o equivalentemente
$\operatorname{Spec}V^{\mathbf Z/p\mathbf Z}$. Entonces $G=H\times_S H'$
describe un grupo $(\mathbf Z/p\mathbf Z)^2$ en característica cero, que
degenera en característica $p$ en el producto de un grupo infinitesimal por
$\mathbf Z/p\mathbf Z$. Aquí $G^\rho$ es la unión de la sección neutra y la
fibra especial, de modo que $G^\rho$ no es abierto, mientras que $G^\sigma$ es
la unión de la sección neutra y la fibra genérica, de modo que no es cerrado,
contrariamente a lo afirmado en los casos de característica igual de (1.1),
(iii) y (iv). Se trata, naturalmente, de fenómenos relacionados con la
característica $p>0$. Los resultados anteriores dan:

\result{COROLARIO}{1.10} Supongamos que todas las características residuales de
$S$ son cero, de modo que $G^\tau=G^\sigma$ y $G^\circ=G^\rho$. Entonces
$G^\tau=G^\sigma$ es abierto; es incluso abierto y cerrado si $G$ es separado
sobre $S$ y los $G_s^\circ$ son propios. Bajo la misma hipótesis,
$G^\circ=G^\rho$ es propio sobre $S$ y, por tanto, cerrado en $G$. Supongamos,
por último, que para todo entero $n>1$ el homomorfismo de potencia $n$-ésima de
$G$ es universalmente abierto. Entonces $G^\circ$ es abierto; y si, además,
los $G_s^\circ$ no tienen componente aditiva (por ejemplo, si los
$G_s^\circ$ son propios, como antes), entonces $G^\tau\longrightarrow S$ es
universalmente abierto, e incluso liso si $S$ es reducido.

Consignemos finalmente el siguiente resultado fácil:

\result{PROPOSICIÓN}{1.11} Existe un abierto $U$ de $S$ tal que el conjunto
$G'$ de los puntos de $G$ en los que $G$ es liso (respectivamente, plano) sobre
$S$ es el conjunto subyacente de un subesquema abierto en grupos inducido por
$G|U$. Además, toda sección de $G$ sobre $U$ es una sección de $G'$ sobre $U$.

\src{231}
\result{COROLARIO}{1.12} Si $G$ es liso (respectivamente, plano) sobre $S$ en
los puntos de la sección neutra, también lo es en los puntos de toda sección de
$G$ sobre $S$ y en todo punto de $G^\circ$. Si, además, para todo entero $n>0$
el homomorfismo de potencia $n$-ésima
$\varphi_n:G\longrightarrow G$ es étale en los puntos de característica prima
con $n$, entonces $G$ es liso (respectivamente, plano) sobre $S$ en todo punto
de $G^\sigma$.

\fsection{2}{Aplicación a las propiedades locales de los esquemas de Picard}

\result{TEOREMA}{2.1}

\noindent(i) Sea $f:X\longrightarrow S$ un morfismo propio y liso, y
supongamos que existe $\operatorname{Pic}_{X/S}$ (por ejemplo, supongamos que
$f$ es proyectivo). Entonces $\operatorname{Pic}_{X/S}$ es separado sobre $S$
y todo subconjunto cerrado $Z$ de $\operatorname{Pic}_{X/S}$ que sea de tipo
finito sobre $S$ es propio sobre $S$.

\noindent(ii) Sea $X$ un preesquema propio geométricamente normal sobre un
cuerpo $k$. Entonces $\operatorname{Pic}_{X/k}^\circ$\footnote{El texto
impreso dice $\operatorname{Pic}_{X/S}^\circ$, aunque $S$ no aparece en este
enunciado sobre el cuerpo $k$.} es propio sobre $k$.

\begin{proof}
(i) Los criterios valuativos ([2], II, §7) nos reducen a demostrar lo
siguiente. Sea $S$ el espectro de un anillo de valoración discreta completo y
sea $U$ el abierto formado por el punto genérico de $S$. Toda sección racional
de $\operatorname{Pic}_{X/S}$ sobre $S$, es decir, toda sección sobre $U$, se
prolonga de manera única a una sección sobre $S$. Por la definición de
$\operatorname{Pic}_{X/S}$, esto equivale al enunciado siguiente: para todo
Módulo invertible $\mathcal L$ sobre $V=f^{-1}(U)$ existe un Módulo invertible sobre $X$
que prolonga $\mathcal L$, único salvo isomorfismo. Esto se deduce fácilmente
de la descripción de los Módulos invertibles sobre $V$ y, respectivamente, sobre
$X$, en términos de clases de divisores de Cartier. En efecto, los anillos
locales de $X$ son regulares ($X$ es liso sobre el esquema regular $S$) y, por
tanto, factoriales por el teorema de Auslander; así, todo divisor sobre
$X$\footnote{El texto impreso dice «sobre $S$». La factorialidad invocada es
la de los anillos locales de $X$.} es un divisor de Cartier. Todo divisor sobre
$V$ se prolonga a un divisor sobre $X$ tomando su clausura.
\end{proof}

\result{OBSERVACIÓN}{2.2} La demostración sigue siendo válida si se supone
solamente que $f$ es plano y que sus fibras $X_s$ son intersecciones completas
locales y lisas sobre $k(s)$ en codimensión $\leq2$. Esto se deduce del teorema
demostrado en [5]: un anillo local noetheriano de intersección completa que sea
regular en codimensión $\leq3$ es factorial («conjetura de Samuel»). Obsérvese
que el

\src{232}
resultado pasa a ser falso si se reemplaza «codimensión $\leq2$» por
«codimensión $\leq1$», es decir, por la hipótesis «normal», como muestra una
familia de cuádricas no singulares que degenera en un cono cuadrático.

(ii) El lema de Chow nos reduce al caso en que $X$ es proyectivo y, por tanto,
está inmerso en $\mathbf P_k^n$; podemos suponer que $X$ es conexo. Si
$\dim X=1$, entonces $X$ es liso sobre $k$ y se aplica (i). Si $\dim X\geq2$,
una variante de los «criterios de equivalencia» conocidos muestra que existe
un número finito de curvas lisas $Y_i$ contenidas en $X$ (obtenidas como
intersecciones de $X$ con subespacios lineales adecuados de $\mathbf P_k^n$),
tales que
\[
\operatorname{Pic}_{X/k}^\tau\longrightarrow
\prod_i\operatorname{Pic}_{Y_i/k}^\tau
\]
es un monomorfismo y, por tanto, induce un monomorfismo sobre las componentes
conexas. La componente conexa del segundo miembro es propia sobre $k$ por el
resultado anterior, y el homomorfismo de esquemas en grupos considerado es
necesariamente una inmersión cerrada. Se sigue que
$\operatorname{Pic}_{X/k}^\circ$ es también propio sobre $k$. Los delicados
criterios de equivalencia pueden evitarse utilizando la estructura de los
grupos algebraicos conmutativos sobre un cuerpo algebraicamente cerrado
(debida a Chevalley--Borel). Se reduce a demostrar que todo morfismo de la
recta afín privada del origen\footnote{El texto impreso dice «de la recta
afín». La fe de erratas vinculada, p.~303, prescribe «recta afín privada del
origen».} a $\operatorname{Pic}_{X/k}^\tau$ es constante. Esto equivale a decir
que todo Módulo invertible sobre $X[t,t^{-1}]$\footnote{El texto impreso dice
$X[t]$; la fe de erratas vinculada, p.~303, prescribe $X[t,t^{-1}]$.} procede
de un Módulo invertible sobre $X$. Este conocido resultado elemental ni siquiera
utiliza que $X$ sea propio sobre $k$ (la normalidad de $X$ reduce
inmediatamente al caso en que $X$ es regular).

\result{COROLARIO}{2.3} Sea $f:X\longrightarrow S$ un morfismo propio normal
(es decir, plano con fibras geométricamente normales), y supongamos que existe
$\operatorname{Pic}_{X/S}$. Entonces $\operatorname{Pic}_{X/S}^\circ$ es propio
sobre $S$ y, por tanto, cerrado en $\operatorname{Pic}_{X/S}$. Además,
$\operatorname{Pic}_{X/S}^\tau$ y $\operatorname{Pic}_{X/S}^\rho$ son también
cerrados, al igual que $\operatorname{Pic}_{X/S}^\sigma$ en «característica
igual».

Aplíquense (1.1) y (2.1), (ii).

\result{COROLARIO}{2.4} Sea $f:X\longrightarrow S$ un morfismo propio y liso
tal que exista $\operatorname{Pic}_{X/S}$ y que este sea una unión disjunta de
esquemas $P^{(i)}$ de tipo finito sobre $S$ (cf. [3], V, 4.1). Entonces cada
$P^{(i)}$ es propio sobre $S$.

Esto se deduce de (2.1), (i). Como se señaló en (2.2), el resultado puede
generalizarse imponiendo hipótesis menos restrictivas sobre las fibras de $f$,
pero pasa a ser falso si se supone meramente que $f$ es normal. En ese caso,
no sé si $\operatorname{Pic}_{X/S}^\tau$ es, no obstante, propio sobre $S$,
incluso suponiendo que sea de tipo finito sobre $S$.

\src{233}
\result{TEOREMA}{2.5} Sea $f:X\longrightarrow S$ un morfismo propio y plano tal
que exista $\operatorname{Pic}_{X/S}$, y para cada entero $n$ sea $\varphi_n$
el homomorfismo de potencia $n$-ésima de este preesquema en grupos. Entonces
$\varphi_n$ es étale en todo punto
$x\in\operatorname{Pic}_{X/S}$\footnote{El texto impreso dice «$x\in X$». La
aplicación $\varphi_n$ es el endomorfismo de potencia del preesquema de
Picard.} cuya característica residual sea prima con $n$.

Por la caracterización infinitesimal de los morfismos étales, esta afirmación
equivale a la siguiente:

\result{LEMA}{2.6} Supongamos que $S$ es el espectro de un anillo artiniano
local $A$ cuyo ideal maximal $\mathfrak m$ tiene potencia $(\nu+1)$-ésima nula,
y pongamos
\[
A_{\nu-1}=A/\mathfrak m^\nu,
\qquad X_{\nu-1}=X\otimes_A A_{\nu-1}.
\]
Sea $\mathcal L$ un Módulo invertible sobre $X$ y sea
$\mathcal L'_{\nu-1}$ un Módulo invertible sobre $X_{\nu-1}$ cuya potencia
tensorial $n$-ésima sea isomorfa a
\[
\mathcal L_{\nu-1}=\mathcal L\otimes_A A_{\nu-1}.
\]
Entonces existe un Módulo invertible $\mathcal L'$ sobre $X$ cuya potencia
tensorial $n$-ésima es isomorfa a $\mathcal L$, siempre que $n$ sea primo con
la característica residual de $k=k(A)$.

Demostremos el lema. Pongamos
\[
V=\mathfrak m^\nu=\mathfrak m^\nu/\mathfrak m^{\nu+1};
\]
este es un espacio vectorial sobre $k=k(A)$. Intentemos primero prolongar
$\mathcal L'_{\nu-1}$ a un Módulo invertible arbitrario $\mathcal L'$ sobre $X$.
La obstrucción pertenece a
\[
H^2(X_0,\mathcal O_{X_0})\otimes_k V.
\]
La hipótesis sobre $\mathcal L'_{\nu-1}$ y el hecho de que
$\mathcal L_{\nu-1}$ pueda prolongarse muestran que $n$ veces esta obstrucción
es cero; la propia obstrucción es, por tanto, nula porque $n$ es primo con la
característica. La ambigüedad en la prolongación elegida pertenece a
\[
H^1(X_0,\mathcal O_{X_0})\otimes_k V,
\]
y la diferencia $\xi$ entre $\mathcal L'^{\otimes n}$ y $\mathcal L$ pertenece
al mismo Módulo. Corregir $\mathcal L'$ para anular esta diferencia equivale a
encontrar un $\eta$ en ese Módulo tal que $n\eta=\xi$. Esto vuelve a ser
posible porque $n$ es primo con la característica.

\result{COROLARIO}{} Bajo las hipótesis de (2.5), supongamos además que los
esquemas de Picard de las fibras $X_s$ no contienen ninguna componente aditiva
(por ejemplo, que los $X_s$ son geométricamente normales; cf. (2.1), (ii)).
Entonces $\operatorname{Pic}_{X/S}\longrightarrow S$ es universalmente abierto
en los puntos de $\operatorname{Pic}_{X/S}^\circ$. Si
$\operatorname{Pic}_{X/S}^\circ$ es cerrado (por ejemplo, si los $X_s$ son
geométricamente normales; cf. (2.3)), entonces
$\operatorname{Pic}_{X/S}^\sigma$ es universalmente abierto sobre $S$. Por
último, en característica igual,
$\operatorname{Pic}_{X/S}^\rho\longrightarrow S$ es universalmente abierto.

Aplíquense (1.5) y (1.1).

\src{234}
\result{COROLARIO}{2.7} Sea $f:X\longrightarrow S$\footnote{El texto impreso
dice «$f:X\longrightarrow Y$», aunque tanto el preesquema de Picard como la
función siguiente tienen base $S$.} un morfismo propio y plano tal que exista
$\operatorname{Pic}_{X/S}$. Entonces la función
\[
s\longmapsto\dim\operatorname{Pic}_{X_s/k(s)}
\]
sobre $S$ es semicontinua superiormente (es decir, puede saltar hacia arriba,
pero no hacia abajo), e incluso continua (es decir, localmente constante) si
los $\operatorname{Pic}_{X_s/k(s)}$ no contienen ninguna componente aditiva.

La primera afirmación es trivial, o casi, para todo preesquema en grupos
localmente de tipo finito sobre una base localmente noetheriana, pues basta
mirar a lo largo de la sección neutra. La segunda afirmación se deduce de
(2.5).

\result{OBSERVACIÓN}{2.8} Sean $s,s'\in S$, con $s$ una especialización de
$s'$. Entonces (2.7) equivale a una desigualdad (respectivamente, una igualdad)
entre las dimensiones de los esquemas de Picard de $X_{s'}$ y de su
«especialización» $X_s$. Incluso antes de que se construyeran los esquemas de
Picard, Serre había observado que, para un morfismo liso
$f:X\longrightarrow S$, la invariancia de la dimensión de los Picard de los
$X_s$ es una consecuencia formal de la teoría de la especialización del grupo
fundamental ([4], X) y de las relaciones clásicas de Kummer entre los puntos de
orden finito de la variedad de Picard y el grupo fundamental abelianizado
([4], XI). En general, sean $\alpha,\mu,\lambda$ las dimensiones de las partes
abeliana, multiplicativa y aditiva de
$\operatorname{Pic}_{X_s/k(s)}^\circ$, y definamos de modo análogo
$\alpha',\mu',\lambda'$. Las relaciones conocidas se expresan mediante las
desigualdades siguientes:
\[
\tag{*}\alpha+\mu+\lambda\geq\alpha'+\mu'+\lambda'
\]
(válida siempre que exista $\operatorname{Pic}_{X/S}$ y, por tanto,
probablemente en todos los casos). Cuando $\lambda=\lambda'=0$, esta
desigualdad se convierte en una igualdad, válida bajo las mismas hipótesis de
existencia:
\[
\alpha+\mu=\alpha'+\mu';
\]
por otra parte,
\[
\tag{**}2\alpha+\mu\leq2\alpha'+\mu'
\]

\src{235}
por el argumento de Serre si los $X_s$ son separables (sin suponer siquiera la
existencia de $\operatorname{Pic}_{X/S}^\tau$), o si los
$ {}_n\operatorname{Pic}_{X/S}$ (los núcleos de $\varphi_n$ en el funtor de
Picard) son separados sobre el abierto de $S$ donde $n$ es invertible (pues son étales sobre ese abierto por (2.5)). Se
inclina uno a conjeturar que (*) es siempre una igualdad, o al menos cuando los
$X_s$ son separables, y que
\[
\tag{***}\alpha\leq\alpha',\qquad\lambda\geq\lambda'
\]
se verifica siempre que exista un preesquema en grupos, localmente de tipo
finito sobre un $S$ localmente noetheriano, cuyas fibras tengan dimensión
constante (véase (1.3) para un resultado positivo en esta dirección).

\result{OBSERVACIÓN}{2.9} En todos los casos conocidos,
$\operatorname{Pic}_{X/S}^\tau$ es universalmente abierto sobre $S$, pero
probablemente no conviene dejarse llevar por ilusiones excesivas, ni siquiera
cuando $f:X\longrightarrow S$ es liso. En todo caso, Mumford ha construido un
ejemplo (ciertamente con $S$ no reducido, de hecho el espectro de un anillo
artiniano) en el que $\operatorname{Pic}_{X/S}^\tau$ no es plano sobre $S$,
haciendo variar infinitesimalmente la superficie de Igusa. El punto considerado
por Mumford pertenece a $\operatorname{Pic}_{X/S}^\rho$, y sigue siendo posible
que, para un $f:X\longrightarrow S$ liso, $\operatorname{Pic}_{X/S}^\tau$ sea
plano sobre $S$ en los puntos de $\operatorname{Pic}_{X/S}^\sigma$. El ponente
duda, sin embargo, que esto sea siempre así, incluso si se restringe a los
puntos de $\operatorname{Pic}_{X/S}^\circ$ y se supone que $S$ es el espectro
de un anillo de valoración discreta. La cuestión está también relacionada con
el estudio de los puntos fijos de un esquema abeliano bajo un grupo finito de
automorfismos, situación para la que parecen faltar ejemplos.

Parecería que, incluso restringiéndose a un morfismo proyectivo liso
$f:X\longrightarrow S$, los resultados de regularidad local para
$\operatorname{Pic}_{X/S}$ enunciados en este parágrafo, junto con las
conjeturas de (2.8), agotan casi todo lo que puede decirse sin hipótesis más
específicas sobre la naturaleza de las fibras de $f$. Recuérdese, sin embargo,
que si las fibras geométricas de $\operatorname{Pic}_{X/S}$ son reducidas y no
tienen componente aditiva, entonces (1.8) y (2.5) implican que
$\operatorname{Pic}_{X/S}$ es liso sobre $S$ en los puntos de
$\operatorname{Pic}_{X/S}^\sigma$ cuando $S$ es reducido. El mismo resultado se
verifica, sin ninguna hipótesis sobre $S$, si $f:X\longrightarrow S$ es normal,
como veremos en (3.5). Consignemos a este respecto:

\result{PROPOSICIÓN}{2.10}

\noindent(i) Si $\operatorname{Pic}_{X/S}$ es liso (respectivamente, plano)
sobre $S$ en los puntos de la sección neutra, también lo es en todo punto de
$\operatorname{Pic}_{X/S}^\sigma$ y en los puntos de toda sección de
$\operatorname{Pic}_{X/S}$ sobre $S$.

\src{236}
\noindent(ii) Sea $s\in S$ tal que
\[
H^2(X_s,\mathcal O_{X_s})=0.
\]
Entonces existe un entorno abierto $U$ de $s$ tal que
$\operatorname{Pic}_{X/S}|U$ es liso sobre $U$.

\noindent(iii) Si $X$ es un esquema propio sobre un cuerpo, entonces
\[
\dim\operatorname{Pic}_{X/k}\leq\dim H^1(X,\mathcal O_X),
\]
con igualdad si y solo si $\operatorname{Pic}_{X/k}$ es liso sobre $k$; este
es siempre el caso cuando $k$ tiene característica cero.

La parte (i) se deduce de (2.5) y (1.12), y (ii) se deduce del criterio
infinitesimal de lisura y de un cálculo conocido de obstrucciones, pues el
«teorema de semicontinuidad» muestra que la hipótesis en $s$ sigue siendo
válida en los puntos próximos. Por último, (iii) se deduce de que
$H^1(X,\mathcal O_X)$ es isomorfo al espacio tangente de Zariski de
$\operatorname{Pic}_{X/k}$ en el elemento neutro. La última afirmación es un
caso particular de un teorema de Cartier que dice que un «grupo formal» en
característica cero es formalmente liso sobre $k$.

\fsection{3}{El subesquema abeliano canónico de
$\operatorname{Pic}_{X/S}$ y el esquema de Albanese}

\result{PROPOSICIÓN}{3.1} Sea $k$ un cuerpo y sea $G$ un esquema conmutativo en
grupos de tipo finito sobre $k$, sin componente aditiva. Entonces
$G_{\mathrm{red}}^\circ$ es separable sobre $k$ y, por consiguiente, es un
subesquema en grupos liso sobre $k$.

Esto es trivial cuando $k$ es perfecto, en particular para $G_{\bar k}$, donde
$\bar k$ es una clausura algebraica de $k$. Basta, por tanto, ver que
$(G_{\bar k}^\circ)_{\mathrm{red}}$ desciende de un subesquema de $G$. De la
hipótesis sobre $G_{\bar k}$ (ausencia de componente aditiva)\footnote{El texto
impreso contiene una construcción mal formada equivalente a «La hipótesis ...
resulta». La fe de erratas vinculada, p.~303, inserta la preposición necesaria
para «De la hipótesis ... resulta».} se deduce fácilmente que existe un entero
$m$ tal que $(G_{\bar k}^\circ)_{\mathrm{red}}$ es la imagen esquemática del
homomorfismo de potencia $m$-ésima de $G_{\bar k}$. Como este homomorfismo
procede del homomorfismo análogo de $G$, la imagen esquemática de este último
da el subesquema buscado.

\result{COROLARIO}{3.2} Sea $X$ un esquema geométricamente normal propio sobre $k$ tal que
exista $\operatorname{Pic}_{X/k}$. Entonces existe un subesquema abeliano $A$
de $\operatorname{Pic}_{X/k}$ cuyo conjunto subyacente es
$\operatorname{Pic}_{X/k}^\circ$.

En efecto, por (2.1), (ii), $\operatorname{Pic}_{X/k}^\circ$ es propio sobre
$k$ y, por tanto, satisface la condición de (3.1).

\src{237}
Así, el resultado anterior muestra que, en ciertos casos, la «variedad de
Picard» clásica (que en la presente teoría se denotaría
$(\operatorname{Pic}_{X/k})_{\mathrm{red}}^\circ$) «está definida sobre $k$»,
sin suponer que $k$ sea perfecto.

Sea ahora $f:X\longrightarrow S$ un esquema relativo propio y plano, con
\[
\mathcal O_S\xrightarrow{\sim}f_*(\mathcal O_X)
\]
para simplificar, tal que exista $\operatorname{Pic}_{X/S}$ y
$\operatorname{Pic}_{X/S}^\circ$ sea propio sobre $S$. Para (3.3), (ii),
suponemos además que existe un abierto de $\operatorname{Pic}_{X/S}$ que
contiene $\operatorname{Pic}_{X/S}^\circ$ y es cuasiproyectivo sobre $S$. Como
se vio anteriormente, esta condición se verifica cuando $f$ es proyectivo con
fibras geométricamente separables e irreducibles. Recuérdese que un esquema
abeliano sobre $S$ es un esquema en grupos sobre $S$ propio y liso sobre $S$,
con fibras geométricamente conexas. Preguntamos si existe un subesquema en
grupos $A$ de $\operatorname{Pic}_{X/S}$ que sea un esquema abeliano y cuyo
conjunto subyacente sea $\operatorname{Pic}_{X/S}^\circ$. Acabamos de ver que
siempre existe uno cuando $S$ es el espectro de un cuerpo. He aquí lo que puede
decirse en el caso general considerado:

\result{TEOREMA}{3.3} Bajo las hipótesis anteriores:

\noindent(i) Si existe un subesquema abeliano de
$\operatorname{Pic}_{X/S}$ cuyo conjunto subyacente sea
$\operatorname{Pic}_{X/S}^\circ$, es único. Por consiguiente, su formación es
compatible con el cambio de base.

\noindent(ii) Para que exista tal subesquema abeliano, basta que exista después
de todo cambio de base $S'\longrightarrow S$ con $S'$ artiniano local. Si $S$
es el espectro de un anillo local, basta comprobar los
$S'=\operatorname{Spec}(A_n)$, donde $A_n=A/\mathfrak m^{n+1}$. Si $S$ es
reducido, basta asimismo comprobar los cambios de base
$S'\longrightarrow S$ en los que $S'$ es el espectro de un anillo de
valoración discreta (completo y con cuerpo residual algebraicamente cerrado,
si se desea).

\noindent(iii) Supongamos que existe $A$, y sea
$B=\operatorname{Alb}^0(X/S)$ el esquema abeliano dual (es decir,
$B=\operatorname{Pic}_{A/S}^\circ$ [7]). Se puede construir entonces,
canónicamente, un espacio homogéneo principal
$P=\operatorname{Alb}^1(X/S)$ bajo $B$, junto con un $S$-morfismo
$X\longrightarrow P$ universal para los $S$-morfismos de $X$ a esquemas
para-abelianos (es decir, espacios homogéneos principales bajo esquemas
abelianos). La formación de $\operatorname{Alb}^0(X/S)$,
$\operatorname{Alb}^1(X/S)$ y
$X\longrightarrow\operatorname{Alb}^1(X/S)$ conmuta con el cambio de base.

Esbocemos la demostración.

La parte (i) es una propiedad general de rigidez de los subesquemas abelianos
de los esquemas conmutativos en grupos: si dos de esos subesquemas coinciden
como conjuntos en un punto $s\in S$, coinciden sobre toda la componente conexa

\src{238}
de $s$ [7]. (Este resultado generaliza un teorema clásico de Chow.)

(ii) Los esquemas de Hilbert muestran que el funtor que asocia a cada $S'$ sobre
$S$ el conjunto (con uno o cero elementos) de los subesquemas abelianos
canónicos de $(\operatorname{Pic}_{X/S})\times_S S'$ está representado por un
esquema $T$ de tipo finito sobre $S$. Por (i), $T\longrightarrow S$ es un
monomorfismo y, por (3.2), es sobreyectivo. La existencia de un subesquema
abeliano canónico de $\operatorname{Pic}_{X/S}$ significa que $T$ posee una
sección sobre $S$, o equivalentemente que $T\longrightarrow S$ es un
isomorfismo. Esto equivale también a que $T\longrightarrow S$ sea étale o,
cuando $S$ es reducido, a que $T\longrightarrow S$ sea propio. Esto demuestra
(ii).

(iii) Directamente de la definición de $\operatorname{Pic}_{X/S}$, para todo
esquema abeliano $C$ sobre $S$, dar un $S$-morfismo de $X$ a un espacio
homogéneo principal bajo $C$ equivale a dar un homomorfismo de grupos
$C'\longrightarrow\operatorname{Pic}_{X/S}$, donde $C'$ es el esquema abeliano
dual de $C$. Si existe el subesquema abeliano canónico $A$ de
$\operatorname{Pic}_{X/S}$, estos homomorfismos se factorizan necesariamente a
través de $A$ (como puede verse también utilizando los puntos de orden finito).
Esto demuestra (iii).

\result{OBSERVACIONES}{3.4} Denotaremos por
$\operatorname{Pic}_{X/S}^{\circ\circ}$ el subesquema abeliano canónico de
$\operatorname{Pic}_{X/S}$, cuando exista. Por desgracia, no siempre existe,
como se ve haciendo variar infinitesimalmente la superficie de Igusa (mediante
la deformación modular de primer orden). Es posible, no obstante, que
$\operatorname{Pic}_{X/S}^{\circ\circ}$ exista al menos cuando $S$ sea reducido
o, equivalentemente por (ii), cuando $S$ sea el espectro de un anillo de
valoración discreta. Sean $X_0,X_1$ las fibras especial y genérica de $X$, y
pongamos
\[
A_1=\operatorname{Pic}_{X_1/K}^\circ,
\]
donde $K$ es el cuerpo de fracciones del anillo de valoración $V$. Por el
teorema de Koizumi, existe un esquema abeliano $A$ sobre $S$, esencialmente
único, cuya fibra genérica es $A_1$. Como en (2.1), (i), se ve fácilmente
(suponiendo ahora que $X$ es liso sobre $S$) que el morfismo identidad de $A_1$
se prolonga a un morfismo
\[
A\longrightarrow\operatorname{Pic}_{X/S}.
\]
De ahí se obtiene un homomorfismo
\[
\tag{*}A_0\longrightarrow\operatorname{Pic}_{X_0/k}^{\circ\circ},
\]
que se demuestra fácilmente que es sobreyectivo y tiene por núcleo un
$p$-grupo finito, donde $p$ es el exponente característico del cuerpo residual
$k$

\src{239}
(de nuevo utilizando los puntos de orden finito). Sentado esto, las condiciones
siguientes sobre $X/S$ son equivalentes:

\medskip
\noindent a. El homomorfismo anterior (*) es un isomorfismo (lo que Shimura
expresaría diciendo que la formación de la «variedad de Picard» es «compatible
con la especialización»).

\noindent b. Existe $\operatorname{Pic}_{X/S}^{\circ\circ}$ (en cuyo caso no es
otro que $A$).

\noindent c. (Para referencia) Existen los
$\operatorname{Pic}_{X_n/S}^{\circ\circ}$.

Por la observación hecha sobre el núcleo de (*), la condición (a) se verifica
cuando la característica residual es cero, pero este resultado se generalizará
considerablemente en (3.5).

Por supuesto, si $\operatorname{Pic}_{X/S}$ es liso sobre $S$ en los puntos de
$\operatorname{Pic}_{X/S}^\circ$, este último es abierto en
$\operatorname{Pic}_{X/S}$ (cf. (1.7)). Provisto de la estructura inducida, es
por tanto un subesquema abeliano de $\operatorname{Pic}_{X/S}$, luego igual a
$\operatorname{Pic}_{X/S}^{\circ\circ}$, que existe en este caso. Pero se tiene
algo mucho más fuerte:

\result{TEOREMA}{3.5} Bajo las hipótesis de (3.3), sea $s\in S$ tal que
$\operatorname{Pic}_{X_s/k(s)}$ sea liso sobre $k(s)$ (o, equivalentemente,
\[
\dim\operatorname{Pic}_{X_s/k(s)}=\dim H^1(X_s,\mathcal O_{X_s})).
\]
Entonces existe un entorno abierto $U$ de $s$ tal que
$\operatorname{Pic}_{X/S}$ es liso sobre $S$ en los puntos de
$\operatorname{Pic}_{X/S}^\circ|U$, que es, por consiguiente, un subesquema
abeliano abierto de $\operatorname{Pic}_{X/S}|U$. A fortiori, existe
$\operatorname{Pic}_{X|U/U}^{\circ\circ}$.

Indiquemos el principio de la demostración. La discusión anterior nos reduce
al caso en que $S$ es el espectro de un anillo artiniano local $A$, y se razona
por inducción sobre el orden infinitesimal de $A$. Podemos suponer, por tanto,
que $\operatorname{Pic}_{X_n/A_n}^\circ$ es liso sobre $A_n$ y basta demostrar
que $\operatorname{Pic}_{X_{n+1}/A_{n+1}}^\circ$ es liso sobre $A_{n+1}$. Para
ello basta construir un esquema abeliano $P_{n+1}$ sobre $A_{n+1}$ que prolongue
\[
P_n=\operatorname{Pic}_{X_n/A_n}^\circ,
\]
junto con un Módulo invertible $\mathcal L_{n+1}$ sobre
$X_{n+1}\times_{A_{n+1}}P_{n+1}$ que prolongue el Módulo invertible
$\mathcal L_n$ sobre $X_n\times_{A_n}P_n$ que interviene en la definición del
esquema de Picard $\operatorname{Pic}_{X_n/A_n}$ como solución de un problema
universal. (N.B. --- Podemos suponer que $X$ posee una sección sobre $S$~...)
La construcción utiliza el resultado clave siguiente: todo esquema abeliano
definido sobre un cociente de un anillo artiniano local puede prolongarse (en
otros términos, el «esquema formal de Módulos» absoluto ([3], II) de un esquema
abeliano sobre un cuerpo algebraicamente cerrado es liso sobre el anillo de

\src{240}
vectores de Witt de $k$). Este resultado se deduce, a su vez, de las propiedades
formales generales de la obstrucción al levantamiento y de las operaciones de
grupo. Admitiéndolo, se prolonga primero $P_n$ arbitrariamente a $P_{n+1}$. Se
encuentra entonces una obstrucción al levantamiento de $\mathcal L_n$, que
pertenece a
\[
H^2(X_0\times P_0,\mathcal O_{X_0\times P_0})\otimes_k V
\qquad(\text{donde }V=\mathfrak m^{n+1}/\mathfrak m^{n+2}),
\]
y, de hecho, al subespacio
\[
H^1(X_0,\mathcal O_{X_0})\otimes
H^1(P_0,\mathcal O_{P_0})\otimes_k V
\]
(porque la restricción de $\mathcal L_n$ a cada uno de los dos factores $X_n$
y $P_n$ es trivial). Este último espacio es precisamente
\[
H^1(P_0,\mathcal T_{P_0/k})\otimes V,
\]
donde $\mathcal T_{P_0/k}$ es el haz tangente de $P_0/k$, y es también, por
tanto, el espacio que mide la indeterminación en el levantamiento de $P_n$ a
$P_{n+1}$ ([3], II). Por consiguiente, este levantamiento puede corregirse (de
una y solo una manera, como debe ser) para anular la obstrucción al
levantamiento de $\mathcal L_n$.

\result{COROLARIO}{3.6} Bajo las hipótesis de (3.5),
$R^1f_*(\mathcal O_X)$ es un Módulo localmente libre sobre $S$ en un entorno de
$s$, y su formación conmuta con el cambio de base.

En efecto, este Módulo es el Módulo tangente de $\operatorname{Pic}_{X/S}$ a lo largo
de la sección neutra.

\result{OBSERVACIÓN}{3.7} El mismo argumento que en (3.5) muestra que, si $S'$
es un subesquema de $S$ definido por un Ideal coherente nilpotente y se supone
solamente que existe $\operatorname{Pic}_{X'/S'}^\circ$ y que tiene fibras
lisas, entonces existe necesariamente $\operatorname{Pic}_{X/S}^\circ$ y es un
esquema abeliano sobre $S$. Esto permite construir esquemas de Picard en
algunos casos pese a la ausencia de una hipótesis de proyectividad; por
ejemplo, siempre existe el esquema abeliano dual de un esquema abeliano sobre
un anillo artiniano. Por otra parte, aplicando (3.6) cuando $X$ es un esquema
abeliano sobre $S$ y utilizando la estructura conocida de álgebra exterior de
$H^*(X_s,\mathcal O_{X_s})$ sobre $H^1(X_s,\mathcal O_{X_s})$
(Rosenlicht--Serre), se obtiene que, para todo $i$, $R^if_*(\mathcal O_X)$ es
localmente libre y, más precisamente, isomorfo a la potencia exterior $i$-ésima
de $R^1f_*(\mathcal O_X)$.

\result{OBSERVACIÓN}{3.8} Para un morfismo proyectivo liso
$f:X\longrightarrow S$, con $S$ reducido y todas las características
residuales nulas, (3.6) era conocido por métodos trascendentes como consecuencia
de la teoría de Hodge. De hecho, todos los $R^pf_*(\Omega^q_{X/S})$ son
entonces localmente libres. Por otra parte, existen contraejemplos

\src{241}
en característica mixta a la libertad local de $R^1f_*(\mathcal O_X)$,
obtenidos a partir de variedades de Serre ([8], p.~50). No parece haber
contraejemplos en característica igual.

\result{COROLARIO}{3.8} Bajo las hipótesis de (3.5),
$\operatorname{Pic}_{X/S}|U$ es liso sobre $U$ en todo punto de
$\operatorname{Pic}_{X/S}^\sigma|U$.

Aplíquese (2.10), (i). En particular, teniendo en cuenta (2.10), (iii):

\result{COROLARIO}{3.9} Supongamos que todas las características residuales de
$S$ son nulas. Entonces $\operatorname{Pic}_{X/S}^\tau$ es liso sobre $S$.

\result{OBSERVACIÓN}{3.10} Si $\operatorname{Pic}_{X/S}^\tau$ es además propio sobre
$S$, se sigue, por ejemplo, que el subgrupo de torsión de los grupos de
Néron--Severi de las fibras geométricas de $f$ es constante sobre cada
componente conexa de $S$ (como también resulta claro por métodos trascendentes
cuando $f$ es liso y proyectivo). Una aplicación directa de (2.5) muestra, más
generalmente, que si $\operatorname{Pic}_{X/S}^\tau$ es propio sobre $S$ (por
ejemplo, si $f:X\longrightarrow S$ es liso y proyectivo) y $q$ es un primo
distinto de todas las características residuales de $S$, entonces la
componente $q$-primaria de los subgrupos de torsión de los grupos de
Néron--Severi de las fibras geométricas de $X/S$ es constante sobre cada
componente conexa de $S$. En característica $p>0$, esto deja de ser cierto para
la componente $p$-primaria de la torsión. Sigue siendo posible, sin embargo,
que el rango sobre $k(s)$ de
\[
\operatorname{Pic}_{X_s/k(s)}^\tau/
\operatorname{Pic}_{X_s/k(s)}^{\circ\circ}=T_{X_s/k(s)}
\]
sea localmente constante. Cuando $S$ es reducido, se puede demostrar que esto
equivale a la existencia de $\operatorname{Pic}_{X/S}^{\circ\circ}$ junto con
la planitud de $\operatorname{Pic}_{X/S}^\tau$ sobre $S$, y basta comprobar la
cuestión cuando $S$ es el espectro de un anillo de valoración discreta. Esto es
lo que he comprobado en los pocos ejemplos que he examinado; pero como el
enunciado correspondiente con $S$ artiniano es falso (cf. Observaciones (2.9)
y (3.4)), probablemente no convenga entregarse a un optimismo excesivo.

\fsection{4}{El teorema de finitud para el esquema de Picard}

Sea $f:X\longrightarrow S$ un morfismo proyectivo y plano tal que exista
$\operatorname{Pic}_{X/S}$. La consideración de los «polinomios de Hilbert»
$Q\in\mathbf Q[t]$ descompone $\operatorname{Pic}_{X/S}$ en una unión disjunta
de abiertos $\operatorname{Pic}_{X/S}^Q$. Sin hipótesis adicionales, por
ejemplo que aseguren que $\operatorname{Pic}_{X/S}$ sea separado sobre $S$,
estos abiertos no son necesariamente de tipo finito sobre $S$; se obtiene un
contraejemplo cuando $X$ es «una cónica que degenera en dos rectas». Es

\src{242}
posible, sin embargo, que sean de tipo finito cuando $f$ sea separable y tenga
fibras geométricamente irreducibles. La cuestión está relacionada con saber si
$\operatorname{Pic}_{X/S}^\tau$ es de tipo finito sobre $S$, lo que podría ser
cierto sin ninguna hipótesis sobre las fibras de $X/S$. Cuando
$f:X\longrightarrow S$ es liso, obsérvese que
$\operatorname{Pic}_{X/S}^\tau$ está contenido en uno de los
$\operatorname{Pic}_{X/S}^Q$ (pues sobre una variedad proyectiva no singular la
equivalencia por torsión es más fina que la equivalencia numérica para los
divisores y, de hecho, por Matsusaka, coincide con ella), y es, por tanto, de
tipo finito sobre $S$ si los $\operatorname{Pic}_{X/S}^Q$ lo son. Las
cuestiones de finitud que acabamos de plantear conservan un sentido evidente
incluso sin suponer la existencia de $\operatorname{Pic}_{X/S}$: dicen que
ciertas familias de Módulos invertibles son «acotadas» en el sentido de [3], IV.
Para todo cuerpo algebraicamente cerrado $k$, considérense los subesquemas
íntegros de $\mathbf P_k^r$ de dimensión $n$ y grado $d$, y los Módulos
invertibles sobre esos preesquemas con polinomio de Hilbert $Q$, donde
$r,n,d,Q$ están fijados. El problema consiste en demostrar que la familia de
estos Módulos (considerados como Módulos coherentes sobre los $\mathbf P_k^r$) es
acotada, es decir, que puede parametrizarse mediante un esquema de tipo finito
sobre $\mathbf Z$.

Utilizando el método de Matsusaka [6], una demostración bastante técnica (que
emplea los «criterios de equivalencia» en una forma adecuada) da una respuesta
afirmativa cuando uno se restringe a las subvariedades no singulares de
$\mathbf P_k^r$. Más precisamente:

\result{TEOREMA}{4.1} Sea $f:X\longrightarrow S$\footnote{El texto impreso
dice «Sea: $X\longrightarrow S$»; el morfismo utilizado más adelante es $f$.}
un morfismo proyectivo liso con fibras geométricamente conexas, siendo $S$
noetheriano. Sea $\mathcal O_X(1)$ un Módulo relativamente muy amplio sobre $X$,
y sea $E$ un subconjunto de $\operatorname{Pic}_{X/S}$ correspondiente a una
familia $(\mathcal L_i)$ de Módulos invertibles sobre las fibras geométricas
$\overline X_{s_i}$ de $X/S$. Sea $D_i$ un divisor (no necesariamente
efectivo) sobre $\overline X_{s_i}$ que defina $\mathcal L_i$, sea
\[
a_n^{(i)}t^n+\cdots+a_0^{(i)}
\]
el polinomio de Hilbert de $\mathcal L_i$, y sea $\xi=\xi_i$ un divisor que
defina $\mathcal O_X(1)$, es decir, una sección hiperplana. Las condiciones
siguientes son equivalentes:

\noindent a. $E$\footnote{El texto impreso dice «$Q$». El subconjunto
introducido en el enunciado, cuya cuasicompacidad expresa la acotación de la
familia, es $E$.} es cuasicompacto, es decir, está contenido en un abierto de
tipo finito sobre $S$; equivalentemente, la familia $(\mathcal L_i)$ es
acotada.

\noindent b. $E$ está contenido en la unión de un número finito de los
conjuntos $\operatorname{Pic}_{X/S}^Q$, $Q\in\mathbf Q[t]$; equivalentemente,
solo aparece un número finito de polinomios de Hilbert entre los
$\mathcal L_i$.

\noindent b'. (Si todas las fibras de $X/S$ tienen la misma dimensión $n$.)
Los dos coeficientes $a_{n-1}^{(i)}$ y $a_{n-2}^{(i)}$ de los polinomios de
Hilbert de los $\mathcal L_i$ recorren conjuntos finitos.

\src{243}
\noindent b''. (Si todas las fibras de $X/S$ tienen la misma dimensión $n$.)
Los enteros
\[
\xi^{n-1}D_i\qquad\text{y}\qquad\xi^{n-2}D_i^2
\]
(los grados proyectivos de $D_i$ y $D_i^2$) permanecen acotados superiormente.

\result{COROLARIO}{4.2} Sea $f:X\longrightarrow S$ un morfismo proyectivo liso
con fibras geométricamente conexas. Entonces los esquemas
$\operatorname{Pic}_{X/S}^Q$ y $\operatorname{Pic}_{X/S}^\tau$
($Q\in\mathbf Q[t]$) son proyectivos sobre $S$.

Son de tipo finito sobre $S$ por (4.1), por lo que se aplica (2.4).

\section*{BIBLIOGRAFÍA}
\addcontentsline{toc}{section}{Bibliografía}

\begin{description}
\item[{[1]}] CHEVALLEY (Claude). -- Sur la théorie de Picard,
\emph{Amer. J. of Math.}, vol. 82, 1960, pp. 435--490.

\item[{[2]}] GROTHENDIECK (A.) y DIEUDONNÉ (J.). -- Éléments de géométrie
algébrique, capítulo I y siguientes. -- París, Presses universitaires de
France, 1960--1961 (Institut des Hautes Études Scientifiques, 4, 8, 11,
\ldots).

\item[{[3]}] GROTHENDIECK (Alexander). -- Technique de descente et théorèmes
d'existence en géométrie algébrique, I--V, \emph{Séminaire Bourbaki}, t. 12,
1959/60, n.º 190, 29 pp., y n.º 195, 22 pp.; t. 13, 1960/61, n.º 212,
20 pp., y n.º 221, 28 pp.; t. 14, 1961/62, n.º 232, 19 pp.

\item[{[4]}] GROTHENDIECK (Alexander). -- Séminaire de géométrie algébrique,
1960/61. -- París, Institut des Hautes Études Scientifiques, 1961.

\item[{[5]}] GROTHENDIECK (Alexander). -- Séminaire de géométrie algébrique,
1961/62, redactado por un grupo de oyentes. -- París, Institut des Hautes Études
Scientifiques (en preparación).

\item[{[6]}] MATSUSAKA (T.). -- The criteria for algebraic equivalence and the
torsion group, \emph{Amer. J. Math.}, vol. 79, 1957, pp. 53--66.

\item[{[7]}] MUMFORD (D.) y TATE (J.). -- Séminaire de géométrie algébrique,
Harvard University, semestre de primavera de 1962 (en preparación).

\item[{[8]}] SERRE (Jean-Pierre). -- Sur la topologie des variétés algébriques
en caractéristique $p$, \emph{Symposium internacional de topología algebraica}
[1956, México], pp. 24--53. -- México, Universidad Nacional Autónoma, 1958.
\end{description}

\def\FgaUnitKey{e236}\clearpage
\chapter*{Fe de erratas y complemento a la exposición n.º 236}
\addcontentsline{toc}{chapter}{Fe de erratas y complemento a la exposición n.º 236}
\src{303}

\noindent\textbf{8. TDTE VI: Los esquemas de Picard. Propiedades generales.}

\begin{center}
[Seminario Bourbaki, t. 14, 1961/62, n.º 236, 23 pp.]
\end{center}

\noindent\textbf{Página 12, línea 16.} --- En lugar de «recta afín», léase
«recta afín privada del origen».

\noindent\textbf{Línea 18.} --- En lugar de «$X[t]$», léase
«$X[t,t^{-1}]$».

\noindent\textbf{Página 16, proposición 3.1, 6.ª línea.} --- En lugar de «Ahora
la hipótesis ...», léase «Ahora bien, de la hipótesis ...».

\noindent\textbf{Página 21, § 4: El teorema de finitud para el esquema de
Picard.} --- Las cuestiones de finitud del género de las planteadas en este
parágrafo han quedado resueltas casi por completo desde que se redactó la
presente exposición. Indiquemos los principales hechos conocidos actualmente
en esta dirección. Para simplificar los enunciados, suponemos implícitamente
que existen todos los preesquemas de Picard que intervienen en ellos, aunque
una modificación evidente de estos enunciados permite prescindir de toda
hipótesis de existencia. En lo sucesivo, $S$ designa un esquema noetheriano y
$X,Y$ esquemas propios sobre $S$.

\medskip
\noindent(i) Sea $f:X\longrightarrow Y$ un morfismo sobreyectivo. Entonces
\[
f^*: \operatorname{Pic}_{Y/S}\longrightarrow\operatorname{Pic}_{X/S}
\]
es un morfismo de tipo finito.

\src{304}
\noindent(ii) Supongamos que $Y$ es proyectivo sobre $S$ y está provisto de un
Módulo invertible relativamente amplio sobre $S$, y sea $X$ el esquema de ceros
de una sección cualquiera de este Módulo, con $f:X\longrightarrow Y$ la inmersión
canónica. Supongamos, por último, que las componentes irreducibles de las
fibras de $Y/S$ tienen dimensión $\geq3$. Entonces
\[
f^*: \operatorname{Pic}_{Y/S}\longrightarrow\operatorname{Pic}_{X/S}
\]
es de tipo finito.

\noindent(iii) Supongamos que $X$ es proyectivo sobre $S$ y que todas sus
fibras geométricas son íntegras y de dimensión $n$. Sea $\mathcal O_X(1)$ un
Módulo invertible amplio sobre $X$ que permita definir polinomios de Hilbert. Sea
$M$ una parte de $\operatorname{Pic}_{X/S}$; entonces $M$ es cuasicompacta si
y solo si, en los polinomios de Hilbert
\[
a_0x^n+a_1x^{n-1}+a_2x^{n-2}+\cdots
\]
de los elementos de $M$, los coeficientes $a_1$ y $a_2$ permanecen acotados, omitiendo todo coeficiente cuyo monomio tendría grado negativo; así, solo interviene $a_1$ si $n=1$, y ninguno de los dos si $n=0$.

\noindent(iv) Para todo entero $n\neq0$, el homomorfismo de potencia $n$-ésima
del preesquema en grupos $\operatorname{Pic}_{X/S}$ es un morfismo de tipo
finito.

Obsérvese que (i) y (ii) significan también que, bajo las hipótesis
consideradas respectivamente en esos enunciados, una parte $M$ de
$\operatorname{Pic}_{Y/S}$ es cuasicompacta si y solo si lo es su imagen en
$\operatorname{Pic}_{X/S}$. Se deduce que un Módulo invertible $\mathcal L$ sobre
$Y$ es $\tau$-equivalente a cero si y solo si lo es su imagen inversa sobre
$X$; en otros términos, $\operatorname{Pic}_{Y/S}^\tau$ es la imagen inversa
de $\operatorname{Pic}_{X/S}^\tau$. En particular, para demostrar que el
primer preesquema es de tipo finito sobre $S$, basta demostrarlo para el
segundo, puesto que $\operatorname{Pic}_{Y/S}\longrightarrow
\operatorname{Pic}_{X/S}$ es de tipo finito. Utilizando entonces (i), el lema
de Chow y (iii), se obtiene:

\medskip
\noindent(v) $\operatorname{Pic}_{X/S}^\tau$ es de tipo finito sobre $S$.

En general, la conjunción de (i) para un morfismo finito y de (ii) permite
reducir la mayoría de las cuestiones de finitud al caso en que $X/S$ tiene
fibras geométricas íntegras normales de dimensión $\leq2$; a menudo, aplicando
(i) a un morfismo sobreyectivo no necesariamente finito y la resolución de
singularidades de las superficies algebraicas (debida, en característica
arbitraria, a Abhyankar), se reduce incluso al caso en que $X/S$ es liso y,
por tanto, tiene fibras geométricas no singulares de dimensión 2. Esto permite,
por ejemplo, teniendo en cuenta (v) y la desigualdad de Picard--Igusa que
acota el rango del grupo de Néron--Severi de una superficie proyectiva no
singular, demostrar la siguiente generalización del teorema de finitud de
Néron:

\medskip
\noindent(vi) Sea $X/S$ propio sobre $S$ y por lo demás arbitrario. Entonces
los grupos de Néron--Severi
\[
\operatorname{Pic}_{X_i/k_i}/\operatorname{Pic}_{X_i/k_i}^\circ
\]
de las fibras geométricas $X_i/k_i$ de $X/S$

\src{305}
son de tipo finito, y sus rangos y los órdenes de sus subgrupos de torsión
permanecen acotados.

El mismo método de reducción al caso de superficies no singulares, junto con
los teoremas conocidos para este caso (a saber, el teorema de finitud de Néron
y el hecho de que la forma de intersección sobre el grupo de Néron--Severi es
no degenerada), implica:

\medskip
\noindent(vii) Sea $X$ un esquema propio sobre un cuerpo algebraicamente
cerrado. Existe un número finito de curvas cerradas íntegras $C_i$
($1\leq i\leq r$) en $X$ con la propiedad siguiente: para toda parte $M$ de
$\operatorname{Pic}_{X/k}$, $M$ es cuasicompacta si y solo si los enteros
\[
\deg_{C_i'}\mathcal L\qquad(\mathcal L\in M)
\]
permanecen acotados (donde $C_i'$ designa la normalización de $C_i$).

Aquí se puede tomar como $r$ el rango del grupo de Néron--Severi. Una vez
sabido que este último es de tipo finito, (vii) se reduce al hecho de que las
formas lineales sobre el grupo de Néron--Severi definidas por curvas $C$ en
$X$ solo se anulan simultáneamente sobre los elementos de torsión del grupo de
Néron--Severi. En el caso en que $X$ es proyectivo y no singular, este
resultado, así como (v), se debía a Matsusaka. Utilizando (vii), se obtiene
fácilmente la siguiente caracterización de los Módulos invertibles
$\tau$-equivalentes a cero sobre $X$:

\medskip
\noindent(viii) Sean $X/k$ un esquema propio sobre un cuerpo y $\mathcal L$ un
Módulo invertible sobre $X$. Las condiciones siguientes son equivalentes:

\noindent a. $\mathcal L$ es $\tau$-equivalente a cero;

\noindent b. Para todo Módulo coherente $F$, se tiene
\[
\chi(F\otimes\mathcal L)=\chi(F),
\]
donde $\chi$ designa la característica de Euler--Poincaré;

\noindent b'. Como (b), pero con $F=\mathcal O_Y$, donde $Y$ es un
subesquema cerrado íntegro de dimensión 1 en $X$;

\noindent c. Para todo $Y$ como el anterior, si $Y'$ designa la curva
normalizada, se tiene
\[
\deg\mathcal L_{Y'}=0.
\]

Si $X/k$ es proyectivo y está provisto de un Módulo invertible amplio
$\mathcal O_X(1)$, las condiciones precedentes son también equivalentes a las
siguientes:

\noindent d. Para todo entero $m$, $\mathcal L^{\otimes m}(1)$ es amplio.

\noindent e. (Si $X$ es íntegro.) Para todo par de enteros $m,n$, se tiene
\[
\chi(\mathcal L^{\otimes m}(n))=\chi(\mathcal O_X(n))
\]

\src{306}
(es decir, (b) es cierto tomando $F=\mathcal L^{\otimes m}(n)$).

Para la suficiencia de esta última condición, obsérvese que significa que los
polinomios de Hilbert de los $\mathcal L^{\otimes m}$ son todos iguales; por
consiguiente, en virtud del criterio de Mumford (iii), los
$\mathcal L^{\otimes m}$ permanecen en una parte cuasicompacta de
$\operatorname{Pic}_{X/k}$, es decir, se tiene (a). Las condiciones (b), (b')
y (c) deben considerarse variantes, sobre un esquema propio arbitrario, de la
noción de equivalencia numérica definida habitualmente sobre variedades
proyectivas no singulares. Para estas últimas, la equivalencia de (a) y (c)
era, naturalmente, bien conocida (Matsusaka).

El criterio (e) implica asimismo el resultado siguiente:

\medskip
\noindent(ix) Sea $f:X\longrightarrow S$ un morfismo proyectivo y plano cuyas
fibras geométricas son íntegras. Entonces $\operatorname{Pic}_{X/S}^\tau$ es
abierto y cerrado en $\operatorname{Pic}_{X/S}$.

Nos limitaremos a algunos comentarios sobre la demostración de los resultados
clave (i), (ii) y (iii) (el resultado (iv) se aparta un poco de los demás y se
demuestra utilizando únicamente (i) para los morfismos finitos sobreyectivos
radiciales, concretamente para un morfismo de Frobenius). Para (i), se
utilizan de manera esencial las ideas del descenso no plano (véase en
particular TDTE I, página 9). Resulta que, dado que solo se buscan resultados
de finitud, la falta de criterios de efectividad para datos de descenso es
inofensiva. Por otra parte, Mumford ha demostrado recientemente una forma un
poco más débil de (iii), a saber, el criterio que hace intervenir todos los
coeficientes del polinomio de Hilbert. Su argumento es extremadamente simple
y está inspirado en la demostración de un criterio de amplitud de Nakai
(enunciado por este último para las superficies no singulares y generalizado
por Mumford a morfismos proyectivos arbitrarios). Me parece, sin embargo, que
este argumento solo es válido bajo una ligera restricción adicional sobre las
fibras de $X/S$, la propiedad $(S_2)$ de Serre, que se verifica cuando las
fibras geométricas son normales. Se utiliza entonces este criterio restringido
en la demostración de (ii): el criterio (i) permite, en efecto, reducirse al
caso en que $Y/S$ es plano con fibras geométricas íntegras y normales y,
aplicando el criterio de Mumford, se reduce fácilmente al caso en que $X/S$
satisface las mismas condiciones. La hipótesis de dimensión implica entonces
que las fibras geométricas de $Y$ y de $X$ tienen profundidad $\geq2$ en sus
puntos cerrados, lo que permite aplicar los «criterios de equivalencia» en la
forma que se les da en SGA 1962 (exposición XII) y concluir la demostración de
(ii). Una vez que se dispone de (i) y (ii), no es difícil eliminar del criterio
de Mumford toda hipótesis de normalidad sobre las fibras y demostrarlo también
en la forma más fuerte dada en (iii).

\src{307}
Obsérvese, por último, que la demostración de (i) y (ii) muestra asimismo que,
cuando $S$ es el espectro de un cuerpo $k$, el morfismo
\[
\operatorname{Pic}_{Y/k}\longrightarrow\operatorname{Pic}_{X/k}
\]
es afín (y no solamente de tipo finito).

\def\FgaUnitKey{com}\unit{Comentarios de 1962}
\sid{com}
\src{297}
\seminares{14.º año, 1961/62, Suplemento\\Mayo de 1962}
\paperhead{FUNDAMENTOS DE LA GEOMETRÍA ALGEBRAICA\\COMENTARIOS}%
  {por Alexander GROTHENDIECK}

\section*{Advertencia}
\addcontentsline{toc}{section}{Advertencia}

Entre 1957 y 1962 presentamos en el Seminario Bourbaki ocho exposiciones
relativas a los Fundamentos de la geometría algebraica. Salvo la primera,
estas exposiciones están redactadas en el marco de los esquemas. Todos los
resultados enunciados encontrarán su lugar en los «Éléments de géométrie
algébrique» de Jean DIEUDONNÉ y Alexander GROTHENDIECK. Sin embargo, el
contenido de ninguna de estas exposiciones está cubierto actualmente ni por
alguno de los capítulos ya publicados o en preparación, ni por ningún otro
libro o artículo, y sin duda no lo estará todavía durante algunos años. Esta
es la razón principal que nos ha llevado a reunir estas exposiciones, que
pondrán a disposición de los lectores cierto número de nociones y resultados
clave de la teoría de esquemas mientras se espera una redacción sistemática.
Por otra parte, la lectura de estos textos permitirá familiarizarse
rápidamente con dichos resultados y nociones, sin verse entorpecido por el
detalle, necesariamente prolijo, de un tratado sistemático, y proporcionará
las motivaciones indispensables para el estudio de tal tratado.

Para comodidad del lector, hemos reunido aquí cierto número de comentarios y
erratas, agrupados por exposiciones, que señalarán en particular, de paso, los
progresos realizados desde la redacción del texto e indicarán algunas
referencias adicionales.

Diversos resultados que figuran en estas exposiciones han sido tratados en
detalle en el Seminario de geometría algebraica del IHES, así como en dos
seminarios consecutivos celebrados en Harvard en 1961/62, impartidos el
primero por mí mismo y el segundo por MUMFORD--TATE, para los cuales se están
preparando actualmente notas mimeografiadas (redactadas por el Sr.
LICHTENBAUM).

\section*{Siglas empleadas}
\addcontentsline{toc}{section}{Siglas empleadas}

\begin{description}
\item[SGA :] Seminario de geometría algebraica, Institut des Hautes Études
Scientifiques, 1960--1962.

\src{298}
\item[SHG :] Seminario GROTHENDIECK, Harvard University, 1961/62.

\item[SHMT :] Seminario MUMFORD--TATE, Harvard University, 1961/62.

\item[EGA :] GROTHENDIECK (A.) y DIEUDONNÉ (J.). --- Éléments de géométrie
algébrique. --- París, Presses universitaires de France, 1960, 1961, …
(Institut des Hautes Études Scientifiques, Publications mathématiques, 4, 8,
11, …).

\item[TDTE :] GROTHENDIECK (Alexander). --- Technique de descente et théorèmes
d'existence en géométrie algébrique, I--VI, \emph{Séminaire Bourbaki}, t. 12,
1959/60, n.os 190 y 195; t. 13, 1960/61, n.os 212 y 221; t. 14, 1961/62,
n.os 232 y 236.
\end{description}

\end{document}
